\documentclass[11pt,oneside]{article}

% =====================================================================
%  Exact Dyadic Extraction of Rvachev's Up-Function
%  from Finite Sinc-Product Splines
%  Canonical consolidated volume.  Editorial merge (2026-09-03) of the six
%  independently written reports filed on 2026-09-02; see the provenance
%  appendix for the absorbed sources and their receipts.
%
%  House style follows the canonical corpus volumes: A4, Libertinus with an
%  lmodern fallback, article class carrying \part/\section/\subsection.
%  All mathematical notation is normative from
%  docs/FabiusFunction_Mathematical_Notation_Catalogue.
% =====================================================================

\usepackage[a4paper,margin=26mm,headheight=15pt]{geometry}
\usepackage[T1]{fontenc}
\usepackage[utf8]{inputenc}
\IfFileExists{libertinus.sty}{\usepackage{libertinus}}{\usepackage{lmodern}}
\usepackage{microtype}
\usepackage{amsmath,amssymb,amsthm,mathtools,mathrsfs,bm}
% Canonical mathematical notation for active FabiusFunction LaTeX documents.
%
% This file is intentionally a plain TeX input rather than a LaTeX package.
% Every standalone document loads it by an explicit path relative to that
% document's directory; included fragments inherit the definitions from their
% root document.  Keep this file independent of document layout and metadata.
%
% Contract:
%   * one command has one mathematical meaning throughout the active corpus;
%   * names encode semantic roles rather than a document-local choice of letter;
%   * visually similar but differently normalized objects have different print;
%   * every command below is defined normatively in the notation catalogue.
%
% The active corpus excludes docs/archive/.  Historical sources may retain
% their original notation only inside an explicitly labelled quotation.

\makeatletter
\@ifundefined{mathbb}{%
  \PackageError{fabius-notation}{The amssymb package must be loaded first}%
    {Load amsmath and amssymb before inputting fabius-notation.tex.}%
}{}
\@ifundefined{operatorname}{%
  \PackageError{fabius-notation}{The amsmath package must be loaded first}%
    {Load amsmath before inputting fabius-notation.tex.}%
}{}
\makeatother

% Version stamp used by the catalogue and by static audits.
\newcommand{\FabiusNotationVersion}{2026-08-31}

% Number systems.  NaturalNumbers is explicitly the nonnegative convention.
\newcommand{\NaturalNumbers}{\mathbb{N}_{0}}
\newcommand{\PositiveIntegers}{\mathbb{N}_{>0}}
\newcommand{\IntegerNumbers}{\mathbb{Z}}
\newcommand{\RationalNumbers}{\mathbb{Q}}
\newcommand{\RealNumbers}{\mathbb{R}}
\newcommand{\ComplexNumbers}{\mathbb{C}}
\newcommand{\ExtendedRealNumbers}{\overline{\mathbb{R}}}

% The three principal Fabius--Rvachev functions.  Subscripts are deliberate:
% a bare F and a calligraphic F have several unrelated uses in the corpus.
\newcommand{\FabiusBounded}{F_{\mathrm{bd}}}
\newcommand{\FabiusGlobal}{F_{\mathrm{gl}}}
\newcommand{\FabiusQuantile}{Q_{F}}
\newcommand{\FabiusClampedQuantile}{Q_{F,\mathrm{cl}}}
\newcommand{\RvachevUp}{\operatorname{up}}

% Constants, elementary operators, and delimiters.
\newcommand{\EulerE}{\mathrm{e}}
\newcommand{\ImaginaryUnit}{\mathrm{i}}
\newcommand{\Differential}{\mathop{}\!\mathrm{d}}
\newcommand{\AbsoluteValue}[1]{\left\lvert #1\right\rvert}
\newcommand{\Norm}[1]{\left\lVert #1\right\rVert}
\newcommand{\Floor}[1]{\left\lfloor #1\right\rfloor}
\newcommand{\Ceiling}[1]{\left\lceil #1\right\rceil}
\newcommand{\FractionalPart}[1]{\left\{#1\right\}}
\newcommand{\PositivePart}[1]{\left(#1\right)_{+}}
\newcommand{\IndicatorOf}[1]{\mathbf{1}_{#1}}
\newcommand{\SetOf}[1]{\left\{#1\right\}}
\newcommand{\InnerProduct}[1]{\left\langle #1\right\rangle}
\newcommand{\CardinalityOf}[1]{\left\lvert #1\right\rvert}
\newcommand{\DefinitionEquals}{\mathrel{\coloneqq}}
\newcommand{\Sign}{\operatorname{sgn}}
\newcommand{\IdentityOperator}{\operatorname{id}}
\newcommand{\SpanOperator}{\operatorname{span}}
\newcommand{\TraceOperator}{\operatorname{tr}}
\newcommand{\RankOperator}{\operatorname{rank}}
\newcommand{\DiagonalOperator}{\operatorname{diag}}
\newcommand{\ResidueOperator}{\operatorname*{Res}}
\newcommand{\LeastCommonMultiple}{\operatorname{lcm}}
\newcommand{\OrderOperator}{\operatorname{ord}}
\newcommand{\PrincipalLogarithm}{\operatorname{Log}}
\newcommand{\FallingFactorial}[2]{\left(#1\right)^{\underline{#2}}}
\newcommand{\RisingFactorial}[2]{\left(#1\right)^{\overline{#2}}}

% Support, topology, convolution, and metric notation.
\newcommand{\SupportOperator}{\operatorname{supp}}
\newcommand{\SupportOf}[1]{\SupportOperator\!\left(#1\right)}
\newcommand{\TopologicalSupportOf}[1]{\operatorname{tsupp}\!\left(#1\right)}
\newcommand{\Convolution}{\mathbin{*}}
\newcommand{\MetricDistance}{\operatorname{dist}}
\newcommand{\TotalVariationDistance}{d_{\mathrm{TV}}}

% Probability.  Equality and convergence in law are not metric distance.
\newcommand{\Probability}{\mathbb{P}}
\newcommand{\Expectation}{\mathbb{E}}
\newcommand{\Variance}{\operatorname{Var}}
\newcommand{\Covariance}{\operatorname{Cov}}
\newcommand{\LawOperator}{\operatorname{Law}}
\newcommand{\LawOf}[1]{\LawOperator\!\left(#1\right)}
\newcommand{\EqualInLaw}{\mathrel{\overset{\mathrm{law}}{=}}}
\newcommand{\ConvergesInLaw}{\mathrel{\xrightarrow{\mathrm{law}}}}

% Fourier and integral transforms.  The printed labels expose normalization.
% FourierTwoPi uses exp(-2*pi*i*x*xi); FourierAngular uses exp(-i*x*omega).
\newcommand{\FourierTwoPi}[1]{\widehat{#1}_{\,2\pi}}
\newcommand{\FourierAngular}[1]{\widehat{#1}_{\,\mathrm{ang}}}
\newcommand{\LaplaceTransformSymbol}{\mathcal{L}}
\newcommand{\MellinTransformSymbol}{\mathcal{M}}
\newcommand{\LaplaceTransformOf}[1]{\LaplaceTransformSymbol\!\left[#1\right]}
\newcommand{\MellinTransformOf}[1]{\MellinTransformSymbol\!\left[#1\right]}
\newcommand{\MomentGeneratingFunction}[1]{M_{#1}}
\newcommand{\CharacteristicFunction}[1]{\varphi_{#1}}

% The two standard sinc functions must never share one symbol.
\newcommand{\SincRad}{\operatorname{sinc}_{\mathrm{rad}}}
\newcommand{\SincPi}{\operatorname{sinc}_{\pi}}

% Binary arithmetic and Thue--Morse notation.
\newcommand{\BinaryDigitSum}{\operatorname{wt}_{2}}
\newcommand{\ThueMorseSignSymbol}{\varepsilon}
\newcommand{\ThueMorseSign}[1]{\varepsilon_{#1}}
\newcommand{\PadicValuation}[1]{\nu_{#1}}
\newcommand{\TwoAdicValuation}{\nu_{2}}

% q-series notation.  Every base is an explicit argument.
\newcommand{\QPochhammer}[3]{\left(#1;#2\right)_{#3}}
\newcommand{\GaussianBinomial}[3]{%
  \genfrac{[}{]}{0pt}{}{#1}{#2}_{#3}%
}
\newcommand{\QInteger}[2]{\left[#1\right]_{#2}}

% Named special functions.
\newcommand{\Polylogarithm}{\operatorname{Li}}
\newcommand{\LambertW}{W}
\newcommand{\GammaFunction}{\Gamma}
\newcommand{\RiemannZeta}{\zeta}

% Asymptotics.  BigO and LittleO are operator tokens for existing formula
% grammar; prefer the At forms so that the limiting regime is printed.
\newcommand{\BigO}{\mathcal{O}}
\newcommand{\LittleO}{o}
\newcommand{\BigOAt}[2]{\mathcal{O}_{#2}\!\left(#1\right)}
\newcommand{\LittleOAt}[2]{o_{#2}\!\left(#1\right)}

\usepackage{booktabs,longtable,tabularx,array,multirow}
\usepackage{graphicx}
\usepackage{xcolor}
\usepackage{enumitem}
\usepackage{fancyhdr}
\usepackage{titlesec}
\usepackage{float}
\usepackage{xurl}
\usepackage{listings}
\usepackage[most]{tcolorbox}
\usepackage{hyperref}
\usepackage[nameinlink,noabbrev,capitalise]{cleveref}

\definecolor{linkblue}{RGB}{24,72,124}
\definecolor{deepgreen}{RGB}{23,102,69}
\definecolor{deepred}{RGB}{140,42,42}
\definecolor{deepgold}{RGB}{122,91,19}
\definecolor{shadegray}{RGB}{246,247,249}
\definecolor{rulegray}{RGB}{195,201,208}
\definecolor{palered}{RGB}{251,236,236}
\definecolor{palegold}{RGB}{251,245,230}
\definecolor{palegreen}{RGB}{233,245,243}

\hypersetup{
  hypertexnames=false,% ed.: avoid duplicate destinations across parts
  colorlinks=true,
  linkcolor=linkblue,
  citecolor=deepgreen,
  urlcolor=linkblue,
  pdftitle={Exact Dyadic Extraction of Rvachev's Up-Function from Finite Sinc-Product Splines},
  pdfauthor={Vladimir Reshetnikov; ProveIt formal-development synthesis},
  pdfsubject={Consolidated volume: exact eventual geometric tails of finite sinc-product splines at dyadic points and the quarter-base Gaussian-binomial extraction of up-function values},
  pdfkeywords={Rvachev up-function, Fabius function, finite sinc product, box spline, dyadic point, geometric tail, q-binomial, q-Pochhammer, Richardson extrapolation, Thue-Morse}
}

\pagestyle{fancy}
\fancyhf{}
\fancyhead[L]{\small Exact dyadic extraction of the up-function}
\fancyhead[R]{\small\nouppercase{\leftmark}}
\fancyfoot[C]{\thepage}
\renewcommand{\headrulewidth}{0.4pt}

\setcounter{tocdepth}{2}
\setcounter{secnumdepth}{3}
\setlist[itemize]{leftmargin=1.45em,itemsep=0.25ex,topsep=0.55ex}
\setlist[enumerate]{leftmargin=1.85em,itemsep=0.35ex,topsep=0.55ex}
\allowdisplaybreaks[2]

\theoremstyle{plain}
\newtheorem{theorem}{Theorem}[section]
\newtheorem{proposition}[theorem]{Proposition}
\newtheorem{lemma}[theorem]{Lemma}
\newtheorem{corollary}[theorem]{Corollary}
\theoremstyle{definition}
\newtheorem{definition}[theorem]{Definition}
\newtheorem{example}[theorem]{Example}
\newtheorem{algorithm}[theorem]{Algorithm}
\newtheorem{exercise}[theorem]{Exercise}
\theoremstyle{remark}
\newtheorem{remark}[theorem]{Remark}

\crefname{exercise}{Exercise}{Exercises}
\Crefname{exercise}{Exercise}{Exercises}
\crefname{algorithm}{Algorithm}{Algorithms}
\Crefname{algorithm}{Algorithm}{Algorithms}

\newtcolorbox{warningbox}[1][]{%
  enhanced,breakable,colback=palered,colframe=deepred,
  boxrule=0.6pt,arc=1.2mm,left=1.8mm,right=1.8mm,top=1.2mm,bottom=1.2mm,
  title={A common trap},fonttitle=\bfseries,#1}
\newtcolorbox{summarybox}[1][]{%
  enhanced,breakable,colback=palegold,colframe=deepgold,
  boxrule=0.6pt,arc=1.2mm,left=1.8mm,right=1.8mm,top=1.2mm,bottom=1.2mm,
  title={Section summary},fonttitle=\bfseries,#1}
\newtcolorbox{checkpointbox}[1][]{%
  enhanced,breakable,colback=palegreen,colframe=deepgreen,
  boxrule=0.6pt,arc=1.2mm,left=1.8mm,right=1.8mm,top=1.2mm,bottom=1.2mm,
  title={Structural checkpoint},fonttitle=\bfseries,#1}

\lstdefinestyle{pseudo}{%
  basicstyle=\ttfamily\small,columns=fullflexible,frame=single,
  rulecolor=\color{rulegray},backgroundcolor=\color{shadegray},
  showstringspaces=false,keepspaces=true,
  xleftmargin=0.4em,xrightmargin=0.4em,aboveskip=0.8em,belowskip=0.8em}
\lstset{style=pseudo}

% ---------------------------------------------------------------------
%  Document-local shorthands.  Each expands to catalogue notation or to an
%  object this volume defines in its setting section; none introduces a
%  meaning the text does not define.  Recorded in the notation appendix.
% ---------------------------------------------------------------------
\newcommand{\Kk}{\RationalNumbers}                          % coefficient field
% LOCAL-PREFIX-SPLINE: \upn{n} is the finite sinc-product spline with n+1
% factors k = 0..n (the six sources' shared convention); \upn{n} = p_{n+1} in
% the Exponents volume and corresponds to fabiusUniformSpline n in Lean.
\newcommand{\upn}[1]{\RvachevUp_{#1}}
\newcommand{\Phin}[1]{\Phi_{#1}}                            % its Fourier transform
\newcommand{\PhiUp}{\Phi}                                   % the infinite product
\newcommand{\dep}[1]{s\!\left(#1\right)}                    % dyadic depth
\newcommand{\Qq}{Q}                                         % the ratio base 1/4
\newcommand{\am}[1]{a_{#1}}                                 % reciprocal-product coefficients
\newcommand{\Cr}[2]{C_{#1}\!\left(#2\right)}                % r-th mode coefficient at x
\newcommand{\qn}[1]{q_{#1}}                                 % the sample \upn{n}(x)
\newcommand{\Wrow}[2]{w^{(#1)}_{#2}}                        % extraction weight
\newcommand{\Trunc}[2]{\operatorname{Trunc}_{<#2}#1}

\begin{document}
\title{\textbf{Exact Dyadic Extraction of Rvachev's Up-Function}\\[0.35em]
  \Large from Finite Sinc-Product Splines}
\author{Vladimir Reshetnikov\\
  \small ProveIt formal-development synthesis}
\date{3 September 2026}

\maketitle

\begin{abstract}
Let \(\PhiUp(t)=\prod_{k\ge0}\SincRad(\pi t/2^{k})\) be the Fourier transform of
Rvachev's up-function under the \(2\pi\) convention, and let \(\upn{n}\) be the
inverse transform of the prefix product with factors \(k=0,\ldots,n\): a
compactly supported piecewise-polynomial box spline of degree at most \(n\)
whose values at dyadic rationals are rational.  Fix a dyadic rational
\(x\in[-1,1]\), let \(\dep{x}\) be the least \(s\ge0\) with \(2^{s}x\in
\IntegerNumbers\), and put \(d=\Floor{\dep{x}/2}\).

The volume proves that the convergence \(\upn{n}(x)\to\RvachevUp(x)\) is not
merely asymptotic: for every \(n\ge\dep{x}\),
\[
 \upn{n}(x)=\RvachevUp(x)+\sum_{r=1}^{d}\Cr{r}{x}\,4^{-rn},
 \qquad
 \Cr{r}{x}=(-1)^{r}\am{r}4^{-r}\RvachevUp^{(2r)}(x),
\]
exactly, with no remainder, where the positive rationals \(\am{r}\) are the
coefficients of \(1/\PhiUp\).  Three finite mechanisms drive the proof: the
omitted Fourier tail is a scaled copy of \(\RvachevUp\) supported in a single
polynomial cell of the finite spline; on that cell the all-orders
deconvolution series is a polynomial in a nilpotent operator and terminates;
and every derivative \(\RvachevUp^{(r)}(x)\) with \(r>\dep{x}\) vanishes.  At
the last level before the onset, where \(x\) is a knot of the spline, the
defect against the law is computed in closed form through Rvachev's
functional-differential equation: it is \(-\ThueMorseSign{k}\,2^{-\binom{s}{2}}\,B_{s}/s!\)
with \(s=\dep{x}\), \(\ThueMorseSign{k}\) the Thue--Morse sign of the knot and
\(B_{s}\) the Bernoulli number, so the onset \(n\ge\dep{x}\) is exact at every
even depth and improves by exactly one level at odd depth; all \(d\) modes are
nonzero when \(\dep{x}\ge2\), so \(d\) is the exact number of modes.

Geometric Lagrange interpolation at the nodes \(1,\Qq,\ldots,\Qq^{d}\),
\(\Qq=1/4\), turns the law into a single extraction row: for any
\(N\ge\dep{x}\),
\[
 \RvachevUp(x)=\frac{1}{\QPochhammer{\Qq}{\Qq}{d}}
 \sum_{j=0}^{d}(-1)^{d-j}\Qq^{\binom{d-j+1}{2}}
 \GaussianBinomial{d}{j}{\Qq}\,\upn{N+j}(x),
\]
recovering the exact value from \(d+1\) consecutive rational samples.  Every
mode coefficient, and hence every even derivative of \(\RvachevUp\) at \(x\),
is recovered by the same interpolation; the row's \(\ell^{\infty}\)
amplification is bounded by \((-\Qq;\Qq)_{\infty}/(\Qq;\Qq)_{\infty}
=1.96926\ldots\).  Against the corpus's Exponents volume, whose signed
quarter-base Richardson combination cancels these modes asymptotically for
general \(x\), the contribution here is exactness at dyadic points after the
depth.  Worked examples, an exact-arithmetic algorithm with its verification
specification, and a formalization register --- the row's algebra and the cell
identity at \(x=1/4\) are kernel-verified in Lean; the general theorem is not
--- complete the volume.
\end{abstract}

\tableofcontents
\clearpage

% ---- fragment 01-setting ----
% ---------------------------------------------------------------------
%  Cluster: Setting -- finite sinc-product splines, dyadic depth, the
%  result in one page.  Label slug: set.
%  Sources: S1 sections 1-3, S2 sections 1-3, S3 sections 1-3,
%           S4 sections 1-3, S5 sections 1-3, S6 sections 1-2.
% ---------------------------------------------------------------------

\section{The result in one page}
\label{due:sec:result}

Rvachev's up-function $\RvachevUp$ \cite{rvachev1990,arias1982} is the
even, nonnegative $C^\infty$ probability density with support $[-1,1]$
whose $2\pi$-frequency Fourier transform is the dyadic sinc product
\begin{equation}
 \PhiUp(\xi)\DefinitionEquals\FourierTwoPi{\RvachevUp}(\xi)
 =\prod_{k=0}^{\infty}\SincRad\!\left(\frac{\pi\xi}{2^{k}}\right),
 \qquad \xi\in\RealNumbers
 \label{due:eq:set-infinite-product}
\end{equation}
(catalogue entries \texttt{rvachev-up} and \texttt{up-fourier-product};
both facts are Lean-proved in the repository \cite{proveit-primary}).
Keeping only the factors $k=0,\ldots,n$ produces the finite product
$\Phin{n}$ and its inverse transform $\upn{n}$, a compactly supported
piecewise polynomial of degree $n$ with knots on a dyadic lattice
(\cref{due:sec:setting}).  At a rational point every value $\upn{n}(x)$ is
a rational number that can be written down exactly, and
$\upn{n}(x)\to\RvachevUp(x)$ as $n\to\infty$.

\paragraph{The problem.}
Fix a dyadic rational $x\in[-1,1]$.  Recover $\RvachevUp(x)$
\emph{exactly}, as an element of $\Kk=\RationalNumbers$, from finitely
many of the finite-spline values $\qn{n}\DefinitionEquals\upn{n}(x)$, by
one closed formula whose weights do not depend on $x$.

\paragraph{The two integers and the two constants.}
Throughout, $\dep{x}$ is the \emph{dyadic depth} of $x$, the least
$s\ge0$ with $2^{s}x\in\IntegerNumbers$ (\cref{due:def:set-depth}), and
\begin{equation}
 d\DefinitionEquals\Floor{\dep{x}/2},
 \qquad
 \Qq\DefinitionEquals\frac14 .
 \label{due:eq:set-d-and-Q}
\end{equation}
The \emph{reciprocal-product coefficients} $\am{m}\in\Kk$ are defined by
the Taylor expansion at the origin
\begin{equation}
 \frac{1}{\PhiUp(z)}=\sum_{m=0}^{\infty}\am{m}\,(2\pi z)^{2m},
 \qquad
 \am{0}=1,\ \am{1}=\frac1{18},\ \am{2}=\frac{31}{16200},\
 \am{3}=\frac{2347}{42865200};
 \label{due:eq:set-am-definition}
\end{equation}
they are positive rationals, and the section that introduces them proves
this, identifies them with the catalogue's Rvachev--Appell coefficients,
and gives the Bernoulli--Bell formula behind the listed values.

\begin{theorem}[Exact eventual geometric law at a dyadic point]
\label{due:thm:set-main}
Let $x\in[-1,1]$ be a dyadic rational with depth $s=\dep{x}$ and let
$d=\Floor{s/2}$.  Then for every integer $n\ge s$,
\begin{equation}
 \boxed{\;
 \upn{n}(x)=\RvachevUp(x)+\sum_{r=1}^{d}\Cr{r}{x}\,4^{-rn},
 \qquad
 \Cr{r}{x}=(-1)^{r}\,\am{r}\,4^{-r}\,\RvachevUp^{(2r)}(x),\;}
 \label{due:eq:set-main-law}
\end{equation}
equivalently
\begin{equation}
 \upn{n}(x)=\sum_{r=0}^{d}(-1)^{r}\,\am{r}\,\RvachevUp^{(2r)}(x)\,
 4^{-r(n+1)}
 \qquad(n\ge s).
 \label{due:eq:set-main-law-shifted}
\end{equation}
Every quantity in \cref{due:eq:set-main-law} lies in $\Kk$.  The identity
is exact: there is no remainder term once $n\ge s$.
\end{theorem}

The two displays are the same statement, because
$\am{r}4^{-r(n+1)}=(\am{r}4^{-r})\,4^{-rn}$; the sources split on which of
the two factors $4^{-r}$ is absorbed into the coefficient, and this volume
always prints $\Cr{r}{x}$ with the factor $4^{-r}$ absorbed, so that the
modes of the \emph{sequence} $(\qn{n})_n$ are literally
$4^{-n},4^{-2n},\ldots,4^{-dn}$
(\cref{due:rmk:set-exponent-conventions}).
% ed.: the proof of the main theorem is distributed over the volume; the
% ed.: section references below were attached at assembly.
The proof of \cref{due:thm:set-main} occupies the body of the volume.  Its
three ingredients are geometric and finite: the omitted Fourier tail is an
exactly rescaled copy of $\RvachevUp$ (\cref{due:prop:set-head-tail}); for
$n\ge s$ that copy is supported inside one polynomial cell of $\upn{n}$
centred at $x$ (\cref{due:lem:set-midpoint,due:lem:set-tail-one-piece}),
so that deconvolution becomes finite linear algebra on a polynomial jet;
and the Taylor jet of $\RvachevUp$ at a dyadic point terminates at order
$s$, which is why only $d=\Floor{s/2}$ even orders survive.  The reciprocal-product coefficients and finite local deconvolution
(\cref{due:sec:coefficients}), the termination of the dyadic jet
(\cref{due:sec:jet}), and the eventual geometric law itself
(\cref{due:sec:main}) assemble these into \cref{due:eq:set-main-law};
its sharpness --- the exact onset and the exact number of modes --- is
\cref{due:sec:sharpness}.

The Exponents volume \cite{proveit-exponents} proves, for its prefix
product $p_n$ (defined in its equation
\texttt{p3:eq:prefix-\allowbreak product}), the sharp $C^{r}$ error law
\[
 \Norm{p_n^{(r)}-\RvachevUp^{(r)}}_\infty=\frac{2^{\binom{r+3}{2}-1}}{9\cdot4^{n}}
 \qquad(\text{its }\texttt{p3:thm:sharp-prefix-law}),
\]
and a signed geometric Richardson combination with quarter-base Lagrange
weights whose leading term is only \emph{asymptotic}, with an
$\BigOAt{\delta_n^{2M+2}}{n\to\infty}$ remainder at a general point (its
equations \texttt{p3:eq:richardson-\allowbreak combination} and, with the same stem,
\texttt{-weights}, \texttt{-moments}, \texttt{-leading},
\texttt{-stability}).  The contribution of the present volume is that
\emph{at a dyadic point of depth $s$ the expansion is exact for every
$n\ge s$}, so the very same Richardson row returns the exact value rather
than an accelerated approximation.  That is the content of the next
theorem.

\begin{theorem}[One-row extraction]
\label{due:thm:set-extraction}
Let $x$, $s$, $d$ be as in \cref{due:thm:set-main} and let $N\ge s$ be any
integer.  Then
\begin{equation}
 \boxed{\;
 \RvachevUp(x)=\frac{1}{\QPochhammer{\Qq}{\Qq}{d}}
 \sum_{j=0}^{d}(-1)^{d-j}\,\Qq^{\binom{d-j+1}{2}}
 \GaussianBinomial{d}{j}{\Qq}\,\upn{N+j}(x),
 \qquad \Qq=\frac14 .\;}
 \label{due:eq:set-extraction-row}
\end{equation}
Reindexing by $i=d-j$ gives the same row in the order in which S6 prints
it,
\begin{equation}
 \RvachevUp(x)=\frac{1}{\QPochhammer{\Qq}{\Qq}{d}}
 \sum_{i=0}^{d}(-1)^{i}\,\Qq^{\binom{i+1}{2}}
 \GaussianBinomial{d}{i}{\Qq}\,\qn{N+d-i},
 \label{due:eq:set-extraction-row-reversed}
\end{equation}
since $\GaussianBinomial{d}{d-i}{\Qq}=\GaussianBinomial{d}{i}{\Qq}$.  The
$d+1$ weights are rational, sum to $1$, and depend on neither $x$ nor $N$;
they depend on $x$ only through $d=\Floor{\dep{x}/2}$.
\end{theorem}

% ed.: this section proved the row in full by geometric Lagrange interpolation,
% ed.: and the extraction section proves the identical theorem by the identical
% ed.: argument.  One proof: the statement stays here, the proof lives there.
\begin{proof}
Deferred to \cref{due:thm:ext-extraction}: the row is geometric Lagrange
interpolation at the nodes $\Qq^{N},\ldots,\Qq^{N+d}$ extrapolated to
$0$, applied to the degree-$\le d$ polynomial that \cref{due:thm:set-main}
supplies; the closed form of the weights is \cref{due:thm:ext-weights}, and
the reversed order is the substitution $i=d-j$ together with the symmetry
$\GaussianBinomial{d}{d-i}{\Qq}=\GaussianBinomial{d}{i}{\Qq}$.
\end{proof}

\begin{remark}[Empty products at $d=0$]
\label{due:rmk:set-d-zero}
For $s\in\{0,1\}$ one has $d=0$; the sum in \cref{due:eq:set-main-law} is
empty, $\QPochhammer{\Qq}{\Qq}{0}=1$ and
$\GaussianBinomial{0}{0}{\Qq}=1$ by the catalogue's empty-product
conventions, and the row \cref{due:eq:set-extraction-row} reads
$\RvachevUp(x)=\upn{N}(x)$ for every $N\ge s$.  This is not vacuous: it
asserts that the finite splines are already exact at the five points
$0,\pm\tfrac12,\pm1$ from level $s$ on.  \Cref{due:ex:set-low-depth}
verifies it directly, with a convention-dependent exception at the single
level $n=0$ for $x=\pm\frac12$.
\end{remark}

\begin{remark}[The two exponent conventions, reconciled]
\label{due:rmk:set-exponent-conventions}
S1, S2 and S6 write the eventual law with modes $4^{-r(n+1)}$ and
coefficients $(-1)^{r}\am{r}\RvachevUp^{(2r)}(x)$ (S1 in the equivalent
form $\gamma_{2r}\RvachevUp^{(2r)}(x)/(2r)!$); S3, S4 and S5 write modes
$4^{-rn}$ with coefficients $\Cr{r}{x}$.  Both are correct and differ by
the factor $4^{-r}$, as \cref{due:eq:set-main-law,due:eq:set-main-law-shifted}
show.  The factor has a geometric meaning: $4^{-(n+1)}=\delta_n^{2}$ is the
square of the half-width $\delta_n=2^{-(n+1)}$ of the omitted tail
(\cref{due:prop:set-head-tail}), so the shifted form is the one that
deconvolution produces, while the unshifted form is the one that names the
ratios of the sequence.  A reader comparing a coefficient table across
sources must first decide which of the two is tabulated: the mode
coefficient $4/9$ at $x=\frac14$ in S1 and S2 is the coefficient
$\Cr{1}{\frac14}=\frac19$ of this volume.  The same shift appears once
more, disguised, in the indexing bridge to the Exponents volume
(\cref{due:sec:setting}): there $\delta_N=2^{-N}$ and the modes are
$4^{-rN}$ with $N=n+1$.
\end{remark}

\begin{remark}[Generality]
\label{due:rmk:set-generality}
Nothing in the proof of \cref{due:thm:set-extraction} used $\Qq=\frac14$.
For any base $q\in(0,1)$, any polynomial $P$ of degree at most $d$, and any
$N$,
\[
 P(0)=\frac{1}{\QPochhammer{q}{q}{d}}\sum_{j=0}^{d}(-1)^{d-j}q^{\binom{d-j+1}{2}}
 \GaussianBinomial{d}{j}{q}\,P(q^{N+j}).
\]
What is special to
$\RvachevUp$ at a dyadic point is \cref{due:thm:set-main}: that the
sequence $\qn{n}$ \emph{is} such a polynomial in $q^{n}$ with $q=\frac14$,
exactly and from level $s$ on.
\end{remark}

\begin{example}[The worked instance $x=\frac14$]
\label{due:ex:set-quarter}
Here $2^{2}\cdot\frac14=1$ and $2\cdot\frac14\notin\IntegerNumbers$, so
$s=2$ and $d=1$: one geometric mode, two samples, onset $n=2$.  The
truncated-power formula of \cref{due:thm:set-finite-spline} gives the
samples at the levels $n=2,3$ in exact arithmetic,
\[
 \qn{2}=\upn{2}\!\left(\tfrac14\right)=\frac{15}{16},
 \qquad
 \qn{3}=\upn{3}\!\left(\tfrac14\right)=\frac{179}{192}
\]
(the level-$2$ value is the five-term sum
$4\,[(\tfrac98)^2-(\tfrac78)^2-(\tfrac58)^2+(\tfrac38)^2-(\tfrac18)^2]
=4\cdot\tfrac{15}{64}$; the level-$3$ value is the ten-term sum
$\tfrac1{48}\sum_{k\le9}\ThueMorseSign{k}(\tfrac{19}{2}-k)^{3}
=\tfrac1{48}\cdot\tfrac{179}{4}$ in the scaled form
\cref{due:eq:set-scaled-cell}).  With $d=1$ the row
\cref{due:eq:set-extraction-row} has $\QPochhammer{\Qq}{\Qq}{1}=\frac34$
and weights $(-\Qq,\,1)/\frac34=\left(-\tfrac13,\ \tfrac43\right)$, so
\[
 \RvachevUp\!\left(\tfrac14\right)
 =-\frac13\cdot\frac{15}{16}+\frac43\cdot\frac{179}{192}
 =-\frac{45}{144}+\frac{179}{144}
 =\frac{134}{144}
 =\frac{67}{72},
\]
which is the classical value $\FabiusBounded(\tfrac34)=\tfrac{67}{72}$
\cite{haugland2016,arias2018}.  The single mode is
$\Cr{1}{\frac14}=\qn{2}\cdot16-\tfrac{67}{72}\cdot16=\tfrac19$, that is
\[
 \upn{n}\!\left(\tfrac14\right)=\frac{67}{72}+\frac19\,4^{-n}
 =\frac{67}{72}+\frac49\,4^{-(n+1)}
 \qquad(n\ge2),
\]
and the analytic formula for the coefficient agrees: the differential law
$\RvachevUp'(x)=2\RvachevUp(2x+1)-2\RvachevUp(2x-1)$ differentiated once
more gives
$\RvachevUp''(\tfrac14)=8[\RvachevUp(4)-\RvachevUp(2)]
-8[\RvachevUp(0)-\RvachevUp(-2)]=-8$, and
$\Cr{1}{\frac14}=(-1)^{1}\am{1}4^{-1}\RvachevUp''(\tfrac14)
=\tfrac1{18}\cdot\tfrac14\cdot8=\tfrac19$.  Two independent checks: the
level $n=4$, which the row did not use, gives
$\upn{4}(\tfrac14)=\tfrac{715}{768}=\tfrac{67}{72}+\tfrac19\cdot\tfrac1{256}$
exactly; and the level $n=1=s-1$, at which $\frac14$ is a knot of
$\upn{1}$ (\cref{due:lem:set-midpoint}), gives $\upn{1}(\tfrac14)=1$,
whereas the law would predict $\tfrac{67}{72}+\tfrac1{36}=\tfrac{23}{24}$.
The onset $n\ge s$ is therefore sharp at this point.  All six packages'
verifiers reproduce these rationals (S1, S2 and S5 print them; S3 and S4
include $x=\frac14$ in their exhaustive sweeps).
\end{example}

\paragraph{A map of the volume.}
% ed.: section titles below are descriptive; the orchestrator should
% ed.: replace them by \cref's to the assembled parts.
\Cref{due:sec:setting} fixes the Fourier normalization, defines
$\Phin{n}$ and $\upn{n}$, proves the finite convolution model, the exact
Thue--Morse truncated-power formula, rationality and strict positivity,
and states the indexing bridge to the Exponents volume and to Lean once
and for all.  \Cref{due:sec:depth} defines the depth, proves that a dyadic
point is the centre of one polynomial cell of $\upn{n}$ exactly when
$n\ge\dep{x}$ and is a knot at $n=\dep{x}-1$, and proves that the omitted
tail sees a single polynomial piece.  The sections that follow introduce
the coefficients $\am{m}$ (with the identification
$\am{m}=(-1)^{m}\gamma_{2m}/(2m)!$ of S1's moment-generating-function
normalization and the Bernoulli--Bell formula), invert the tail
convolution on a finite polynomial jet, prove that the dyadic Taylor jet
of $\RvachevUp$ terminates at order $\dep{x}$, and combine the three to
obtain \cref{due:thm:set-main}, including the statement that every one of
the $d$ modes is nonzero for $s\ge2$ and the one-stage sharpening at odd
depth.  The extraction part proves \cref{due:thm:set-extraction} a second
time as the first row of an inverse geometric Vandermonde system,
recovers every $\Cr{r}{x}$ and every even derivative from the same
samples, gives the integer base-$4$ form, the annihilating $q$-difference
recurrence, the Romberg-style triangular table, the minimality of $d+1$
samples and the uniform stability bound on the weights.  The remaining
parts give worked exact examples through depth seven, an exact algorithm
with its verification suite, structural consequences, and the appendices
with the finite $q$-binomial identities, the notation concordance, the
repair log and the formalization register.

\section{Fourier normalization and the finite splines}
\label{due:sec:setting}

\subsection{The \texorpdfstring{$2\pi$}{2pi} convention and the two sinc normalizations}
\label{due:sec:set-fourier}

All transforms in this volume are $2\pi$-frequency transforms in the sense
of the catalogue entry \texttt{fourier-two-pi}: for $f\in L^{1}(\RealNumbers)$,
\begin{equation}
 \FourierTwoPi{f}(\xi)\DefinitionEquals
 \int_{\RealNumbers}f(x)\,\EulerE^{-2\pi\ImaginaryUnit x\xi}\,\Differential x,
 \qquad
 f(x)=\int_{\RealNumbers}\FourierTwoPi{f}(\xi)\,
 \EulerE^{2\pi\ImaginaryUnit x\xi}\,\Differential\xi,
 \label{due:eq:set-fourier-convention}
\end{equation}
the inversion formula being asserted only where it holds classically
(here: whenever $\FourierTwoPi{f}\in L^{1}$ and $f$ is continuous).  There
is no scalar prefactor in either direction, convolution goes to the
pointwise product, $\FourierTwoPi{f\Convolution g}=\FourierTwoPi{f}\,
\FourierTwoPi{g}$, and dilation obeys
$\FourierTwoPi{\delta^{-1}f(\cdot/\delta)}(\xi)=\FourierTwoPi{f}(\delta\xi)$
for $\delta>0$.  The sinc function is the \emph{radian} sinc of the
catalogue entry \texttt{sinc},
\begin{equation}
 \SincRad(z)\DefinitionEquals
 \begin{cases}\sin z/z,&z\ne0,\\ 1,&z=0,\end{cases}
 \qquad
 \SincPi(z)=\SincRad(\pi z),
 \label{due:eq:set-sinc}
\end{equation}
so that the product \cref{due:eq:set-infinite-product} may equally be
written $\prod_{k\ge0}\SincPi(\xi/2^{k})$.  Under
\cref{due:eq:set-fourier-convention}, $\PhiUp=\FourierTwoPi{\RvachevUp}$ is
the product \cref{due:eq:set-infinite-product}; this is the normalization of
Rvachev and of Arias de Reyna \cite{rvachev1990,arias1982}, and the
catalogue records that the product converges locally uniformly on
$\ComplexNumbers$ to an entire function with $\PhiUp(0)=1$.

\subsection{Boxes, the finite product, and the finite spline}
\label{due:sec:set-boxes}

\begin{definition}[Dyadic boxes, finite products, finite splines]
\label{due:def:set-finite-spline}
For $k\in\NaturalNumbers$ let
\begin{equation}
 b_k(x)\DefinitionEquals 2^{k}\,\IndicatorOf{[-2^{-k-1},\,2^{-k-1}]}(x),
 \qquad x\in\RealNumbers,
 \label{due:eq:set-box}
\end{equation}
the centred uniform density of total mass one and half-width $2^{-k-1}$.
For $n\in\NaturalNumbers$ define
\begin{equation}
 \Phin{n}(\xi)\DefinitionEquals\prod_{k=0}^{n}
 \SincRad\!\left(\frac{\pi\xi}{2^{k}}\right),
 \qquad
 \upn{n}\DefinitionEquals b_0\Convolution b_1\Convolution\cdots\Convolution b_n,
 \qquad
 \delta_n\DefinitionEquals 2^{-(n+1)} .
 \label{due:eq:set-finite-product}
\end{equation}
Thus $\Phin{n}$ has exactly $n+1$ factors, indexed $k=0,\ldots,n$, and
$\delta_n$ is the half-width of the last box $b_n$.
\end{definition}

\begin{lemma}[Transform of a box]
\label{due:lem:set-box-transform}
$\FourierTwoPi{b_k}(\xi)=\SincRad(\pi\xi/2^{k})$ for every
$k\in\NaturalNumbers$ and $\xi\in\RealNumbers$.
\end{lemma}

\begin{proof}
With $h=2^{-k-1}$ and $\xi\ne0$,
$\FourierTwoPi{b_k}(\xi)=2^{k}\int_{-h}^{h}\EulerE^{-2\pi\ImaginaryUnit x\xi}
\Differential x=2^{k}\,\dfrac{\sin(2\pi h\xi)}{\pi\xi}
=\dfrac{\sin(\pi\xi/2^{k})}{\pi\xi/2^{k}}$, since $2^{k}\cdot2h=1$; at
$\xi=0$ both sides equal $1$.
\end{proof}

\begin{proposition}[Finite convolution model]
\label{due:prop:set-convolution-model}
For every $n\in\NaturalNumbers$:
\begin{enumerate}
\item $\FourierTwoPi{\upn{n}}=\Phin{n}$.  For $n\ge1$ the function
      $\Phin{n}$ is integrable, $\upn{n}$ is continuous, and
      $\upn{n}=\InverseFourierTwoPiOperator[\Phin{n}]$ holds pointwise;
      for $n=0$, $\upn{0}=b_0$ and inversion holds at every point except
      the two jumps $\pm\frac12$, where the inversion integral returns the
      half-height $\frac12$.
\item $\upn{n}$ is an even, nonnegative probability density with
      \begin{equation}
       \SupportOf{\upn{n}}=[-1+\delta_n,\ 1-\delta_n].
       \label{due:eq:set-finite-support}
      \end{equation}
\item (Probabilistic model.)  Let $V_0,V_1,\ldots$ be independent random
      variables, each uniform on $[-\frac12,\frac12]$, and put
      \begin{equation}
       Z_n\DefinitionEquals\sum_{k=0}^{n}2^{-k}V_k,
       \qquad
       Z\DefinitionEquals\sum_{k=0}^{\infty}2^{-k}V_k .
       \label{due:eq:set-random-sums}
      \end{equation}
      The series converges absolutely and surely, $\AbsoluteValue{Z}\le1$,
      $\upn{n}$ is the density of $Z_n$, and $\RvachevUp$ is the density of
      $Z$.  In particular $\upn{n}$ is the density of a sum of $n+1$
      independent centred uniforms of half-widths
      $\frac12,\frac14,\ldots,2^{-n-1}$.
\end{enumerate}
\end{proposition}

\begin{proof}
(1) Each $b_k\in L^{1}$, so the convolution theorem and
\cref{due:lem:set-box-transform} give
$\FourierTwoPi{\upn{n}}=\prod_{k=0}^{n}\FourierTwoPi{b_k}=\Phin{n}$.  Since
$\AbsoluteValue{\SincRad(\pi\xi/2^{k})}\le 2^{k}/(\pi\AbsoluteValue{\xi})$
and also $\le1$, one has
$\AbsoluteValue{\Phin{n}(\xi)}\le\min\{1,\,C_n\AbsoluteValue{\xi}^{-(n+1)}\}$
with $C_n=2^{\binom{n+1}{2}}\pi^{-(n+1)}$, which is integrable for
$n\ge1$.  A convolution of a bounded function with an $L^{1}$ function is
continuous, so $\upn{n}=b_0\Convolution(b_1\Convolution\cdots\Convolution b_n)$
is continuous for $n\ge1$, and the classical inversion theorem applies.
For $n=0$ the statement is the Dirichlet--Jordan theorem for the box.

(2) Convolution of even, nonnegative functions of integral one is even,
nonnegative, and of integral one (Fubini).  The support of a convolution of
compactly supported nonnegative functions is the Minkowski sum of the
supports, so the support radius is $\sum_{k=0}^{n}2^{-k-1}=1-2^{-(n+1)}$.

(3) $\AbsoluteValue{2^{-k}V_k}\le2^{-k-1}$, so the series converges
absolutely for every sample point and $\AbsoluteValue{Z}\le\sum_k2^{-k-1}=1$.
The density of $2^{-k}V_k$ is $b_k$, and the density of a sum of independent
variables is the convolution of the densities, so $Z_n$ has density
$\upn{n}$.  The characteristic function of $Z$ at $-2\pi\xi$, that is the
$2\pi$-transform of its law, is
$\Expectation\,\EulerE^{-2\pi\ImaginaryUnit\xi Z}
=\lim_{n}\Expectation\,\EulerE^{-2\pi\ImaginaryUnit\xi Z_n}
=\lim_{n}\Phin{n}(\xi)=\PhiUp(\xi)$ by dominated convergence and the
locally uniform convergence of the product.  The finite measure
$\RvachevUp(x)\,\Differential x$ has the same transform
\cref{due:eq:set-infinite-product}, and a finite Borel measure is determined
by its Fourier transform; hence $Z$ has density $\RvachevUp$.
\end{proof}

\begin{proposition}[Head--tail factorization]
\label{due:prop:set-head-tail}
For $\delta>0$ write $T_\delta(y)\DefinitionEquals\delta^{-1}\RvachevUp(y/\delta)$
for the copy of $\RvachevUp$ rescaled to the support $[-\delta,\delta]$; it
is an even probability density with $\FourierTwoPi{T_\delta}(\xi)=\PhiUp(\delta\xi)$.
For every $n\in\NaturalNumbers$,
\begin{equation}
 \PhiUp(\xi)=\Phin{n}(\xi)\,\PhiUp(\delta_n\xi),
 \qquad
 \boxed{\;\RvachevUp=\upn{n}\Convolution T_{\delta_n},\;}
 \qquad
 Z\EqualInLaw Z_n+\delta_n Z',
 \label{due:eq:set-head-tail}
\end{equation}
where $Z'$ is a copy of $Z$ independent of $Z_n$.  Equivalently,
$\Phin{n}(\xi)=\PhiUp(\xi)/\PhiUp(\delta_n\xi)$ wherever the denominator is
nonzero.
\end{proposition}

\begin{proof}
Reindex the tail of the product by $k=n+1+j$:
\[
 \prod_{k=n+1}^{\infty}\SincRad\!\left(\frac{\pi\xi}{2^{k}}\right)
 =\prod_{j=0}^{\infty}\SincRad\!\left(\frac{\pi(2^{-(n+1)}\xi)}{2^{j}}\right)
 =\PhiUp(\delta_n\xi),
\]
which is the first identity.  The dilation rule gives
$\FourierTwoPi{T_{\delta_n}}(\xi)=\PhiUp(\delta_n\xi)$, so the first identity
says $\FourierTwoPi{\RvachevUp}=\FourierTwoPi{\upn{n}}\,\FourierTwoPi{T_{\delta_n}}
=\FourierTwoPi{\upn{n}\Convolution T_{\delta_n}}$.  Both $\RvachevUp$ and
$\upn{n}\Convolution T_{\delta_n}$ are in $L^{1}$, and an $L^{1}$ function is
determined almost everywhere by its transform; both sides are continuous
(the right side as a convolution of a bounded function with an $L^{1}$
function), so they agree everywhere.  For the random-variable form, the
omitted sum is $\sum_{k\ge n+1}2^{-k}V_k=2^{-(n+1)}\sum_{j\ge0}2^{-j}V_{n+1+j}$,
and $Z'\DefinitionEquals\sum_{j\ge0}2^{-j}V_{n+1+j}$ is built from the
variables $V_{n+1},V_{n+2},\ldots$, hence is independent of $Z_n$ and has
the law of $Z$.
\end{proof}

The factorization \cref{due:eq:set-head-tail} is exact for every $n$; no
limit is taken.  It is the source of the base $\Qq=\frac14$: the omitted
tail lives at the spatial scale $\delta_n$, its law is even so only even
moments enter, and $\delta_{n+1}^{2}/\delta_n^{2}=\frac14$.

\subsection{The exact truncated-power formula}
\label{due:sec:set-truncated-power}

Write $\ThueMorseSign{k}=(-1)^{\BinaryDigitSum(k)}$ for the signed
Thue--Morse sequence, where $\BinaryDigitSum(k)$ is the number of ones in
the binary expansion of $k$, so that
$(\ThueMorseSign{k})_{k\ge0}=1,-1,-1,1,-1,1,1,-1,\ldots$; and write
$\PositivePart{y}^{n}=y^{n}$ for $y>0$ and $0$ for $y\le0$, with the
convention $\PositivePart{y}^{0}=\IndicatorOf{(0,\infty)}(y)$ at $n=0$.

\begin{theorem}[Exact finite spline with Thue--Morse knots]
\label{due:thm:set-finite-spline}
For every $n\ge1$ and every $x\in\RealNumbers$,
\begin{equation}
 \boxed{\;
 \upn{n}(x)=\frac{2^{\binom{n+1}{2}}}{n!}
 \sum_{k=0}^{2^{n+1}-1}\ThueMorseSign{k}\,
 \PositivePart{x+1-\delta_n-\frac{k}{2^{n}}}^{\,n}.\;}
 \label{due:eq:set-truncated-power}
\end{equation}
Consequently:
\begin{enumerate}
\item The knots are the $2^{n+1}$ points
      \begin{equation}
       t_{n,k}\DefinitionEquals-1+\delta_n+\frac{k}{2^{n}}
       =-1+(2k+1)\,\delta_n
       =\bigl(2k+1-2^{n+1}\bigr)\delta_n,
       \qquad 0\le k<2^{n+1},
       \label{due:eq:set-knots}
      \end{equation}
      the odd multiples of $\delta_n$ in $[-1+\delta_n,1-\delta_n]$;
      consecutive knots are $2\delta_n$ apart, the first and last knots are
      the endpoints of the support.
\item $\upn{n}\in C^{\,n-1}(\RealNumbers)$, and on each open interval free of
      knots $\upn{n}$ coincides with a polynomial of degree at most $n$.
\item Off the knots the $n$-th derivative is the step function
      \begin{equation}
       {\upn{n}}^{(n)}(x)=2^{\binom{n+1}{2}}
       \sum_{k=0}^{2^{n+1}-1}\ThueMorseSign{k}\,\IndicatorOf{(t_{n,k},\infty)}(x),
       \label{due:eq:set-nth-derivative}
      \end{equation}
      which jumps by $2^{\binom{n+1}{2}}\ThueMorseSign{k}\ne0$ at $t_{n,k}$.
      Hence every $t_{n,k}$ is a genuine knot: $\upn{n}$ agrees with no
      polynomial on any neighbourhood of $t_{n,k}$.
\end{enumerate}
For $n=0$ the right-hand side of \cref{due:eq:set-truncated-power} equals
$\IndicatorOf{[-\frac12,\frac12)}$, which is $b_0=\upn{0}$ except at the
jump $x=\frac12$.
\end{theorem}

% ed.: the sources say "integrate n+1 times from -infinity"; the closure
% ed.: (a polynomial vanishing on a half-line is zero) is supplied here, as
% ed.: is the genuine-knot statement (3), which the sources do not state
% ed.: and which \cref{due:lem:set-midpoint} needs for "x is a knot".
\begin{proof}
\emph{Step 1: the top derivative is atomic.}
With $h_k=2^{-k-1}$ one has
$b_k=2^{k}\bigl(\IndicatorOf{(-h_k,\infty)}-\IndicatorOf{(h_k,\infty)}\bigr)$
almost everywhere, hence in the sense of distributions
$Db_k=2^{k}(\delta_{-h_k}-\delta_{h_k})$, a difference of two unit point
masses.  Convolution of compactly supported distributions is associative and
commutative and satisfies $D(u\Convolution v)=(Du)\Convolution v$, so
\begin{align*}
 D^{n+1}\upn{n}
 &=(Db_0)\Convolution(Db_1)\Convolution\cdots\Convolution(Db_n)
 =\prod_{k=0}^{n}2^{k}\sum_{\beta\in\{0,1\}^{n+1}}
 (-1)^{\beta_0+\cdots+\beta_n}\,\delta_{c(\beta)},\\
 c(\beta)&=\sum_{k=0}^{n}(-1)^{1-\beta_k}h_k
 =-(1-\delta_n)+\sum_{k=0}^{n}\beta_k2^{-k},
\end{align*}
where $\beta_k=1$ selects the right endpoint $+h_k$ of the $k$-th box (sign
$-1$) and $\beta_k=0$ the left endpoint $-h_k$ (sign $+1$), and the second
form of $c(\beta)$ uses $\sum_k h_k=1-\delta_n$ and $2h_k=2^{-k}$.  Put
\[
 k(\beta)\DefinitionEquals\sum_{j=0}^{n}\beta_j2^{\,n-j}.
\]
As $\beta$ runs over the binary cube, $k(\beta)$ runs exactly once over
$0,1,\ldots,2^{n+1}-1$ (binary expansion with the digits read in reverse
order), $\sum_j\beta_j2^{-j}=k(\beta)/2^{n}$, and
$(-1)^{\sum_j\beta_j}=\ThueMorseSign{k(\beta)}$ because reversing the
binary digits does not change their sum.  With
$\prod_{k=0}^{n}2^{k}=2^{\binom{n+1}{2}}$ this gives
\begin{equation}
 D^{n+1}\upn{n}=2^{\binom{n+1}{2}}\sum_{k=0}^{2^{n+1}-1}
 \ThueMorseSign{k}\,\delta_{t_{n,k}},
 \qquad t_{n,k}=-1+\delta_n+\frac{k}{2^{n}} .
 \label{due:eq:set-atomic-derivative}
\end{equation}

\emph{Step 2: integrate back.}
Let $R(x)$ denote the right-hand side of \cref{due:eq:set-truncated-power}.
For $n\ge1$ the function $y\mapsto\PositivePart{y}^{n}/n!$ is of class
$C^{\,n-1}$, its $n$-th derivative is $\IndicatorOf{(0,\infty)}$ off the
origin, and its $(n+1)$-st distributional derivative is $\delta_0$.  Hence
$D^{n+1}R$ equals the right-hand side of
\cref{due:eq:set-atomic-derivative}, so $D^{n+1}(\upn{n}-R)=0$ and
$\upn{n}-R$ is a polynomial of degree at most $n$ (the only distributions
annihilated by $D^{n+1}$).  For $x<-1+\delta_n=t_{n,0}$ every truncated
power in $R$ vanishes, and $\upn{n}$ vanishes by
\cref{due:eq:set-finite-support}; a polynomial vanishing on a half-line is
identically zero.  This proves \cref{due:eq:set-truncated-power}.

\emph{Consequences.}
(1) is read off the formula: $t_{n,k}=-1+(2k+1)2^{-(n+1)}$, and
$2k+1-2^{n+1}$ runs over the odd integers from $1-2^{n+1}$ to $2^{n+1}-1$.
(2) Each summand is $C^{\,n-1}$ and polynomial of degree at most $n$ off its
own knot.  (3) Differentiate \cref{due:eq:set-truncated-power} $n$ times off
the knots, using $D^{n}\PositivePart{y}^{n}=n!\,\IndicatorOf{(0,\infty)}(y)$;
the knots are pairwise distinct, so exactly one indicator switches at each
$t_{n,k}$, and the jump is $2^{\binom{n+1}{2}}\ThueMorseSign{k}$.  A
polynomial has a continuous $n$-th derivative, so no polynomial agrees with
$\upn{n}$ on a neighbourhood of a knot.
For $n=0$, $\PositivePart{x+\tfrac12}^{0}-\PositivePart{x-\tfrac12}^{0}
=\IndicatorOf{[-\frac12,\frac12)}(x)$ with the stated convention.
\end{proof}

\begin{corollary}[Scaled cell formula and rationality]
\label{due:cor:set-scaled-cell}
Let $n\ge1$ and $x\in[-1,1]$, and put
\begin{equation}
 y\DefinitionEquals 2^{n}(x+1-\delta_n)=2^{n}(x+1)-\tfrac12,
 \qquad 0\le\Floor{y}\le2^{n+1}-1 .
 \label{due:eq:set-scaled-coordinate}
\end{equation}
Then
\begin{equation}
 \upn{n}(x)=\frac{2^{-\binom{n}{2}}}{n!}
 \sum_{0\le k\le\Floor{y}}\ThueMorseSign{k}\,(y-k)^{n},
 \label{due:eq:set-scaled-cell}
\end{equation}
where the sum is empty if $y<0$.  In particular:
\begin{enumerate}
\item if $x\in\RationalNumbers$ then $\upn{n}(x)\in\Kk$;
\item if $x=m/2^{s}$ with $m\in\IntegerNumbers$ and $n\ge s$, then
      $y\in\IntegerNumbers+\frac12$ and
      $2^{\binom{n+1}{2}}\,n!\,\upn{n}(x)\in\IntegerNumbers$.
\end{enumerate}
\end{corollary}

\begin{proof}
$x+1-\delta_n-k/2^{n}=2^{-n}(y-k)$, and the term is active exactly when
$k<y$, that is $k\le\Floor{y}$ when $y\notin\IntegerNumbers$ and
$k\le y-1$ when $y\in\IntegerNumbers$; in the latter case the term $k=y$
contributes $0$ anyway.  Factor $2^{-n^2}$ out of the $n$-th powers:
$2^{\binom{n+1}{2}-n^{2}}=2^{-\binom n2}$.  For $x\le1$ one has
$y\le2^{n+1}-\frac12$, so no index beyond $2^{n+1}-1$ is ever active, and
for $x\ge-1$ one has $y\ge-\frac12$.  Rationality is now immediate.  If
$2^{s}x\in\IntegerNumbers$ and $n\ge s$, then $2^{n}(x+1)\in\IntegerNumbers$,
so $y-k\in\IntegerNumbers+\frac12$, $(y-k)^{n}\in2^{-n}\IntegerNumbers$, and
the sum lies in $2^{-n}\IntegerNumbers$; multiplying by
$2^{\binom n2}n!\cdot2^{n}=2^{\binom{n+1}{2}}n!$ gives an integer.
\end{proof}

The scaled formula \cref{due:eq:set-scaled-cell} is the form in which the
sources compute (S3, S4, S6) and is, term for term, the definition of
\texttt{fabiusUniformSpline} in the repository
(\cref{due:sec:set-bridge}).  The denominator bound in (2) is sharp in the
examples: $\upn{2}(\tfrac14)=\tfrac{15}{16}$ with
$2^{3}\cdot2!=16$, and $\upn{3}(\tfrac14)=\tfrac{179}{192}$ with
$2^{6}\cdot3!=384$.
% ed.: the denominator bound is new; it is what makes the phrase "exact
% ed.: rational arithmetic" quantitative for the verifiers.

\begin{remark}[Prouhet cancellation, in advance]
\label{due:rmk:set-prouhet}
Summing \cref{due:eq:set-atomic-derivative} over all $k$ gives
$\sum_{k<2^{n+1}}\ThueMorseSign{k}=0$ for $n\ge0$, which is the reason the
step function \cref{due:eq:set-nth-derivative} returns to $0$ beyond the
last knot.  More is true: $\sum_{k<2^{m}}\ThueMorseSign{k}\,k^{j}=0$ for all
$0\le j<m$ (the generating polynomial
$\sum_{k<2^{m}}\ThueMorseSign{k}z^{k}=\prod_{h<m}(1-z^{2^{h}})$ has a zero of
order $m$ at $z=1$; catalogue entry \texttt{thue-morse-block-sum}).  This is
the mechanism by which, at a dyadic point, the local degree of $\upn{n}$
collapses from $n$ to $\dep{x}$; the section on local degree collapse
develops it.
\end{remark}

\subsection{Basic analytic properties and strict positivity}
\label{due:sec:set-basic}

\begin{proposition}[Basic properties of $\RvachevUp$ and $\upn{n}$]
\label{due:prop:set-basic}
\begin{enumerate}
\item $\RvachevUp$ is an even, nonnegative $C^{\infty}$ probability density
      with topological support $[-1,1]$, $\RvachevUp(0)=1$,
      $\RvachevUp(\pm1)=0$, all derivatives vanishing at $\pm1$, and
      \begin{equation}
       \RvachevUp'(x)=2\RvachevUp(2x+1)-2\RvachevUp(2x-1),
       \qquad x\in\RealNumbers .
       \label{due:eq:set-differential-law}
      \end{equation}
\item $\upn{n}(x)>0$ for $\AbsoluteValue{x}<1-\delta_n$, for every $n\ge0$.
\item $\RvachevUp(x)>0$ for $\AbsoluteValue{x}<1$.
\item $\upn{n}(0)=1$ for every $n\ge0$.
\end{enumerate}
\end{proposition}

\begin{proof}
(1) The regularity, support, normalization and the differential law are the
content of the catalogue entry \texttt{rvachev-up}, Lean-proved in the
repository \cite{proveit-primary}; we recall the one-line derivations that
are used again below.  Smoothness: by
\cref{due:prop:set-convolution-model}(1),
$\AbsoluteValue{\PhiUp(\xi)}\le\AbsoluteValue{\Phin{K}(\xi)}
\le C_K\AbsoluteValue{\xi}^{-(K+1)}$ for every $K$, so
$\xi^{m}\PhiUp(\xi)\in L^{1}$ for every $m$ and the inversion integral may
be differentiated under the sign arbitrarily often.  The value at $0$:
by \cref{due:eq:set-head-tail} with $n=0$, $\RvachevUp=b_0\Convolution T_{1/2}$,
and $T_{1/2}$ is a probability density supported in
$[-\frac12,\frac12]$, on which $b_0\equiv1$; hence
$\RvachevUp(0)=\int b_0(-y)T_{1/2}(y)\Differential y=1$.  The same identity
reads $\RvachevUp(x)=\int_{x-1/2}^{x+1/2}2\RvachevUp(2y)\Differential y
=\int_{2x-1}^{2x+1}\RvachevUp(z)\Differential z$, and differentiating
gives \cref{due:eq:set-differential-law}.

(2) Let $f,g\ge0$ be integrable with $f>0$ on $(-\alpha,\alpha)$ and $g>0$
on $(-\beta,\beta)$, $f$ vanishing outside $[-\alpha,\alpha]$.  For
$\AbsoluteValue{x}<\alpha+\beta$ the set
$\{y:\AbsoluteValue{y}<\beta,\ \AbsoluteValue{x-y}<\alpha\}
=(-\beta,\beta)\cap(x-\alpha,x+\alpha)$ is a nonempty open interval (its
endpoints $\max(-\beta,x-\alpha)<\min(\beta,x+\alpha)$ precisely because
$\AbsoluteValue{x}<\alpha+\beta$), and on it the integrand of
$(f\Convolution g)(x)=\int f(x-y)g(y)\Differential y$ is strictly positive;
hence $(f\Convolution g)(x)>0$.  Apply this inductively to
$\upn{n}=\upn{n-1}\Convolution b_n$ with $\alpha=1-\delta_{n-1}$ and
$\beta=\delta_n$, noting $\alpha+\beta=1-\delta_n$, starting from
$\upn{0}=b_0>0$ on $(-\frac12,\frac12)$.

(3) Fix $\AbsoluteValue{x}<1$ and choose $n$ with $2\delta_n<1-\AbsoluteValue{x}$.
For $\AbsoluteValue{y}\le\delta_n$ one has
$\AbsoluteValue{x-y}\le\AbsoluteValue{x}+\delta_n<1-\delta_n$, so by (2) and
continuity ($n\ge1$) the function $y\mapsto\upn{n}(x-y)$ is bounded below
by a positive constant $c$ on $[-\delta_n,\delta_n]$, which carries the
whole mass of $T_{\delta_n}$.  By \cref{due:eq:set-head-tail},
$\RvachevUp(x)=\int\upn{n}(x-y)T_{\delta_n}(y)\Differential y\ge c>0$.

(4) $\upn{n}=b_0\Convolution g_n$ with
$g_n=b_1\Convolution\cdots\Convolution b_n$ a probability density supported
in $[-\frac12+\delta_n,\frac12-\delta_n]\subset[-\frac12,\frac12]$
($g_0=\delta_0$), on which $b_0\equiv1$; hence
$\upn{n}(0)=\int b_0(-y)g_n(y)\Differential y=1$.
\end{proof}

\subsection{The indexing bridge}
\label{due:sec:set-bridge}

The finite spline of this volume has appeared under three indexings in the
corpus.  The box below fixes the dictionary once; every later use of an
outside result is to be read through it.
% ed.: a breakable tcolorbox placed directly after a heading inherits the
% ed.: heading's \nobreak and is moved whole to the next page, leaving the
% ed.: heading orphaned; the lead-in paragraph above prevents that.

% ed.: consolidates S2 "Index conversion to the repository convention",
% ed.: S4 "Factor-count and user indexing" and S6 "Exact indexing bridge";
% ed.: the Lean correspondence was checked against the definition of
% ed.: fabiusUniformSpline and the module docstring of
% ed.: QuarterSplineLocalPolynomial.lean.
\begin{warningbox}[title={One spline, three indexings}]
All six sources, and this volume, index the finite spline by the
\emph{largest factor index}: $\Phin{n}=\prod_{k=0}^{n}\SincRad(\pi\xi/2^{k})$
has $n+1$ factors, $\upn{n}$ has degree $n$, and the omitted tail has
half-width $\delta_n=2^{-(n+1)}$.  Two other conventions are in use in the
corpus, and every formula transported from elsewhere must be shifted once.
\begin{itemize}
\item \textbf{The Exponents volume} \cite{proveit-exponents} indexes by the
      \emph{number of factors}: its prefix product (equation
      \texttt{p3:eq:prefix-product}) is
      $\Phi_n(\xi)=\prod_{j=0}^{n-1}\SincRad(\pi\xi/2^{j})$ with $n$ factors,
      $\FourierTwoPi{p_n}=\Phi_n$, $p_n=u_{2^{-1}}\Convolution\cdots\Convolution u_{2^{-n}}$
      with $u_a=(2a)^{-1}\IndicatorOf{[-a,a]}$, and its tail scale is
      $\delta_N=2^{-N}$.  The bridge is
      \begin{equation}
       \boxed{\;\upn{n}=p_{n+1},\qquad \Phin{n}=\Phi_{n+1},\qquad
       \delta_n^{\text{(here)}}=2^{-(n+1)}=\delta_{n+1}^{\text{(there)}},\;}
       \label{due:eq:set-bridge-exponents}
      \end{equation}
      so the sharp error law $2^{\binom{r+3}{2}-1}/(9\cdot4^{N})$ of
      \texttt{p3:thm:sharp-prefix-law} and the exact CDF law
      $1/(9\cdot4^{N})$ of \texttt{p3:cor:kolmogorov} read, for $\upn{n}$,
      with $4^{N}=4^{n+1}$; and a mode $4^{-rN}$ there is a mode
      $4^{-r(n+1)}=4^{-r}\cdot4^{-rn}$ here
      (\cref{due:rmk:set-exponent-conventions}).
\item \textbf{The Lean development} \cite{proveit-primary} works on the
      Fabius interval.  Its centred spline
      \texttt{fabiusUniformSpline p y} is, for $y\in[0,1]$, the sum
      $\bigl(2^{\binom p2}p!\bigr)^{-1}\sum_{r<\Floor{2^{p}y+1/2}}
      \ThueMorseSign{r}\,(2^{p}y-\tfrac12-r)^{p}$, that is
      \cref{due:eq:set-scaled-cell} with $y=2^{p}(x+1)-\frac12$ replaced
      by $2^{p}y-\frac12$.  Hence, for $n\ge1$ and $x\in[-1,1]$,
      \begin{equation}
       \boxed{\;\upn{n}(x)=\mathtt{fabiusUniformSpline}\;n\;(1+x)
       =\mathtt{fabiusUniformSpline}\;n\;(1-\AbsoluteValue{x}),\;}
       \label{due:eq:set-bridge-lean}
      \end{equation}
      the second form by evenness of $\upn{n}$.  The degree matches the
      index: $\upn{n}$ \emph{is} \texttt{fabiusUniformSpline n}, and the
      Lean theorem
      \texttt{fabiusUniformSpline\_\allowbreak eq\_\allowbreak centeredPartialCDF}
      identifies it on $[0,1]$ with the distribution function of
      $\sum_{j=1}^{n}2^{-j}U_j+2^{-n-1}$, $U_j$ uniform on $[0,1]$ (the
      finite-level form of the fold $\RvachevUp(x)=\FabiusBounded(1-\AbsoluteValue{x})$
      of the catalogue).  The module
      \texttt{QuarterSplineLocalPolynomial.\allowbreak lean} writes $P_n$
      for the density of the first $n$ centred uniforms, the
      report convention of the inverse frontier, and records
      $P_n=\mathtt{fabiusUniformSpline}\,(n-1)$: thus
      \begin{equation}
       P_n=\upn{n-1}=p_n,
       \label{due:eq:set-bridge-report}
      \end{equation}
      and its kernel-checked local formula
      $P_n(\tfrac14+z)=\tfrac{5}{72}+z+4z^{2}-\tfrac49\,4^{-n}$ for
      $n\ge3$, $\AbsoluteValue{z}\le2^{-n}$, is in this volume's terms
      $\upn{m}(\tfrac34-z)=\tfrac5{72}+z+4z^{2}-\tfrac19\,4^{-m}$ for
      $m=n-1\ge2$, the depth-two instance $x=\frac34$ of
      \cref{due:thm:set-main} with $\Cr{1}{\frac34}=-\frac19$ (compare
      $\Cr{1}{\frac14}=+\frac19$ in \cref{due:ex:set-quarter}; the sign
      flips with $\RvachevUp''$).
\end{itemize}
The rule of this volume: \emph{every} formula keeps the sources' $n$;
$p_n$, $\Phi_n$ (Exponents), $P_n$ (report) and
\texttt{fabiusUniformSpline} are mentioned only in this box and in the
formalization register, always through
\cref{due:eq:set-bridge-exponents,due:eq:set-bridge-lean,due:eq:set-bridge-report}.
The most common error the shift produces is the confusion between
$4^{-n}$ and $4^{-(n+1)}$ in the coefficient of a mode; the second most
common is an onset stated as $N\ge s+1$ (factor count) being copied as
$n\ge s+1$ (largest index).  The onset in this volume's indexing is
$n\ge s$, and \cref{due:ex:set-quarter} shows it cannot be lowered.
\end{warningbox}

\begin{summarybox}[title={Setting, in five lines}]
$2\pi$-frequency transforms, radian sinc; $\PhiUp=\FourierTwoPi{\RvachevUp}
=\prod_{k\ge0}\SincRad(\pi\xi/2^{k})$; $\Phin{n}$ keeps $k=0,\ldots,n$
($n+1$ factors); $\upn{n}=b_0\Convolution\cdots\Convolution b_n$ is its
inverse transform, the density of $Z_n=\sum_{k\le n}2^{-k}V_k$, an even
$C^{\,n-1}$ box spline of degree $n$ on $[-1+\delta_n,1-\delta_n]$,
$\delta_n=2^{-(n+1)}$, with the exact Thue--Morse truncated-power formula
\cref{due:eq:set-truncated-power} and knots at the odd multiples of
$\delta_n$; $\RvachevUp=\upn{n}\Convolution T_{\delta_n}$ exactly, for every
$n$; $\upn{n}(x)\in\Kk$ at rational $x$, with denominator dividing
$2^{\binom{n+1}{2}}n!$ once $n\ge\dep{x}$; $\upn{n}=p_{n+1}$ in the
Exponents volume and $\upn{n}=\mathtt{fabiusUniformSpline}\;n$ in Lean.
\end{summarybox}

\section{Dyadic depth and cell geometry}
\label{due:sec:depth}

\subsection{The depth}
\label{due:sec:set-depth-def}

\begin{definition}[Dyadic depth]
\label{due:def:set-depth}
A dyadic rational is a number $x=m/2^{n}$ with $m\in\IntegerNumbers$ and
$n\in\NaturalNumbers$ (catalogue entry \texttt{fabius-dyadic-cells}).  Its
\emph{depth} is
\begin{equation}
 \dep{x}\DefinitionEquals\min\SetOf{s\in\NaturalNumbers:\ 2^{s}x\in\IntegerNumbers}.
 \label{due:eq:set-depth}
\end{equation}
Thus $\dep{x}=0$ exactly when $x\in\IntegerNumbers$, and for $s=\dep{x}\ge1$
the integer $a\DefinitionEquals2^{s}x$ is odd, so that $x=a/2^{s}$ is the
unique representation of $x$ with an odd numerator; for $0<x<1$ this is the
catalogue's ``reduced at level $s$''.  The four letters $r$ (S1), $d$
(S2), $m$ (S3) and $s$ (S4, S5, S6) of the sources all denote $\dep{x}$.
\end{definition}

If $a=2^{s}x$ were even for $s\ge1$, then $2^{s-1}x=a/2$ would be an integer,
contradicting minimality; this is the whole content of ``odd''.

\begin{lemma}[Representation independence and invariances]
\label{due:lem:set-depth-representation}
Let $x=m/2^{n}$ with $m\in\IntegerNumbers$, $n\in\NaturalNumbers$, not
necessarily reduced, and let $\TwoAdicValuation(m)$ be the $2$-adic
valuation of $m$, with $\TwoAdicValuation(0)=+\infty$.  Then
\begin{equation}
 \dep{x}=\max\SetOf{0,\ n-\TwoAdicValuation(m)},
 \label{due:eq:set-depth-valuation}
\end{equation}
which is invariant under $(m,n)\mapsto(2m,n+1)$.  Moreover, for every
dyadic $x$ and every $c\in\IntegerNumbers$,
\begin{equation}
 \dep{-x}=\dep{x},\qquad \dep{x+c}=\dep{x},\qquad
 \dep{1-\AbsoluteValue{x}}=\dep{x}.
 \label{due:eq:set-depth-invariances}
\end{equation}
\end{lemma}

% ed.: the catalogue's audit rule demands invariance under (m,n) ->
% ed.: (2m,n+1) for any exact evaluator; none of the sources states it.
\begin{proof}
If $m=0$ then $x=0$, $\dep{x}=0$, and the right-hand side of
\cref{due:eq:set-depth-valuation} is $\max\{0,-\infty\}=0$.  Otherwise
write $m=2^{v}u$ with $v=\TwoAdicValuation(m)$ and $u$ odd.  Then
$2^{s}x=2^{s+v-n}u$ is an integer if and only if $s+v-n\ge0$, because $u$
is odd; the least such $s\ge0$ is $\max\{0,n-v\}$.  Replacing $(m,n)$ by
$(2m,n+1)$ increases both $v$ and $n$ by one.  The invariances: $2^{s}x$
is an integer if and only if $-2^{s}x$ is; if and only if
$2^{s}x+2^{s}c$ is; and $1-\AbsoluteValue{x}$ is $1-x$ or $1+x$, both
covered by the first two.
\end{proof}

The third invariance is what allows S6 to work with $a=1-\AbsoluteValue{x}\in[0,1]$
throughout: the depth, and hence $d$, is the same for $x$ and for its fold
$1-\AbsoluteValue{x}$, and \cref{due:eq:set-bridge-lean} transports the
samples.

\subsection{Knots and cells}
\label{due:sec:set-cells}

\begin{lemma}[The knot lattice of $\upn{n}$]
\label{due:lem:set-knots}
Let $n\ge1$.  The knot set of $\upn{n}$ is
\begin{equation}
 K_n\DefinitionEquals\SetOf{t_{n,k}:0\le k<2^{n+1}}
 =\SetOf{(2j+1)\delta_n:\ j\in\IntegerNumbers,\ \AbsoluteValue{2j+1}\le2^{n+1}-1},
 \label{due:eq:set-knot-set}
\end{equation}
the odd multiples of $\delta_n$ of absolute value at most $1-\delta_n$.
The complement of $K_n$ in $(-1,1)$ is the disjoint union of the
$2^{n+1}-1$ open cells
\begin{equation}
 \mathcal C_{n,j}\DefinitionEquals\bigl(2j\delta_n-\delta_n,\ 2j\delta_n+\delta_n\bigr),
 \qquad -2^{n}<j<2^{n},
 \label{due:eq:set-cells}
\end{equation}
each of width $2\delta_n$ and centred at the even multiple $2j\delta_n$ of
$\delta_n$; on each cell $\upn{n}$ is a polynomial of degree at most $n$,
and on no neighbourhood of a knot is it a polynomial.  Outside
$[-1+\delta_n,1-\delta_n]$ the function vanishes identically.  The cell
centres are exactly the even multiples of $\delta_n$ in $(-1,1)$; $0$ is
always a cell centre.
\end{lemma}

\begin{proof}
By \cref{due:thm:set-finite-spline}(1), $t_{n,k}=(2k+1-2^{n+1})\delta_n$
with $0\le k<2^{n+1}$, and $2k+1-2^{n+1}$ runs over the odd integers of
absolute value at most $2^{n+1}-1$.  Two consecutive odd multiples
$(2j-1)\delta_n<(2j+1)\delta_n$ of $\delta_n$ bound the open interval
$\mathcal C_{n,j}$, whose midpoint is $2j\delta_n$; both bounding points
are knots exactly when $\AbsoluteValue{2j\pm1}\le2^{n+1}-1$, that is
$\AbsoluteValue{j}\le2^{n}-1$.  The polynomial and genuine-knot statements
are \cref{due:thm:set-finite-spline}(2),(3), and the vanishing outside the
support is \cref{due:eq:set-finite-support}.  Since $1=2^{n+1}\delta_n$ is
an even multiple of $\delta_n$, the even multiples strictly between $-1$
and $1$ are $2j\delta_n$ with $\AbsoluteValue{j}<2^{n}$, and $j=0$ gives
$0$.
\end{proof}

\subsection{A dyadic point becomes a cell midpoint}
\label{due:sec:set-midpoint}

For $x\in\RealNumbers$ and $n\in\NaturalNumbers$ call
\begin{equation}
 W_n(x)\DefinitionEquals(x-\delta_n,\ x+\delta_n)
 \label{due:eq:set-tail-window}
\end{equation}
the \emph{tail window} at $x$: it is the open interval of half-width
$\delta_n$ around $x$, which is exactly the set of points $x-y$ with $y$
in the interior of the support $[-\delta_n,\delta_n]$ of the tail kernel
$T_{\delta_n}$ of \cref{due:prop:set-head-tail}.  Its width $2\delta_n$
equals the width of a cell of $\upn{n}$.  Whether the window meets a knot
is therefore decided by the position of $x$ relative to the lattice
$\delta_n\IntegerNumbers$, and that position is governed by the depth.

\begin{lemma}[A dyadic point is a cell midpoint exactly from level $\dep{x}$ on]
\label{due:lem:set-midpoint}
Let $x\in[-1,1]$ be dyadic with depth $s=\dep{x}$, and let $n\ge1$.  Write
$\theta_n\DefinitionEquals x/\delta_n=2^{n+1}x$.
\begin{enumerate}
\item[(i)] If $n\ge s$, then $\theta_n$ is an even integer.  If moreover
      $\AbsoluteValue{x}<1$, the point $x$ is the centre of the cell
      $\mathcal C_{n,\theta_n/2}$, the two knots adjacent to $x$ are
      $x\pm\delta_n$, and the tail window $W_n(x)$ is that cell: it
      contains no knot.  If $x=\pm1$, then $\upn{n}\equiv0$ on $W_n(x)$.
      In all cases there is a polynomial $P_{n,x}$ of degree at most $n$,
      the \emph{cell polynomial}, with
      \begin{equation}
       \upn{n}(x+h)=P_{n,x}(h)\qquad\text{for }\AbsoluteValue{h}<\delta_n,
       \label{due:eq:set-cell-polynomial}
      \end{equation}
      and $P_{n,\pm1}=0$.
\item[(ii)] If $n=s-1$ (so $s\ge2$), then $\theta_n=2^{s}x$ is an odd
      integer of absolute value at most $2^{s}-1$: the point $x$ is a knot
      of $\upn{n}$, the only knot in $W_n(x)$, and $\upn{n}$ agrees with no
      polynomial on $W_n(x)$.
\item[(iii)] If $1\le n\le s-2$, then $\theta_n\notin\IntegerNumbers$: the
      point $x$ lies in the interior of a cell but is not its centre, and
      $W_n(x)$ contains exactly one knot, which is different from $x$; again
      $\upn{n}$ agrees with no polynomial on $W_n(x)$.
\end{enumerate}
Consequently $W_n(x)$ is free of knots if and only if $n\ge s$.  For the
degree-zero spline the same trichotomy holds with $\delta_0=\frac12$ and
the two jumps $\pm\frac12$ of $b_0$ as knots.
\end{lemma}

% ed.: (i) merges S1 "Cell-center stabilization", S2 "A dyadic point
% ed.: becomes a cell midpoint", S4 "Exact midpoint alignment" and S5
% ed.: "Dyadic midpoint lemma" (the latter two in factor-count indexing,
% ed.: N >= s+1, which is n >= s here); (ii) is S1's item 3; (iii) and the
% ed.: "if and only if" are added so that the onset is explained for every
% ed.: level below s, not only for s-1.
\begin{proof}
For $s=0$, $x\in\{-1,0,1\}$ and $\theta_n=2^{n+1}x$ is an even integer for
every $n\ge0$; cases (ii) and (iii) are empty.  For $s\ge1$ write
$x=a/2^{s}$ with $a$ odd; then $\theta_n=a\,2^{\,n+1-s}$.  This is an
integer if and only if $n+1\ge s$, and it is then even if and only if
$n+1-s\ge1$, that is $n\ge s$; at $n=s-1$ it equals the odd integer $a$;
and for $n\le s-2$ it is a rational with denominator $2^{\,s-1-n}\ge2$.

(i) For $n\ge s$ and $\AbsoluteValue{x}<1$, $x=2j\delta_n$ with
$j=\theta_n/2$ and $\AbsoluteValue{j}<2^{n}$, so by
\cref{due:lem:set-knots} $x$ is the centre of $\mathcal C_{n,j}=W_n(x)$,
whose endpoints $(2j\pm1)\delta_n=x\pm\delta_n$ are the adjacent knots and
whose interior contains none.  The cell polynomial is the polynomial that
$\upn{n}$ equals on $\mathcal C_{n,j}$, shifted by $x$.  For $x=1$ the
window is $(1-\delta_n,1+\delta_n)$, on which $\upn{n}=0$ by
\cref{due:eq:set-finite-support}; the case $x=-1$ is symmetric; so
$P_{n,\pm1}=0$.

(ii) At $n=s-1$, $\theta_n=a$ is odd with $\AbsoluteValue{a}=2^{s}\AbsoluteValue{x}<2^{s}$
(note $\AbsoluteValue{x}<1$ because $s\ge1$), hence
$\AbsoluteValue{a}\le2^{s}-1=2^{n+1}-1$ and $x=a\delta_n\in K_n$ by
\cref{due:eq:set-knot-set}.  The neighbouring knots $x\pm2\delta_n$ lie
outside the open window $W_n(x)$, so $x$ is the only knot in it, and
\cref{due:thm:set-finite-spline}(3) says $\upn{n}$ is not a polynomial
near $x$.

(iii) In units of $\delta_n$ the window is the open interval
$(\theta_n-1,\theta_n+1)$ of length $2$ about a non-integer, which contains
exactly the two integers $\Floor{\theta_n}$ and $\Floor{\theta_n}+1$,
exactly one of which is odd; call it $2j+1$.  Both integers are strictly
inside the window, and $\AbsoluteValue{2j+1}\le2^{n+1}-1$ because
$\AbsoluteValue{\theta_n}<2^{n+1}$ and $\theta_n\notin\IntegerNumbers$
force $\AbsoluteValue{\Floor{\theta_n}},\AbsoluteValue{\Floor{\theta_n}+1}\le2^{n+1}$
with the value $\pm2^{n+1}$ even.  So $(2j+1)\delta_n$ is a knot of
$\upn{n}$ lying in $W_n(x)$, different from $x$ since $\theta_n$ is not an
integer, and it is the only one since the next odd integers are outside.
The point $x$ itself is not a knot and lies in the open cell between the
odd integers $2j\pm1$, at distance $\AbsoluteValue{\theta_n-2j}\delta_n\ne0$
from its centre.  Non-polynomiality on $W_n(x)$ follows again from
\cref{due:thm:set-finite-spline}(3).

The ``if and only if'' collects (i)--(iii).  For $n=0$, $\theta_0=2x$,
$\delta_0=\frac12$, and the knots of $b_0$ are $\pm\frac12=\pm\delta_0$,
which are the odd multiples of $\delta_0$ in $[-\frac12,\frac12]$; the
same case analysis applies verbatim.
\end{proof}

\Cref{due:lem:set-midpoint} is the knot-geometry explanation of the onset
$n\ge s$ in \cref{due:thm:set-main}.  The main theorem is obtained by
deconvolving $\RvachevUp=\upn{n}\Convolution T_{\delta_n}$ on the tail
window; that inversion is finite linear algebra only when the window
carries a single polynomial, which by the lemma happens for all $n\ge s$
and for no $n<s$.  At the levels $n<s$ the window straddles a knot, and
the convolution mixes two polynomial pieces; the resulting values may
still coincide with the eventual law by accident (the section on the
odd-level sharpening isolates the one systematic instance, $n=s-1$ with
$s$ odd), but no such coincidence is needed, and \cref{due:ex:set-quarter}
shows that at $n=s-1$ with $s$ even it fails.

\begin{example}[Depth zero and depth one, exactly]
\label{due:ex:set-low-depth}
The five dyadic points with $d=0$ are $0,\pm1$ (depth $0$) and
$\pm\frac12$ (depth $1$), and \cref{due:thm:set-main} asserts
$\upn{n}(x)=\RvachevUp(x)$ from level $n=\dep{x}$ on.  This holds, and can
be seen directly:
\begin{itemize}
\item $\upn{n}(0)=1=\RvachevUp(0)$ for all $n\ge0$ and
      $\upn{n}(\pm1)=0=\RvachevUp(\pm1)$ for all $n\ge0$, by
      \cref{due:prop:set-basic}(1),(4) and \cref{due:eq:set-finite-support}.
\item $\upn{n}(\pm\tfrac12)=\tfrac12=\RvachevUp(\pm\tfrac12)$ for all
      $n\ge1$.  Indeed $\upn{n}=b_0\Convolution g_n$ with
      $g_n=b_1\Convolution\cdots\Convolution b_n$ an even probability
      density supported in $[-\frac12+\delta_n,\frac12-\delta_n]$, and
      $b_0(\tfrac12-y)=1$ exactly for $y\in[0,1]$, so
      $\upn{n}(\tfrac12)=\int_{0}^{1}g_n=\tfrac12$ by evenness (the point
      $0$ carries no mass).  The same computation with the tail
      $T_{1/2}$ in place of $g_n$ gives $\RvachevUp(\tfrac12)=\tfrac12$.
\end{itemize}
At $n=0=\dep{\frac12}-1$ the point $\frac12$ is a jump of $b_0$, in
accordance with \cref{due:lem:set-midpoint}(ii): $\upn{0}(\tfrac12)$ is
$0$ with the right-continuous convention of
\cref{due:thm:set-finite-spline}, $1$ with the closed box
\cref{due:eq:set-box}, and $\tfrac12$ with the half-height inversion
convention of \cref{due:prop:set-convolution-model}(1).  Only the last
convention makes the law hold one level early; S3 calls it the
symmetric-rectangle convention at $s=1$.  This is the single
convention-dependent value in the volume, and the onset $n\ge s$ avoids it.
\end{example}

\subsection{The tail sees only one polynomial piece}
\label{due:sec:set-tail-one-piece}

\begin{lemma}[Jet transfer through the tail]
\label{due:lem:set-tail-one-piece}
Let $x\in[-1,1]$ be dyadic with depth $s$, let $n\ge\max\{s,1\}$, and let
$P_{n,x}$ be the cell polynomial of \cref{due:eq:set-cell-polynomial}.
Then, with $Z'$ a random variable of law $\RvachevUp$,
\begin{equation}
 \boxed{\;
 \RvachevUp^{(r)}(x)=\int_{-\delta_n}^{\delta_n}P_{n,x}^{(r)}(-y)\,
 T_{\delta_n}(y)\,\Differential y
 =\Expectation\bigl[P_{n,x}^{(r)}(-\delta_nZ')\bigr],
 \qquad 0\le r\le n .\;}
 \label{due:eq:set-jet-transfer}
\end{equation}
In particular $\RvachevUp(x)=\Expectation[P_{n,x}(-\delta_nZ')]$: the value
of $\RvachevUp$ at $x$ is a moment functional of $Z'$ applied to a single
polynomial of degree at most $n$, and the omitted tail never sees a second
polynomial piece.
\end{lemma}

% ed.: S1's "Exact self-similar tail" and S4's "the tail sees only one
% ed.: polynomial piece" state the r = 0 case; S6 states the jet form
% ed.: "in the distributional sense".  The differentiation under the
% ed.: integral for r <= n-1 and the absolute-continuity argument at r = n
% ed.: are supplied here, since the deconvolution section uses r = n.
\begin{proof}
Start from $\RvachevUp(x+h)=\int\upn{n}(x+h-y)T_{\delta_n}(y)\Differential y$,
valid for all $h$ by \cref{due:eq:set-head-tail}.

\emph{Orders $r\le n-1$.}  By \cref{due:thm:set-finite-spline}(2) the
derivatives ${\upn{n}}^{(r)}$, $r\le n-1$, exist everywhere, are continuous
and bounded (compact support).  Differentiating under the integral sign
$r$ times is justified by dominated convergence, the difference quotients
of ${\upn{n}}^{(r-1)}$ being bounded by $\sup\AbsoluteValue{{\upn{n}}^{(r)}}$
and $T_{\delta_n}$ being integrable; hence
$\RvachevUp^{(r)}(x)=\int{\upn{n}}^{(r)}(x-y)T_{\delta_n}(y)\Differential y$.

\emph{Order $r=n$.}  The function ${\upn{n}}^{(n-1)}$ is continuous and
piecewise polynomial, hence Lipschitz, hence absolutely continuous with
almost-everywhere derivative $g\DefinitionEquals{\upn{n}}^{(n)}$, the
bounded step function \cref{due:eq:set-nth-derivative} defined off the
finite knot set.  Then
\[
 \frac{\RvachevUp^{(n-1)}(x+h)-\RvachevUp^{(n-1)}(x)}{h}
 =\int\frac1h\int_{0}^{h}g(x+t-y)\Differential t\;T_{\delta_n}(y)\Differential y .
\]
For every $y$ outside the finite set $x-K_n$ the inner average tends to
$g(x-y)$ as $h\to0$, because $g$ is continuous at $x-y$; the inner average
is bounded by $\sup\AbsoluteValue{g}$; and $T_{\delta_n}$ is integrable.
Dominated convergence gives
$\RvachevUp^{(n)}(x)=\int g(x-y)T_{\delta_n}(y)\Differential y$.

\emph{Localization.}  $T_{\delta_n}$ vanishes outside $[-\delta_n,\delta_n]$,
and for $\AbsoluteValue{y}<\delta_n$ the point $x-y$ lies in the tail window
$W_n(x)$, on which $\upn{n}(x-y)=P_{n,x}(-y)$ by
\cref{due:lem:set-midpoint}(i) (this is where $n\ge s$ enters), so that
${\upn{n}}^{(r)}(x-y)=P_{n,x}^{(r)}(-y)$ for $r\le n-1$ and
$g(x-y)=P_{n,x}^{(n)}(-y)$ for $y$ off the two endpoints.  The endpoints
$y=\pm\delta_n$ are a null set.  Substituting yields the integral in
\cref{due:eq:set-jet-transfer}, and the substitution $y=\delta_n z$ with
$T_{\delta_n}(\delta_nz)\Differential y=\RvachevUp(z)\Differential z$
turns it into the expectation.
\end{proof}

\begin{remark}[What the lemma says and what it does not]
\label{due:rmk:set-tail-scope}
\Cref{due:eq:set-jet-transfer} holds for the $n$-jet only.  For $r>n$ the
right-hand side would be $0$, and $\RvachevUp^{(r)}(x)$ is not in general
$0$; the correct statement for higher orders is the termination theorem
for the dyadic jet, proved later, which says $\RvachevUp^{(r)}(x)=0$ for
$r>\dep{x}$, a bound independent of $n$.  The two facts combine in the
deconvolution section: since $\RvachevUp$ is even, the law of $Z'$ is even
and only the even moments $\Expectation[(Z')^{2m}]$ enter; the finite
Taylor expansion of $P_{n,x}^{(r)}$ then expresses the $n$-jet of
$\RvachevUp$ at $x$ through the $n$-jet of $P_{n,x}$ by a triangular
system, whose inverse is the reciprocal moment series truncated at degree
$n$, with the coefficients $\am{m}$ of \cref{due:eq:set-am-definition}.
For $n<s$ the lemma is false as stated, because the window contains a
knot and $\upn{n}(x-y)$ is given by two different polynomials on the two
sides of it; for irrational $x$ no level $n$ makes $x$ a cell centre,
which is why the eventual law is a genuinely dyadic phenomenon.
\end{remark}

\begin{summarybox}[title={Depth and cells, in four lines}]
$\dep{x}=\min\{s:2^{s}x\in\IntegerNumbers\}$, computed from any
representation $m/2^{n}$ as $\max\{0,n-\TwoAdicValuation(m)\}$, invariant
under $x\mapsto-x$, $x+c$, $1-\AbsoluteValue{x}$; the knots of $\upn{n}$
are the odd multiples of $\delta_n=2^{-(n+1)}$, all genuine, and the cell
centres are the even multiples; a dyadic $x$ is a cell centre of $\upn{n}$
if and only if $n\ge\dep{x}$, is a knot at $n=\dep{x}-1$, and sits off
centre with one knot inside its tail window at every lower level; for
$n\ge\dep{x}$ the whole $n$-jet of $\RvachevUp$ at $x$ is the expectation
of the jet of one cell polynomial, $\RvachevUp^{(r)}(x)
=\Expectation[P_{n,x}^{(r)}(-\delta_nZ')]$, $0\le r\le n$.
\end{summarybox}

% ---- fragment 02-mechanisms ----
% ---------------------------------------------------------------------
%  Cluster "mech": the two finite mechanisms.
%  Absorbs S1 sections 4-6, S2 sections 4-5, S3 sections 5-7, S4 sections
%  4-6, S5 sections 4-6, S6 sections 3-4 (tags as in the provenance table).
% ---------------------------------------------------------------------

\section{The reciprocal-product coefficients and finite deconvolution}
\label{due:sec:coefficients}

Two finite mechanisms turn the eventual law of \cref{due:sec:main} from an
asymptotic expansion into an identity.  The first, proved in this section,
is algebraic: on the one polynomial cell of \(\upn{n}\) centred at a dyadic
point, division by the omitted Fourier factor is a polynomial in a nilpotent
operator, so the all-orders deconvolution series is a finite sum with
rational coefficients.  The second, proved in \cref{due:sec:jet}, is
arithmetic: the derivative jet of \(\RvachevUp\) at a dyadic point of depth
\(s\) stops at order \(s\), and its last even entry is not zero.  This
section introduces the coefficients the first mechanism produces, fixes
their normalization against the three normalizations used by the sources and
by the catalogue, proves that every one of them is a positive rational, and
proves the deconvolution identity.

\paragraph{Standing conventions.}
\(\PhiUp(z)=\FourierTwoPi{\RvachevUp}(z)=\prod_{k\ge0}\SincRad(\pi z/2^{k})\)
is the \(2\pi\)-convention transform of the up-function (catalogue entries
\texttt{rvachev-up}, \texttt{up-fourier-product}, \texttt{sinc}); it is
entire, even, and \(\PhiUp(0)=1\).  \(\Phin{n}\) is its prefix with the
\(n+1\) factors \(k=0,\ldots,n\), \(\upn{n}=\FourierTwoPiOperator^{-1}[\Phin{n}]\),
and
\[
 \delta_n\DefinitionEquals2^{-(n+1)} .
\]
From \cref{due:sec:setting} we use two facts.  First, \(\upn{n}\) is the
density of \(\sum_{k=0}^{n}2^{-k-1}U_k\) with \(U_0,U_1,\ldots\) independent
and uniform on \([-1,1]\); equivalently it is the convolution of the box
densities \(b_k=2^{k}\IndicatorOf{[-2^{-k-1},2^{-k-1}]}\), \(k=0,\ldots,n\).
Hence it is even, nonnegative, supported on \([-(1-\delta_n),1-\delta_n]\),
of class \(C^{n-1}\) when \(n\ge1\), and on each open interval between
consecutive knots it is one polynomial of degree at most \(n\); its knots are
the odd multiples of \(\delta_n\) (the sums \(\sum_{k=0}^{n}\pm2^{-k-1}\)).
Second, the head--tail identity: with the dilated kernel
\begin{equation}
 \rho_\delta(y)\DefinitionEquals\delta^{-1}\RvachevUp(y/\delta)
 \qquad(\delta>0),
 \label{due:eq:mech-dilated-kernel}
\end{equation}
which is a smooth even probability density supported on \([-\delta,\delta]\)
with \(\FourierTwoPi{\rho_\delta}(t)=\PhiUp(\delta t)\),
\begin{equation}
 \RvachevUp=\upn{n}\Convolution\rho_{\delta_n}\qquad(n\ge0).
 \label{due:eq:mech-head-tail}
\end{equation}
Indeed the factors of \(\PhiUp\) omitted from \(\Phin{n}\) are
\(\prod_{k\ge n+1}\SincRad(\pi t/2^{k})=\prod_{j\ge0}\SincRad(\pi\delta_nt/2^{j})
=\PhiUp(\delta_nt)\), so \(\PhiUp(t)=\Phin{n}(t)\,\PhiUp(\delta_nt)\), and
\cref{due:eq:mech-head-tail} is Fourier inversion of this product.  The
scale of the omitted tail is therefore \(\delta_n\), and the ratio base is
\(\Qq=\delta_{n+1}^{2}/\delta_n^{2}=1/4\).

\subsection{Definition, and the three normalizations of the same sequence}

\begin{lemma}[The reciprocal product near the origin]
\label{due:lem:mech-zero-free}
\(\PhiUp\) has no zero in the open unit disc \(\AbsoluteValue{z}<1\), and
\(\pm1\) are simple zeros.  Consequently \(1/\PhiUp\) is holomorphic and even
on \(\AbsoluteValue{z}<1\), equals \(1\) at \(0\), and its Taylor series at
the origin has radius of convergence exactly \(1\).
\end{lemma}

\begin{proof}
The factor \(\SincRad(\pi z/2^{k})=\sin(\pi z/2^{k})/(\pi z/2^{k})\) vanishes
exactly at \(z=2^{k}j\), \(j\in\IntegerNumbers\setminus\SetOf{0}\), each zero
being simple.  Every such point has modulus at least \(1\), and \(z=\pm1\) is
attained only by the pair \((k,j)=(0,\pm1)\); since the product converges
locally uniformly (catalogue entry \texttt{up-fourier-product}), its zero set
is the union of the zero sets of the factors with multiplicities added.
Evenness is inherited from the factors.  The radius statement is
Cauchy--Hadamard together with the simple pole of \(1/\PhiUp\) at \(\pm1\).
\end{proof}

\begin{definition}[Reciprocal-product coefficients]
\label{due:def:mech-am}
The sequence \((\am{m})_{m\ge0}\) is defined by
\begin{equation}
 \frac{1}{\PhiUp(z)}=\sum_{m=0}^{\infty}\am{m}\,(2\pi z)^{2m},
 \qquad\AbsoluteValue{z}<1 .
 \label{due:eq:mech-am-def}
\end{equation}
By \cref{due:lem:mech-zero-free} the series exists, has only even powers, and
\(\am{0}=1\).
\end{definition}

The factor \((2\pi)^{2m}\) is part of the definition: it makes \(\am{m}\) the
coefficient of the angular variable \(\omega=2\pi z\), in which the omitted
tail acts on a polynomial cell as the differential operator
\(\sum_m(-1)^m\am{m}\delta_n^{2m}D^{2m}\)
(\cref{due:lem:mech-finite-inverse}).  Three other descriptions of the same
sequence occur in the sources and in the catalogue; the next proposition
identifies them exactly.

Let \(X\) be a random variable with density \(\RvachevUp\).  Its
moment-generating function and its raw moments are
\begin{equation}
 M(z)\DefinitionEquals\Expectation\,\EulerE^{zX}
 =\int_{-1}^{1}\EulerE^{zx}\RvachevUp(x)\,\Differential x,
 \qquad
 m_n^{\mathrm{Rv}}\DefinitionEquals\Expectation X^{n}
 =\int_{-1}^{1}x^{n}\RvachevUp(x)\,\Differential x ,
 \label{due:eq:mech-mgf}
\end{equation}
so that \(m^{\mathrm{Rv}}_{2j}=c_j\) is the \(j\)-th even full moment of the
catalogue entry \texttt{full-moments} and \(m^{\mathrm{Rv}}_{2j+1}=0\)
(entry \texttt{rvachev-\allowbreak appell-\allowbreak deconvolution}).  That entry defines the
binomial-convolution reciprocal \((\gamma^{\mathrm{Rv}}_n)_{n\ge0}\) as the
unique rational sequence with
\begin{equation}
 \sum_{k=0}^{n}\binom nk\,m_k^{\mathrm{Rv}}\gamma^{\mathrm{Rv}}_{n-k}
 =\delta_{n,0}\qquad(n\in\NaturalNumbers).
 \label{due:eq:mech-gamma-rv-def}
\end{equation}
S1 expands the reciprocal moment-generating function \(1/M(z)\) in the form
\(\sum_j\gamma_j\,z^{j}/j!\); S6 expands it in the form
\(\sum_k\gamma_k\,z^{2k}/(2k)!\), with the odd indices suppressed.

\begin{proposition}[The three normalizations]
\label{due:prop:mech-normalizations}
\begin{enumerate}[label=\textup{(\roman*)}]
\item \(M\) is entire and even, \(M(z)=\sum_{n\ge0}m_n^{\mathrm{Rv}}z^{n}/n!\),
and
\begin{equation}
 M(z)=\PhiUp\!\left(\frac{\ImaginaryUnit z}{2\pi}\right)
 =\prod_{k=0}^{\infty}\frac{\sinh(z/2^{k+1})}{z/2^{k+1}},
 \qquad
 \PhiUp(t)=M(2\pi\ImaginaryUnit t).
 \label{due:eq:mech-mgf-product}
\end{equation}
\item For \(\AbsoluteValue{z}<2\pi\),
\begin{equation}
 \frac1{M(z)}=\sum_{n=0}^{\infty}\gamma^{\mathrm{Rv}}_n\frac{z^{n}}{n!}
 =\sum_{m=0}^{\infty}(-1)^{m}\am{m}z^{2m},
 \label{due:eq:mech-gamma-egf}
\end{equation}
so that \(\gamma^{\mathrm{Rv}}_{2m+1}=0\) and
\begin{equation}
 \boxed{\;
 \am{m}=(-1)^{m}\,\frac{\gamma^{\mathrm{Rv}}_{2m}}{(2m)!}
 \qquad(m\ge0).\;}
 \label{due:eq:mech-am-gamma}
\end{equation}
In particular S1's \(\gamma_j\) and S6's \(\gamma_k\) are
\(\gamma^{\mathrm{Rv}}_j\) and \(\gamma^{\mathrm{Rv}}_{2k}\), and S1's mode
coefficient \(\gamma_{2\ell}/(2\ell)!\) is \((-1)^{\ell}\am{\ell}\).
\item The even moments and the coefficients are reciprocal under the two
equivalent finite convolutions
\begin{equation}
 \sum_{j=0}^{m}(-1)^{m-j}\am{m-j}\frac{c_j}{(2j)!}=\delta_{m,0},
 \qquad
 \sum_{k=0}^{m}\binom{2m}{2k}\gamma^{\mathrm{Rv}}_{2k}\,c_{m-k}=\delta_{m,0}
 \qquad(m\in\NaturalNumbers).
 \label{due:eq:mech-moment-convolution}
\end{equation}
\end{enumerate}
\end{proposition}

\begin{proof}
(i) \(X\) is bounded, so \(M\) is entire with the displayed Taylor series,
and even because \(\RvachevUp\) is even.  Substituting
\(\xi=\ImaginaryUnit z/(2\pi)\) into
\(\PhiUp(\xi)=\int\RvachevUp(x)\EulerE^{-2\pi\ImaginaryUnit\xi x}\,\Differential x\)
gives \(\EulerE^{-2\pi\ImaginaryUnit\xi x}=\EulerE^{zx}\), hence
\(M(z)=\PhiUp(\ImaginaryUnit z/2\pi)\); inverting, \(\PhiUp(t)=M(-2\pi\ImaginaryUnit t)
=M(2\pi\ImaginaryUnit t)\).  Termwise,
\(\SincRad(\pi\cdot\ImaginaryUnit z/(2\pi2^{k}))=\SincRad(\ImaginaryUnit z/2^{k+1})
=\sinh(z/2^{k+1})/(z/2^{k+1})\), which is the product.

(ii) By \cref{due:lem:mech-zero-free} and (i), \(M\) has no zero in
\(\AbsoluteValue{z}<2\pi\) (its zeros are \(2\pi\ImaginaryUnit\) times the
zeros of \(\PhiUp\)), so \(1/M\) is holomorphic there.  Putting
\(z\mapsto\ImaginaryUnit z/(2\pi)\) in \cref{due:eq:mech-am-def} gives, for
\(\AbsoluteValue{z}<2\pi\),
\[
 \frac1{M(z)}=\sum_{m\ge0}\am{m}(\ImaginaryUnit z)^{2m}
 =\sum_{m\ge0}(-1)^{m}\am{m}z^{2m}.
\]
Write the same function as an exponential generating series
\(\sum_n g_nz^{n}/n!\); then \(g_{2m+1}=0\) and
\(g_{2m}/(2m)!=(-1)^m\am{m}\).  Multiplying the two convergent series
\(M(z)\cdot M(z)^{-1}=1\) and comparing the coefficient of \(z^{n}/n!\) gives
\(\sum_{k}\binom nk m_k^{\mathrm{Rv}}g_{n-k}=\delta_{n,0}\), which is
\cref{due:eq:mech-gamma-rv-def}.  Since \(m_0^{\mathrm{Rv}}=1\), the
relation \cref{due:eq:mech-gamma-rv-def} determines the sequence uniquely by
triangular recursion, so \(g_n=\gamma^{\mathrm{Rv}}_n\), which proves
\cref{due:eq:mech-gamma-egf,due:eq:mech-am-gamma}.

(iii) Multiply the catalogue's expansion
\(\PhiUp(z)=\sum_j(-1)^{j}c_j(2\pi z)^{2j}/(2j)!\) (entry
\texttt{full-moments}; it is \cref{due:eq:mech-mgf-product} read at
\(z=2\pi\ImaginaryUnit t\)) by \cref{due:eq:mech-am-def} and compare the
coefficients of \((2\pi z)^{2m}\): this is the first identity.  The second is
\cref{due:eq:mech-gamma-rv-def} at \(n=2m\) after discarding the odd terms,
which vanish.
\end{proof}

\emph{Formal versions.}  In module \texttt{RvachevMomentAppell}:
\begin{sloppypar}
\begin{itemize}
\item \texttt{Fabius.\allowbreak rvachevRawMomentRat} is \(m^{\mathrm{Rv}}_n\), with even
entries the executable rational moments \texttt{Fabius.moment}~\(j=c_j\) of
\texttt{Arithmetic};
\item \texttt{Fabius.\allowbreak rvachevReciprocalMomentRat} is \(\gamma^{\mathrm{Rv}}_n\),
and \texttt{Fabius.\allowbreak binomialConv\_\allowbreak rvachevRawMomentRat\_\allowbreak reciprocal} is
\cref{due:eq:mech-gamma-rv-def};
\item \texttt{Fabius.\allowbreak rvachevReciprocalMomentRat\_\allowbreak eq\_\allowbreak completeBellPolynomial}
is the Bell-polynomial form of that reciprocal.
\end{itemize}
\end{sloppypar}
Thus \cref{due:eq:mech-am-gamma} expresses every \(\am{m}\) as a
Lean-computable rational, \((-1)^m/(2m)!\) times the \(2m\)-th entry of
\texttt{rvachevReciprocalMomentRat}.  The analytic identification
\cref{due:eq:mech-gamma-egf} of the reciprocal sequence with the Taylor
coefficients of \(1/M\) is not formalized; the polynomial identities that use
it are (see \cref{due:rmk:mech-appell}).

\subsection{Rationality and positivity}

\begin{proposition}[Logarithm of the reciprocal product]
\label{due:prop:mech-log}
For \(\AbsoluteValue{z}<1\), with the branch of the logarithm that vanishes
at \(z=0\),
\begin{equation}
 \log\frac1{\PhiUp(z)}=\sum_{j=1}^{\infty}\alpha_j\,(2\pi z)^{2j},
 \qquad
 \boxed{\;
 \alpha_j=\frac{\RiemannZeta(2j)}{j\,(2\pi)^{2j}\,(1-\Qq^{j})}
 =\frac{\AbsoluteValue{B_{2j}}}{2j\,(2j)!\,(1-\Qq^{j})}
 \in\RationalNumbers_{>0},\;}
 \label{due:eq:mech-alpha}
\end{equation}
where \(B_{2j}\) are the Bernoulli numbers and \(\Qq=1/4\).  Equivalently,
for \(\AbsoluteValue{z}<2\pi\), the cumulants \(\kappa_j\) of \(X\) satisfy
\begin{equation}
 \log M(z)=\sum_{j\ge1}\kappa_{2j}\frac{z^{2j}}{(2j)!},
 \qquad
 \kappa_{2j}=\frac{B_{2j}}{2j\,(1-\Qq^{j})},
 \qquad
 \alpha_j=(-1)^{j+1}\frac{\kappa_{2j}}{(2j)!},
 \qquad
 \kappa_{2j+1}=0 .
 \label{due:eq:mech-kappa}
\end{equation}
\end{proposition}

\begin{proof}
For \(\AbsoluteValue{w}<\pi\), the product \(\sin w/w=\prod_{h\ge1}(1-w^{2}/(\pi^{2}h^{2}))\)
has all factors in the disc \(\AbsoluteValue{u-1}<1\), so
\[
 \log\frac{\sin w}{w}=\sum_{h\ge1}\log\Bigl(1-\frac{w^{2}}{\pi^{2}h^{2}}\Bigr)
 =-\sum_{h\ge1}\sum_{j\ge1}\frac{w^{2j}}{j\,\pi^{2j}h^{2j}}
 =-\sum_{j\ge1}\frac{\RiemannZeta(2j)}{j\,\pi^{2j}}\,w^{2j},
\]
the interchange being justified by absolute convergence
(\(\sum_h\sum_j\AbsoluteValue{w}^{2j}/(j\pi^{2j}h^{2j})
=-\sum_h\log(1-\AbsoluteValue{w}^{2}/(\pi^{2}h^{2}))<\infty\)).
Put \(w=\pi z/2^{k}\), \(k\ge0\), which satisfies \(\AbsoluteValue{w}<\pi\)
when \(\AbsoluteValue{z}<1\), and sum over \(k\):
\[
 \sum_{k\ge0}\log\SincRad\!\Bigl(\frac{\pi z}{2^{k}}\Bigr)
 =-\sum_{k\ge0}\sum_{j\ge1}\frac{\RiemannZeta(2j)}{j}\,z^{2j}\Qq^{jk}
 =-\sum_{j\ge1}\frac{\RiemannZeta(2j)}{j\,(1-\Qq^{j})}\,z^{2j},
\]
again by absolute convergence
(\(\RiemannZeta(2j)\le\RiemannZeta(2)\) and
\(\sum_k\sum_j\AbsoluteValue{z}^{2j}\Qq^{jk}/j=-\sum_k\log(1-\AbsoluteValue{z}^{2}\Qq^{k})<\infty\)).
The left side is a holomorphic function on the disc whose exponential is
\(\PhiUp(z)\) and which vanishes at \(0\); by \cref{due:lem:mech-zero-free}
the disc is a simply connected zero-free domain of \(\PhiUp\), so this
function is the branch of \(\log\PhiUp\) fixed in the statement.  Negating and
writing \(z^{2j}=(2\pi z)^{2j}/(2\pi)^{2j}\) gives the first form of
\(\alpha_j\); Euler's evaluation
\(\RiemannZeta(2j)=\AbsoluteValue{B_{2j}}(2\pi)^{2j}/(2\,(2j)!)\) gives the
second, which is visibly a positive rational.  Finally
\(\log M(z)=\log\PhiUp(\ImaginaryUnit z/2\pi)
=-\sum_j\alpha_j(\ImaginaryUnit z)^{2j}=\sum_j(-1)^{j+1}\alpha_jz^{2j}\)
for \(\AbsoluteValue{z}<2\pi\), and since \(\operatorname{sign}B_{2j}=(-1)^{j+1}\),
\((-1)^{j+1}(2j)!\,\alpha_j=B_{2j}/(2j(1-\Qq^{j}))\).  Odd cumulants vanish
because \(M\) is even.
\end{proof}

% ed.: S1 derives the cumulants probabilistically (cumulants add under
% convolution and scale homogeneously); that argument is recorded here as a
% remark and the analytic derivation above is the proof, because it also
% supplies the domain of validity of every series used later.
\begin{remark}[The probabilistic reading]
\label{due:rmk:mech-cumulants-probabilistic}
S1 obtains \cref{due:eq:mech-kappa} by writing \(X=\sum_{k\ge0}2^{-k-1}U_k\):
the \(2j\)-th cumulant of a uniform variable on \([-1,1]\) is
\(2^{2j-1}B_{2j}/j\) (from \(\log(\sinh z/z)\)), cumulants of independent
summands add, and the \(2j\)-th cumulant of \(cU\) is \(c^{2j}\) times that of
\(U\); summing \(\sum_{k\ge0}2^{-2j(k+1)}=\Qq^{j}/(1-\Qq^{j})\) gives
\(\kappa_{2j}=2^{2j-1}B_{2j}\Qq^{j}/(j(1-\Qq^{j}))=B_{2j}/(2j(1-\Qq^{j}))\).
The two derivations agree.
\end{remark}

\begin{theorem}[Rationality, positivity, recurrence, closed form, growth]
\label{due:thm:mech-am-rational-positive}
\begin{enumerate}[label=\textup{(\alph*)}]
\item \(\am{0}=1\) and, for \(m\ge1\),
\begin{equation}
 \boxed{\;m\,\am{m}=\sum_{j=1}^{m}j\,\alpha_j\,\am{m-j}.\;}
 \label{due:eq:mech-recurrence}
\end{equation}
\item In terms of the complete exponential Bell polynomials,
\begin{equation}
 \am{m}=\frac1{m!}\,
 \ExponentialCompleteBellPolynomial{m}\bigl(1!\,\alpha_1,\,2!\,\alpha_2,\ldots,m!\,\alpha_m\bigr)
 =\sum_{k_1+2k_2+\cdots+mk_m=m}\;\prod_{j=1}^{m}\frac{\alpha_j^{k_j}}{k_j!}.
 \label{due:eq:mech-bell}
\end{equation}
\item Every \(\am{m}\) is a positive rational number.  Equivalently
\((-1)^{m}\gamma^{\mathrm{Rv}}_{2m}>0\) for every \(m\); no even entry of the
reciprocal moment sequence vanishes.
\item \((2\pi)^{2m}\am{m}\to2/\PhiUp(1/2)=3.6115993860280564\ldots\) as
\(m\to\infty\), with
\((2\pi)^{2m}\am{m}-2/\PhiUp(1/2)=\BigOAt{r^{-2m}}{m\to\infty}\) for every
\(r<2\); here \(\PhiUp(1/2)=\prod_{k\ge1}\SincRad(\pi/2^{k})
=0.5537712758888100\ldots\).
\end{enumerate}
\end{theorem}

\begin{proof}
Put \(w=(2\pi z)^{2}\) and \(A(w)=\sum_{j\ge1}\alpha_jw^{j}\), so that
\cref{due:eq:mech-alpha,due:eq:mech-am-def} read
\(\sum_m\am{m}w^{m}=\exp A(w)\) for \(\AbsoluteValue{w}<4\pi^{2}\).  Since
\(A(0)=0\), the composition \(\exp\circ A\) is also a well-defined formal
power series and the analytic and formal coefficients agree.

(a) Differentiating \(F=\exp A\) gives \(F'=A'F\); the coefficient of
\(w^{m-1}\) is \cref{due:eq:mech-recurrence}.

(b) The exponential formula
\(\exp\bigl(\sum_{j\ge1}x_jw^{j}/j!\bigr)
=\sum_{m\ge0}\ExponentialCompleteBellPolynomial{m}(x_1,\ldots,x_m)\,w^{m}/m!\)
with \(x_j=j!\,\alpha_j\) gives the first form; expanding
\(\exp A=\prod_j\exp(\alpha_jw^{j})=\prod_j\sum_{k_j}\alpha_j^{k_j}w^{jk_j}/k_j!\)
and collecting \(w^{m}\) gives the second (a finite sum, because
\(k_j=0\) for \(j>m\)).

(c) By \cref{due:prop:mech-log} every \(\alpha_j\) is a positive rational.
Induction on \(m\) in \cref{due:eq:mech-recurrence}: if
\(\am{0},\ldots,\am{m-1}\) are positive rationals, the right side is a
positive rational and so is \(\am{m}\).  The reformulation is
\cref{due:eq:mech-am-gamma}.

% ed.: (d) is not in any source; it quantifies the size of the modes and is
% the analytic reason the extraction weights of the later sections need no
% growth hypothesis on the coefficients.
(d) By \cref{due:lem:mech-zero-free}, \(1/\PhiUp\) is meromorphic on
\(\AbsoluteValue{z}<2\) with simple poles at \(\pm1\) and no other pole there
(the zeros \(2^{k}j\) of \(\PhiUp\) other than \(\pm1\) have modulus at least
\(2\)).  Near \(z=1\), \(\SincRad(\pi z)=\sin(\pi z)/(\pi z)=(1-z)+\BigOAt{(1-z)^{2}}{z\to1}\)
and \(\prod_{k\ge1}\SincRad(\pi z/2^{k})=\PhiUp(z/2)\to\PhiUp(1/2)\), so the
residue of \(1/\PhiUp\) at \(1\) is \(-1/\PhiUp(1/2)\); by evenness the
residue at \(-1\) is \(1/\PhiUp(1/2)\).  Hence
\[
 g(z)\DefinitionEquals\frac1{\PhiUp(z)}-\frac1{\PhiUp(1/2)}
 \Bigl(\frac1{1-z}+\frac1{1+z}\Bigr)
 =\frac1{\PhiUp(z)}-\frac{2}{\PhiUp(1/2)}\,\frac1{1-z^{2}}
\]
is holomorphic on \(\AbsoluteValue{z}<2\), and Cauchy's estimate on the
circle \(\AbsoluteValue{z}=r<2\) bounds its \(z^{2m}\)-coefficient by
\(\max_{\AbsoluteValue{z}=r}\AbsoluteValue{g}\cdot r^{-2m}\).  The
\(z^{2m}\)-coefficient of \(1/\PhiUp\) is \((2\pi)^{2m}\am{m}\) and that of
\(2/(\PhiUp(1/2)(1-z^{2}))\) is \(2/\PhiUp(1/2)\).  Every factor
\(\SincRad(\pi/2^{k})\), \(k\ge1\), is positive, so \(\PhiUp(1/2)>0\); the
numerical values are the convergent product evaluated to the digits shown.
\end{proof}

The values \((2\pi)^{2m}\am{m}\) for \(m=1,\ldots,5\) are
\(2.193\), \(2.982\), \(3.369\), \(3.527\), \(3.584\); the next three are
\(3.603\), \(3.609\), \(3.611\), approaching the limit \(3.6116\) roughly at
the rate \(\Qq^{m}\) that the double pole of \(1/\PhiUp\) at \(\pm2\) (from
the factors \(k=0\) and \(k=1\)) dictates.

\begin{proposition}[\(q\)-Pochhammer coefficient formula]
\label{due:prop:mech-qpoch}
\(\PhiUp(z)=\prod_{h\ge1}\QPochhammer{z^{2}/h^{2}}{\Qq}{\infty}\) for all
\(z\in\ComplexNumbers\), and for every \(m\ge0\)
\begin{equation}
 \boxed{\;
 (2\pi)^{2m}\am{m}
 =\sum_{\substack{\nu_1,\nu_2,\ldots\ge0\\ \nu_1+\nu_2+\cdots=m}}
 \;\prod_{h=1}^{\infty}
 \frac{h^{-2\nu_h}}{\QPochhammer{\Qq}{\Qq}{\nu_h}},\;}
 \label{due:eq:mech-qpoch}
\end{equation}
the sum running over the finitely supported sequences \((\nu_h)_{h\ge1}\) of
total weight \(m\); each summand is positive and the sum converges.  In
S1's normalization the same statement reads
\(\gamma_{2m}/(2m)!=(-1)^{m}(4\pi^{2})^{-m}\sum_\nu\prod_hh^{-2\nu_h}/\QPochhammer{\Qq}{\Qq}{\nu_h}\).
\end{proposition}

\begin{proof}
Expanding each sinc factor by its Euler product,
\[
 \PhiUp(z)=\prod_{k\ge0}\prod_{h\ge1}\Bigl(1-\frac{z^{2}}{h^{2}4^{k}}\Bigr)
 =\prod_{h\ge1}\prod_{k\ge0}\Bigl(1-\frac{z^{2}}{h^{2}}\Qq^{k}\Bigr)
 =\prod_{h\ge1}\QPochhammer{z^{2}/h^{2}}{\Qq}{\infty};
\]
the rearrangement is legitimate because
\(\sum_{k,h}\AbsoluteValue{z}^{2}/(h^{2}4^{k})<\infty\) makes the double
product absolutely convergent.  For \(\AbsoluteValue{a}<1\) Euler's
\(q\)-exponential identity (\cite[(1.3.15)]{gasper-rahman}, \cite{dlmf-q})
\[
 \frac1{\QPochhammer{a}{\Qq}{\infty}}
 =\sum_{\nu=0}^{\infty}\frac{a^{\nu}}{\QPochhammer{\Qq}{\Qq}{\nu}}
\]
applies to \(a=z^{2}/h^{2}\) as soon as \(\AbsoluteValue{z}<1\).  Fix
\(H\ge1\) and let \(P_H(z)=\prod_{h\le H}\QPochhammer{z^{2}/h^{2}}{\Qq}{\infty}^{-1}\).
Multiplying \(H\) absolutely convergent power series,
\[
 P_H(z)=\sum_{m\ge0}S_{H,m}\,z^{2m},
 \qquad
 S_{H,m}=\sum_{\substack{\nu_1,\ldots,\nu_H\ge0\\ \nu_1+\cdots+\nu_H=m}}
 \prod_{h=1}^{H}\frac{h^{-2\nu_h}}{\QPochhammer{\Qq}{\Qq}{\nu_h}},
\]
a finite sum of positive terms that increases with \(H\).  As
\(H\to\infty\), \(P_H\to1/\PhiUp\) locally uniformly on
\(\AbsoluteValue{z}<1\), because the tails
\(\prod_{h>H}\QPochhammer{z^{2}/h^{2}}{\Qq}{\infty}\to1\) locally uniformly
and \(\PhiUp\) is zero-free there; by Cauchy's formula the coefficients
converge, \(S_{H,m}\to(2\pi)^{2m}\am{m}\).  A monotone sequence of finite
sums of positive terms converges to the sum of the infinite series of those
terms, which is the right side of \cref{due:eq:mech-qpoch}.  The S1 form is
\cref{due:eq:mech-am-gamma}.
\end{proof}

For \(m=1\) the formula reads \((2\pi)^{2}\am{1}=\RiemannZeta(2)/(1-\Qq)
=\tfrac43\cdot\tfrac{\pi^{2}}{6}=\tfrac{2\pi^{2}}9\), i.e.\ \(\am{1}=1/18\);
for \(m=2\) it reads
\((2\pi)^{4}\am{2}=\RiemannZeta(4)/\QPochhammer{\Qq}{\Qq}{2}
+\tfrac12(\RiemannZeta(2)^{2}-\RiemannZeta(4))/\QPochhammer{\Qq}{\Qq}{1}^{2}
=\tfrac{32}{2025}\pi^{4}+\tfrac{2}{135}\pi^{4}=\tfrac{62}{2025}\pi^{4}\), i.e.\
\(\am{2}=62/(2025\cdot16)=31/16200\), in agreement with
\cref{due:tab:mech-am}.  The same quarter-base \(q\)-calculus governs the
extraction weights of the later sections; that is the conceptual content of
the formula.  For computation, \cref{due:eq:mech-recurrence} is the tool.

% ed.: the sources disagree on which argument proves positivity; the
% orchestrating question was "which proof establishes positivity, not just
% nonvanishing".  This remark records the answer exactly.
\begin{remark}[Which argument proves what]
\label{due:rmk:mech-positivity-routes}
Three of the sources' arguments prove strict positivity and one does not.
\begin{itemize}
\item S1's proposition ``Rational cumulants and reciprocal coefficients''
proves \emph{rationality only}: its recurrence
\(\gamma_n=-\sum_{j=1}^{n}\binom{n-1}{j-1}\kappa_j\gamma_{n-j}\) (the
cumulant form of \cref{due:eq:mech-gamma-rv-def}) has terms of both signs,
so no sign information follows from it.  Positivity in S1 comes from its
\(q\)-Pochhammer formula, \cref{due:prop:mech-qpoch} here, whose summands are
all positive: this is the second proof of
\cref{due:thm:mech-am-rational-positive}(c).
\item S2, S3 and S4 use the exponential recurrence
\cref{due:eq:mech-recurrence} with \(\alpha_j>0\): the proof adopted above.
\item S6's lemma ``Rationality and nonvanishing of the reciprocal moments''
proves, despite its title, the strict inequality \((-1)^{k}\gamma_{k}>0\) for
its \(\gamma_k=\gamma^{\mathrm{Rv}}_{2k}\), that is, \(\am{k}>0\).  Its
argument: from
\cref{due:eq:mech-mgf-product}, \(1/M(z)=\prod_{k\ge0}(z/2^{k+1})/\sinh(z/2^{k+1})\),
and \(w/\sinh w=\sum_{r\ge0}(2-2^{2r})B_{2r}w^{2r}/(2r)!=\sum_r(-1)^{r}b_rw^{2r}\)
with \(b_0=1\) and \(b_r=2(2^{2r-1}-1)\AbsoluteValue{B_{2r}}/(2r)!>0\) for
\(r\ge1\).  The \(z^{2m}\)-coefficient of the product of the first \(K\)
factors is \((-1)^{m}\) times a finite sum of positive terms containing the
term \(b_m4^{-m}\) (all of \(2m\) taken from the factor \(k=0\)); the partial
products converge locally uniformly on \(\AbsoluteValue{z}<2\pi\) to
\(1/M\), so the coefficients converge and
\begin{equation}
 \am{m}=(-1)^{m}[z^{2m}]\frac1{M(z)}\;\ge\;\frac{b_m}{4^{m}}
 =\frac{2(2^{2m-1}-1)\AbsoluteValue{B_{2m}}}{4^{m}\,(2m)!}\qquad(m\ge1),
 \label{due:eq:mech-lower-bound}
\end{equation}
a third proof, with an explicit lower bound (\(1/24\le1/18\),
\(7/5760\le31/16200\), \ldots).
\end{itemize}
The justification of coefficient extraction from an infinite product of
series, which S1 and S6 pass over, is the locally uniform convergence used
in \cref{due:prop:mech-qpoch} and in the third bullet.
\end{remark}

\begin{table}[htbp]
\centering
\caption{The first reciprocal-product coefficients, their logarithmic
coefficients \(\alpha_m\) and cumulants \(\kappa_{2m}\)
(\cref{due:prop:mech-log}), and the catalogue's reciprocal moments
\(\gamma^{\mathrm{Rv}}_{2m}=(-1)^{m}(2m)!\,\am{m}\).  The values
\(\am{0},\ldots,\am{5}\) are the ground truth of the verification outputs of
S4 and S5; the rest are computed by \cref{due:eq:mech-recurrence} in exact
arithmetic.}
\label{due:tab:mech-am}
\footnotesize
\renewcommand{\arraystretch}{1.25}
\begin{tabular}{@{}c l l@{}}
\toprule
\(m\) & \(\am{m}\) & \(\gamma^{\mathrm{Rv}}_{2m}=(-1)^{m}(2m)!\,\am{m}\)\\
\midrule
0 & \(1\) & \(1\)\\
1 & \(1/18\) & \(-1/9\)\\
2 & \(31/16200\) & \(31/675\)\\
3 & \(2347/42865200\) & \(-2347/59535\)\\
4 & \(1904369/1311675120000\) & \(1904369/32531625\)\\
5 & \(3309760579/88561680752160000\) & \(-3309760579/24405225075\)\\
6 & \(1884491755011469/1980124051433319792000000\) & \(1884491755011469/4133856862760625\)\\
\bottomrule
\end{tabular}
\par\medskip
\begin{tabular}{@{}c l l@{}}
\toprule
\(m\) & \(\alpha_m\) & \(\kappa_{2m}\)\\
\midrule
1 & \(1/18\) & \(1/9\)\\
2 & \(1/2700\) & \(-2/225\)\\
3 & \(1/178605\) & \(16/3969\)\\
4 & \(1/9639000\) & \(-16/3825\)\\
5 & \(1/478533825\) & \(256/33759\)\\
6 & \(691/15688261338750\) & \(-353792/16769025\)\\
\bottomrule
\end{tabular}
\end{table}

\subsection{Finite deconvolution on a polynomial jet}

For \(\delta>0\) and a polynomial \(P\) define the tail-averaging operator
\begin{equation}
 (\mathcal A_\delta P)(t)\DefinitionEquals
 \int_{-1}^{1}P(t-\delta y)\,\RvachevUp(y)\,\Differential y
 =(P\Convolution\rho_\delta)(t),
 \label{due:eq:mech-averaging}
\end{equation}
the second form by the substitution \(u=\delta y\) in
\cref{due:eq:mech-dilated-kernel}.  Writing \(D\) for \(\Differential/\Differential t\):

\begin{lemma}[Symbol of the tail average]
\label{due:lem:mech-symbol}
On the space of polynomials of degree at most \(n\),
\begin{equation}
 \mathcal A_\delta=M(\delta D)=\sum_{j=0}^{\Floor{n/2}}\frac{c_j}{(2j)!}\,\delta^{2j}D^{2j},
 \label{due:eq:mech-symbol}
\end{equation}
so \(\mathcal A_\delta\) preserves degree and leading coefficient, commutes
with \(D\), and \((\mathcal A_\delta P)^{(r)}=\mathcal A_\delta(P^{(r)})\) for all
\(r\ge0\).
\end{lemma}

\begin{proof}
Taylor's formula for a polynomial is finite:
\(P(t-\delta y)=\sum_{k=0}^{n}P^{(k)}(t)(-\delta y)^{k}/k!\).  Integrating
against \(\RvachevUp(y)\,\Differential y\) and using
\(m^{\mathrm{Rv}}_k=0\) for odd \(k\), \(m^{\mathrm{Rv}}_{2j}=c_j\), gives
\cref{due:eq:mech-symbol}.  The \(j=0\) term is the identity and every other
term lowers degree.  Differentiation under the integral sign in
\cref{due:eq:mech-averaging} gives the last assertion.
\end{proof}

\begin{lemma}[Finite polynomial deconvolution]
\label{due:lem:mech-finite-inverse}
Let \(P\) be a polynomial of degree at most \(n\) and \(\delta>0\).  Then
\begin{equation}
 \boxed{\;
 P=\sum_{m=0}^{\Floor{n/2}}(-1)^{m}\am{m}\,\delta^{2m}D^{2m}(\mathcal A_\delta P),\;}
 \label{due:eq:mech-finite-inverse}
\end{equation}
and for every \(r\ge0\) and every \(t\),
\begin{equation}
 P^{(r)}(t)=\sum_{m=0}^{\Floor{(n-r)/2}}(-1)^{m}\am{m}\,\delta^{2m}
 (\mathcal A_\delta P)^{(r+2m)}(t),
 \label{due:eq:mech-finite-inverse-jet}
\end{equation}
where an empty sum (\(r>n\)) is \(0\).  Both identities hold in the
polynomial ring \(\Kk[t]\) with \(\Kk=\RationalNumbers\) whenever
\(P\in\RationalNumbers[t]\) and \(\delta\in\RationalNumbers\); no limiting
process is involved.
\end{lemma}

\begin{proof}
On the space \(V_n\) of polynomials of degree at most \(n\), \(D\) is
nilpotent, \(D^{n+1}=0\), so every power series in \(\delta D\) is a
polynomial in \(D\) of degree at most \(n\), and products of such series are
computed by finite coefficient convolutions.  Put
\(R\DefinitionEquals\sum_{m\le\Floor{n/2}}(-1)^{m}\am{m}\delta^{2m}D^{2m}\).
Composing with \cref{due:eq:mech-symbol},
\[
 R\,\mathcal A_\delta
 =\sum_{s\ge0}\delta^{2s}D^{2s}
 \sum_{m+j=s}(-1)^{m}\am{m}\frac{c_j}{(2j)!}
 =\sum_{s\ge0}\delta_{s,0}\,\delta^{2s}D^{2s}
 =\mathrm{id}_{V_n},
\]
where the terms with \(2s>n\) vanish on \(V_n\) regardless of their
coefficients and the terms with \(2s\le n\) are evaluated by the first
identity of \cref{due:eq:mech-moment-convolution}.  This is
\cref{due:eq:mech-finite-inverse}.  Apply \(D^{r}\) to it; since \(D\)
commutes with \(R\) and with \(\mathcal A_\delta\) and
\(D^{r+2m}\mathcal A_\delta P=0\) once \(r+2m>n\), only \(m\le\Floor{(n-r)/2}\)
survive, which is \cref{due:eq:mech-finite-inverse-jet}.
\end{proof}

\begin{remark}[The catalogue's Appell form of the same inverse]
\label{due:rmk:mech-appell}
With \(A_n^{\mathrm{Rv}}(x)=\sum_k\binom nk\gamma^{\mathrm{Rv}}_{n-k}x^{k}\) the
Rvachev--Appell polynomials of the entry
\texttt{rvachev-\allowbreak appell-\allowbreak deconvolution}, one has
\(A_n^{\mathrm{Rv}}=\sum_{j}\gamma^{\mathrm{Rv}}_j D^{j}(x^{n})/j!
=\sum_m(-1)^m\am{m}D^{2m}(x^{n})\) by
\cref{due:eq:mech-am-gamma}, so the catalogue's deconvolution operator
\(\mathscr D_{\RvachevUp}\colon x^{n}\mapsto A^{\mathrm{Rv}}_n\) is exactly the
operator \(R\) of the proof at \(\delta=1\), and the catalogue's smoothing
identity \(\int(\mathscr D_{\RvachevUp}P)(x-y)\RvachevUp(y)\,\Differential y=P(x)\)
is \(\mathcal A_1\mathscr D_{\RvachevUp}=\mathrm{id}\), the other order of
\cref{due:eq:mech-finite-inverse}.  At scale \(\delta\), conjugating by the
dilation \(t\mapsto t/\delta\) gives
\(\mathcal A_\delta^{-1}P=\sum_kp_k\delta^{k}A^{\mathrm{Rv}}_k(t/\delta)\)
for \(P=\sum_kp_kt^{k}\).  \emph{Formal versions} (all in \texttt{RvachevMomentAppell}):
\begin{itemize}
\item \texttt{Fabius.\allowbreak rvachevAppellPolynomialRat}, the rational Appell family;
\item \texttt{Fabius.\allowbreak rvachevDeconvolvedPolynomial}, the operator
\(\mathscr D_{\RvachevUp}\);
\item \texttt{Fabius.\allowbreak integral\_\allowbreak eval\_\allowbreak rvachevAppellPolynomial\_\allowbreak sub\_\allowbreak mul\_\allowbreak rvachev},
the smoothing identity with \(x-y\).
\end{itemize}
Consequently \cref{due:lem:mech-finite-inverse} at \(\delta=1\) is
Lean-proved in the Appell normalization; the dilation to general \(\delta\) is
a change of variable.
\end{remark}

The lemma is applied on the cell that the tail window
\([x-\delta_n,x+\delta_n]\) occupies.  We recall the geometric input from
\cref{due:sec:setting} in the precise form needed, with its short proof.

\begin{lemma}[A dyadic point is the centre of one cell]
\label{due:lem:mech-cell}
Let \(x\in[-1,1]\) be dyadic with \(\dep{x}=s\) and let \(n\ge s\).  Then
there is one polynomial \(P_{n,x}\) of degree at most \(n\) with
\begin{equation}
 \upn{n}(z)=P_{n,x}(z)\qquad\text{for all }z\in(x-\delta_n,x+\delta_n);
 \label{due:eq:mech-cell-polynomial}
\end{equation}
for \(x=\pm1\) it is the zero polynomial.  For \(n\ge1\) and \(0\le r\le n-1\),
\({\upn{n}}^{(r)}(x)=P_{n,x}^{(r)}(x)\); for \(r=n\), \(P_{n,x}^{(n)}\) is the
constant value of the piecewise-constant \(n\)-th derivative of \(\upn{n}\) on
the open cell.  If \(x\in(-1,1)\) and \(n=s-1\ge0\), then \(x\) is a knot of
\(\upn{n}\).
\end{lemma}

\begin{proof}
Write \(x=a/2^{s}\) with \(a\) odd if \(s\ge1\) and \(a\in\SetOf{-1,0,1}\)
if \(s=0\).  Since \(\delta_n=2^{-n-1}\), \(x/\delta_n=a\,2^{\,n+1-s}\), an
even integer when \(n\ge s\) and the odd integer \(a\) when \(n=s-1\).  The
knots of \(\upn{n}\) are the odd multiples of \(\delta_n\) in
\([-(1-\delta_n),1-\delta_n]\).  If \(\AbsoluteValue{x}<1\), then
\(\AbsoluteValue{a}\le2^{s}-1\), so
\(\AbsoluteValue{x/\delta_n}\le2^{n+1}-2^{n+1-s}\le2^{n+1}-2\): \(x\) is an
even multiple of \(\delta_n\) at distance at least \(\delta_n\) from the two
extreme knots \(\pm(2^{n+1}-1)\delta_n\), hence the midpoint of the open cell
between the two neighbouring odd multiples, on which \(\upn{n}\) is one
polynomial.  If \(x=\pm1\), the interval \((x-\delta_n,x+\delta_n)\) lies
outside the support \([-(1-\delta_n),1-\delta_n]\) except for its endpoint,
so \(\upn{n}=0\) there.  The derivative statements follow from
\(\upn{n}\in C^{n-1}\) and from the polynomial form on the open cell.  When
\(n=s-1\) and \(\AbsoluteValue{x}<1\), \(x\) is an odd multiple of
\(\delta_n\) of modulus at most \((2^{n+1}-1)\delta_n\), a knot.
\end{proof}

\begin{theorem}[Exact local deconvolution on a centred cell]
\label{due:thm:mech-local-deconvolution}
Let \(x\in[-1,1]\) be dyadic with \(\dep{x}=s\), let \(n\ge s\), and let
\(P=P_{n,x}\) be the cell polynomial of \cref{due:lem:mech-cell}.  Then, for
\emph{every} \(r\ge0\),
\begin{align}
 \RvachevUp^{(r)}(x)
 &=(\mathcal A_{\delta_n}P)^{(r)}(x)
 =\sum_{j=0}^{\Floor{(n-r)/2}}\frac{c_j}{(2j)!}\,\delta_n^{2j}\,P^{(r+2j)}(x),
 \label{due:eq:mech-forward-jet}\\
 P^{(r)}(x)
 &=\sum_{m=0}^{\Floor{(n-r)/2}}(-1)^{m}\am{m}\,\delta_n^{2m}\,
 \RvachevUp^{(r+2m)}(x).
 \label{due:eq:mech-local-deconvolution}
\end{align}
In particular, at \(r=0\), since \(P(x)=\upn{n}(x)\) and
\(\delta_n^{2m}=4^{-m(n+1)}=\Qq^{m}\,4^{-mn}\),
\begin{equation}
 \boxed{\;
 \upn{n}(x)=\sum_{m=0}^{\Floor{n/2}}(-1)^{m}\am{m}\,4^{-m(n+1)}\,
 \RvachevUp^{(2m)}(x)
 =\RvachevUp(x)+\sum_{m=1}^{\Floor{n/2}}
 \Cr{m}{x}\,4^{-mn},\;}
 \label{due:eq:mech-local-value}
\end{equation}
with \(\Cr{m}{x}=(-1)^{m}\am{m}4^{-m}\RvachevUp^{(2m)}(x)\): an identity for
each \(n\ge s\), with no remainder.
\end{theorem}

\begin{proof}
% ed.: S1, S2 and S6 differentiate the spline distributionally and restrict
% to r <= n; S4 and S5 differentiate the smooth kernel instead.  The kernel
% route is adopted because it needs no distribution theory and holds for all
% r, which is what makes the second proof of jet termination in
% \cref{due:rmk:mech-second-proof} available.
By \cref{due:eq:mech-head-tail}, \(\RvachevUp=\upn{n}\Convolution\rho_{\delta_n}\)
with \(\upn{n}\) bounded and integrable and \(\rho_{\delta_n}\) smooth with
compact support.  Differentiation under the integral sign therefore gives,
for every \(r\ge0\),
\[
 \RvachevUp^{(r)}(x)=\int_{\RealNumbers}\upn{n}(z)\,\rho_{\delta_n}^{(r)}(x-z)\,\Differential z .
\]
The integrand vanishes unless \(\AbsoluteValue{x-z}\le\delta_n\); on the
open interval \((x-\delta_n,x+\delta_n)\) we have \(\upn{n}=P\) by
\cref{due:lem:mech-cell}, and the two endpoints have measure zero.  Hence
\[
 \RvachevUp^{(r)}(x)=\int_{\RealNumbers}P(z)\,\rho_{\delta_n}^{(r)}(x-z)\,\Differential z
 =(P\Convolution\rho_{\delta_n})^{(r)}(x)=(\mathcal A_{\delta_n}P)^{(r)}(x),
\]
the middle equality because \(P\Convolution\rho_{\delta_n}\) is a polynomial
whose \(r\)-th derivative is obtained by differentiating the kernel, and the
last by \cref{due:eq:mech-averaging}.  Expanding with
\cref{due:eq:mech-symbol} (applied to \(P^{(r)}\), of degree at most
\(n-r\)) gives \cref{due:eq:mech-forward-jet}.  Now
\cref{due:eq:mech-finite-inverse-jet} with \(\delta=\delta_n\), evaluated at
\(t=x\), reads
\(P^{(r)}(x)=\sum_m(-1)^m\am{m}\delta_n^{2m}(\mathcal A_{\delta_n}P)^{(r+2m)}(x)\),
and each \((\mathcal A_{\delta_n}P)^{(r+2m)}(x)\) equals
\(\RvachevUp^{(r+2m)}(x)\) by the identity just proved (at order \(r+2m\)).
This is \cref{due:eq:mech-local-deconvolution}; \(r=0\) and
\(P(x)=\upn{n}(x)\) give \cref{due:eq:mech-local-value}.
\end{proof}

\begin{warningbox}[title={An all-orders asymptotic expansion is not an identity}]
For an arbitrary \(x\), dividing \(\PhiUp\) by \(\PhiUp(\delta_nt)\) on the
Fourier side yields, for every truncation order \(K\),
\(\upn{n}=\sum_{m\le K}(-1)^m\am{m}\delta_n^{2m}\RvachevUp^{(2m)}+\BigOAt{\delta_n^{2K+2}}{n\to\infty}\)
in the uniform norm; this is the asymptotic Richardson calculus of the
Exponents volume \cite{proveit-exponents} (its \texttt{p3:eq:richardson-\allowbreak moments}
and \texttt{p3:eq:richardson-\allowbreak leading}).  At a dyadic point the displayed
derivatives eventually vanish (\cref{due:thm:mech-jet-termination}), but that
alone does not force the remainder to vanish: a remainder smaller than every
power of \(4^{-n}\) is not excluded by any finite number of vanishing terms.
What excludes it is \cref{due:thm:mech-local-deconvolution}: on the centred
cell the reciprocal series acts on a nilpotent operator, so the identity is a
finite algebraic one and there is no remainder to estimate.  The onset
\(n\ge\dep{x}\) is exactly the range in which the tail window
\([x-\delta_n,x+\delta_n]\) crosses no knot.
\end{warningbox}

\begin{remark}[Six forms of one lemma]
\label{due:rmk:mech-six-forms}
\cref{due:lem:mech-finite-inverse,due:thm:mech-local-deconvolution} absorb
S1's ``Finite jet inversion'' and ``Exact local deconvolution of a sinc
prefix'' (stated with \(\gamma_k\delta^{k}/k!\) over all \(k\), odd terms
vanishing), S2's ``Finite local deconvolution'' (the \(a_m\) form with
\(0\le r\le n\), inverting the triangular relation
\cref{due:eq:mech-forward-jet}), S3's ``Finite polynomial deconvolution''
(the operator \(M(\delta D)\) on polynomials, applied to a cell polynomial of
degree at most \(s\), see \cref{due:rmk:mech-second-proof}), S4's ``Exact
inverse on polynomials'' and ``The tail sees only one polynomial piece'' (the
kernel-differentiation route adopted here, in the \(p_N\) indexing with
\(4^{-mN}=4^{-m(n+1)}\)), S5's ``Finite polynomial inverse'' and ``Jet
transfer at a cell centre'' (the angular normalization
\(\Phi(-h^{2}D^{2})\) with \(h_N=2^{-N}=\delta_n\)), and S6's ``Exact
reciprocal moment deconvolution on a matched cell'' (the form
\(\sum_k\gamma_k4^{-k(n+1)}F^{(2k)}(a)/(2k)!\) for the centred CDF spline
\(S_n\) at \(a=1-\AbsoluteValue{x}\), where \(F^{(2k)}(a)=\RvachevUp^{(2k)}(x)\)
and \(S_n(a)=\upn{n}(x)\)).  Under \cref{due:prop:mech-normalizations} and the
index bridge of \cref{due:sec:setting} all six are the identity
\cref{due:eq:mech-local-value}.
\end{remark}

% ed.: S3's odd-level onset is the one genuinely stronger statement among
% the six; the endpoint case s = 1 is separated from it because it depends on
% a convention for the value of a jump.
\begin{proposition}[One stage earlier at odd depth]
\label{due:prop:mech-odd-onset}
Let \(x\in(-1,1)\) be dyadic with \(\dep{x}=s=2d+1\) odd, \(s\ge3\), and put
\(n=s-1=2d\).  Then \cref{due:eq:mech-local-value} also holds at this
\(n\):
\begin{equation}
 \upn{s-1}(x)=\sum_{m=0}^{d}(-1)^{m}\am{m}\,4^{-ms}\,\RvachevUp^{(2m)}(x).
 \label{due:eq:mech-odd-onset}
\end{equation}
For \(s=1\), \(x=\pm1/2\), the same formula at \(n=0\) reads
\(\upn{0}(\pm1/2)=\RvachevUp(\pm1/2)=1/2\) and holds if and only if the jump
of the box \(\upn{0}=\IndicatorOf{[-1/2,1/2]}\) is assigned its Fourier-inversion
value \(1/2\).
\end{proposition}

\begin{proof}
By \cref{due:lem:mech-cell}, \(x\) is a knot of \(\upn{n}\); the tail window
\([x-\delta_n,x+\delta_n]\) covers the right half of the cell
\((x-2\delta_n,x)\) and the left half of the cell \((x,x+2\delta_n)\).  Let
\(P_-\) and \(P_+\) be the cell polynomials of \(\upn{n}\) on these two cells.
Since \(\upn{n}\in C^{n-1}\) (\(n=2d\ge2\)), the difference \(P_+-P_-\)
vanishes to order \(n\) at \(x\), so \(P_+-P_-=c\,(z-x)^{n}\) for a constant
\(c\).  Put \(P_*\DefinitionEquals\tfrac12(P_-+P_+)\), so that
\(P_\pm=P_*\pm\tfrac c2(z-x)^{n}\) and \(P_*(x)=P_\pm(x)=\upn{n}(x)\).  As in
the proof of \cref{due:thm:mech-local-deconvolution}, for every \(r\ge0\),
\begin{align*}
 \RvachevUp^{(r)}(x)
 &=\int_{x-\delta_n}^{x}P_-(z)\rho_{\delta_n}^{(r)}(x-z)\,\Differential z
  +\int_{x}^{x+\delta_n}P_+(z)\rho_{\delta_n}^{(r)}(x-z)\,\Differential z\\
 &=(\mathcal A_{\delta_n}P_*)^{(r)}(x)
  +\frac c2\int_{0}^{\delta_n}u^{n}\Bigl[\rho_{\delta_n}^{(r)}(-u)-(-1)^{n}\rho_{\delta_n}^{(r)}(u)\Bigr]\Differential u,
\end{align*}
by the substitutions \(z=x+u\) and \(z=x-u\) in the two half-windows.  The
kernel \(\rho_{\delta_n}\) is even, so \(\rho_{\delta_n}^{(r)}\) has parity
\((-1)^{r}\), and the bracket vanishes whenever \(r\equiv n\pmod2\).  Since
\(n\) is even, \(\RvachevUp^{(r)}(x)=(\mathcal A_{\delta_n}P_*)^{(r)}(x)\)
for every even \(r\).  Now \cref{due:eq:mech-finite-inverse} for \(P_*\)
(degree at most \(n=2d\)), evaluated at \(x\), uses only the even-order
derivatives \((\mathcal A_{\delta_n}P_*)^{(2m)}(x)\), \(m\le d\):
\[
 \upn{n}(x)=P_*(x)=\sum_{m=0}^{d}(-1)^{m}\am{m}\delta_n^{2m}
 (\mathcal A_{\delta_n}P_*)^{(2m)}(x)
 =\sum_{m=0}^{d}(-1)^{m}\am{m}\delta_n^{2m}\RvachevUp^{(2m)}(x),
\]
and \(\delta_n^{2m}=4^{-m(n+1)}=4^{-ms}\).  For \(s=1\), \(n=0\): the two
cell polynomials of the box at \(x=1/2\) are the constants \(1\) and \(0\),
the same computation gives \(\RvachevUp(1/2)=\int_0^{\delta_0}\rho_{\delta_0}(u)\,\Differential u
=\tfrac12\), and \cref{due:eq:mech-odd-onset} with \(d=0\) asserts
\(\upn{0}(1/2)=1/2\), the mean of the two one-sided limits, which is the
value the inverse Fourier integral of \(\SincRad(\pi t)\) takes at a jump.
\end{proof}

The parity computation also shows why the even-depth case admits no such
sharpening in general: for \(n=s-1\) odd the bracket does not vanish for even
\(r\), and the exact examples of \cref{due:sec:main} (for instance
\(x=1/4\), where \(\upn{1}(1/4)=1\) while the law gives
\(67/72+\tfrac19\cdot\tfrac14=23/24\)) show that
\cref{due:eq:mech-local-value} fails at \(n=s-1\) for even \(s\).

\begin{summarybox}[title={What the first mechanism delivers}]
The coefficients \(\am{m}\) of \(1/\PhiUp(z)=\sum_m\am{m}(2\pi z)^{2m}\) are
positive rationals, computable by the Bernoulli recurrence
\cref{due:eq:mech-recurrence}, equal to \((-1)^{m}\gamma^{\mathrm{Rv}}_{2m}/(2m)!\)
in the catalogue's reciprocal-moment normalization, and given in closed
\(q\)-Pochhammer form by \cref{due:eq:mech-qpoch}.  At a dyadic point of
depth \(s\) and every \(n\ge s\) (and at \(n=s-1\) when \(s\ge3\) is odd),
\(\upn{n}(x)=\sum_{m\le\Floor{n/2}}(-1)^{m}\am{m}4^{-m(n+1)}\RvachevUp^{(2m)}(x)\)
exactly, because the deconvolution series is a polynomial in the nilpotent
operator \(D\) on the cell polynomial.  The sum still runs to
\(\Floor{n/2}\); cutting it to \(\Floor{s/2}\), independently of \(n\), is
the second mechanism.
\end{summarybox}

\section{Termination of the dyadic derivative jet}
\label{due:sec:jet}

The second mechanism is arithmetic.  Rvachev's functional-differential
equation (catalogue entry \texttt{rvachev-up}, ``Differential law'')
\begin{equation}
 \RvachevUp'(x)=2\RvachevUp(2x+1)-2\RvachevUp(2x-1)\qquad(x\in\RealNumbers)
 \label{due:eq:mech-refinement}
\end{equation}
expresses each derivative of \(\RvachevUp\) as a signed sum of values of
\(\RvachevUp\) at dilated arguments; at a dyadic point the dilated arguments
eventually all land on odd integers, where \(\RvachevUp\) vanishes.

\subsection{The Thue--Morse derivative stencil}

Recall \(\ThueMorseSign{k}=(-1)^{s_2(k)}\), \(s_2(k)\) the binary digit sum, so
that \(\ThueMorseSign{2k}=\ThueMorseSign{k}\) and
\(\ThueMorseSign{2k+1}=-\ThueMorseSign{k}\).

\begin{theorem}[Thue--Morse derivative stencil]
\label{due:thm:mech-stencil}
For every \(r\ge0\) and every \(x\in\RealNumbers\),
\begin{equation}
 \boxed{\;
 \RvachevUp^{(r)}(x)=2^{\binom{r+1}{2}}
 \sum_{k=0}^{2^{r}-1}\ThueMorseSign{k}\,
 \RvachevUp\bigl(2^{r}x+2^{r}-1-2k\bigr).\;}
 \label{due:eq:mech-stencil}
\end{equation}
\end{theorem}

\begin{proof}
Induction on \(r\).  For \(r=0\) the sum has the single term
\(\RvachevUp(x)\).  Assume \cref{due:eq:mech-stencil} at order \(r\) and
differentiate; the chain rule contributes \(2^{r}\), and
\cref{due:eq:mech-refinement} applied to each term gives
\begin{align*}
 \RvachevUp^{(r+1)}(x)
 &=2^{\binom{r+1}{2}+r+1}\sum_{k<2^{r}}\ThueMorseSign{k}
 \Bigl[\RvachevUp\bigl(2^{r+1}x+2^{r+1}-1-2(2k)\bigr)\\
 &\hspace{9em}-\RvachevUp\bigl(2^{r+1}x+2^{r+1}-1-2(2k+1)\bigr)\Bigr].
\end{align*}
The first bracketed term carries the index \(2k\) with sign
\(\ThueMorseSign{k}=\ThueMorseSign{2k}\), the second the index \(2k+1\) with sign
\(-\ThueMorseSign{k}=\ThueMorseSign{2k+1}\); together they enumerate
\(0\le k'<2^{r+1}\).  Finally \(\binom{r+1}{2}+r+1=\binom{r+2}{2}\).
\end{proof}

\begin{corollary}[Exact derivative tiling; global form]
\label{due:cor:mech-tiling}
For \(r\ge0\), \(0\le k<2^{r}\) and \(y\in[-1,1]\),
\begin{equation}
 \boxed{\;
 \RvachevUp^{(r)}\Bigl(-1+\frac{2k+1+y}{2^{r}}\Bigr)
 =2^{\binom{r+1}{2}}\,\ThueMorseSign{k}\,\RvachevUp(y).\;}
 \label{due:eq:mech-tiling}
\end{equation}
Equivalently, with the signed global extension
\(\FabiusGlobal(t)=\sum_{b\ge0}\ThueMorseSign{b}\RvachevUp(t-2b-1)\) of the
catalogue entry \texttt{fabius-global},
\begin{equation}
 \RvachevUp^{(r)}(x)=2^{\binom{r+1}{2}}\,\FabiusGlobal\bigl(2^{r}(x+1)\bigr)
 \qquad(x\in[-1,1],\ r\ge0).
 \label{due:eq:mech-global-form}
\end{equation}
\end{corollary}

\begin{proof}
At \(x=-1+(2k+1+y)/2^{r}\) the argument of the \(j\)-th term of
\cref{due:eq:mech-stencil} is \(2(k-j)+y\).  For \(j\ne k\) and
\(\AbsoluteValue{y}\le1\) it has modulus at least \(1\), where
\(\RvachevUp\) vanishes (\(\RvachevUp(\pm1)=0\) and the support is
\([-1,1]\)); the term \(j=k\) is \(\ThueMorseSign{k}\RvachevUp(y)\).  For
\cref{due:eq:mech-global-form}: \(t=2^{r}(x+1)\in[0,2^{r+1}]\) lies in a block
\([2k,2k+2]\) with \(0\le k<2^{r}\) (for \(t=2^{r+1}\) take \(k=2^{r}-1\)),
say \(t=2k+1+y\) with \(y\in[-1,1]\); in the locally finite sum defining
\(\FabiusGlobal(t)\) only \(b=k\) can contribute, the neighbouring translates
being evaluated at \(\pm1\) or beyond, so \(\FabiusGlobal(t)=\ThueMorseSign{k}\RvachevUp(y)\),
and \cref{due:eq:mech-tiling} applies.
\end{proof}

\emph{Formal versions.}
\begin{sloppypar}
\begin{itemize}
\item \texttt{Fabius.\allowbreak iteratedDeriv\_\allowbreak rvachevUp\_\allowbreak eq\_\allowbreak thueMorse\_\allowbreak sum}
(\texttt{Differential}) is \cref{due:thm:mech-stencil}, term for term;
\item \texttt{Fabius.\allowbreak iteratedDeriv\_\allowbreak rvachev} (\texttt{GlobalExtension}) is
the global form \cref{due:eq:mech-global-form};
\item \texttt{Fabius.\allowbreak iteratedDeriv\_\allowbreak extendedFabius} (same module) is the
scaling law \(\FabiusGlobal^{(r)}(t)=2^{\binom{r+1}{2}}\FabiusGlobal(2^{r}t)\)
that the catalogue's \texttt{fabius-dyadic-cells} entry (field ``Derivative
scaling'') invokes.
\end{itemize}
\end{sloppypar}
S5 and S6 derive the jet results from that scaling law; S1, S2, S3 and S4
from the stencil.  The two routes are the same computation read through
\cref{due:cor:mech-tiling}.

\subsection{Values at the two shallowest points, and strict positivity}

\begin{lemma}[Values and positivity]
\label{due:lem:mech-values}
\(\RvachevUp(0)=1\), \(\RvachevUp(\pm1/2)=1/2\), and \(\RvachevUp(y)>0\) for
every \(y\in(-1,1)\).
\end{lemma}

\begin{proof}
% ed.: the two values are derived from the first mechanism at d = 0, as S3
% does, rather than from a separate argument; positivity follows S4.
The point \(0\) has depth \(0\) and \(\pm1/2\) depth \(1\).
\cref{due:eq:mech-local-value} at \(x=0\), \(n=0\), has the single term
\(m=0\) (\(\Floor{0/2}=0\)): \(\RvachevUp(0)=\upn{0}(0)=1\), the box
\(b_0\) at its centre.  At \(x=\pm1/2\), \(n=1\), again only \(m=0\):
\(\RvachevUp(\pm1/2)=\upn{1}(\pm1/2)=(b_0\Convolution b_1)(\pm1/2)
=2\int_{1/4}^{3/4}\IndicatorOf{[-1/2,1/2]}(y)\,\Differential y=\tfrac12\).
(Both agree with the catalogue: \(\RvachevUp(0)=1\) is normative, and
\(\RvachevUp(1/2)=\FabiusBounded(1/2)=1/2\) by the reflection
\(\FabiusBounded(t)+\FabiusBounded(1-t)=1\).)

Positivity.  First, \(\upn{n}(y)>0\) for \(\AbsoluteValue{y}<1-\delta_n\), by
induction on \(n\): \(\upn{0}=b_0\) is positive on \((-1/2,1/2)\); and if
\(\upn{n-1}\) is nonnegative and positive on the open interval
\((-(1-\delta_{n-1}),1-\delta_{n-1})\), then for
\(\AbsoluteValue{y}<1-\delta_n=(1-\delta_{n-1})+2^{-n-1}\),
\[
 \upn{n}(y)=(\upn{n-1}\Convolution b_n)(y)
 =2^{n}\int_{y-2^{-n-1}}^{y+2^{-n-1}}\upn{n-1}(u)\,\Differential u>0,
\]
because the interval of integration meets the positivity interval of
\(\upn{n-1}\) in a nonempty open interval.  Now fix \(y\in(-1,1)\) and choose
\(n\ge1\) with \(2\delta_n<1-\AbsoluteValue{y}\).  For
\(\AbsoluteValue{u}\le\delta_n\), \(\AbsoluteValue{y-u}<1-\delta_n\), so the
continuous function \(u\mapsto\upn{n}(y-u)\) is positive on the compact
support \([-\delta_n,\delta_n]\) of \(\rho_{\delta_n}\), and
\cref{due:eq:mech-head-tail} gives
\(\RvachevUp(y)=\int\upn{n}(y-u)\rho_{\delta_n}(u)\,\Differential u
\ge\min_{\AbsoluteValue{u}\le\delta_n}\upn{n}(y-u)>0\).
\end{proof}

\emph{Formal versions.}
\begin{sloppypar}
\begin{itemize}
\item \texttt{Fabius.rvachevUp\_zero} (\texttt{Basic}): \(\RvachevUp(0)=1\);
\item \texttt{Fabius.\allowbreak rvachevUp\_\allowbreak pos\_\allowbreak of\_\allowbreak mem\_\allowbreak Ioo} (\texttt{Monotonicity}):
positivity on \((-1,1)\).
\end{itemize}
\end{sloppypar}

\subsection{The jet stops at the depth, and its last even entry is not zero}

\begin{theorem}[Dyadic jet termination]
\label{due:thm:mech-jet-termination}
Let \(x\in[-1,1]\) be dyadic with \(\dep{x}=s\), and write \(x=a/2^{s}\)
(\(a\) odd if \(s\ge1\)).
\begin{enumerate}[label=\textup{(\roman*)}]
\item \(\RvachevUp^{(r)}(x)=0\) for every \(r>s\).
\item If \(s\ge1\), the derivative of order exactly \(s\) is a signed power
of two,
\begin{equation}
 \RvachevUp^{(s)}(x)=2^{\binom{s+1}{2}}\,\ThueMorseSign{k},
 \qquad k=\frac{2^{s}(x+1)-1}{2}=\frac{a+2^{s}-1}{2}\in\SetOf{0,\ldots,2^{s}-1}.
 \label{due:eq:mech-top-derivative}
\end{equation}
\item If \(s\ge1\) and \(0\le r<s\), then \(\RvachevUp^{(r)}(x)\ne0\); in
particular, for odd \(s\),
\begin{equation}
 \RvachevUp^{(s-1)}(x)=2^{\binom{s}{2}-1}\,\ThueMorseSign{k'},
 \qquad k'=\Floor{\frac{2^{s}(x+1)-1}{4}}.
 \label{due:eq:mech-subtop-derivative}
\end{equation}
\end{enumerate}
Thus, for \(s\ge1\), \(s=\max\SetOf{r:\RvachevUp^{(r)}(x)\ne0}\) and the
set of orders with a nonzero derivative is exactly \(\SetOf{0,1,\ldots,s}\);
for \(s=0\) (the points \(0,\pm1\)) every derivative of positive order
vanishes.
\end{theorem}

\begin{proof}
Put \(t_r\DefinitionEquals2^{r}(x+1)=2^{r-s}(a+2^{s})\).

(i) For \(r>s\), \(t_r\) is an even integer (for \(s=0\) directly, since
\(a+1\in\SetOf{0,1,2}\) and \(r\ge1\); for \(s\ge1\) because
\(r-s\ge1\)).  Then every argument \(2^{r}x+2^{r}-1-2j=t_r-1-2j\) in
\cref{due:eq:mech-stencil} is an odd integer, of modulus at least \(1\),
where \(\RvachevUp\) vanishes.  So the whole sum is zero.  (Equivalently,
\(t_r\) is a block endpoint in \cref{due:eq:mech-tiling}, \(y=\pm1\).)

(ii) For \(r=s\ge1\), \(t_s=a+2^{s}\) is odd, and \(0<t_s<2^{s+1}\) because
\(\AbsoluteValue{a}<2^{s}\); write \(t_s=2k+1\) with \(0\le k<2^{s}\).  This
is \cref{due:eq:mech-tiling} with \(y=0\), and \(\RvachevUp(0)=1\).

(iii) For \(0\le r<s\), \(t_r=(a+2^{s})/2^{s-r}\) is not an integer, so
\(t_r\in(2k,2k+2)\) for a unique \(k\), \(0\le k<2^{r}\) (as
\(0<t_r<2^{r+1}\)), and \(t_r=2k+1+y\) with \(y\in(-1,1)\).  By
\cref{due:eq:mech-tiling} and \cref{due:lem:mech-values},
\(\RvachevUp^{(r)}(x)=2^{\binom{r+1}{2}}\ThueMorseSign{k}\RvachevUp(y)\ne0\).
For odd \(s\) and \(r=s-1\), \(t_{s-1}=(a+2^{s})/2\) is a half-integer:
writing the odd number \(A=a+2^{s}=2^{s}(x+1)\) as \(4k'+1\) or \(4k'+3\)
gives \(t_{s-1}=2k'+1\mp\tfrac12\), so \(y=\mp\tfrac12\) with
\(k'=\Floor{(A-1)/4}\), and \(\RvachevUp(\pm1/2)=1/2\) gives
\cref{due:eq:mech-subtop-derivative}.
\end{proof}

\emph{Formal versions.}
\begin{sloppypar}
\begin{itemize}
\item \texttt{Fabius.\allowbreak iteratedDeriv\_\allowbreak rvachevUp\_\allowbreak dyadic\_\allowbreak eq\_\allowbreak zero}
(\texttt{DyadicDerivativeFiltration}) is (i);
\item \texttt{Fabius.\allowbreak iteratedDeriv\_\allowbreak rvachevUp\_\allowbreak dyadic\_\allowbreak critical} (same
module) is (ii), in the form
\(\RvachevUp^{(s)}((2c+1)/2^{s})=-\ThueMorseSign{c}2^{\binom{s+1}{2}}\) for
\(0\le c<2^{s-1}\), which is \cref{due:eq:mech-top-derivative} with
\(k=c+2^{s-1}\) and \(\ThueMorseSign{c+2^{s-1}}=-\ThueMorseSign{c}\);
\item \texttt{Fabius.\allowbreak dyadic\_\allowbreak depth\_\allowbreak eq\_\allowbreak max\_\allowbreak nonzero\_\allowbreak iteratedDeriv} (same
module) is their conjunction, ``the depth is the largest order with a
nonzero derivative'';
\item \texttt{Fabius.\allowbreak iteratedDeriv\_\allowbreak rvachev\_\allowbreak centeredDyadic\_\allowbreak eq\_\allowbreak zero} and
\texttt{Fabius.\allowbreak iteratedDeriv\_\allowbreak rvachev\_\allowbreak centeredDyadic\_\allowbreak ne\_\allowbreak zero}
(\texttt{OriginalPaperSupplement}) are the same two facts in the
centred-cell coordinates of the original paper;
\item \texttt{Fabius.\allowbreak iteratedDeriv\_\allowbreak fabiusReal\_\allowbreak dyadicGrid\_\allowbreak eq\_\allowbreak zero\_\allowbreak iff}
(\texttt{BoundedDerivatives}) is the ``zero iff \(m\) even'' form on the
matched grid \(m/2^{k}\), \(m<2^{k}\).
\end{itemize}
\end{sloppypar}
Part (iii) for \(0<r<s\) is not formalized as a separate statement; it is
the positivity lemma composed with the tiling.

% ed.: S6 states the cutoff for even derivatives in terms of half the
% depth; this corollary is the exact reconciliation, and it also carries
% S1's "highest even derivative is nonzero" with its magnitude.
\begin{corollary}[The even jet: half the depth]
\label{due:cor:mech-even-jet}
Let \(x\in(-1,1)\) be dyadic with \(\dep{x}=s\) and \(d=\Floor{s/2}\).  Then
\begin{equation}
 \RvachevUp^{(2m)}(x)\ne0\iff m\le d;
 \label{due:eq:mech-even-jet}
\end{equation}
at the two endpoints \(x=\pm1\) (depth \(0\)) the whole jet vanishes,
\(\RvachevUp^{(r)}(\pm1)=0\) for all \(r\ge0\).  For
\(s\ge2\) the last nonzero even entry has magnitude
\begin{equation}
 \AbsoluteValue{\RvachevUp^{(2d)}(x)}=
 \begin{cases}
 2^{\binom{s+1}{2}}, & s=2d,\\[0.3em]
 2^{\binom{s}{2}-1}, & s=2d+1,
 \end{cases}
 \label{due:eq:mech-top-even}
\end{equation}
with the signs \(\ThueMorseSign{k}\) and \(\ThueMorseSign{k'}\) of
\cref{due:thm:mech-jet-termination}.  In S6's terms, with
\(a=1-\AbsoluteValue{x}\) and \(F=\FabiusBounded\), \(F^{(2m)}(a)=\RvachevUp^{(2m)}(x)\)
vanishes for \(m>d\) and is nonzero for \(m\le d\): ``beyond half the dyadic
depth'' for the even derivatives is the same cutoff as ``beyond the depth''
for all derivatives, because \(2m>s\iff m>\Floor{s/2}\).
\end{corollary}

\begin{proof}
\(2m\le s\iff m\le d\).  For \(2m\le s\) the derivative is nonzero by
\cref{due:thm:mech-jet-termination}(ii)--(iii) (and by
\cref{due:lem:mech-values} when \(m=0\)); for \(2m>s\) it vanishes by (i).
At \(x=\pm1\), (i) with \(s=0\) gives the vanishing for \(r\ge1\), and
\(\RvachevUp(\pm1)=0\).  The magnitudes are
\cref{due:eq:mech-top-derivative} at \(2d=s\) and
\cref{due:eq:mech-subtop-derivative} at \(2d=s-1\).  For the CDF form: for
\(x\ge0\), \(\RvachevUp(x)=\FabiusBounded(1-x)\) (catalogue entry
\texttt{fabius-relations}), so \(\RvachevUp^{(2m)}(x)=(-1)^{2m}F^{(2m)}(1-x)=F^{(2m)}(a)\);
for \(x<0\) use evenness.
\end{proof}

\begin{remark}[A second proof of termination, from the first mechanism; local
degree collapse]
\label{due:rmk:mech-second-proof}
% ed.: this equivalence is editorial.  It is what makes S3's Prouhet
% cancellation and the other five sources' jet-termination arguments two
% proofs of one fact rather than two different theorems.
\cref{due:thm:mech-local-deconvolution} was proved for every \(r\ge0\), and
its right side \cref{due:eq:mech-forward-jet} is empty when \(r>n\).  Taking
\(n=s\) gives, without any Thue--Morse computation,
\(\RvachevUp^{(r)}(x)=(\mathcal A_{\delta_s}P_{s,x})^{(r)}(x)=0\) for
\(r>s\), because \(\mathcal A_{\delta_s}P_{s,x}\) is a polynomial of degree at
most \(s\): the jet of \(\RvachevUp\) at \(x\) terminates because
\(\RvachevUp\) is, at \(x\), the average of a degree-\(s\) polynomial over a
window that crosses no knot.  Conversely, for every \(n\ge s\),
\cref{due:eq:mech-local-deconvolution} at orders \(r>s\) gives
\(P_{n,x}^{(r)}(x)=0\), so
\begin{equation}
 \deg P_{n,x}\le s\qquad(n\ge s),
 \label{due:eq:mech-degree-collapse}
\end{equation}
and at order \(r=s\) it gives \(P_{n,x}^{(s)}(x)=\RvachevUp^{(s)}(x)\)
(the terms \(m\ge1\) vanish by (i)).  Thus the cell polynomial of
\emph{every} \(\upn{n}\), \(n\ge s\), at a point of depth \(s\) has degree
exactly \(s\) with the \(n\)-independent leading coefficient
\(\RvachevUp^{(s)}(x)/s!=2^{\binom{s+1}{2}}\ThueMorseSign{k}/s!\), although
\(\upn{n}\) has global degree \(n\).  This is S3's theorem ``Local degree
collapse'', proved there by Prouhet cancellation in the truncated-power
formula (blocks of \(2^{n-s}\) consecutive Thue--Morse signs annihilate
polynomials of degree below \(n-s\)).  The two statements, jet termination
and degree collapse, are equivalent through
\cref{due:thm:mech-local-deconvolution}: each implies the other by one of
\cref{due:eq:mech-forward-jet,due:eq:mech-local-deconvolution}.  By the catalogue's caveat in \texttt{fabius-dyadic-cells}, no local
analyticity is inferred from this finite jet, and none is needed anywhere in
this volume.

\emph{Formal versions.}
\begin{sloppypar}
\begin{itemize}
\item \texttt{Fabius.\allowbreak iteratedDeriv\_\allowbreak fabiusUniformSpline\_\allowbreak dyadic\_\allowbreak center}
(\texttt{PlateauLimit}) is the constancy of the leading coefficient, the
``plateau'' of the Lean corpus: in the centred-CDF picture at the anchor
\(2^{-r}\) (depth \(r\)), the \(r\)-th derivative of
\texttt{fabiusUniformSpline}~\(p\) equals \(2^{\binom{r+1}{2}}\) for every
\(p\ge r\), the value of \(\RvachevUp^{(r)}\) there;
\item \texttt{Fabius.\allowbreak rvachevDyadicTaylorPolynomial\_\allowbreak natDegree\_\allowbreak eq}
(\texttt{OriginalPaperSupplement}): the Taylor polynomial of \(\RvachevUp\)
at a centred dyadic of level \(s\) has degree exactly \(s\).
\end{itemize}
\end{sloppypar}
\end{remark}

\subsection{The two mechanisms together: the exact number of modes}

\begin{corollary}[Finite geometric tail with all modes present]
\label{due:cor:mech-exact-modes}
Let \(x\in[-1,1]\) be dyadic with \(\dep{x}=s\) and \(d=\Floor{s/2}\).  For
every \(n\ge s\), and also for \(n=s-1\) when \(s\ge3\) is odd,
\begin{equation}
 \boxed{\;
 \upn{n}(x)=\sum_{m=0}^{d}(-1)^{m}\am{m}\,4^{-m(n+1)}\,\RvachevUp^{(2m)}(x)
 =\RvachevUp(x)+\sum_{m=1}^{d}\Cr{m}{x}\,4^{-mn},\;}
 \label{due:eq:mech-exact-modes}
\end{equation}
with \(\Cr{m}{x}=(-1)^{m}\am{m}4^{-m}\RvachevUp^{(2m)}(x)\ne0\) for every
\(1\le m\le d\).  In particular:
\begin{itemize}
\item for \(s\in\SetOf{0,1}\) (the points \(0,\pm1,\pm1/2\)) there is no
mode at all, \(\upn{n}(x)=\RvachevUp(x)\) for every \(n\ge s\);
\item for \(s\ge2\) the tail has exactly \(d\ge1\) nonconstant geometric
modes \(4^{-mn}\), \(1\le m\le d\), with the top coefficient
\begin{equation}
 \AbsoluteValue{\Cr{d}{x}}=\am{d}\,4^{-d}\,\AbsoluteValue{\RvachevUp^{(2d)}(x)}
 =\begin{cases}
 \am{d}\,2^{\binom{s+1}{2}-2d}, & s=2d,\\[0.3em]
 \am{d}\,2^{\binom{s}{2}-2d-1}, & s=2d+1,
 \end{cases}
 \label{due:eq:mech-top-mode}
\end{equation}
so the mode count \(d\) is sharp and the sample count \(d+1\) of the
extraction in \cref{due:sec:main} cannot be lowered within the mode model.
\end{itemize}
All quantities in \cref{due:eq:mech-exact-modes} lie in
\(\Kk=\RationalNumbers\).
\end{corollary}

\begin{proof}
\cref{due:eq:mech-local-value} holds for \(n\ge s\), and
\cref{due:eq:mech-odd-onset} for \(n=s-1\) when \(s\ge3\) is odd; in both the
sum runs to \(\Floor{n/2}\ge d\) (for \(n=s-1=2d\) it runs exactly to \(d\)).
By \cref{due:cor:mech-even-jet} (and \cref{due:thm:mech-jet-termination}(i)
at \(x=\pm1\)), \(\RvachevUp^{(2m)}(x)=0\) for \(m>d\), so
the sum stops at \(d\); and for \(1\le m\le d\), \(\am{m}>0\)
(\cref{due:thm:mech-am-rational-positive}) and \(\RvachevUp^{(2m)}(x)\ne0\)
give \(\Cr{m}{x}\ne0\).  The rewriting
\(4^{-m(n+1)}=4^{-m}4^{-mn}\) produces \(\Cr{m}{x}\), and
\cref{due:eq:mech-top-even} gives \cref{due:eq:mech-top-mode}.  For \(s\le1\),
\(d=0\).

% ed.: rationality by the Vandermonde argument of S6, which needs no
% separate arithmetic theorem about up at dyadic points.
Rationality.  The samples \(\upn{n}(x)\), \(n=s,\ldots,s+d\), are rational by
the truncated-power formula of \cref{due:sec:setting}.  Read at these
\(d+1\) values of \(n\), \cref{due:eq:mech-exact-modes} is a square linear
system for the \(d+1\) unknowns \(\RvachevUp(x),\Cr{1}{x},\ldots,\Cr{d}{x}\)
whose matrix \(\bigl(4^{-mn}\bigr)_{n,m}\) is the Vandermonde matrix of the
distinct rational nodes \(4^{-s},\ldots,4^{-(s+d)}\); it is invertible over
\(\RationalNumbers\), so the unknowns are rational, and then so is
\(\RvachevUp^{(2m)}(x)=(-1)^{m}4^{m}\Cr{m}{x}/\am{m}\).  (This recovers the
classical rationality of \(\RvachevUp\) at dyadic points from the finite
splines alone.)  The sample count \(d+1\) cannot be lowered within the model:
the kernel of the \(d\times(d+1)\) system obtained from any \(d\) of the
levels is spanned by the coefficient vector of
\(\prod_{n}(t-4^{-n})\) over those levels, whose constant term
\(\pm\prod_n4^{-n}\) is not zero, so \(d\) samples leave the constant term
\(\RvachevUp(x)\) undetermined.
\end{proof}

The exact examples of the later sections confirm the top coefficient
literally: at \(x=1/4\) (\(s=2\), \(d=1\)) it is
\(\am{1}2^{3-2}=1/9\); at \(x=1/8\) (\(s=3\), \(d=1\)) it is
\(\am{1}2^{3-3}=1/18\); at \(x=1/16\) (\(s=4\), \(d=2\)) it is
\(\am{2}2^{10-4}=64\cdot31/16200=248/2025\); at \(x=1/32\) (\(s=5\),
\(d=2\)) it is \(\am{2}2^{10-5}=124/2025\); the signs are
\(\ThueMorseSign{k}\) or \(\ThueMorseSign{k'}\) times \((-1)^{d}\).

\begin{checkpointbox}[title={The shape of the argument so far}]
\begin{enumerate}
\item \emph{Cell geometry} (\cref{due:sec:setting}, recalled in
\cref{due:lem:mech-cell}): for \(n\ge\dep{x}\) the tail window
\([x-\delta_n,x+\delta_n]\) is one polynomial cell of \(\upn{n}\).
\item \emph{First mechanism} (\cref{due:thm:mech-local-deconvolution}):
on that cell the deconvolution identity \cref{due:eq:mech-local-value}
holds exactly, a finite sum to \(\Floor{n/2}\) with the positive rational
coefficients \(\am{m}\) of \cref{due:thm:mech-am-rational-positive}.
\item \emph{Second mechanism} (\cref{due:thm:mech-jet-termination},
\cref{due:cor:mech-even-jet}): \(\RvachevUp^{(2m)}(x)=0\) for \(m>\Floor{s/2}\)
and \(\ne0\) for \(m\le\Floor{s/2}\).
\item \emph{Together} (\cref{due:cor:mech-exact-modes}): an exact
\(d\)-mode geometric law from \(n=\dep{x}\) on, every mode present.  The
extraction of \(\RvachevUp(x)\) from \(d+1\) consecutive samples by the
quarter-base Gaussian-binomial row, and the recovery of every mode, are the
subject of \cref{due:sec:main} and the sections that follow it.
\end{enumerate}
\end{checkpointbox}

% ---- fragment 03-main-theorem ----
% ==========================================================================
% Cluster 03: the exact eventual geometric law and its sharpness
% ==========================================================================

\section{The exact eventual geometric law}
\label{due:sec:main}

The three mechanisms are now in place: a dyadic point is the centre of one
polynomial cell of $\upn{n}$ from level $n=\dep{x}$ on
(\cref{due:sec:depth}); on such a cell the omitted tail is deconvolved
exactly by the finite reciprocal-product operator with the rational
coefficients $\am{m}$ (\cref{due:sec:coefficients}); and the derivative jet
of $\RvachevUp$ at the point terminates at the order $\dep{x}$
(\cref{due:sec:jet}).  This section assembles them into the proof of
\cref{due:thm:set-main}, deferred from \cref{due:sec:result}, and records
the immediate consequences that the six sources state in various forms: the
derivative version of the law, the rationality of everything in sight, the
rational tail generating function, and the eventually annihilating
$q$-difference recurrence.  \Cref{due:sec:sharpness} then shows that
nothing in the statement can be improved.

Throughout, $x\in[-1,1]$ is dyadic, $s=\dep{x}$, $d=\Floor{s/2}$,
$\Qq=1/4$, $\delta_n=2^{-(n+1)}$, and $\qn{n}=\upn{n}(x)$.

\subsection{Proof of the main theorem}
\label{due:sec:main-proof}

\begin{proof}[Proof of \cref{due:thm:set-main}]
Fix $n\ge s$.

\emph{Step 1: one polynomial on the whole tail window.}
By \cref{due:lem:set-midpoint} (equivalently \cref{due:lem:mech-cell}),
$x$ is the midpoint of a knot cell of $\upn{n}$: the closed window
$[x-\delta_n,x+\delta_n]$ contains no knot in its interior, and
\cref{due:thm:set-finite-spline} makes $\upn{n}$ a single polynomial
$P=P_{n,x}$ of degree at most $n$ on it, with $P(x)=\upn{n}(x)$.

\emph{Step 2: exact local deconvolution.}
By \cref{due:prop:set-head-tail}, $\RvachevUp=\upn{n}\Convolution T_{\delta_n}$
with $T_{\delta}=\delta^{-1}\RvachevUp(\cdot/\delta)$ supported in
$[-\delta,\delta]$.  The convolution at $x$, and all its derivatives at
$x$, therefore see only the polynomial $P$, and
\cref{due:thm:mech-local-deconvolution} gives, at $r=0$,
\begin{equation}
 \upn{n}(x)=P(x)
 =\sum_{m=0}^{\Floor{n/2}}(-1)^{m}\am{m}\,\delta_n^{2m}\,\RvachevUp^{(2m)}(x),
 \qquad \delta_n^{2m}=4^{-m(n+1)} .
 \label{due:eq:main-step2}
\end{equation}
This is an identity, with no remainder: the reciprocal series acts on a
polynomial of degree at most $n$ and stops by itself.

\emph{Step 3: termination.}
By \cref{due:thm:mech-jet-termination}(i), $\RvachevUp^{(r)}(x)=0$ for every
$r>s$, so the terms of \cref{due:eq:main-step2} with $2m>s$, that is with
$m>d$, vanish; and $n\ge s$ gives $\Floor{n/2}\ge d$, so all terms with
$m\le d$ are present.  Hence
\begin{equation}
 \upn{n}(x)=\sum_{m=0}^{d}(-1)^{m}\am{m}\,4^{-m(n+1)}\,\RvachevUp^{(2m)}(x)
 \qquad(n\ge s),
 \label{due:eq:main-step3}
\end{equation}
which is \cref{due:eq:set-main-law-shifted}.  Writing
$4^{-m(n+1)}=4^{-m}\cdot4^{-mn}$ and separating the term $m=0$ gives
\cref{due:eq:set-main-law} with
$\Cr{m}{x}=(-1)^{m}\am{m}4^{-m}\RvachevUp^{(2m)}(x)$.

\emph{Step 4: rationality.}
The samples $\qn{n}$ are rational by \cref{due:cor:set-scaled-cell}, and
the $\am{m}$ are positive rationals by
\cref{due:thm:mech-am-rational-positive}.  Read \cref{due:eq:main-step3} at
the $d+1$ levels $n=s,\ldots,s+d$ as a linear system for the $d+1$ unknowns
$\RvachevUp^{(2m)}(x)$, $0\le m\le d$: its matrix
$\bigl((-1)^{m}\am{m}4^{-m(n+1)}\bigr)_{n,m}$ is a Vandermonde matrix in
the distinct nodes $4^{-s},\ldots,4^{-(s+d)}$ times an invertible diagonal
matrix, hence invertible over $\RationalNumbers$
(\cref{due:lem:rec-det} gives the determinant).  So every
$\RvachevUp^{(2m)}(x)$, $0\le m\le d$, and every $\Cr{m}{x}$ is a rational
combination of rational samples, and all quantities in
\cref{due:eq:set-main-law} lie in $\Kk=\RationalNumbers$.
\end{proof}

% ed.: none of the six sources isolates the rationality step from the
% ed.: classical arithmetic of Arias de Reyna; the volume derives it from the
% ed.: law itself, so the volume is self-contained on this point.
\begin{remark}[What the proof used, and what it did not]
\label{due:rmk:main-ingredients}
The proof uses no property of $\RvachevUp$ beyond the head--tail
factorization, the knot geometry of the finite splines, the evenness of the
tail, and the termination of the dyadic jet.  In particular it does not use
the classical rationality of $\RvachevUp$ at dyadic points
\cite{arias1982,arias2018}, which it re-derives (\cref{due:cor:main-rational}),
nor any convergence statement: the only limit in the volume is the one that
defines $\RvachevUp$ itself.  The same four ingredients, with the level
$n=s-1$ in place of $n\ge s$, are what \cref{due:sec:sharp-onset} analyses;
the difference is that at $n=s-1$ the window contains two polynomial pieces.
\end{remark}

\subsection{The derivative version}
\label{due:sec:main-derivative}

The identity \cref{due:eq:mech-local-deconvolution} holds for every
derivative order, and the jet terminates at the same order $s$ for all of
them, so the law has a derivative version with fewer modes.  For $n\ge s$
and $0\le r\le n$ write $\upn{n}^{(r)}(x)$ for the $r$-th derivative at $x$
of the cell polynomial $P_{n,x}$; for $r\le n-1$ this is the ordinary
$r$-th derivative of $\upn{n}\in C^{\,n-1}$, and for $r=n$ it is the
common one-sided value of the $n$-th derivative on the open cell.

\begin{corollary}[Eventual law for the derivatives]
\label{due:cor:main-derivative}
Let $0\le r\le s$ and $d_r=\Floor{(s-r)/2}$.  For every $n\ge s$,
\begin{equation}
 \boxed{\;
 \upn{n}^{(r)}(x)
 =\sum_{m=0}^{d_r}(-1)^{m}\am{m}\,4^{-m(n+1)}\,\RvachevUp^{(r+2m)}(x)
 =\RvachevUp^{(r)}(x)+\sum_{m=1}^{d_r}\Cr{r,m}{x}\,4^{-mn},\;}
 \label{due:eq:main-derivative-law}
\end{equation}
with $\Cr{r,m}{x}=(-1)^{m}\am{m}4^{-m}\RvachevUp^{(r+2m)}(x)\ne0$ for
$1\le m\le d_r$; for $r>s$ and $n\ge s$ both $\upn{n}^{(r)}(x)$ and
$\RvachevUp^{(r)}(x)$ vanish.  Consequently the extraction row of order
$d_r$ (\cref{due:thm:ext-extraction} with $d_r$ in place of $d$) recovers
$\RvachevUp^{(r)}(x)$ exactly from the $d_r+1$ consecutive derivative
samples $\upn{N}^{(r)}(x),\ldots,\upn{N+d_r}^{(r)}(x)$, for any $N\ge s$.
\end{corollary}

\begin{proof}
\Cref{due:eq:mech-local-deconvolution} at order $r$ reads
\[
 P^{(r)}(x)=\sum_{m\le\Floor{(n-r)/2}}(-1)^{m}\am{m}\delta_n^{2m}
 \RvachevUp^{(r+2m)}(x),
\]
and $\RvachevUp^{(r+2m)}(x)=0$ once $r+2m>s$
(\cref{due:thm:mech-jet-termination}(i)), which cuts the sum at $m=d_r$;
since $n\ge s$, $\Floor{(n-r)/2}\ge\Floor{(s-r)/2}=d_r$, so no term is
lost.  The nonvanishing of the modes is
\cref{due:thm:mech-jet-termination}(iii): $\RvachevUp^{(r+2m)}(x)\ne0$ for
$r+2m\le s$.  For $r>s$ the cell polynomial has degree at most $n$, and
the same identity with an empty right side gives $P^{(r)}(x)=0$, since
every $\RvachevUp^{(r+2m)}(x)$ vanishes.  The extraction statement is the
moment identity \cref{due:eq:ext-moments} of the row of order $d_r$
applied to the finite sum \cref{due:eq:main-derivative-law}, exactly as in
the proof of \cref{due:thm:ext-extraction}.
\end{proof}

% ed.: S2 states the derivative law with the depth letter in the summation
% ed.: bound (its d is this volume's s) and the onset "n >= d"; S4 states it
% ed.: for p_{n+1}; S3 states it for the centred splines.  All three are the
% ed.: same statement; the mode count floor((s-r)/2) is S4's.
\begin{remark}[Parity of the derivative order]
\label{due:rmk:main-derivative-parity}
For odd $r$ the derivatives on the right of
\cref{due:eq:main-derivative-law} are the odd-order derivatives
$\RvachevUp^{(r)},\RvachevUp^{(r+2)},\ldots$, so the derivative law is not a
consequence of the value law at other points; it is the value law for the
odd part of the local jet.  The mode ratios are the same powers of $1/4$
in every case, because they come from the even tail and not from the order
of differentiation.
\end{remark}

\subsection{Rationality}
\label{due:sec:main-rational}

\begin{corollary}[Exact rationality]
\label{due:cor:main-rational}
For every dyadic $x\in[-1,1]$, the value $\RvachevUp(x)$, every derivative
$\RvachevUp^{(r)}(x)$, and every mode coefficient $\Cr{m}{x}$ are rational
numbers, and each is a fixed rational combination of $d+1$ rational
finite-spline values:
\begin{equation}
 \RvachevUp(x)=\sum_{j=0}^{d}\Wrow{d}{j}\,\qn{N+j},
 \qquad
 \RvachevUp^{(2m)}(x)
 =\frac{(-1)^{m}4^{m(N+1)}}{\am{m}}
  \sum_{j=0}^{d}\qn{N+j}\,\frac{(-1)^{d-m}e^{(j)}_{d-m}}{D_{j}}
 \qquad(N\ge s),
 \label{due:eq:main-rational}
\end{equation}
with the weights $\Wrow{d}{j}$ of \cref{due:thm:ext-weights} and the
quantities $D_{j}$, $e^{(j)}_{r}$ of \cref{due:thm:rec-inverse}.  The odd
derivatives are rational by the Thue--Morse tiling
\cref{due:cor:mech-tiling}, and every derivative of order above $s$ is
zero.
\end{corollary}

\begin{proof}
The first identity is \cref{due:thm:ext-extraction}, the second is
\cref{due:cor:rec-jet}; both rest only on \cref{due:thm:set-main} and on the
rationality of the samples (\cref{due:cor:set-scaled-cell}).  The
statements about odd and high orders are
\cref{due:cor:mech-tiling,due:thm:mech-jet-termination}.
\end{proof}

The classical proof of rationality \cite{arias1982,arias2018} runs
through the functional-differential equation and explicit recurrences for
the values at dyadic points; the present one runs through finite
convolution data alone, and it returns the value rather than a rationality
certificate.  The size of the denominators is the subject of
\cref{due:prop:stab-denominator}: the extraction adds to the denominators
of the samples only the odd factor $P_{d}=\prod_{k=1}^{d}(4^{k}-1)$.

\subsection{The tail generating function and the annihilating recurrence}
\label{due:sec:main-generating}

The finite geometric law is equivalent to the rationality of a generating
function, and its coefficient form is the annihilating recurrence of
\cref{due:sec:recovery}.

\begin{proposition}[Rational tail generating function]
\label{due:prop:main-generating}
Let $N\ge s$.  Then, for $\AbsoluteValue{z}<1$,
\begin{equation}
 \sum_{k=0}^{\infty}\qn{N+k}\,z^{k}
 =\frac{\RvachevUp(x)}{1-z}
 +\sum_{m=1}^{d}\frac{\Cr{m}{x}\,\Qq^{mN}}{1-\Qq^{m}z},
 \label{due:eq:main-generating}
\end{equation}
a rational function of $z$ whose denominator divides
$(1-z)\prod_{m=1}^{d}(1-\Qq^{m}z)$ and whose numerator has degree at most
$d$; for $s\ge2$ the denominator is exactly this product, since every
partial fraction has a nonzero numerator.  Conversely, a sequence whose
generating function has this form for some $N$ satisfies the law
\cref{due:eq:set-main-law} for all $n\ge N$.
\end{proposition}

\begin{proof}
Insert \cref{due:eq:set-main-law} at $n=N+k$ and sum the $d+1$ geometric
series; the ratios $1$ and $\Qq^{m}$ are all less than $1/\AbsoluteValue{z}$
when $\AbsoluteValue{z}<1$.  Distinct poles $1,\Qq^{-1},\ldots,\Qq^{-d}$
with nonzero residues (\cref{due:thm:sharp-modes}) give the exact
denominator.  The converse is the uniqueness of partial-fraction
expansions.
\end{proof}

\begin{corollary}[Eventually annihilating recurrence]
\label{due:cor:main-recurrence}
The sequence $(\qn{n})_{n\ge s}$ is annihilated by the product of shifts
$\prod_{m=0}^{d}(\mathsf E-\Qq^{m})$, that is, by the recurrence
\cref{due:eq:rec-recurrence} of \cref{due:thm:rec-recurrence}, with the
integer form \cref{due:eq:rec-recurrence-int}; for $s\ge2$ no recurrence of
lower order with constant coefficients annihilates it.  The recurrence is
\emph{eventual}: \cref{due:sec:sharp-onset} shows that at even depth its
residual at $n=s-1$ is a nonzero rational whose value is known exactly.
\end{corollary}

\begin{proof}
Multiplying \cref{due:eq:main-generating} by the denominator gives a
polynomial of degree at most $d$, which is the statement that the
coefficients of $z^{k}$, $k\ge d+1$, of the product vanish; those
coefficients are the residuals $\rho_{N+k-d-1}$ of \cref{due:eq:rec-residual}.
Minimality is \cref{due:prop:rec-diagnostics}\ref{due:it:diag-order}
together with $\Cr{d}{x}\ne0$.
\end{proof}

\begin{summarybox}[title={Section summary}]
\Cref{due:thm:set-main} is proved: for $n\ge s$ the sample $\qn{n}$ is
$\RvachevUp(x)$ plus exactly $d=\Floor{s/2}$ geometric modes
$\Cr{m}{x}4^{-mn}$, with no remainder, because the window seen by the
omitted tail is one polynomial cell (\cref{due:sec:depth}), the tail is
deconvolved by a finite operator (\cref{due:sec:coefficients}), and the jet
stops at order $s$ (\cref{due:sec:jet}).  The same argument gives the law
for every derivative order $r\le s$ with $\Floor{(s-r)/2}$ modes
(\cref{due:cor:main-derivative}), rationality of the whole jet from finite
convolution data (\cref{due:cor:main-rational}), a rational tail generating
function with the exact denominator $(1-z)\prod_{m\le d}(1-\Qq^{m}z)$
(\cref{due:prop:main-generating}), and the eventually annihilating
recurrence (\cref{due:cor:main-recurrence}).
\end{summarybox}

% ==========================================================================
\section{Sharpness: the number of modes and the onset}
\label{due:sec:sharpness}

\Cref{due:thm:set-main} makes three quantitative claims: the tail has
$d=\Floor{s/2}$ modes, the law holds from level $s$ on, and (through
\cref{due:thm:ext-extraction}) $d+1$ samples suffice.  This section shows
that the first is exact (every mode is present), settles the second
completely at even depth by computing the defect at the last level before
the onset in closed form, records what is proved and what is verified at
odd depth, shows that the third is minimal within the mode model, and
makes precise the relation to the asymptotic statements of the Exponents
volume \cite{proveit-exponents}.

\subsection{The number of modes is exactly \texorpdfstring{$d$}{d}}
\label{due:sec:sharp-modes}

\begin{theorem}[All $d$ modes are present]
\label{due:thm:sharp-modes}
Let $x\in(-1,1)$ be dyadic with $s=\dep{x}\ge2$ and $d=\Floor{s/2}$.  Then
$\Cr{m}{x}\ne0$ for every $1\le m\le d$; the polynomial
$P_{N}(z)=\sum_{m=0}^{d}B_{m}z^{m}$ of \cref{due:thm:rec-inverse}, whose
values at $z=\Qq^{k}$ are the samples $\qn{N+k}$, has exact degree $d$; the
least order of a constant-coefficient recurrence annihilating
$(\qn{n})_{n\ge s}$ is $d$; and the top coefficient has the magnitude
\begin{equation}
 \AbsoluteValue{\Cr{d}{x}}
 =\begin{cases}
  \am{d}\,2^{\binom{s+1}{2}-2d}, & s=2d,\\[0.3em]
  \am{d}\,2^{\binom{s}{2}-2d-1}, & s=2d+1.
 \end{cases}
 \label{due:eq:sharp-top-mode}
\end{equation}
For $s\in\SetOf{0,1}$ (the points $0,\pm1,\pm1/2$) there is no mode:
$\qn{n}=\RvachevUp(x)$ for every $n\ge s$.
\end{theorem}

\begin{proof}
$\Cr{m}{x}=(-1)^{m}\am{m}4^{-m}\RvachevUp^{(2m)}(x)$ with $\am{m}>0$
(\cref{due:thm:mech-am-rational-positive}) and $\RvachevUp^{(2m)}(x)\ne0$
for $2m\le s$ (\cref{due:cor:mech-even-jet}); the magnitude is
\cref{due:eq:mech-top-mode}.  The degree statement is $B_{d}=\Cr{d}{x}\Qq^{dN}\ne0$
(\cref{due:eq:rec-modes}), and the recurrence order is
\cref{due:prop:rec-diagnostics}\ref{due:it:diag-order}.  For $s\le1$ the
sum in \cref{due:eq:set-main-law} is empty.
\end{proof}

% ed.: S6 is the source that proves all modes nonzero; S1's "sharp maximal
% ed.: mode" proves only the top one, and S3's "sharp mode" corollary is the
% ed.: same top-mode statement.  The volume's proof (cor:mech-even-jet) is
% ed.: S6's argument, made exact by the tiling identity of cor:mech-tiling.
\begin{remark}[Which source proved what]
\label{due:rmk:sharp-modes-sources}
Sources S1 and S3 prove only that the top mode $\Cr{d}{x}$ is nonzero;
S6 proves that all $d$ modes are nonzero, from the vanishing pattern of the
even jet, and \cref{due:cor:mech-even-jet} is its argument.  The exhaustive
exact verification of the volume (\cref{due:sec:algorithm}) found all
$684$ mode coefficients at the $252$ interior points of depth $2$ to $7$
nonzero, as the theorem requires.
\end{remark}

\subsection{The onset: the exact defect at the knot level}
\label{due:sec:sharp-onset}

At the level $n=s-1$ the point $x$ is a knot of $\upn{n}$, not the centre
of a cell (\cref{due:lem:set-midpoint}(ii)), and the proof of the law does
not apply.  Whether the law nevertheless holds there was, in the sources, a
matter of computation: S3 proved that it does at odd depth
(\cref{due:prop:mech-odd-onset}), S4's verification found that it fails
at every point of even depth through depth $7$, and no source proved the
failure.  The following theorem settles the question at every depth by
computing the defect in closed form: it is a Thue--Morse sign times a
Bernoulli number.

Write, for $s\ge1$ and $n=s-1$,
\begin{equation}
 \Delta_{s}(x)\DefinitionEquals
 \upn{s-1}(x)-\sum_{m=0}^{d}(-1)^{m}\am{m}\,4^{-ms}\,\RvachevUp^{(2m)}(x)
 \label{due:eq:sharp-defect-def}
\end{equation}
for the difference between the sample at the knot level and the value the
law \cref{due:eq:set-main-law-shifted} would predict there (note
$4^{-m(n+1)}=4^{-ms}$).  For $s\ge2$ the spline $\upn{s-1}$ is continuous,
so $\Delta_{s}(x)$ is unambiguous; for $s=1$ it depends on the value
assigned to the box $\upn{0}$ at its jumps, as in
\cref{due:prop:mech-odd-onset}.

We use the moment generating function $M(z)=\sum_{n}m^{\mathrm{Rv}}_{n}z^{n}/n!$
of \cref{due:prop:mech-normalizations}, $1/M(z)=\sum_{m}(-1)^{m}\am{m}z^{2m}$
(\cref{due:eq:mech-gamma-egf}), and its half-range companion
\begin{equation}
 M_{+}(z)\DefinitionEquals\int_{0}^{1}\EulerE^{zu}\,\RvachevUp(u)\,\Differential u
 =\sum_{j=0}^{\infty}\frac{M_{j}}{2}\,\frac{z^{j}}{j!},
 \qquad
 M_{j}\DefinitionEquals\int_{-1}^{1}\AbsoluteValue{u}^{j}\,\RvachevUp(u)\,\Differential u,
 \label{due:eq:sharp-half-mgf-def}
\end{equation}
so that $M(z)=M_{+}(z)+M_{+}(-z)$, $M_{0}=1$, and $M_{j}=m^{\mathrm{Rv}}_{j}$
for even $j$.

% ed.: this lemma is new to the volume; it is the one-line consequence of
% ed.: Rvachev's equation up'(x) = 2up(2x+1) - 2up(2x-1) that the sources'
% ed.: numerics were tracking without naming.
\begin{lemma}[Half-range moment generating function]
\label{due:lem:sharp-half-mgf}
$M_{+}$ is entire and
\begin{equation}
 \boxed{\;z\,M_{+}(z)=\EulerE^{z/2}\,M(z/2)-1,\qquad
 M_{+}(z)+\frac1z=\frac{\EulerE^{z/2}M(z/2)}{z}.\;}
 \label{due:eq:sharp-half-mgf}
\end{equation}
Consequently
\begin{equation}
 \frac{M_{+}(z)+z^{-1}}{M(z)}
 =\frac{\EulerE^{z/2}}{2\sinh(z/2)}
 =\frac12\Bigl(1+\coth\frac z2\Bigr)
 =\frac12+\frac1z+\sum_{j=1}^{\infty}\frac{B_{2j}}{(2j)!}\,z^{2j-1},
 \label{due:eq:sharp-coth}
\end{equation}
where $B_{2j}$ are the Bernoulli numbers.
\end{lemma}

\begin{proof}
Integrating by parts, with $\RvachevUp(0)=1$ and $\RvachevUp(1)=0$
(\cref{due:prop:set-basic}),
\[
 z\,M_{+}(z)=\bigl[\EulerE^{zu}\RvachevUp(u)\bigr]_{0}^{1}
 -\int_{0}^{1}\EulerE^{zu}\,\RvachevUp'(u)\,\Differential u
 =-1-\int_{0}^{1}\EulerE^{zu}\,\RvachevUp'(u)\,\Differential u .
\]
Rvachev's equation $\RvachevUp'(u)=2\RvachevUp(2u+1)-2\RvachevUp(2u-1)$
(\cref{due:prop:set-basic}) reduces on $0<u<1$ to
$\RvachevUp'(u)=-2\RvachevUp(2u-1)$, because $2u+1>1$ lies outside the
support.  Substituting $v=2u-1$,
\[
 -\int_{0}^{1}\EulerE^{zu}\,\RvachevUp'(u)\,\Differential u
 =2\int_{0}^{1}\EulerE^{zu}\,\RvachevUp(2u-1)\,\Differential u
 =\int_{-1}^{1}\EulerE^{z(v+1)/2}\,\RvachevUp(v)\,\Differential v
 =\EulerE^{z/2}\,M(z/2),
\]
which is the first identity; the second is the first divided by $z$.
For \cref{due:eq:sharp-coth}, the product formula
\cref{due:eq:mech-mgf-product} gives
$M(z)=\dfrac{\sinh(z/2)}{z/2}\,M(z/2)$, so
\[
 \frac{M_{+}(z)+z^{-1}}{M(z)}
 =\frac{\EulerE^{z/2}M(z/2)}{z}\cdot\frac{z/2}{\sinh(z/2)\,M(z/2)}
 =\frac{\EulerE^{z/2}}{2\sinh(z/2)}
 =\frac{\cosh(z/2)+\sinh(z/2)}{2\sinh(z/2)},
\]
and $\frac z2\coth\frac z2=\sum_{j\ge0}B_{2j}z^{2j}/(2j)!$ is the
generating function of the even Bernoulli numbers.
\end{proof}

\begin{theorem}[Exact defect at the knot level]
\label{due:thm:sharp-knot-defect}
Let $x\in(-1,1)$ be dyadic with depth $s\ge2$, $d=\Floor{s/2}$, and let
$k=\bigl(2^{s}(x+1)-1\bigr)/2\in\SetOf{0,\ldots,2^{s}-1}$ be the index of
$x$ in the knot lattice $t_{s-1,k}=-1+(2k+1)2^{-s}$ of $\upn{s-1}$
(\cref{due:lem:set-knots}).  Then
\begin{equation}
 \boxed{\;
 \Delta_{s}(x)=-\ThueMorseSign{k}\;2^{-\binom{s}{2}}\;\beta_{s},
 \qquad
 \beta_{s}=\begin{cases}
  \dfrac{B_{s}}{s!}, & s\ \text{even},\\[0.7em]
  0, & s\ \text{odd}.
 \end{cases}\;}
 \label{due:eq:sharp-knot-defect}
\end{equation}
In particular, at every point of even depth $s$ the law fails at
$n=s-1$ by the universal amount
$\AbsoluteValue{\Delta_{s}(x)}=2^{-\binom s2}\AbsoluteValue{B_{s}}/s!$, which
is $\tfrac1{24}$, $\tfrac1{46080}$, $\tfrac1{990904320}$,
$\tfrac1{324699527577600}$ for $s=2,4,6,8$; and at every point of odd
depth $s\ge3$ the law already holds at $n=s-1$.
\end{theorem}

\begin{proof}
Put $n=s-1\ge1$ and $\delta=\delta_{n}=2^{-s}$, so that $4^{-ms}=\delta^{2m}$
and $\Delta_{s}(x)=\upn{n}(x)-\sum_{m=0}^{d}(-1)^{m}\am{m}\delta^{2m}
\RvachevUp^{(2m)}(x)$.

\emph{Two pieces on the window.}
By \cref{due:lem:set-midpoint}(ii), $x=t_{n,k}$ is the only knot of
$\upn{n}$ in the window $[x-\delta,x+\delta]$.  By
\cref{due:thm:set-finite-spline} the $n$-th derivative of $\upn{n}$ is a
step function whose jump at $t_{n,k}$, right value minus left value, is
$J\DefinitionEquals2^{\binom{n+1}{2}}\ThueMorseSign{k}=2^{\binom s2}\ThueMorseSign{k}$,
and $\upn{n}$ has degree at most $n$ on each of the two cells.  Hence on
the window
\begin{equation}
 \upn{n}(t)=P(t)+\frac{J}{n!}\,(t-x)_{+}^{\,n},
 \label{due:eq:sharp-two-pieces}
\end{equation}
where $P$ is the polynomial piece on the left cell, and, since
$n\ge1$ makes $\upn{n}$ continuous, $\upn{n}(x)=P(x)$.

\emph{Transfer of the extra piece through the tail.}
Write $\mathcal A_{\delta}P=P\Convolution T_{\delta}$ for the tail average
of \cref{due:sec:coefficients}.  By \cref{due:prop:set-head-tail},
$\RvachevUp=\upn{n}\Convolution T_{\delta}$, so for every $r\ge0$
\begin{equation}
 \RvachevUp^{(r)}(x)=(\mathcal A_{\delta}P)^{(r)}(x)+J\,X_{r},
 \qquad
 X_{r}\DefinitionEquals\frac{1}{n!}\,\frac{\Differential^{r}}{\Differential t^{r}}
 \Bigl[(t-x)_{+}^{\,n}\Convolution T_{\delta}\Bigr](x),
 \label{due:eq:sharp-transfer}
\end{equation}
and we compute $X_{r}$ in the three ranges of $r$.
\begin{enumerate}[label=\textup{(\alph*)}]
\item $0\le r\le n-1$.  The truncated power is $C^{\,n-1}$, so the
 derivative falls on it:
 \[
  X_{r}=\frac{1}{(n-r)!}\int_{-\delta}^{0}(-y)^{\,n-r}\,T_{\delta}(y)\,\Differential y
  =\frac{\delta^{\,n-r}}{(n-r)!}\int_{-1}^{0}\AbsoluteValue{u}^{\,n-r}\RvachevUp(u)\,\Differential u
  =\frac{\delta^{\,n-r}\,M_{n-r}}{2\,(n-r)!},
 \]
 by the substitution $y=\delta u$ and the evenness of $\RvachevUp$.
\item $r=n$.  The $(n-1)$-st derivative of the truncated power is
 absolutely continuous, so the $n$-th derivative of its convolution with the
 smooth $T_{\delta}$ is the convolution of its a.e.\ derivative
 $\IndicatorOf{t>x}$ with $T_{\delta}$, and
 $X_{n}=\int_{-\delta}^{0}T_{\delta}(y)\,\Differential y=\tfrac12
 =\delta^{0}M_{0}/(2\cdot0!)$, the same formula as in (a).
\item $r\ge n+1$.  Now the derivatives fall on $T_{\delta}$:
 $X_{r}=\frac1{n!}\int_{-\delta}^{0}(-y)^{n}\,T_{\delta}^{(r)}(y)\,\Differential y$.
 Integrate by parts $n+1$ times.  Every boundary term at $y=-\delta$
 vanishes because all derivatives of $T_{\delta}$ vanish at $-\delta$
 (\cref{due:thm:mech-jet-termination} at depth $0$), and at $y=0$ the
 derivatives of $(-y)^{n}$ of order below $n$ vanish while the
 $n$-th is the constant $(-1)^{n}n!$; the $(n+1)$-st derivative of
 $(-y)^{n}$ is zero, so no integral remains, and
 \[
  X_{r}=\frac{1}{n!}\cdot(-1)^{n}\cdot(-1)^{n}n!\;T_{\delta}^{(r-n-1)}(0)
  =T_{\delta}^{(r-n-1)}(0)
  =\delta^{\,n-r}\,\RvachevUp^{(r-n-1)}(0)
  =\delta^{-1}\,[\,r=n+1\,],
 \]
 because $\RvachevUp(0)=1$ and every derivative of positive order vanishes
 at the origin (\cref{due:thm:mech-jet-termination} at depth $0$).
\end{enumerate}
In all three ranges,
\begin{equation}
 X_{r}=\delta^{\,n-r}\,c_{n-r},
 \qquad
 c_{j}\DefinitionEquals\frac{M_{j}}{2\,j!}\ (j\ge0),\qquad c_{-1}\DefinitionEquals1,\qquad
 c_{j}\DefinitionEquals0\ (j\le-2).
 \label{due:eq:sharp-Xr}
\end{equation}

\emph{The defect.}
Insert \cref{due:eq:sharp-transfer} at $r=2m$ into the definition of
$\Delta_{s}(x)$:
\[
 \Delta_{s}(x)=\Bigl[P(x)-\sum_{m=0}^{d}(-1)^{m}\am{m}\delta^{2m}
 (\mathcal A_{\delta}P)^{(2m)}(x)\Bigr]
 -J\sum_{m=0}^{d}(-1)^{m}\am{m}\,\delta^{2m}X_{2m}.
\]
The bracket vanishes: by \cref{due:lem:mech-finite-inverse} the finite
deconvolution $\sum_{m}(-1)^{m}\am{m}\delta^{2m}(\mathcal A_{\delta}P)^{(2m)}$
returns $P$ for every polynomial $P$, the sum extending over
$2m\le\deg P\le n$; the terms with $2m\le n$ are all among $m\le d$, and
the possible extra term $m=d$ with $2d=n+1$ (even $s$) is zero because
$\deg P\le n<2d$.  By \cref{due:eq:sharp-Xr},
\[
 \Delta_{s}(x)=-J\,\delta^{\,n}\sum_{m=0}^{d}(-1)^{m}\am{m}\,c_{n-2m},
\]
and the sum runs over exactly those $m$ with $n-2m\ge-1$, that is
$m\le\Floor{(n+1)/2}=\Floor{s/2}=d$, so it is the full coefficient of
$z^{n}$ in the product of the two series $\sum_{m}(-1)^{m}\am{m}z^{2m}=1/M(z)$
and $\sum_{j\ge-1}c_{j}z^{j}=M_{+}(z)+z^{-1}$
(\cref{due:eq:sharp-half-mgf-def}):
\[
 \Delta_{s}(x)=-J\,\delta^{\,n}\,[z^{n}]\;\frac{M_{+}(z)+z^{-1}}{M(z)}
 =-J\,\delta^{\,n}\,[z^{n}]\;\frac12\Bigl(1+\coth\frac z2\Bigr)
\]
by \cref{due:lem:sharp-half-mgf}.  For $n\ge1$ the coefficient of $z^{n}$
in \cref{due:eq:sharp-coth} is $B_{n+1}/(n+1)!$ when $n$ is odd and $0$ when
$n$ is even; with $n=s-1$ this is $\beta_{s}$.  Finally
$J\delta^{n}=2^{\binom s2}\ThueMorseSign{k}\,2^{-s(s-1)}
=\ThueMorseSign{k}\,2^{-\binom s2}$.

The numerical values: $B_{2}/2!=\tfrac1{12}$, $B_{4}/4!=-\tfrac1{720}$,
$B_{6}/6!=\tfrac1{30240}$, $B_{8}/8!=-\tfrac1{1209600}$, multiplied by
$2^{-1},2^{-6},2^{-15},2^{-28}$.
\end{proof}

% ed.: the s = 1 case is the convention question of prop:mech-odd-onset
% ed.: seen once more: the same computation at n = 0 gives -J/2.
\begin{remark}[Depth one, and a check]
\label{due:rmk:sharp-depth-one}
The computation applies verbatim at $s=1$, $n=0$, $x=\pm1/2$, if
$\upn{0}(x)$ is read as the left-cell value $P(x)$: the coefficient of $z^{0}$
in \cref{due:eq:sharp-coth} is $\tfrac12$, so $\Delta_{1}=-J/2$ with
$J=\ThueMorseSign{k}$ the jump of the box, that is $\Delta_{1}(1/2)=\tfrac12$
and $\Delta_{1}(-1/2)=-\tfrac12$; with the Fourier-inversion value
$\upn{0}(\pm1/2)=\tfrac12$ the defect is $0$, and with the right-cell value it
is $+J/2$.  This is \cref{due:prop:mech-odd-onset}'s remark on the
symmetric-rectangle convention.  At $x=1/4$ ($s=2$, $k=2$,
$\ThueMorseSign{2}=-1$) the theorem gives $\Delta_{2}=+\tfrac1{24}$, and indeed
$\upn{1}(1/4)=1$ while the law predicts $\tfrac{23}{24}$
(\cref{due:ex:set-quarter}); at $x=3/4$ ($k=3$, $\ThueMorseSign{3}=+1$) it
gives $-\tfrac1{24}$, and $\upn{1}(3/4)=0$ against a predicted
$\tfrac1{24}$.  The formula \cref{due:eq:sharp-knot-defect} was checked in
exact rational arithmetic at all $508$ interior dyadic points of depth $2$
to $8$ (\cref{due:sec:algorithm}).
\end{remark}

\begin{corollary}[The exact onset]
\label{due:cor:sharp-onset}
Let $n_{0}(x)$ be the least level from which the law
\cref{due:eq:set-main-law} holds for all $n\ge n_{0}(x)$.
\begin{enumerate}[label=\textup{(\roman*)}]
\item If $s$ is even, $s\ge2$, then $n_{0}(x)=s$: the onset stated in
 \cref{due:thm:set-main} is exact at every point of even depth.
\item If $s$ is odd, $s\ge3$, then $n_{0}(x)\le s-1$, and $n_{0}(x)=s-1$
 exactly when the law fails at $n=s-2$.  This failure holds at every
 interior point of depth $3$, $5$ and $7$, with
 $\AbsoluteValue{\Delta}=\tfrac1{96}$, $\tfrac{7}{737280}$,
 $\tfrac{31}{63417876480}$ respectively, independent of the point
 (verified in exact arithmetic); a proof for all odd depths is open.
\item For $s\in\SetOf{0,1}$, $n_{0}(x)=s$ with the Fourier-inversion
 convention at $s=1$, and $n_{0}(\pm1/2)=1$ with either one-sided
 convention.
\end{enumerate}
\end{corollary}

\begin{proof}
(i) is \cref{due:thm:sharp-knot-defect} with $B_{s}\ne0$ for even $s$.
(ii) is \cref{due:thm:sharp-knot-defect} with $\beta_{s}=0$
(equivalently \cref{due:prop:mech-odd-onset}); at $n=s-2$ the window
contains one knot strictly inside, at distance $2^{-s}$ from $x$
(\cref{due:lem:set-midpoint}(iii)), and the two-piece computation above
does not apply as it stands, because the extra truncated power is centred
off $x$.  The listed values are from the exhaustive verification.  (iii)
is \cref{due:rmk:sharp-depth-one} and \cref{due:ex:set-low-depth}.
\end{proof}

% ed.: the sources' onset statements (S1 "cell-center level", S2/S4/S5
% ed.: n >= s or N >= s+1, S3 the odd-level sharpening) are all consistent
% ed.: with the corollary; none of them claimed exactness at even depth,
% ed.: which is new here.
\begin{remark}[The two-knot levels]
\label{due:rmk:sharp-two-knots}
The defects at $n=s-2$ recorded in \cref{due:cor:sharp-onset}(ii), and the
analogous even-depth values $\tfrac1{15360}$ ($s=4$) and
$\tfrac{1}{198180864}$ ($s=6$) at $n=s-2$, are again independent of the
point within a depth, which suggests that a closed form like
\cref{due:eq:sharp-knot-defect} exists at every level below the onset, with
the half-range moments of \cref{due:lem:sharp-half-mgf} replaced by partial
moments $\int_{-1}^{\pm1/2}(\pm\tfrac12-u)^{j}\RvachevUp(u)\,\Differential u$ of
the tail over the off-centre knot.  The numerators $1,7,31=2^{s-2}-1$ at
odd depth and the denominators containing $\prod_{k\le d}(4^{k}-1)$ point
to such a formula.  It is not needed for the results of this volume and is
left open.
\end{remark}

\subsection{Minimal number of samples}
\label{due:sec:sharp-samples}

\begin{theorem}[$d+1$ samples are necessary and sufficient within the mode model]
\label{due:thm:sharp-samples}
Let $s\ge2$.
\begin{enumerate}[label=\textup{(\roman*)}]
\item The row of order $d$ returns $\RvachevUp(x)$ exactly from any $d+1$
 consecutive samples at levels $\ge s$ (\cref{due:thm:ext-extraction}), and
 any row of order $D\ge d$ does the same from $D+1$ samples
 (\cref{due:cor:ext-overorder}).
\item The sequence $(\qn{n})_{n\ge s}$ has exact degree $d$ as a polynomial
 in $\Qq^{n}$, and for any $k\le d$ levels $n_{0}<\cdots<n_{k-1}$ (all
 $\ge s$) there are two sequences of the model class
 $\mathcal M_{d}=\SetOf{L+\sum_{r=1}^{d}A_{r}\Qq^{rn}}$ that agree at these
 levels and have different constants $L$.  Hence no rule, linear or not,
 that reads only $k\le d$ samples returns the constant on all of
 $\mathcal M_{d}$, and $d$ cannot be replaced by $d-1$ for the up-function
 sequence at $x$.
\end{enumerate}
\end{theorem}

\begin{proof}
(i) is quoted.  (ii): the exact degree is \cref{due:thm:sharp-modes}; the
two sequences and the conclusion are
\cref{due:prop:rec-minimal}, whose hypothesis $\Cr{d}{x}\ne0$ is
\cref{due:thm:sharp-modes}.
\end{proof}

\begin{remark}[Outside the model class]
\label{due:rmk:sharp-samples-outside}
The theorem is a statement about rules that see only the samples.  A rule
that also knows $x$ may exploit relations among the coefficients, for
instance the fact that the top mode is determined up to sign by $s$
(\cref{due:eq:sharp-top-mode}), and S1's remark that fewer samples ``can
succeed only by using additional point-specific information'' is exactly
this caveat.  Within the volume no such rule is used: the extraction is
universal in $x$ and in the window.
\end{remark}

\subsection{Relation to the Exponents volume: exact versus asymptotic}
\label{due:sec:sharp-exponents}

The Exponents volume \cite{proveit-exponents} treats the same finite
splines, under the bridge $\upn{n}=p_{n+1}$ of
\cref{due:eq:set-bridge-exponents}, at a \emph{general} point.  Three of its
statements are relevant here, and each becomes exact at a dyadic point.

\begin{proposition}[The sharp $C^{r}$ law is attained at dyadic points]
\label{due:prop:sharp-attained}
Let $r\ge0$ and let $x\in(-1,1)$ be dyadic of depth $s=r+2$.  Then for every
$n\ge s$,
\begin{equation}
 \upn{n}^{(r)}(x)-\RvachevUp^{(r)}(x)
 =-\frac{\am{1}}{4^{\,n+1}}\,\RvachevUp^{(r+2)}(x)
 =-\ThueMorseSign{k}\;\frac{2^{\binom{r+3}{2}-1}}{9\cdot4^{\,n+1}},
 \label{due:eq:sharp-attained}
\end{equation}
with $k$ the Thue--Morse index of \cref{due:eq:mech-top-derivative}.  The
right-hand side has absolute value $2^{\binom{r+3}{2}-1}/(9\cdot4^{N})$ with
$N=n+1$, which is the value of $\Norm{p_{N}^{(r)}-\RvachevUp^{(r)}}_{\infty}$
in the sharp ordinary-prefix law of the Exponents volume (its
\texttt{p3:thm:sharp-prefix-law}, valid for $N\ge r+3$): the supremum in
that law is attained at every dyadic point of depth $r+2$, for every level
from the onset on, and the one-mode law \cref{due:eq:sharp-attained} is the
exact form of the extremal error there.
\end{proposition}

\begin{proof}
\Cref{due:cor:main-derivative} with $s=r+2$ has $d_{r}=1$ and a single
mode,
\[
 \upn{n}^{(r)}(x)-\RvachevUp^{(r)}(x)=-\am{1}4^{-(n+1)}\RvachevUp^{(r+2)}(x);
\]
here $\am{1}=\tfrac1{18}$ and, by \cref{due:thm:mech-jet-termination}(ii),
$\RvachevUp^{(s)}(x)=2^{\binom{s+1}{2}}\ThueMorseSign{k}$ with
$s+1=r+3$.  Then $2^{\binom{r+3}{2}}/18=2^{\binom{r+3}{2}-1}/9$.  The
hypothesis $N\ge r+3$ of the cited law is $n\ge r+2=s$ under the bridge,
which is exactly the onset.
\end{proof}

\begin{proposition}[The exact CDF law is the depth-two density law]
\label{due:prop:sharp-cdf}
For every $n\ge0$ and every real $x$,
\begin{equation}
 \upn{n+1}(x)=\int_{2x-1}^{2x+1}\upn{n}(t)\,\Differential t,
 \qquad
 \RvachevUp(x)=\int_{2x-1}^{2x+1}\RvachevUp(t)\,\Differential t .
 \label{due:eq:sharp-finite-fde}
\end{equation}
Consequently, with $P_{n}(x)=\int_{-\infty}^{x}\upn{n}$ and
$P(x)=\int_{-\infty}^{x}\RvachevUp$,
\begin{equation}
 P_{n}(\tfrac12)=\upn{n+1}(\tfrac14),\qquad
 P(\tfrac12)=\RvachevUp(\tfrac14)=\tfrac{67}{72},\qquad
 P_{n}(\tfrac12)-P(\tfrac12)=\frac{1}{9\cdot4^{\,n+1}}
 \quad(n\ge1).
 \label{due:eq:sharp-cdf-half}
\end{equation}
The last identity is the exact CDF law $1/(9\cdot4^{N})$ of the Exponents
volume (its \texttt{p3:cor:kolmogorov}, at its extremizer $x=\tfrac12$)
under $N=n+1$: it is the case $s=2$, $d=1$ of \cref{due:thm:set-main} at
$x=\tfrac14$, transported to the distribution function by
\cref{due:eq:sharp-finite-fde}.
\end{proposition}

\begin{proof}
On the Fourier side, $\Phin{n+1}(\xi)=\prod_{k=0}^{n+1}\SincRad(\pi\xi/2^{k})
=\SincRad(\pi\xi)\prod_{j=0}^{n}\SincRad\bigl(\pi(\xi/2)/2^{j}\bigr)
=\SincRad(\pi\xi)\,\Phin{n}(\xi/2)$.  The inverse transform of
$\SincRad(\pi\xi)$ is the box $\IndicatorOf{[-1/2,1/2]}$
(\cref{due:lem:set-box-transform}) and that of $\Phin{n}(\xi/2)$ is
$2\upn{n}(2\,\cdot)$, so $\upn{n+1}(x)=\int_{x-1/2}^{x+1/2}2\upn{n}(2t)\,\Differential t
=\int_{2x-1}^{2x+1}\upn{n}(u)\,\Differential u$.  Letting $n\to\infty$ (or
integrating Rvachev's equation) gives the second identity.  For
\cref{due:eq:sharp-cdf-half}, evenness and total mass $1$ give
$P_{n}(\tfrac12)=1-P_{n}(-\tfrac12)=\int_{-1/2}^{1}\upn{n}
=\int_{-1/2}^{3/2}\upn{n}=\upn{n+1}(\tfrac14)$, using the support
$[-1,1]$; likewise for $P$.  Then \cref{due:thm:set-main} at $x=\tfrac14$,
where $s=2$, $d=1$, $\Cr{1}{\tfrac14}=\tfrac19$
(\cref{due:ex:set-quarter}), gives
$\upn{n+1}(\tfrac14)-\RvachevUp(\tfrac14)=\tfrac19\cdot4^{-(n+1)}$ for
$n+1\ge2$.
\end{proof}

% ed.: the cluster brief asked whether the CDF law 1/(9 4^n) and the density
% ed.: law at x = 1/2 are the same object; they are not (at x = 1/2 the
% ed.: density law has no mode at all), but the CDF law at 1/2 IS the density
% ed.: law at 1/4 through the finite functional equation, which is the
% ed.: precise statement.
\begin{remark}[Two laws at the point $\tfrac12$]
\label{due:rmk:sharp-half}
At $x=\tfrac12$ itself the density law has no mode: $\dep{\tfrac12}=1$, and
$\upn{n}(\tfrac12)=\RvachevUp(\tfrac12)=\tfrac12$ for every $n\ge1$
(\cref{due:ex:set-low-depth}).  The CDF discrepancy $1/(9\cdot4^{n+1})$ at
the same point is not a contradiction: the distribution function at
$\tfrac12$ is the density at $\tfrac14$, one level up, and $\tfrac14$ has
depth $2$.  In the Lean corpus the CDF form of this one-mode law at the
quarter point is kernel-verified
(\texttt{Fabius.\allowbreak fabiusUniformSpline\_\allowbreak quarter\_\allowbreak eq\_\allowbreak quadratic}; see
\cref{due:sec:set-bridge} and the register in \cref{due:sec:register}).
\end{remark}

\begin{corollary}[Richardson acceleration is exact at dyadic points]
\label{due:cor:sharp-richardson}
In the notation of the Exponents volume, let
$R_{n,M}=\sum_{j=0}^{M-1}w_{M,j}\,p_{n+j}$ be its signed geometric
Richardson combination with the quarter-base Lagrange weights $w_{M,j}$
(its \texttt{p3:eq:richardson-\allowbreak combination} and
\texttt{p3:eq:richardson-\allowbreak weights}).  Let $x$ be dyadic of depth $s$ and
$d=\Floor{s/2}$.  Then for every $n\ge s+1$ and every $M\ge d+1$,
\begin{equation}
 R_{n,M}(x)=\RvachevUp(x)
 \label{due:eq:sharp-richardson-exact}
\end{equation}
exactly: the leading term $-\am{M}q^{\binom M2}\delta_{n}^{2M}\RvachevUp^{(2M)}(x)$
and the $\BigOAt{\delta_{n}^{2M+2}}{n\to\infty}$ remainder of its asymptotic
expansion (its \texttt{p3:eq:richardson-leading}) both vanish identically,
because
$\RvachevUp^{(2M)}(x)=0$ for $2M>s$.  The same holds for the derivative
combinations $R^{(r)}_{n,M}$ with $M\ge\Floor{(s-r)/2}+1$.
\end{corollary}

\begin{proof}
Under the bridge $p_{n+j}=\upn{n+j-1}$, and $w_{M,j}=\Wrow{M-1}{j}$ is the
weight of \cref{due:thm:ext-weights} of order $M-1\ge d$; the levels
$n-1,\ldots,n+M-2$ are all $\ge s$.  \Cref{due:cor:ext-overorder} gives the
value law, and \cref{due:cor:main-derivative} with the same row gives the
derivative law.
\end{proof}

\begin{remark}[Away from dyadic points]
\label{due:rmk:sharp-nondyadic}
At a non-dyadic $x$ two things fail at once: the window
$[x-\delta_n,x+\delta_n]$ contains a knot of $\upn{n}$ for every $n$, so
no single polynomial describes the convolution neighbourhood, and the jet
of $\RvachevUp$ at $x$ does not terminate.  The all-orders expansion of the
Exponents volume and the asymptotic Richardson acceleration remain valid,
and its sharp prefix law is the general statement; the finite exact law is
an arithmetic property of the dyadic points.  Between the two lie the
rationals with odd denominator and the other non-dyadic points, at which
the finite splines are still polynomial pieces but the tail sees a knot;
nothing in this volume applies there.
\end{remark}

\begin{summarybox}[title={Section summary}]
Every one of the $d$ modes is present for $s\ge2$
(\cref{due:thm:sharp-modes}).  At the last level before the onset the
defect is known in closed form,
$\Delta_{s}(x)=-\ThueMorseSign{k}2^{-\binom s2}\beta_{s}$ with
$\beta_{s}=B_{s}/s!$ for even $s$ and $0$ for odd $s$
(\cref{due:thm:sharp-knot-defect}), which proves that the onset $n\ge s$
is exact at even depth and reproves the one-level improvement at odd depth
(\cref{due:cor:sharp-onset}); the failure at $n=s-2$ for odd $s$ is
verified through depth $7$ and left open.  Within the mode model $d+1$
samples are necessary and sufficient (\cref{due:thm:sharp-samples}).  The
sharp $C^{r}$ law of the Exponents volume is attained at every dyadic point
of depth $r+2$ (\cref{due:prop:sharp-attained}), its exact CDF law at
$\tfrac12$ is the $s=2$ density law at $\tfrac14$
(\cref{due:prop:sharp-cdf}), and its signed Richardson combination is exact,
not asymptotic, at dyadic points (\cref{due:cor:sharp-richardson}).
\end{summarybox}

% ---- fragment 04-extraction ----
% =====================================================================
%  Cluster: Extraction -- geometric Lagrange interpolation, the
%  q-binomial row, recovery of every mode, stability.   Slug: ext
%  Sources merged: S1 sec. 8-9 + app. B; S2 sec. 7-9 + app. A;
%  S3 sec. 8-10 + app. A-B; S4 sec. 8-10 + app. A-B; S5 sec. 8-9 + app. A;
%  S6 sec. 5 + app.
% =====================================================================

\section{The quarter-base Gaussian-binomial extraction row}
\label{due:sec:extraction}

The preceding part of this volume proves that the finite spline values at a
dyadic point obey an \emph{exact} finite geometric law once the level has
reached the depth of the point.  From here on that law is the only input:
everything in the present three sections is finite-dimensional linear algebra
over the coefficient field $\Kk=\RationalNumbers$, and every identity is an
equality of rational numbers.  We first fix the standing data.

\paragraph{Standing hypotheses and data.}
Throughout, the Fourier transform is $\FourierTwoPi{f}(\xi)
=\int_{\RealNumbers}f(x)\EulerE^{-2\pi\ImaginaryUnit x\xi}\Differential x$,
$\Phin{n}(\xi)=\prod_{k=0}^{n}\SincRad(\pi\xi/2^{k})$ has $n+1$ factors,
$\upn{n}=\FourierTwoPiOperator^{-1}[\Phin{n}]$, and
$\RvachevUp=\FourierTwoPiOperator^{-1}[\PhiUp]$ with
$\PhiUp=\prod_{k\ge0}\SincRad(\pi\xi/2^{k})$.  Fix a dyadic point
$x\in[-1,1]$ of depth $s=\dep{x}$ (the least $s\ge0$ with
$2^{s}x\in\IntegerNumbers$), put
\begin{equation}
 d\DefinitionEquals\Floor{s/2},\qquad
 \Qq\DefinitionEquals\tfrac14,\qquad
 \qn{n}\DefinitionEquals\upn{n}(x)\quad(n\ge0),
 \label{due:eq:ext-data}
\end{equation}
% ed.: cross-reference to the main theorem attached at assembly.
and recall the exact eventual geometric law, \cref{due:thm:set-main}, proved
in \cref{due:sec:main}: for every $n\ge s$,
\begin{equation}
 \qn{n}
 =\RvachevUp(x)+\sum_{r=1}^{d}\Cr{r}{x}\,\Qq^{rn},
 \qquad
 \Cr{r}{x}=(-1)^{r}\am{r}\,4^{-r}\,\RvachevUp^{(2r)}(x),
 \label{due:eq:ext-model}
\end{equation}
where $\am{r}\in\RationalNumbers_{>0}$ are the reciprocal-product
coefficients, $1/\PhiUp(z)=\sum_{m\ge0}\am{m}(2\pi z)^{2m}$.  For $d=0$ (that
is, $s\in\SetOf{0,1}$, $x\in\SetOf{0,\pm\tfrac12,\pm1}$) the sum is empty and
\cref{due:eq:ext-model} says $\qn{n}=\RvachevUp(x)$ for all $n\ge s$.
Nothing below uses the specific values of the $\Cr{r}{x}$; only the shape of
\cref{due:eq:ext-model} matters.  Consequently every statement is proved for an
arbitrary sequence of the form $\qn{n}=L+\sum_{r=1}^{d}A_{r}\Qq^{rn}$ with
$L,A_{r}$ in a field, and the up-function enters only through
\cref{due:eq:ext-model}.  Where the base $\Qq=1/4$ is inessential we write a
general $q$ and say what is needed of it.

\subsection{Geometric Lagrange interpolation at the nodes
\texorpdfstring{$1,\Qq,\ldots,\Qq^{d}$}{1, Q, ..., Q\^{}d}}

Fix a starting level $N\ge s$ and rewrite \cref{due:eq:ext-model} at the
$d+1$ consecutive levels $N,N+1,\ldots,N+d$:
\begin{equation}
 \qn{N+j}
 =\sum_{r=0}^{d}B_{r}\,\bigl(\Qq^{j}\bigr)^{r},
 \qquad
 B_{0}=\RvachevUp(x),\quad
 B_{r}=\Cr{r}{x}\,\Qq^{rN}\ (1\le r\le d),
 \qquad 0\le j\le d.
 \label{due:eq:ext-samples}
\end{equation}
Thus the samples are the values of the polynomial
$P_{N}(z)\DefinitionEquals\sum_{r=0}^{d}B_{r}z^{r}$, of degree at most $d$,
at the geometric nodes $z=1,\Qq,\Qq^{2},\ldots,\Qq^{d}$, and the number wanted
is $P_{N}(0)=B_{0}=\RvachevUp(x)$.  Extrapolating a degree-$d$ interpolant
from $d+1$ nodes to the origin is Lagrange interpolation; the whole content of
this section is that at geometric nodes the Lagrange weights are
$q$-Pochhammer quotients.

\begin{definition}[Geometric Lagrange weights]
\label{due:def:ext-weights}
Let $K$ be a field and $q\in K$ with $q\neq0$ and
$\QPochhammer{q}{q}{d}=\prod_{h=1}^{d}(1-q^{h})\neq0$; equivalently, the
$d+1$ nodes $q^{0},q^{1},\ldots,q^{d}$ are pairwise distinct.  For
$0\le j\le d$ let
\[
 \ell^{(d)}_{j}(z)
 \DefinitionEquals
 \prod_{\substack{0\le k\le d\\k\neq j}}\frac{z-q^{k}}{q^{j}-q^{k}}
\]
be the Lagrange basis polynomial of the node $q^{j}$, and define the
\emph{geometric Lagrange weight}
\begin{equation}
 \Wrow{d}{j}
 \DefinitionEquals\ell^{(d)}_{j}(0)
 =\prod_{\substack{0\le k\le d\\k\neq j}}\frac{-q^{k}}{q^{j}-q^{k}}.
 \label{due:eq:ext-weight-product}
\end{equation}
When $K=\RationalNumbers$ and $q=\Qq=1/4$ we call
$(\Wrow{d}{0},\ldots,\Wrow{d}{d})$ the \emph{quarter-base row of order $d$}.
For $d=0$ the product is empty and $\Wrow{0}{0}=1$.
\end{definition}

The equivalence claimed in the definition is elementary: $q^{i}=q^{j}$ with
$i<j\le d$ and $q\ne0$ holds exactly when $q^{j-i}=1$ with
$1\le j-i\le d$, that is, exactly when some factor of $\QPochhammer{q}{q}{d}$
vanishes.  For $q=1/4$ the condition holds for every $d$.

\begin{lemma}[Moment identities of a geometric Lagrange row]
\label{due:lem:ext-moments}
Under the hypotheses of \cref{due:def:ext-weights}:
\begin{enumerate}
\item\label{due:it:ext-exact}
 For every polynomial $p\in K[z]$ of degree at most $d$,
 $\sum_{j=0}^{d}\Wrow{d}{j}\,p(q^{j})=p(0)$.  In particular
 \begin{equation}
  \sum_{j=0}^{d}\Wrow{d}{j}\,q^{jm}=\delta_{m0}
  \qquad(0\le m\le d).
  \label{due:eq:ext-moments}
 \end{equation}
\item\label{due:it:ext-unique}
 The row $(\Wrow{d}{0},\ldots,\Wrow{d}{d})$ is the \emph{only} vector in
 $K^{d+1}$ satisfying \cref{due:eq:ext-moments}.
\item\label{due:it:ext-uncancelled}
 For every $r\ge0$, writing $h_{r}$ for the complete homogeneous symmetric
 polynomial of degree $r$,
 \begin{equation}
  \sum_{j=0}^{d}\Wrow{d}{j}\,q^{j(d+1+r)}
  =(-1)^{d}q^{\binom{d+1}{2}}\,h_{r}\bigl(1,q,\ldots,q^{d}\bigr).
  \label{due:eq:ext-uncancelled}
 \end{equation}
 In particular the first uncancelled moment is
 $\sum_{j}\Wrow{d}{j}q^{j(d+1)}=(-1)^{d}q^{\binom{d+1}{2}}$.
\end{enumerate}
\end{lemma}

\begin{proof}
\ref{due:it:ext-exact}: the polynomial $p-\sum_{j}p(q^{j})\ell^{(d)}_{j}$ has
degree at most $d$ and vanishes at the $d+1$ distinct nodes $q^{0},\ldots,q^{d}$
(because $\ell^{(d)}_{j}(q^{k})=\delta_{jk}$), hence is identically zero;
evaluate at $z=0$.  Taking $p=z^{m}$ gives \cref{due:eq:ext-moments}.

\ref{due:it:ext-unique}: the system \cref{due:eq:ext-moments} is
$V^{\mathsf T}w=e_{0}$ with $V=(q^{jm})_{0\le j,m\le d}$ the Vandermonde matrix
of the distinct nodes, and $\det V=\prod_{0\le i<k\le d}(q^{k}-q^{i})\neq0$.

\ref{due:it:ext-uncancelled}: put $\omega(z)=\prod_{k=0}^{d}(z-q^{k})$ and
divide with remainder, $z^{d+1+r}=\omega(z)h(z)+\rho(z)$ with
$\deg\rho\le d$.  Since $\omega(q^{j})=0$, $\rho(q^{j})=q^{j(d+1+r)}$, so
by \ref{due:it:ext-exact}
$\sum_{j}\Wrow{d}{j}q^{j(d+1+r)}=\rho(0)=-\omega(0)h(0)$.  Here
$\omega(0)=\prod_{k=0}^{d}(-q^{k})=(-1)^{d+1}q^{\binom{d+1}{2}}$.  For $h(0)$,
expand at $z=\infty$:
\[
 \frac{z^{d+1+r}}{\omega(z)}
 =z^{r}\prod_{k=0}^{d}\bigl(1-q^{k}z^{-1}\bigr)^{-1}
 =z^{r}\sum_{i\ge0}h_{i}\bigl(1,q,\ldots,q^{d}\bigr)z^{-i},
\]
because $\prod_{k}(1-x_{k}t)^{-1}=\sum_{i}h_{i}(x)t^{i}$ is the generating
function of the complete homogeneous symmetric polynomials.  Since
$\rho/\omega=O(z^{-1})$, the quotient $h$ is the polynomial part of this Laurent
series, $h(z)=\sum_{i=0}^{r}h_{i}z^{r-i}$, whose constant term is $h_{r}$.
Therefore the moment equals $(-1)^{d}q^{\binom{d+1}{2}}h_{r}$.
\end{proof}

\begin{theorem}[Closed form of the geometric Lagrange weights]
\label{due:thm:ext-weights}
Under the hypotheses of \cref{due:def:ext-weights}, for $0\le j\le d$,
\begin{equation}
 \boxed{\;
 \Wrow{d}{j}
 =\frac{(-1)^{d-j}\,q^{\binom{d-j+1}{2}}}
       {\QPochhammer{q}{q}{j}\,\QPochhammer{q}{q}{d-j}}
 =\frac{(-1)^{d-j}\,q^{\binom{d-j+1}{2}}}{\QPochhammer{q}{q}{d}}
  \,\GaussianBinomial{d}{j}{q}.\;}
 \label{due:eq:ext-weights}
\end{equation}
\end{theorem}

% ed.: the six sources prove this identically (S1 App. B, S2 App. A, S3 App. A,
% ed.: S4 App. A, S5 App. A, S5 Thm "Universal q-Lagrange weights"); one proof.
\begin{proof}
Split the product \cref{due:eq:ext-weight-product} into its numerator and the
factors with $k<j$ and $k>j$.  The numerator is
\[
 \prod_{\substack{0\le k\le d\\k\neq j}}(-q^{k})
 =(-1)^{d}q^{\binom{d+1}{2}-j}.
\]
For $k<j$ write $q^{j}-q^{k}=-q^{k}(1-q^{j-k})$; the $j$ factors give
$(-1)^{j}q^{\binom{j}{2}}\prod_{h=1}^{j}(1-q^{h})
=(-1)^{j}q^{\binom{j}{2}}\QPochhammer{q}{q}{j}$.
For $k>j$ write $q^{j}-q^{k}=q^{j}(1-q^{k-j})$; the $d-j$ factors give
$q^{j(d-j)}\QPochhammer{q}{q}{d-j}$.  Hence
\[
 \Wrow{d}{j}
 =\frac{(-1)^{d}q^{\binom{d+1}{2}-j}}
       {(-1)^{j}q^{\binom{j}{2}+j(d-j)}\QPochhammer{q}{q}{j}\QPochhammer{q}{q}{d-j}},
\]
and the exponent is
$\binom{d+1}{2}-j-\binom{j}{2}-j(d-j)
=\tfrac12\bigl(d^{2}+d-j-2dj+j^{2}\bigr)
=\tfrac12(d-j)(d-j+1)=\binom{d-j+1}{2}$.
This is the first form.  The second follows from
$\GaussianBinomial{d}{j}{q}
=\QPochhammer{q}{q}{d}/(\QPochhammer{q}{q}{j}\QPochhammer{q}{q}{d-j})$,
which is the catalogue's definition and is legitimate because
$\QPochhammer{q}{q}{d}\ne0$.
\end{proof}

The row has a compact generating polynomial, from which the moment identities
can be read off a second time.  We first record the finite $q$-binomial
theorem with its short proof, since the same identity is used four times
below (row polynomial, integer form, annihilating recurrence, Lebesgue factor).

\begin{lemma}[Finite $q$-binomial theorem]
\label{due:lem:ext-qbinomial}
In any commutative ring, for $n\ge0$,
\begin{equation}
 \QPochhammer{z}{q}{n}=\prod_{h=0}^{n-1}(1-zq^{h})
 =\sum_{k=0}^{n}(-1)^{k}q^{\binom{k}{2}}\GaussianBinomial{n}{k}{q}\,z^{k},
 \label{due:eq:ext-qbinomial}
\end{equation}
where $\GaussianBinomial{n}{k}{q}$ is the Gaussian polynomial (the catalogue's
polynomial-in-$q$ reading).  In particular $\GaussianBinomial{n}{k}{q}$ has
non-negative integer coefficients, so it takes integer values at every
integer $q$.
\end{lemma}

\begin{proof}
The Gaussian polynomials satisfy the Pascal recurrence
$\GaussianBinomial{n}{k}{q}
=\GaussianBinomial{n-1}{k}{q}+q^{n-k}\GaussianBinomial{n-1}{k-1}{q}$
for $1\le k\le n$: in the quotient form, both terms on the right have the
common denominator $\QPochhammer{q}{q}{k}\QPochhammer{q}{q}{n-k}$ after
multiplying the first by $(1-q^{n-k})$ and the second by $(1-q^{k})$, and
$(1-q^{n-k})+q^{n-k}(1-q^{k})=1-q^{n}$; so the sum is
$\QPochhammer{q}{q}{n-1}(1-q^{n})/(\QPochhammer{q}{q}{k}\QPochhammer{q}{q}{n-k})
=\GaussianBinomial{n}{k}{q}$.  Together with
$\GaussianBinomial{n}{0}{q}=\GaussianBinomial{n}{n}{q}=1$ the recurrence
shows by induction on $n$ that the Gaussian polynomials lie in
$\IntegerNumbers_{\ge0}[q]$.

Now induct on $n$; the case $n=0$ reads $1=1$.  Multiplying the identity for
$n-1$ by $(1-zq^{n-1})$ gives, as coefficient of $z^{k}$,
\[
 (-1)^{k}\Bigl(q^{\binom{k}{2}}\GaussianBinomial{n-1}{k}{q}
 +q^{\binom{k-1}{2}+n-1}\GaussianBinomial{n-1}{k-1}{q}\Bigr)
 =(-1)^{k}q^{\binom{k}{2}}\Bigl(\GaussianBinomial{n-1}{k}{q}
 +q^{n-k}\GaussianBinomial{n-1}{k-1}{q}\Bigr),
\]
because $\binom{k-1}{2}+n-1-\binom{k}{2}=n-k$.  The bracket is
$\GaussianBinomial{n}{k}{q}$ by the Pascal recurrence.
\end{proof}

\begin{proposition}[Row polynomial and mode annihilation]
\label{due:prop:ext-rowpoly}
Under the hypotheses of \cref{due:def:ext-weights}, the generating polynomial
of the row and its reverse are
\begin{equation}
 W_{d}(z)\DefinitionEquals\sum_{j=0}^{d}\Wrow{d}{j}z^{j}
 =\prod_{r=1}^{d}\frac{z-q^{r}}{1-q^{r}}
 =\frac{z^{d}\,\QPochhammer{q/z}{q}{d}}{\QPochhammer{q}{q}{d}},
 \qquad
 \sum_{j=0}^{d}\Wrow{d}{d-j}z^{j}
 =\frac{\QPochhammer{qz}{q}{d}}{\QPochhammer{q}{q}{d}}
 =\prod_{r=1}^{d}\frac{1-q^{r}z}{1-q^{r}}.
 \label{due:eq:ext-rowpoly}
\end{equation}
Consequently $W_{d}(1)=1$ and $W_{d}(q^{r})=0$ for $1\le r\le d$: the roots of
the row polynomial are exactly the ratios $q,q^{2},\ldots,q^{d}$ of the
decaying modes, and the roots of the reversed polynomial are their inverses.
\end{proposition}

\begin{proof}
By \cref{due:lem:ext-qbinomial} with $z$ replaced by $q/z$,
\[
 \prod_{r=1}^{d}(z-q^{r})
 =z^{d}\prod_{h=0}^{d-1}\Bigl(1-\frac{q}{z}q^{h}\Bigr)
 =\sum_{k=0}^{d}(-1)^{k}q^{\binom{k}{2}+k}\GaussianBinomial{d}{k}{q}z^{d-k}
 =\sum_{j=0}^{d}(-1)^{d-j}q^{\binom{d-j+1}{2}}\GaussianBinomial{d}{j}{q}z^{j},
\]
using $\binom{k}{2}+k=\binom{k+1}{2}$, $j=d-k$ and the symmetry
$\GaussianBinomial{d}{d-j}{q}=\GaussianBinomial{d}{j}{q}$.  Dividing by
$\prod_{r=1}^{d}(1-q^{r})=\QPochhammer{q}{q}{d}$ and comparing with
\cref{due:eq:ext-weights} gives the first display.  Reversing, with
$z\mapsto qz$ in \cref{due:eq:ext-qbinomial},
$\sum_{j}\Wrow{d}{d-j}z^{j}
=\QPochhammer{q}{q}{d}^{-1}\sum_{k}(-1)^{k}q^{\binom{k}{2}}
 \GaussianBinomial{d}{k}{q}(qz)^{k}
=\QPochhammer{qz}{q}{d}/\QPochhammer{q}{q}{d}$.
The values $W_{d}(1)=1$ and $W_{d}(q^{r})=0$ are read off the product form.
\end{proof}

\begin{remark}[Two proofs of the moment identities]
\label{due:rmk:ext-tworoutes}
$W_{d}(q^{m})=\sum_{j}\Wrow{d}{j}q^{jm}$, so the root statement of
\cref{due:prop:ext-rowpoly} \emph{is} \cref{due:eq:ext-moments}: once the
closed form \cref{due:eq:ext-weights} is known, the finite $q$-binomial
theorem proves the moment identities without any interpolation.  This is the
route of S6 and of S2's appendix; the interpolation route of
\cref{due:lem:ext-moments} has the advantage of giving the uncancelled
moments \cref{due:eq:ext-uncancelled} as well, which the row polynomial does
not see.  In operator language, with $\mathsf E$ the forward shift
$(\mathsf E\,q)_{n}=\qn{n+1}$, the identity
$\sum_{j}\Wrow{d}{j}\qn{N+j}=\bigl(W_{d}(\mathsf E)q\bigr)_{N}$ and the
factorization $W_{d}(\mathsf E)=\prod_{r=1}^{d}(\mathsf E-q^{r})/(1-q^{r})$
say that the row is the product of the elementary filters
$(\mathsf E-q^{r})/(1-q^{r})$, each of which kills one mode and preserves
constants; this is the form in which the Romberg tableau of
\cref{due:sec:recovery} evaluates the row.
\end{remark}

\subsection{The extraction theorem}

For $\Qq=1/4$ it is convenient to have the finite products as integers.  Put
\begin{equation}
 P_{k}\DefinitionEquals\prod_{r=1}^{k}(4^{r}-1)
 \qquad(P_{0}=1,\ P_{1}=3,\ P_{2}=45,\ P_{3}=2835,\ P_{4}=722925,\
 P_{5}=739552275),
 \label{due:eq:ext-Pk}
\end{equation}
so that $\QPochhammer{\Qq}{\Qq}{k}=\prod_{r=1}^{k}(1-4^{-r})
=4^{-\binom{k+1}{2}}P_{k}$.  Every $P_{k}$ is odd.

\begin{theorem}[Single Gaussian-binomial extraction formula]
\label{due:thm:ext-extraction}
Let $x\in[-1,1]$ be dyadic of depth $s=\dep{x}$, let $d=\Floor{s/2}$,
$\Qq=1/4$, and let $N\ge s$ be any integer.  Then, with
$\qn{n}=\upn{n}(x)$ in the $(n+1)$-factor convention,
\begin{equation}
 \boxed{\;
 \RvachevUp(x)
 =\sum_{j=0}^{d}\frac{(-1)^{d-j}\,\Qq^{\binom{d-j+1}{2}}}
   {\QPochhammer{\Qq}{\Qq}{j}\,\QPochhammer{\Qq}{\Qq}{d-j}}\;\qn{N+j}
 =\frac{1}{\QPochhammer{\Qq}{\Qq}{d}}
  \sum_{j=0}^{d}(-1)^{d-j}\,\Qq^{\binom{d-j+1}{2}}
  \GaussianBinomial{d}{j}{\Qq}\;\upn{N+j}(x).\;}
 \label{due:eq:ext-extraction}
\end{equation}
The same row can be written in any of the following equivalent ways:
\begin{align}
 \RvachevUp(x)
 &=\frac{1}{\QPochhammer{\Qq}{\Qq}{d}}
   \sum_{k=0}^{d}(-1)^{k}\,\Qq^{\binom{k+1}{2}}
   \GaussianBinomial{d}{k}{\Qq}\;\qn{N+d-k}
 \qquad\text{(reversed order, S6)},
 \label{due:eq:ext-reversed}\\
 \RvachevUp(x)
 &=\sum_{j=0}^{d}(-1)^{d-j}\,\frac{4^{\binom{j+1}{2}}}{P_{j}\,P_{d-j}}\;\qn{N+j}
 \qquad\text{(integer-product weights, S3)},
 \label{due:eq:ext-integer-product}\\
 \RvachevUp(x)
 &=\frac{1}{P_{d}}\sum_{j=0}^{d}(-1)^{d-j}\,4^{\binom{j+1}{2}}
   \GaussianBinomial{d}{j}{4}\;\qn{N+j}
 \qquad\text{(common odd denominator, S4)}.
 \label{due:eq:ext-base4}
\end{align}
In \cref{due:eq:ext-base4} every numerator coefficient is an integer.  All
four displays are identities in $\RationalNumbers$: with exact rational inputs
$\qn{N},\ldots,\qn{N+d}$ they return the exact rational number
$\RvachevUp(x)$, with no remainder and no limit.
\end{theorem}

% ed.: S1 Thm 8.1, S2 Thm 7.1, S3 Thm 1.3, S4 Thm 8.1, S5 Thm 8.1, S6 Thm 1.2
% ed.: are all this statement; the forward/reversed/integer/base-4 displays are
% ed.: reconciled here once.
\begin{proof}
By \cref{due:eq:ext-samples} the samples are $\qn{N+j}=P_{N}(\Qq^{j})$ for a
polynomial $P_{N}$ of degree at most $d$ with $P_{N}(0)=\RvachevUp(x)$.
The hypotheses of \cref{due:def:ext-weights} hold for $q=\Qq=1/4$ and every
$d$ (no positive power of $1/4$ equals $1$), so
\cref{due:lem:ext-moments}\ref{due:it:ext-exact} gives
$\RvachevUp(x)=P_{N}(0)=\sum_{j}\Wrow{d}{j}\qn{N+j}$, and
\cref{due:thm:ext-weights} gives both forms in
\cref{due:eq:ext-extraction}.  Explicitly, with
$\sum_{j}\Wrow{d}{j}\Qq^{jr}=\delta_{r0}$:
\[
 \sum_{j=0}^{d}\Wrow{d}{j}\qn{N+j}
 =\sum_{r=0}^{d}B_{r}\sum_{j=0}^{d}\Wrow{d}{j}\Qq^{jr}
 =B_{0}=\RvachevUp(x).
\]

\Cref{due:eq:ext-reversed} is \cref{due:eq:ext-extraction} with the summation
index $k=d-j$ and $\GaussianBinomial{d}{d-k}{\Qq}=\GaussianBinomial{d}{k}{\Qq}$;
S6 writes the samples in the order $\qn{N+d},\qn{N+d-1},\ldots,\qn{N}$, which
is the same formula.

\Cref{due:eq:ext-integer-product}: substitute
$\QPochhammer{\Qq}{\Qq}{k}=4^{-\binom{k+1}{2}}P_{k}$ into the first form of
\cref{due:eq:ext-extraction}; the power of $4$ in the $j$-th weight becomes
$4^{-\binom{d-j+1}{2}+\binom{j+1}{2}+\binom{d-j+1}{2}}=4^{\binom{j+1}{2}}$.

\Cref{due:eq:ext-base4}: the inversion identity
$\GaussianBinomial{d}{j}{q^{-1}}=q^{-j(d-j)}\GaussianBinomial{d}{j}{q}$ holds
because
\[
 \GaussianBinomial{d}{j}{q^{-1}}
 =\prod_{i=1}^{j}\frac{1-q^{-(d-j+i)}}{1-q^{-i}}
 =\prod_{i=1}^{j}\frac{q^{-(d-j+i)}\bigl(q^{d-j+i}-1\bigr)}{q^{-i}\bigl(q^{i}-1\bigr)}
 =q^{-j(d-j)}\GaussianBinomial{d}{j}{q}.
\]
With $q=4$ this gives
$\GaussianBinomial{d}{j}{\Qq}=4^{-j(d-j)}\GaussianBinomial{d}{j}{4}$, and
$\QPochhammer{\Qq}{\Qq}{d}=4^{-\binom{d+1}{2}}P_{d}$, so the $j$-th weight in
the second form of \cref{due:eq:ext-extraction} is
\[
 (-1)^{d-j}\,4^{\binom{d+1}{2}-\binom{d-j+1}{2}-j(d-j)}\,
 \frac{\GaussianBinomial{d}{j}{4}}{P_{d}},
\]
where
\[
 \binom{d+1}{2}-\binom{d-j+1}{2}-j(d-j)
 =\tfrac12\bigl(2dj-j^{2}+j\bigr)-\bigl(dj-j^{2}\bigr)=\binom{j+1}{2}.
\]
The numerators are integers by \cref{due:lem:ext-qbinomial}, and $P_{d}$ is a
product of odd numbers.
\end{proof}

\begin{warningbox}[title={Exact, not asymptotic: the same row as in the Exponents volume}]
The weights \cref{due:eq:ext-weights} at $q=1/4$ are exactly the signed
geometric Richardson weights $w_{M,j}$ of the Exponents volume
\cite{proveit-exponents}, Part~III (labels \texttt{p3:eq:richardson-\allowbreak weights},
\texttt{p3:eq:richardson-\allowbreak moments}), with $M=d+1$ weights on the nodes
$q^{j}$, $0\le j<M$; under the bridge $\upn{n}=p_{n+1}$ our row
$\sum_{j=0}^{d}\Wrow{d}{j}\upn{N+j}$ is their combination $R_{N+1,\,d+1}$
(label \texttt{p3:eq:richardson-\allowbreak combination}).  That volume proves, for
\emph{general} $x$, only the asymptotic statement
$R_{n,M}=\RvachevUp-\am{M}q^{\binom M2}\delta_{n}^{2M}\RvachevUp^{(2M)}
+O(\delta_{n}^{2M+2})$ (label \texttt{p3:eq:richardson-\allowbreak leading}); the sign
and the factor $q^{\binom{M}{2}}$ of its leading term are the first
uncancelled moment $(-1)^{d}q^{\binom{d+1}{2}}$ of
\cref{due:eq:ext-uncancelled} applied to the mode $r=d+1$.  The contribution
of the present volume is that at a dyadic point of depth $s$ the mode
expansion \cref{due:eq:ext-model} \emph{terminates} at $r=d$ as soon as
$n\ge s$, so the very same row is an exact identity
(\cref{due:thm:ext-extraction}), not an acceleration with a remainder.  Do
not read the $O(\delta_{n}^{2M+2})$ of the Exponents volume as a bound on an
error of \cref{due:eq:ext-extraction}: at a dyadic point of depth $s$ and
$N\ge s$ that error is zero.
\end{warningbox}

\begin{corollary}[Over-ordering, other windows, other level sets]
\label{due:cor:ext-overorder}
In the situation of \cref{due:thm:ext-extraction}:
\begin{enumerate}
\item For every $M\ge d$ and every $N\ge s$,
 $\RvachevUp(x)=\sum_{j=0}^{M}\Wrow{M}{j}\qn{N+j}$: a row of higher order
 than necessary is still exact.
\item Any two windows $N,N'\ge s$ (of orders $M,M'\ge d$) return the same
 rational number.
\item For distinct offsets $0\le t_{0}<t_{1}<\cdots<t_{d}$ and $N\ge s$,
 \begin{equation}
  \RvachevUp(x)
  =\sum_{j=0}^{d}\qn{N+t_{j}}
   \prod_{\substack{0\le k\le d\\k\neq j}}
   \frac{-\Qq^{t_{k}}}{\Qq^{t_{j}}-\Qq^{t_{k}}};
  \label{due:eq:ext-nonconsecutive}
 \end{equation}
 for the stride $t_{j}=\sigma j$ this is \cref{due:eq:ext-extraction} with
 $\Qq^{\sigma}$ in place of $\Qq$.
\end{enumerate}
\end{corollary}

\begin{proof}
(1) The samples $\qn{N+j}$, $0\le j\le M$, are still values of the polynomial
$P_{N}$ of degree at most $d\le M$ at the nodes $\Qq^{j}$, and
\cref{due:lem:ext-moments}\ref{due:it:ext-exact} applies with $M$ in place of
$d$.  (2) Both windows return $\RvachevUp(x)$.  (3) The values
$\qn{N+t_{j}}=P_{N}(\Qq^{t_{j}})$ sit at the distinct nodes $\Qq^{t_{j}}$, and
the displayed product is the Lagrange weight of the node $\Qq^{t_{j}}$ for
evaluation at $0$; for $t_{j}=\sigma j$ the nodes are the powers of
$\Qq^{\sigma}$, which satisfy \cref{due:def:ext-weights}, and
\cref{due:thm:ext-weights} applies with $q=\Qq^{\sigma}$.  Only in the
consecutive case does the common factor $\Qq^{N}$ cancel to a $q$-Pochhammer
quotient.
\end{proof}

\begin{corollary}[One completely explicit finite formula]
\label{due:cor:ext-explicit}
% ed.: the Thue--Morse truncated-power representation is proved in the setting
% ed.: part of this volume; it is quoted here in the (n+1)-factor indexing, and
% ed.: was re-validated in exact arithmetic (up(1/2)=1/2, up(3/4)=5/72,
% ed.: up(7/8)=1/288) before being printed.
Insert into \cref{due:eq:ext-extraction} the Thue--Morse truncated-power
representation of the finite spline, \cref{due:thm:set-finite-spline}, which in
the present indexing reads, for $n\ge0$ and $x\in[-1,1]$,
\[
 \upn{n}(x)
 =\frac{2^{\binom{n+1}{2}}}{n!}
  \sum_{k=0}^{2^{n+1}-1}\ThueMorseSign{k}
  \Bigl(x+1-2^{-(n+1)}-\frac{k}{2^{n}}\Bigr)_{+}^{\,n},
\]
with $\ThueMorseSign{k}=(-1)^{\text{(binary digit sum of }k)}$ and
$t_{+}^{n}=\max(t,0)^{n}$.  Then, for every $N\ge s$,
\begin{equation}
 \RvachevUp(x)
 =\sum_{j=0}^{d}
 \frac{(-1)^{d-j}\,4^{\binom{j+1}{2}}}{P_{j}P_{d-j}}\;
 \frac{2^{\binom{N+j+1}{2}}}{(N+j)!}
 \sum_{k=0}^{2^{N+j+1}-1}\ThueMorseSign{k}
 \Bigl(x+1-2^{-(N+j+1)}-\frac{k}{2^{N+j}}\Bigr)_{+}^{\,N+j}.
 \label{due:eq:ext-explicit}
\end{equation}
Every sum and product is finite, every term is rational, and $N=s$ gives the
shortest guaranteed instance, with exactly $d+1$ finite splines.
\end{corollary}

\begin{proof}
Substitute the representation of $\upn{N+j}(x)$ into
\cref{due:eq:ext-integer-product}.  No passage to a limit remains.
\end{proof}

\subsection{The first rows, verified by hand}

\begin{example}[The rows of order $1$, $2$, $3$]
\label{due:ex:ext-hand}
With $\Qq=1/4$: $\QPochhammer{\Qq}{\Qq}{1}=\tfrac34$,
$\QPochhammer{\Qq}{\Qq}{2}=\tfrac34\cdot\tfrac{15}{16}=\tfrac{45}{64}$,
$\QPochhammer{\Qq}{\Qq}{3}=\tfrac{45}{64}\cdot\tfrac{63}{64}=\tfrac{2835}{4096}$.
Evaluating the first form of \cref{due:eq:ext-weights}:
\begin{align*}
 d=1:&\quad
 \Wrow{1}{0}=\frac{-\Qq^{1}}{\QPochhammer{\Qq}{\Qq}{1}}=-\frac{1/4}{3/4}=-\frac13,
 \qquad
 \Wrow{1}{1}=\frac{1}{\QPochhammer{\Qq}{\Qq}{1}}=\frac43;\\
 d=2:&\quad
 \Wrow{2}{0}=\frac{\Qq^{3}}{\QPochhammer{\Qq}{\Qq}{2}}=\frac{1/64}{45/64}=\frac1{45},
 \qquad
 \Wrow{2}{1}=\frac{-\Qq}{\QPochhammer{\Qq}{\Qq}{1}^{2}}=-\frac{1/4}{9/16}=-\frac49,
 \qquad
 \Wrow{2}{2}=\frac{1}{\QPochhammer{\Qq}{\Qq}{2}}=\frac{64}{45};\\
 d=3:&\quad
 \Wrow{3}{0}=\frac{-\Qq^{6}}{\QPochhammer{\Qq}{\Qq}{3}}=-\frac{1/4096}{2835/4096}=-\frac1{2835},\\
 &\quad
 \Wrow{3}{3}=\frac{1}{\QPochhammer{\Qq}{\Qq}{3}}=\frac{4096}{2835},\\
 &\quad
 \Wrow{3}{1}=\frac{\Qq^{3}}{\QPochhammer{\Qq}{\Qq}{1}\QPochhammer{\Qq}{\Qq}{2}}
 =\frac{1/64}{135/256}=\frac4{135},\\
 &\quad
 \Wrow{3}{2}=\frac{-\Qq}{\QPochhammer{\Qq}{\Qq}{2}\QPochhammer{\Qq}{\Qq}{1}}
 =-\frac{1/4}{135/256}=-\frac{64}{135}.
\end{align*}
The rows sum to $1$: $-\tfrac13+\tfrac43=1$;
$\tfrac1{45}-\tfrac{20}{45}+\tfrac{64}{45}=1$;
$\tfrac{-1+84-1344+4096}{2835}=\tfrac{2835}{2835}=1$.  The moments
$1\le m\le d$ vanish: $-\tfrac13+\tfrac43\cdot\tfrac14=0$;
$\tfrac1{45}-\tfrac49\cdot\tfrac14+\tfrac{64}{45}\cdot\tfrac1{16}
=\tfrac{1-5+4}{45}=0$; and for $d=3$
\begin{align*}
 m=1:&\quad
 \frac{-1+84\cdot4^{-1}-1344\cdot4^{-2}+4096\cdot4^{-3}}{2835}
 =\frac{-1+21-84+64}{2835}=0,\\
 m=2:&\quad
 \frac{-1+84\cdot4^{-2}-1344\cdot4^{-4}+4096\cdot4^{-6}}{2835}
 =\frac{-1+\frac{21}{4}-\frac{21}{4}+1}{2835}=0,\\
 m=3:&\quad
 \frac{-1+84\cdot4^{-3}-1344\cdot4^{-6}+4096\cdot4^{-9}}{2835}
 =\frac{-1+\frac{21}{16}-\frac{21}{64}+\frac{1}{64}}{2835}\\
 &\qquad
 =\frac{(-64+84-21+1)/64}{2835}=0.
\end{align*}
The common-denominator form \cref{due:eq:ext-base4} for $d=3$ uses
$\GaussianBinomial{3}{0}{4}=\GaussianBinomial{3}{3}{4}=1$ and
$\GaussianBinomial{3}{1}{4}=\GaussianBinomial{3}{2}{4}=\tfrac{1-4^{3}}{1-4}=21$,
giving the numerators $(-1,\;4\cdot21,\;-4^{3}\cdot21,\;4^{6})
=(-1,84,-1344,4096)$ over $P_{3}=2835$, in agreement with
$\tfrac{4}{135}=\tfrac{84}{2835}$ and $\tfrac{64}{135}=\tfrac{1344}{2835}$.
These vectors coincide with the weight vectors printed in the Exponents volume
\cite{proveit-exponents} after \texttt{p3:eq:richardson-\allowbreak leading} for
$M=2,3,4$, namely $(-\tfrac13,\tfrac43)$,
$(\tfrac1{45},-\tfrac49,\tfrac{64}{45})$ and
$(-\tfrac1{2835},\tfrac4{135},-\tfrac{64}{135},\tfrac{4096}{2835})$; their
$M$ is our $d+1$.
\end{example}

\begin{table}[htbp]
\centering
\caption{The quarter-base rows through order $5$, in the common-denominator
form \cref{due:eq:ext-base4}.  The integer vector multiplies
$(\qn{N},\qn{N+1},\ldots,\qn{N+d})$ and the result is divided by $P_{d}$.
Rows $d\le5$ agree with the verification outputs of S1 (\texttt{m=0..5}) and
S4 (\texttt{d=0..4}); the rightmost numerator is always $4^{\binom{d+1}{2}}$
and the leftmost is $(-1)^{d}$.}
\label{due:tab:ext-rows}
\begin{tabular}{c@{\qquad}l@{\qquad}r}
\toprule
$d$ & integer numerator vector & $P_{d}$\\
\midrule
0 & $(1)$ & $1$\\
1 & $(-1,\ 4)$ & $3$\\
2 & $(1,\ -20,\ 64)$ & $45$\\
3 & $(-1,\ 84,\ -1344,\ 4096)$ & $2835$\\
4 & $(1,\ -340,\ 22848,\ -348160,\ 1048576)$ & $722925$\\
5 & $(-1,\ 1364,\ -371008,\ 23744512,\ -357564416,\ 1073741824)$ & $739552275$\\
\bottomrule
\end{tabular}
\end{table}

The rows depend on $x$ only through $d=\Floor{\dep{x}/2}$: one row serves
every dyadic point of depths $2d$ and $2d+1$, and every admissible starting
level $N\ge s$.  The sample count therefore grows by one only every second
time the depth grows: depths $0,1$ need one sample, depths $2,3$ need two,
depths $4,5$ need three, and so on.

\begin{remark}[Lean crosswalk]
\label{due:rmk:ext-lean}
The row and its algebra are formalized in the repository
\cite{proveit-primary}; in Lean the order $d$ is the parameter \texttt{p} and
the nodes are \texttt{q\,\^{}\,j}, \texttt{j <= p}.
\begin{itemize}
\item The weight of \cref{due:def:ext-weights} is
 \nolinkurl{Fabius.geometricLagrangeWeight} (the evaluation-at-zero Lagrange
 weight on the nodes $q^{0},\ldots,q^{d}$), with its literal product form
 \nolinkurl{Fabius.geometricLagrangeWeight_eq_product}
 (\nolinkurl{GeometricLagrangeWeights.lean}).
\item \Cref{due:thm:ext-weights} is
 \nolinkurl{Fabius.geometricLagrangeWeight_eq_geometricQPochhammer} over any
 field with $q\ne0$; at the quarter base, unconditionally,
 \nolinkurl{Fabius.quarterLagrangeWeight_eq_qPochhammer} and
 \nolinkurl{Fabius.quarterLagrangeWeight_eq_qBinomial}
 (\nolinkurl{GeometricLagrangeQBinomial.lean}).  The Gaussian form carries
 the explicit hypothesis $\QPochhammer{q}{q}{d}\ne0$ in its field-generic
 version, discharged at $q=1/4$ by
 \nolinkurl{Fabius.quarter_qPochhammer_self_ne_zero}.
\item The moment identities \cref{due:eq:ext-moments} are
 \nolinkurl{Fabius.sum_quarterLagrangeWeight_eq_one} and
 \nolinkurl{Fabius.sum_quarterLagrangeWeight_mul_pow_eq_zero}; the
 uncancelled moments \cref{due:eq:ext-uncancelled} are
 \nolinkurl{Fabius.sum_geometricLagrangeWeight_mul_pow_card_add} (with the
 complete homogeneous polynomial) and
 \nolinkurl{Fabius.sum_quarterLagrangeWeight_firstUncancelled}.
\item The row polynomial identity \cref{due:eq:ext-rowpoly} is
 \nolinkurl{Fabius.quarterLagrangeWeightPolynomial_eq_forwardRichardson},
 whose right side is the normalized root polynomial with roots
 $q,\ldots,q^{d}$.
\item The filter's action on the \emph{infinite} product is formalized as
 \nolinkurl{Fabius.quarterLagrange_rvachevFourierProduct_eq_gaussian_tsum}
 (\nolinkurl{RvachevQBinomialFilter.lean}): for the infinite sinc product
 $\Phi$ in the repository's normalization,
 \[
  \sum_{j=0}^{d}\Wrow{d}{j}\,\Phi(2^{-j}z)
  =1+\sum_{r\ge0}(-1)^{d}\Qq^{\binom{d+1}{2}}
  \GaussianBinomial{d+r}{d}{\Qq}\,\mathrm{m}_{d+1+r}(z),
 \]
 where $\mathrm m_{k}(z)$ is the $k$-th even Taylor mode of $\Phi$: the
 quarter-base row cancels the first $d$ nonconstant even moments of the sinc
 product and multiplies every surviving mode by an explicit Gaussian
 coefficient.  This is the Fourier-side counterpart of
 \cref{due:eq:ext-uncancelled}; on the function side it is the asymptotic
 statement of the Exponents volume.
\end{itemize}
What is \emph{not} formalized is \cref{due:thm:ext-extraction} itself,
because its input \cref{due:eq:ext-model} (the exact eventual geometric law
at a dyadic point) has no Lean proof yet; once it has one, the theorem is the
moment identities applied to a finite sum and requires no new analysis.
\end{remark}

\begin{summarybox}[title={Section summary}]
Lagrange extrapolation from the geometric nodes $1,\Qq,\ldots,\Qq^{d}$ to the
origin has the $q$-Pochhammer weights
$\Wrow{d}{j}=(-1)^{d-j}\Qq^{\binom{d-j+1}{2}}/
(\QPochhammer{\Qq}{\Qq}{j}\QPochhammer{\Qq}{\Qq}{d-j})$
(\cref{due:thm:ext-weights}); they annihilate the moments
$\Qq^{jm}$, $1\le m\le d$, preserve constants, and multiply the mode
$\Qq^{j(d+1)}$ by $(-1)^{d}\Qq^{\binom{d+1}{2}}$
(\cref{due:lem:ext-moments}).  Applied to $d+1$ consecutive finite-spline
values at a dyadic point of depth $s\le N$, the row returns $\RvachevUp(x)$
exactly (\cref{due:thm:ext-extraction}), in four equivalent notations, with
an odd common denominator $P_{d}=\prod_{r\le d}(4^{r}-1)$ and integer
numerators $(-1)^{d-j}4^{\binom{j+1}{2}}\GaussianBinomial{d}{j}{4}$.
\end{summarybox}

\section{Recovering every mode and the Romberg tableau}
\label{due:sec:recovery}

The extraction row is the first row of the inverse of a geometric Vandermonde
matrix.  This section writes down the whole inverse in closed form, so that
the same $d+1$ samples return every mode coefficient $\Cr{r}{x}$ and hence
every even derivative $\RvachevUp^{(2r)}(x)$ that occurs in
\cref{due:eq:ext-model}; it then organizes the row as a triangular
Romberg-type tableau, and turns the annihilating recurrence of the sequence
into an exact test for the onset level.

\subsection{The geometric Vandermonde system and its exact inverse}

With $B_{r}=\Cr{r}{x}\Qq^{rN}$ as in \cref{due:eq:ext-samples}, the $d+1$
samples $\qn{N},\ldots,\qn{N+d}$ satisfy
\begin{equation}
 \underbrace{\begin{pmatrix}
  1&1&1&\cdots&1\\
  1&\Qq&\Qq^{2}&\cdots&\Qq^{d}\\
  1&\Qq^{2}&\Qq^{4}&\cdots&\Qq^{2d}\\
  \vdots&\vdots&\vdots&&\vdots\\
  1&\Qq^{d}&\Qq^{2d}&\cdots&\Qq^{d^{2}}
 \end{pmatrix}}_{V=(\Qq^{ij})_{0\le i,j\le d}}
 \begin{pmatrix}B_{0}\\B_{1}\\B_{2}\\\vdots\\B_{d}\end{pmatrix}
 =\begin{pmatrix}\qn{N}\\\qn{N+1}\\\qn{N+2}\\\vdots\\\qn{N+d}\end{pmatrix}.
 \label{due:eq:rec-vandermonde}
\end{equation}
The matrix $V$ is symmetric.

\begin{lemma}[Determinant]
\label{due:lem:rec-det}
For any field element $q$ with distinct powers $q^{0},\ldots,q^{d}$,
\begin{equation}
 \det\bigl(q^{ij}\bigr)_{0\le i,j\le d}
 =\prod_{0\le i<k\le d}\bigl(q^{k}-q^{i}\bigr)
 =(-1)^{\binom{d+1}{2}}\,q^{\binom{d+1}{3}}\prod_{k=1}^{d}\QPochhammer{q}{q}{k}
 \;\neq0.
 \label{due:eq:rec-det}
\end{equation}
\end{lemma}

\begin{proof}
The first equality is the Vandermonde determinant at the nodes $q^{k}$.  For
fixed $k$, $\prod_{i<k}(q^{k}-q^{i})=\prod_{i<k}(-q^{i})(1-q^{k-i})
=(-1)^{k}q^{\binom{k}{2}}\QPochhammer{q}{q}{k}$; multiply over
$k=0,\ldots,d$ and use $\sum_{k=0}^{d}k=\binom{d+1}{2}$ and
$\sum_{k=0}^{d}\binom{k}{2}=\binom{d+1}{3}$.  Each factor
$\QPochhammer{q}{q}{k}$ is nonzero by the distinctness of the powers.
\end{proof}

\begin{theorem}[Closed inverse of the geometric Vandermonde system; mode projectors]
\label{due:thm:rec-inverse}
Let $N\ge s$ and let $\ell^{(d)}_{j}$ be the Lagrange basis polynomials of
\cref{due:def:ext-weights} at $q=\Qq$.  Put
\begin{equation}
\begin{aligned}
 D_{j}&\DefinitionEquals\prod_{\substack{0\le k\le d\\k\ne j}}\bigl(\Qq^{j}-\Qq^{k}\bigr)
 =(-1)^{j}\,\Qq^{\binom{j}{2}+j(d-j)}\,
   \QPochhammer{\Qq}{\Qq}{j}\,\QPochhammer{\Qq}{\Qq}{d-j},\\
 e^{(j)}_{r}&\DefinitionEquals
 e_{r}\bigl(1,\Qq,\ldots,\widehat{\Qq^{j}},\ldots,\Qq^{d}\bigr),
\end{aligned}
 \label{due:eq:rec-Dj}
\end{equation}
where $e_{r}$ is the elementary symmetric polynomial of degree $r$ and the hat
omits the node $\Qq^{j}$.  Then:
\begin{enumerate}
\item\label{due:it:rec-lagrange}
 $P_{N}(z)=\sum_{j=0}^{d}\qn{N+j}\,\ell^{(d)}_{j}(z)$, and for $0\le m\le d$
 \begin{equation}
  B_{m}=\sum_{j=0}^{d}\qn{N+j}\,[z^{m}]\ell^{(d)}_{j}(z),
  \qquad
  \bigl(V^{-1}\bigr)_{mj}=[z^{m}]\ell^{(d)}_{j}(z)=[z^{j}]\ell^{(d)}_{m}(z)
  =\frac{(-1)^{d-m}\,e^{(j)}_{d-m}}{D_{j}}.
  \label{due:eq:rec-inverse-entries}
 \end{equation}
\item\label{due:it:rec-esym}
 The elementary symmetric functions of the punctured node set have the
 Gaussian expansion
 \begin{equation}
  e^{(j)}_{r}=[t^{r}]\,\frac{\QPochhammer{-t}{\Qq}{d+1}}{1+t\Qq^{j}}
  =\sum_{h=0}^{r}(-1)^{h}\,\Qq^{jh+\binom{r-h}{2}}
   \GaussianBinomial{d+1}{r-h}{\Qq}.
  \label{due:eq:rec-esym}
 \end{equation}
\item\label{due:it:rec-projector}
 (Mode projector.)  With $\mathsf E$ the forward shift on the sample sequence,
 for $0\le m\le d$,
 \begin{equation}
  B_{m}=\frac{1}{D_{m}}
  \Bigl[\prod_{\substack{0\le k\le d\\k\ne m}}(\mathsf E-\Qq^{k})\,q\Bigr]_{N},
  \label{due:eq:rec-projector}
 \end{equation}
 and the operator expands as
 \begin{equation}
  \Bigl[\prod_{\substack{0\le k\le d\\k\ne m}}(\mathsf E-\Qq^{k})\,q\Bigr]_{N}
  =\sum_{a=0}^{m}\sum_{b=0}^{d-m}(-1)^{a+b}\,
  \Qq^{\binom{a}{2}+\binom{b}{2}+(m+1)b}\,
  \GaussianBinomial{m}{a}{\Qq}\GaussianBinomial{d-m}{b}{\Qq}\;
  \qn{N+d-a-b}.
  \label{due:eq:rec-projector-expanded}
 \end{equation}
\item\label{due:it:rec-modes}
 Consequently every mode coefficient is a rational combination of the same
 $d+1$ samples,
 \begin{equation}
  \Cr{m}{x}=\Qq^{-mN}B_{m}
  =4^{mN}\sum_{j=0}^{d}\qn{N+j}\,\frac{(-1)^{d-m}e^{(j)}_{d-m}}{D_{j}}
  \qquad(0\le m\le d),
  \label{due:eq:rec-modes}
 \end{equation}
 independent of the admissible window $N\ge s$; for $m=0$ this is
 \cref{due:eq:ext-extraction}.
\end{enumerate}
\end{theorem}

% ed.: S1 Thm "Full coefficient reconstruction", S2 Prop "Closed inverse for all
% ed.: modes", S3 Thm "Mode projector", S4 sec. "The Vandermonde system",
% ed.: S5 sec. 9, S6 sec. 5.1 -- merged into one statement.  The symmetry
% ed.: [z^m] l_j = [z^j] l_m (which reconciles S1/S2's row formula with S3's
% ed.: operator formula) is not remarked in any source; it is proved here.
\begin{proof}
\ref{due:it:rec-lagrange}: $P_{N}$ has degree at most $d$ and
$P_{N}(\Qq^{j})=\qn{N+j}$, so $P_{N}=\sum_{j}\qn{N+j}\ell^{(d)}_{j}$ as in the
proof of \cref{due:lem:ext-moments}; comparing the coefficient of $z^{m}$
gives $B_{m}=\sum_{j}\qn{N+j}[z^{m}]\ell^{(d)}_{j}$, that is
$(V^{-1})_{mj}=[z^{m}]\ell^{(d)}_{j}$, because this expresses the solution of
\cref{due:eq:rec-vandermonde} linearly in the right-hand side for every
right-hand side.  Since $V$ is symmetric, so is $V^{-1}$, whence
$[z^{m}]\ell^{(d)}_{j}=[z^{j}]\ell^{(d)}_{m}$.  Finally
$\ell^{(d)}_{j}(z)=D_{j}^{-1}\prod_{k\ne j}(z-\Qq^{k})$ and, by Vieta's
formulas, $\prod_{k\ne j}(z-\Qq^{k})=\sum_{m=0}^{d}(-1)^{d-m}e^{(j)}_{d-m}z^{m}$.
The closed form of $D_{j}$ is the denominator computation in the proof of
\cref{due:thm:ext-weights}: the factors with $k<j$ give
$(-1)^{j}\Qq^{\binom{j}{2}}\QPochhammer{\Qq}{\Qq}{j}$, those with $k>j$ give
$\Qq^{j(d-j)}\QPochhammer{\Qq}{\Qq}{d-j}$.

\ref{due:it:rec-esym}: $\prod_{k=0}^{d}(1+t\Qq^{k})=\QPochhammer{-t}{\Qq}{d+1}
=\sum_{i=0}^{d+1}\Qq^{\binom{i}{2}}\GaussianBinomial{d+1}{i}{\Qq}t^{i}$ by
\cref{due:lem:ext-qbinomial} with $z=-t$, and
$\prod_{k\ne j}(1+t\Qq^{k})=\QPochhammer{-t}{\Qq}{d+1}\cdot(1+t\Qq^{j})^{-1}$
with $(1+t\Qq^{j})^{-1}=\sum_{h\ge0}(-1)^{h}\Qq^{jh}t^{h}$ as a formal power
series.  The coefficient of $t^{r}$ in the left side is $e^{(j)}_{r}$; in the
right side it is the displayed convolution.

\ref{due:it:rec-projector}: put $u_{i}=\qn{N+i}=\sum_{r}B_{r}(\Qq^{r})^{i}$.
The shift acts on the sequence $i\mapsto(\Qq^{r})^{i}$ by the scalar $\Qq^{r}$,
so
$\bigl[\prod_{k\ne m}(\mathsf E-\Qq^{k})u\bigr]_{0}
=\sum_{r=0}^{d}B_{r}\prod_{k\ne m}(\Qq^{r}-\Qq^{k})$; for $r\ne m$ the product
contains the factor $k=r$ and vanishes, and for $r=m$ it is $D_{m}$.  This is
\cref{due:eq:rec-projector}.  For the expansion, split the operator into the
factors below and above $m$ and apply \cref{due:lem:ext-qbinomial} to each:
\[
 \prod_{k=0}^{m-1}(\mathsf E-\Qq^{k})
 =\mathsf E^{m}\QPochhammer{\mathsf E^{-1}}{\Qq}{m}
 =\sum_{a=0}^{m}(-1)^{a}\Qq^{\binom{a}{2}}\GaussianBinomial{m}{a}{\Qq}\mathsf E^{m-a},
\]
\begin{align*}
 \prod_{k=m+1}^{d}(\mathsf E-\Qq^{k})
 &=\prod_{i=1}^{d-m}\bigl(\mathsf E-\Qq^{m+i}\bigr)
 =\mathsf E^{d-m}\QPochhammer{\Qq^{m+1}\mathsf E^{-1}}{\Qq}{d-m}\\
 &=\sum_{b=0}^{d-m}(-1)^{b}\Qq^{\binom{b}{2}+(m+1)b}
  \GaussianBinomial{d-m}{b}{\Qq}\mathsf E^{d-m-b}.
\end{align*}
(The operators are polynomials in the single commuting symbol $\mathsf E$, so
the formal manipulation with $\mathsf E^{-1}$ is legitimate after multiplying
back by the displayed power of $\mathsf E$.)  Multiplying and applying
$\mathsf E^{d-a-b}$ to $u$ at index $0$ gives $u_{d-a-b}=\qn{N+d-a-b}$.
Comparing \cref{due:eq:rec-projector} with \cref{due:eq:rec-inverse-entries}
is the symmetry $[z^{i}]\ell^{(d)}_{m}=[z^{m}]\ell^{(d)}_{i}$ once more, since
$D_{m}^{-1}\prod_{k\ne m}(\mathsf E-\Qq^{k})=\ell^{(d)}_{m}(\mathsf E)$.

\ref{due:it:rec-modes}: $B_{m}=\Cr{m}{x}\Qq^{mN}$ by definition, and
$\Cr{m}{x}$ does not depend on $N$; the right side of \cref{due:eq:rec-modes}
is therefore the same rational number for every window $N\ge s$.
\end{proof}

\begin{corollary}[Even-jet recovery]
\label{due:cor:rec-jet}
For $1\le m\le d$ and any $N\ge s$,
\begin{equation}
 \RvachevUp^{(2m)}(x)
 =\frac{(-1)^{m}4^{m}}{\am{m}}\,\Cr{m}{x}
 =\frac{(-1)^{m}4^{m(N+1)}}{\am{m}}
  \sum_{j=0}^{d}\qn{N+j}\,\frac{(-1)^{d-m}e^{(j)}_{d-m}}{D_{j}}
 \;\in\RationalNumbers.
 \label{due:eq:rec-jet}
\end{equation}
Thus the $d+1$ consecutive finite-spline values at a dyadic point of depth $s$
determine, besides $\RvachevUp(x)$, every even derivative
$\RvachevUp^{(2)}(x),\ldots,\RvachevUp^{(2d)}(x)$ that occurs in the
terminating expansion, and all of them are rational.
\end{corollary}

\begin{proof}
Solve $\Cr{m}{x}=(-1)^{m}\am{m}4^{-m}\RvachevUp^{(2m)}(x)$ for the derivative;
$\am{m}\neq0$ because the $\am{m}$ are positive.  Insert
\cref{due:eq:rec-modes}.  Rationality: the samples, the nodes, and $\am{m}$
are rational.
\end{proof}

\begin{example}[The two modes at $x=1/16$]
\label{due:ex:rec-sixteenth}
% ed.: all rationals below recomputed exactly from the truncated-power
% ed.: representation and Gauss--Jordan on the 3x3 system; they agree with
% ed.: S1 verification_output.txt lines 16-17, S2 experiment_output.txt lines
% ed.: 25-27, and S5's table (5/648, -248/2025).
For $x=1/16$ one has $s=4$, $d=2$, and the window $N=4$ gives
\[
 \qn{4}=\frac{24575}{24576},\qquad
 \qn{5}=\frac{1965959}{1966080},\qquad
 \qn{6}=\frac{94365509}{94371840}.
\]
Solving \cref{due:eq:rec-vandermonde} exactly,
$(B_{0},B_{1},B_{2})=\bigl(\tfrac{2073457}{2073600},\tfrac{5}{165888},
-\tfrac{31}{16588800}\bigr)$, hence
\[
 \RvachevUp\bigl(\tfrac1{16}\bigr)=\frac{2073457}{2073600},\qquad
 \Cr{1}{\tfrac1{16}}=4^{4}B_{1}=\frac{5}{648},\qquad
 \Cr{2}{\tfrac1{16}}=4^{8}B_{2}=-\frac{248}{2025},
\]
and the window $N=5$ returns the same three numbers.  By
\cref{due:eq:rec-jet} with $\am{1}=\tfrac1{18}$, $\am{2}=\tfrac{31}{16200}$:
$\RvachevUp''(1/16)=-4\cdot18\cdot\tfrac{5}{648}=-\tfrac59$ and
$\RvachevUp^{(4)}(1/16)=16\cdot\tfrac{16200}{31}\cdot\bigl(-\tfrac{248}{2025}\bigr)
=-1024$.  The first value has an independent check from the
functional-differential equation $\RvachevUp'(x)=2\RvachevUp(2x+1)-2\RvachevUp(2x-1)$:
$\RvachevUp''(1/16)=4\bigl(\RvachevUp'(9/8)-\RvachevUp'(-7/8)\bigr)
=-4\RvachevUp'(-7/8)=-4\cdot2\,\RvachevUp(-3/4)=-8\cdot\tfrac{5}{72}=-\tfrac59$,
using $\RvachevUp(3/4)=5/72$ (itself the $d=1$ row at $x=3/4$).
\end{example}

\begin{proposition}[Why $d+1$ samples are needed within the mode model]
\label{due:prop:rec-minimal}
Consider the class $\mathcal M_{d}$ of all sequences
$n\mapsto L+\sum_{r=1}^{d}A_{r}\Qq^{rn}$ with $L,A_{1},\ldots,A_{d}\in\Kk$.
\begin{enumerate}
\item Any $d+1$ samples at distinct levels determine $(L,A_{1},\ldots,A_{d})$
 uniquely (\cref{due:lem:rec-det} for consecutive levels; for arbitrary
 distinct levels $n_{0}<\cdots<n_{d}$ the matrix $(\Qq^{rn_{i}})$ is the
 Vandermonde matrix of the distinct nodes $\Qq^{n_{i}}$).
\item For any $k\le d$ levels $n_{0}<\cdots<n_{k-1}$ there exist two sequences
 in $\mathcal M_{d}$ with identical values at these $k$ levels and different
 constants $L$.  Hence no rule whatsoever (linear or not) that reads only
 $k\le d$ samples returns $L$ on all of $\mathcal M_{d}$.
\item If $\Cr{d}{x}\neq0$, the up-function sequence $\qn{n}$ ($n\ge s$) has
 exact degree $d$ in $\Qq^{n}$; so $\mathcal M_{d}$ cannot be replaced by
 $\mathcal M_{d-1}$ for this $x$, and the order-$d$ row is the minimal
 universal rule.
\end{enumerate}
\end{proposition}

% ed.: S4 Prop "Why d+1 samples are generically necessary", S6 sec. 5.5, S1
% ed.: sec. 11.3.  The sources prove (2) only for consecutive levels and only
% ed.: for linear rules; the polynomial construction below covers both.
\begin{proof}
(1) as stated.  (2) Take a sequence $\qn{n}=\sum_{r}B_{r}\Qq^{r(n-N)}$ of the
class, $N\DefinitionEquals n_{0}$, and let
$v(z)=\prod_{i=0}^{k-1}\bigl(z-\Qq^{n_{i}-N}\bigr)=\sum_{r=0}^{k}v_{r}z^{r}$,
a polynomial of degree $k\le d$.  The sequence
$\qn{n}'=\qn{n}+\sum_{r=0}^{k}v_{r}\Qq^{r(n-N)}$ belongs to $\mathcal M_{d}$,
agrees with $\qn{n}$ at every $n_{i}$ because $v(\Qq^{n_{i}-N})=0$, and its
constant differs from $L$ by $v(0)=(-1)^{k}\Qq^{\sum_{i}(n_{i}-N)}\neq0$.
(3) If $\Cr{d}{x}\ne0$, the polynomial $P_{N}$ of \cref{due:eq:ext-samples}
has degree exactly $d$, so the sequence is not in $\mathcal M_{d-1}$: a
nontrivial combination of the modes $1,\Qq^{n},\ldots,\Qq^{dn}$ with distinct
ratios is not the zero sequence (Vandermonde again).
\end{proof}

% ed.: the nonvanishing of C_d(x) for s >= 2 is S6's theorem (S3 proves the
% ed.: top mode only); it is proved in the jet section and used here only as
% ed.: a hypothesis.
The nonvanishing $\Cr{d}{x}\ne0$ for every dyadic $x$ of depth $s\ge2$ is
\cref{due:cor:mech-even-jet}, restated as \cref{due:thm:sharp-modes}; here it
enters only as a hypothesis.  Statement (2) concerns the abstract model class: it
does not exclude that a rule exploiting point-specific information about the
particular numbers $\Cr{r}{x}$ could use fewer samples, but no such
information is available without first knowing the answer.

\subsection{The Romberg tableau}

The row can be evaluated without forming its large numerators, one mode at a
time, by the classical Romberg pattern (Richardson's deferred approach to the
limit \cite{richardson1911}, in the geometric-node form treated in
\cite{brezinski-redivo-zaglia}).

\begin{algorithm}[Exact quarter-base $q$-Richardson tableau]
\label{due:alg:rec-tableau}
Given exact samples $\qn{n}$ for $n\ge N$, set $T_{0,n}=\qn{n}$ and, for
$m=1,2,\ldots$,
\begin{equation}
 \boxed{\;
 T_{m,n}=\frac{T_{m-1,n+1}-\Qq^{m}\,T_{m-1,n}}{1-\Qq^{m}}
 =\frac{4^{m}\,T_{m-1,n+1}-T_{m-1,n}}{4^{m}-1}.\;}
 \label{due:eq:rec-tableau}
\end{equation}
Entry $T_{m,n}$ uses the samples $\qn{n},\ldots,\qn{n+m}$.  Output $T_{d,N}$.
\end{algorithm}

\begin{theorem}[What the tableau computes]
\label{due:thm:rec-tableau}
For every $m\ge0$ and $n\ge N$:
\begin{enumerate}
\item\label{due:it:rec-tab-row}
 $T_{m,n}=\sum_{j=0}^{m}\Wrow{m}{j}\qn{n+j}
 =\Bigl[\prod_{i=1}^{m}\frac{\mathsf E-\Qq^{i}}{1-\Qq^{i}}\,q\Bigr]_{n}$: the
 $m$-th column of the tableau is the extraction row of order $m$ applied at
 level $n$.
\item\label{due:it:rec-tab-neville}
 $T_{m,n}$ is the value at $z=0$ of the polynomial of degree at most $m$
 interpolating the data $(\Qq^{j},\qn{n+j})$, $0\le j\le m$; the recursion
 \cref{due:eq:rec-tableau} is Neville's recurrence for these nodes.
\item\label{due:it:rec-tab-modes}
 If $n\ge s$, then
 \begin{equation}
  T_{m,n}=\RvachevUp(x)+\sum_{r=m+1}^{d}\Cr{r}{x}\,\Qq^{rn}
  \prod_{i=1}^{m}\frac{\Qq^{r}-\Qq^{i}}{1-\Qq^{i}}.
  \label{due:eq:rec-tableau-modes}
 \end{equation}
 In particular $T_{m,n}=\RvachevUp(x)$ for every $m\ge d$ and every $n\ge s$,
 while for $m<d$ the column $T_{m,\cdot}$ is a nonconstant function of $n$
 whenever $\Cr{m+1}{x}\ne0$, its deviation from $\RvachevUp(x)$ being an
 exact finite sum whose leading term is
 $\Cr{m+1}{x}\prod_{i=1}^{m}\frac{\Qq^{m+1}-\Qq^{i}}{1-\Qq^{i}}\,\Qq^{(m+1)n}$.
\end{enumerate}
\end{theorem}

% ed.: S1 Alg "Exact quarter-base tableau", S5 Alg "Exact q-Richardson table",
% ed.: S6 sec. 5.4 -- one statement; the Neville identification (2) and the
% ed.: explicit mode formula (3) are supplied.
\begin{proof}
\ref{due:it:rec-tab-row}: induction on $m$.  For $m=0$ both sides are
$\qn{n}$.  If the operator identity holds for $m-1$, then
\cref{due:eq:rec-tableau} says
$T_{m,\cdot}=\frac{\mathsf E-\Qq^{m}}{1-\Qq^{m}}T_{m-1,\cdot}$, and composing
with the operator for $m-1$ gives the operator for $m$.  Its expansion in
powers of $\mathsf E$ is the row of order $m$ by
\cref{due:prop:ext-rowpoly}: $W_{m}(\mathsf E)=\sum_{j}\Wrow{m}{j}\mathsf E^{j}$.

\ref{due:it:rec-tab-neville}: by \ref{due:it:rec-tab-row} and
\cref{due:lem:ext-moments}\ref{due:it:ext-exact} (with $m$ in place of $d$),
$T_{m,n}=\sum_{j}\Wrow{m}{j}p(\Qq^{j})=p(0)$ for the interpolant $p$ of
$(\Qq^{j},\qn{n+j})_{0\le j\le m}$, which is the first claim.  For the
recurrence, let $p_{[0,m]}$ be that interpolant, $p_{[0,m-1]}$ the interpolant
of the first $m$ data and $p_{[1,m]}$ that of the last $m$.  Neville's identity
$p_{[0,m]}(z)=\frac{(z-\Qq^{0})p_{[1,m]}(z)-(z-\Qq^{m})p_{[0,m-1]}(z)}{\Qq^{m}-\Qq^{0}}$
(both sides have degree at most $m$ and agree at all $m+1$ nodes) gives at
$z=0$: $p_{[0,m]}(0)=\frac{p_{[1,m]}(0)-\Qq^{m}p_{[0,m-1]}(0)}{1-\Qq^{m}}$.
Here $p_{[0,m-1]}(0)=T_{m-1,n}$ by the first claim, and
$p_{[1,m]}(0)=T_{m-1,n+1}$ because $p_{[1,m]}(\Qq z)$ interpolates
$(\Qq^{j},\qn{n+1+j})_{0\le j\le m-1}$ and has the same value at $0$.

\ref{due:it:rec-tab-modes}: apply the operator of \ref{due:it:rec-tab-row} to
\cref{due:eq:ext-model}, valid for all indices $\ge n\ge s$: the constant is
multiplied by $\prod_{i}(1-\Qq^{i})/(1-\Qq^{i})=1$, and the mode $\Qq^{rn}$ by
$\prod_{i=1}^{m}(\Qq^{r}-\Qq^{i})/(1-\Qq^{i})$, which vanishes for
$1\le r\le m$ and is nonzero for $r>m$.  If $\Cr{m+1}{x}\ne0$ the right side
of \cref{due:eq:rec-tableau-modes} is a nonzero combination of the distinct
geometric sequences $\Qq^{rn}$, $m<r\le d$, hence nonconstant in $n$.
\end{proof}

The tableau costs $\binom{d+1}{2}$ entries, each one multiplication, one
subtraction and one division in $\RationalNumbers$; the intermediate
denominators are products of the odd numbers $4^{m}-1$ times the sample
denominators, never larger than in the closed row.  The columns are also the
natural diagnostic display (\cref{due:prop:rec-diagnostics} below): column
$d$ is constant along $n\ge s$, and column $d-1$ is not.

\begin{example}[Tableau at $x=1/16$]
\label{due:ex:rec-tableau}
% ed.: recomputed exactly; row m=2 was checked constant for n = 4,5,6, which
% ed.: uses samples through q_8.
With the samples of \cref{due:ex:rec-sixteenth} and $\qn{7},\qn{8}$:
\begin{center}
\begin{tabular}{c@{\qquad}l@{\qquad}l@{\qquad}l}
\toprule
$m$ & $T_{m,4}$ & $T_{m,5}$ & $T_{m,6}$\\
\midrule
0 & $24575/24576$ & $1965959/1966080$ & $94365509/94371840$\\
1 & $1474459/1474560$ & $70774001/70778880$ & $226476797/226492416$\\
2 & $2073457/2073600$ & $2073457/2073600$ & $2073457/2073600$\\
\bottomrule
\end{tabular}
\end{center}
Column $2$ is constant, equal to $\RvachevUp(1/16)$; column $1$ is not,
because $\Cr{2}{1/16}=-248/2025\ne0$; for instance
$T_{1,4}-\RvachevUp(1/16)=\Cr{2}{1/16}\,\Qq^{8}\,\frac{\Qq^{2}-\Qq}{1-\Qq}
=-\frac{248}{2025}\cdot\frac{1}{65536}\cdot\Bigl(-\frac14\Bigr)
=\frac{31}{66355200}$, which is $\frac{1474459}{1474560}-\frac{2073457}{2073600}
=\frac{66350655-66350624}{66355200}$ exactly.
\end{example}

\subsection{The annihilating recurrence and onset diagnostics}

\begin{theorem}[Annihilating recurrence]
\label{due:thm:rec-recurrence}
For every $n\ge s$,
\begin{equation}
 \Bigl[\prod_{m=0}^{d}(\mathsf E-\Qq^{m})\,q\Bigr]_{n}
 =\sum_{j=0}^{d+1}(-1)^{d+1-j}\,\Qq^{\binom{d+1-j}{2}}
 \GaussianBinomial{d+1}{j}{\Qq}\,\qn{n+j}=0,
 \label{due:eq:rec-recurrence}
\end{equation}
or, after multiplication by $4^{\binom{d+1}{2}}$, with integer coefficients
\begin{equation}
 \sum_{j=0}^{d+1}(-1)^{d+1-j}\,4^{\binom{j}{2}}\GaussianBinomial{d+1}{j}{4}\,\qn{n+j}=0
 \qquad(n\ge s).
 \label{due:eq:rec-recurrence-int}
\end{equation}
For $d=1$ this reads $\qn{n}-5\qn{n+1}+4\qn{n+2}=0$, for $d=2$ it reads
$-\qn{n}+21\qn{n+1}-84\qn{n+2}+64\qn{n+3}=0$.
\end{theorem}

\begin{proof}
Each factor $\mathsf E-\Qq^{m}$ annihilates the mode $\Qq^{mn}$ of
\cref{due:eq:ext-model} and the factors commute, so the product annihilates
the whole sequence from index $s$ on.  For the expansion,
$\prod_{m=0}^{d}(\mathsf E-\Qq^{m})=\mathsf E^{d+1}\QPochhammer{\mathsf E^{-1}}{\Qq}{d+1}
=\sum_{k=0}^{d+1}(-1)^{k}\Qq^{\binom{k}{2}}\GaussianBinomial{d+1}{k}{\Qq}\mathsf E^{d+1-k}$
by \cref{due:lem:ext-qbinomial}; put $j=d+1-k$ and use the symmetry of the
Gaussian coefficient.  The integer form follows from
$\GaussianBinomial{d+1}{j}{\Qq}=4^{-j(d+1-j)}\GaussianBinomial{d+1}{j}{4}$
(inversion identity in the proof of \cref{due:thm:ext-extraction}) and
$\binom{d+1}{2}-\binom{d+1-j}{2}-j(d+1-j)=\binom{j}{2}$.
The displayed rows are $\GaussianBinomial{2}{1}{4}=5$ and
$\GaussianBinomial{3}{1}{4}=\GaussianBinomial{3}{2}{4}=21$ with
$4^{\binom{2}{2}}=4$, $4^{\binom{3}{2}}=64$.
\end{proof}

Note the difference from the extraction row: the recurrence contains the
factor $\mathsf E-1$ and annihilates \emph{everything}, including the
constant; the row omits that factor and therefore preserves the constant
while killing the rest.  Write
\begin{equation}
 \rho_{n}\DefinitionEquals
 \sum_{j=0}^{d+1}(-1)^{d+1-j}\,4^{\binom{j}{2}}\GaussianBinomial{d+1}{j}{4}\,\qn{n+j}
 \label{due:eq:rec-residual}
\end{equation}
for the residual of \cref{due:eq:rec-recurrence-int} at level $n$.

\begin{proposition}[Onset and model-order diagnostics]
\label{due:prop:rec-diagnostics}
Let $x$ be dyadic of depth $s$ and $d=\Floor{s/2}$.
\begin{enumerate}
\item\label{due:it:diag-onset}
 (Onset detection is a finite test.)  Let $0\le n_{0}\le s$.  If
 $\rho_{n}=0$ for every $n$ with $n_{0}\le n\le s-1$, then the geometric law
 \cref{due:eq:ext-model} already holds for all $n\ge n_{0}$, and every
 extraction row and every mode formula of
 \cref{due:thm:ext-extraction,due:thm:rec-inverse} is exact from any window
 $N\ge n_{0}$.  Conversely, if the law holds for all $n\ge n_{0}$ then
 $\rho_{n}=0$ for all $n\ge n_{0}$.
\item\label{due:it:diag-order}
 (Model order.)  For the sequence restricted to $n\ge s$, the least $k$ such
 that $\prod_{m=0}^{k}(\mathsf E-\Qq^{m})$ annihilates it equals
 $\max\SetOf{r:\Cr{r}{x}\ne0}$ (and $0$ if all modes vanish).  In particular,
 when $\Cr{d}{x}\ne0$ the least annihilating order is exactly
 $d=\Floor{s/2}$, so exact data at an \emph{unreduced} dyadic fraction reveal
 $d$, and $s$ up to parity, without reducing the fraction.
\item\label{due:it:diag-exact}
 (Certificate.)  In exact arithmetic, $\rho_{n}\ne0$ for some $n\ge s$ is
 impossible; a nonzero residual signals a wrong depth, a wrong indexing
 convention (for instance $p_{n}$ in place of $\upn{n}$), or an arithmetic
 error.  Two windows $N\neq N'$ both $\ge s$ must return identical
 rationals.
\item\label{due:it:diag-float}
 (Floating point.)  With rounded samples $\tilde q_{n}=\qn{n}+e_{n}$ and
 $n\ge s$, the computed residual is
 $\tilde\rho_{n}=\sum_{j}c_{j}e_{n+j}$ with the integer coefficients $c_{j}$
 of \cref{due:eq:rec-residual}, so
 $\AbsoluteValue{\tilde\rho_{n}}\le\bigl(\sum_{j}\AbsoluteValue{c_{j}}\bigr)\max_{j}\AbsoluteValue{e_{n+j}}$
 with $\sum_{j}\AbsoluteValue{c_{j}}=4^{\binom{d+1}{2}}\QPochhammer{-1}{\Qq}{d+1}$
 ($=10$ for $d=1$, $170$ for $d=2$), and a residual larger than this
 multiple of the sample accuracy proves that level $n$ is still in the
 transient range.
\end{enumerate}
\end{proposition}

% ed.: S4 sec. "Annihilating recurrences and onset diagnostics", S2 Thm
% ed.: "Eventual annihilating recurrence" + diagnostics, S1 Cor "Annihilating
% ed.: recurrence" + Remark "Why initial values can fail", S3 Thm 10.1.  The
% ed.: sources state onset detection informally ("holds on several successive
% ed.: blocks"); (1) makes it a finite, provably sufficient test.
\begin{proof}
\ref{due:it:diag-onset}: by \cref{due:thm:rec-recurrence}, $\rho_{n}=0$ for
$n\ge s$; with the hypothesis, $\rho_{n}=0$ for all $n\ge n_{0}$.  The
sequence $(\qn{n})_{n\ge n_{0}}$ therefore satisfies a linear recurrence of
order $d+1$ with characteristic polynomial $\prod_{m=0}^{d}(z-\Qq^{m})$, whose
roots $1,\Qq,\ldots,\Qq^{d}$ are distinct.  The solution space of such a
recurrence on $n\ge n_{0}$ is $(d+1)$-dimensional (a solution is determined
by its first $d+1$ values) and contains the $d+1$ linearly independent
geometric sequences $\Qq^{mn}$; so $\qn{n}=L'+\sum_{r=1}^{d}A'_{r}\Qq^{rn}$
for all $n\ge n_{0}$ with some constants.  On $n\ge s$ this representation
and \cref{due:eq:ext-model} both hold; their difference is a combination of
the $d+1$ distinct geometric sequences vanishing on $d+1$ consecutive indices,
hence (\cref{due:lem:rec-det}) all coefficients agree: $L'=\RvachevUp(x)$,
$A'_{r}=\Cr{r}{x}$.  So \cref{due:eq:ext-model} holds for $n\ge n_{0}$, and
the proofs of \cref{due:thm:ext-extraction,due:thm:rec-inverse} go through
verbatim with $n_{0}$ in place of $s$.  The converse is the proof of
\cref{due:thm:rec-recurrence}.

\ref{due:it:diag-order}: $\prod_{m=0}^{k}(\mathsf E-\Qq^{m})$ kills the modes
$r\le k$ and multiplies the mode $\Qq^{rn}$, $r>k$, by the nonzero scalar
$\prod_{m\le k}(\Qq^{r}-\Qq^{m})$; the image is a combination of the distinct
geometric sequences $\Qq^{rn}$, $k<r\le d$, with coefficients
$\Cr{r}{x}\prod_{m\le k}(\Qq^{r}-\Qq^{m})$, which is the zero sequence
exactly when all $\Cr{r}{x}$ with $r>k$ vanish.

\ref{due:it:diag-exact} restates \cref{due:thm:rec-recurrence} and
\cref{due:cor:ext-overorder}.

\ref{due:it:diag-float}: linearity, the triangle inequality, and
\[
 \sum_{j}\AbsoluteValue{c_{j}}
 =4^{\binom{d+1}{2}}\sum_{j=0}^{d+1}\Qq^{\binom{d+1-j}{2}}\GaussianBinomial{d+1}{j}{\Qq}
 =4^{\binom{d+1}{2}}\sum_{k=0}^{d+1}\Qq^{\binom{k}{2}}\GaussianBinomial{d+1}{k}{\Qq}
 =4^{\binom{d+1}{2}}\QPochhammer{-1}{\Qq}{d+1}
\]
by \cref{due:lem:ext-qbinomial} at $z=-1$.  For $d=1$: $1+5+4=10$; for
$d=2$: $1+21+84+64=170$.
\end{proof}

\begin{remark}[What the residual sees below the depth]
\label{due:rmk:rec-below-onset}
% ed.: exact residuals computed for this remark: x = 7/8, 3/8, 1/8 (s = 3),
% ed.: 3/32, 1/32 (s = 5): rho_{s-1} = 0, rho_{s-2} != 0; x = 15/16, 1/16
% ed.: (s = 4): rho_{s-1} != 0.  Cross-references attached at assembly; the
% ed.: closed form of rho_{s-1} is the sharpness section's knot-defect theorem.
At $n=s-1$ the point $x$ is a knot of the spline $\upn{n}$, not the centre of
a polynomial cell, and the mechanism that proves \cref{due:eq:ext-model}
does not apply; the residual $\rho_{s-1}$ need not vanish, and in general it
does not: for $x=1/16$ ($s=4$, $d=2$) exact computation gives
$\rho_{2}=-\tfrac{1}{1920}$, $\rho_{3}=\tfrac{1}{46080}$, $\rho_{4}=0$.  For
odd depth the picture is different: at $x=7/8,3/8,1/8$ ($s=3$) and at
$x=3/32,1/32$ ($s=5$) one finds $\rho_{s-1}=0$ but $\rho_{s-2}\neq0$, in
agreement with \cref{due:prop:mech-odd-onset}, according to which the law
\cref{due:eq:ext-model} holds from $n=s-1$ on when $s\ge3$ is odd.  The
residual at the knot level is in fact known in closed form at every depth:
since the law holds at the levels $s,\ldots,s+d$, only the first term of
\cref{due:eq:rec-residual} survives, and
$\rho_{s-1}=(-1)^{d+1}\Delta_{s}(x)=(-1)^{d}\ThueMorseSign{k}\,2^{-\binom s2}\beta_{s}$
with $\beta_{s}=B_{s}/s!$ for even $s$ and $0$ for odd $s$
(\cref{due:thm:sharp-knot-defect}); at $x=1/16$ this is
$(-1)^{2}\cdot(-1)\cdot2^{-6}\cdot(-\tfrac1{720})=\tfrac1{46080}$, the value
computed above.  By \cref{due:prop:rec-diagnostics}%
\ref{due:it:diag-onset} the residual test detects this earlier onset
automatically, and then the row may start at $N=s-1$; the guaranteed onset
$N\ge s$ of \cref{due:thm:ext-extraction} remains the safe choice when no
test is run.
\end{remark}

\begin{summarybox}[title={Section summary}]
The $d+1$ samples $\qn{N},\ldots,\qn{N+d}$, $N\ge s$, determine every mode:
$(V^{-1})_{mj}=[z^{m}]\ell^{(d)}_{j}(z)=(-1)^{d-m}e^{(j)}_{d-m}/D_{j}$ with
$D_{j}=(-1)^{j}\Qq^{\binom j2+j(d-j)}\QPochhammer{\Qq}{\Qq}{j}
\QPochhammer{\Qq}{\Qq}{d-j}$ (\cref{due:thm:rec-inverse}), hence
$\Cr{m}{x}$ and $\RvachevUp^{(2m)}(x)=(-1)^{m}4^{m}\Cr{m}{x}/\am{m}$ are
rational combinations of the samples (\cref{due:cor:rec-jet}); fewer than
$d+1$ samples cannot do this within the mode model
(\cref{due:prop:rec-minimal}).  The Romberg recursion
$T_{m,n}=(4^{m}T_{m-1,n+1}-T_{m-1,n})/(4^{m}-1)$ evaluates the rows of all
orders at once, and its column $d$ is constant from $n=s$ on
(\cref{due:thm:rec-tableau}).  The order-$(d+1)$ recurrence with integer
coefficients $(-1)^{d+1-j}4^{\binom j2}\GaussianBinomial{d+1}{j}{4}$
annihilates the sequence for $n\ge s$; a finite residual test certifies an
earlier onset and reveals the true number of modes
(\cref{due:thm:rec-recurrence,due:prop:rec-diagnostics}).
\end{summarybox}

\section{Stability of the extraction}
\label{due:sec:stability}

Exact rational input makes the extraction exact, and that is how it should be
used.  This section records what happens when the input is inexact: the row
is uniformly well conditioned in the depth, the full mode recovery is not, and
the only genuine floating-point hazard lies upstream, in the evaluation of the
samples.

\subsection{The Lebesgue factor of the row}

\begin{proposition}[Exact $\ell^{\infty}$ noise amplification]
\label{due:prop:stab-lebesgue}
For $0<q<1$ and $d\ge0$ let
$\Lambda_{d}(q)\DefinitionEquals\sum_{j=0}^{d}\AbsoluteValue{\Wrow{d}{j}}$.  Then
\begin{equation}
 \boxed{\;
 \Lambda_{d}(q)=\frac{\QPochhammer{-q}{q}{d}}{\QPochhammer{q}{q}{d}}
 =\prod_{h=1}^{d}\frac{1+q^{h}}{1-q^{h}},\;}
 \qquad
 \Lambda_{0}<\Lambda_{1}<\Lambda_{2}<\cdots<
 \Lambda_{\infty}(q)\DefinitionEquals\frac{\QPochhammer{-q}{q}{\infty}}{\QPochhammer{q}{q}{\infty}}<\infty,
 \label{due:eq:stab-lebesgue}
\end{equation}
and $\Lambda_{d}(q)\to\Lambda_{\infty}(q)$.  At the quarter base,
\begin{equation}
 \Lambda_{\infty}(\tfrac14)
 =\frac{\QPochhammer{-1/4}{1/4}{\infty}}{\QPochhammer{1/4}{1/4}{\infty}}
 =1.9692603536682694319471936969672656700661899623\ldots
 \label{due:eq:stab-lebesgue-limit}
\end{equation}
Consequently, if the samples are perturbed, $\tilde q_{N+j}=\qn{N+j}+e_{j}$,
the row \cref{due:eq:ext-extraction} returns
$\widetilde{\RvachevUp}(x)$ with
\begin{equation}
 \AbsoluteValue{\widetilde{\RvachevUp}(x)-\RvachevUp(x)}
 \le\Lambda_{d}(\tfrac14)\max_{j}\AbsoluteValue{e_{j}}
 <1.96927\,\max_{j}\AbsoluteValue{e_{j}},
 \label{due:eq:stab-bound}
\end{equation}
uniformly in the depth, and the first inequality is attained by
$e_{j}=\varepsilon\operatorname{sign}\Wrow{d}{j}$.
\end{proposition}

% ed.: S1 Prop "Exact l-infinity noise amplification", S2 sec. 9, S3 Prop
% ed.: "Lebesgue factor", S4 Prop "Uniform weight bound", S5 sec. 9.2, S6
% ed.: sec. 5.5 -- one statement.  Monotonicity and the limit are proved (the
% ed.: sources assert them); the constant is recomputed to 46 digits from the
% ed.: partial products (400 factors; the tail is below 4^{-400}).
\begin{proof}
For $0<q<1$ every $\QPochhammer{q}{q}{k}$ is positive, so by
\cref{due:eq:ext-weights} the sign of $\Wrow{d}{j}$ is $(-1)^{d-j}$ and
$\AbsoluteValue{\Wrow{d}{j}}=q^{\binom{d-j+1}{2}}\GaussianBinomial{d}{j}{q}/\QPochhammer{q}{q}{d}$.
With $k=d-j$ and the symmetry of the Gaussian coefficient,
\[
 \Lambda_{d}(q)
 =\frac{1}{\QPochhammer{q}{q}{d}}\sum_{k=0}^{d}q^{\binom{k+1}{2}}\GaussianBinomial{d}{k}{q}
 =\frac{1}{\QPochhammer{q}{q}{d}}\sum_{k=0}^{d}(-1)^{k}q^{\binom{k}{2}}
  \GaussianBinomial{d}{k}{q}(-q)^{k}
 =\frac{\QPochhammer{-q}{q}{d}}{\QPochhammer{q}{q}{d}},
\]
the last step being \cref{due:lem:ext-qbinomial} at $z=-q$; and
$\QPochhammer{-q}{q}{d}/\QPochhammer{q}{q}{d}=\prod_{h=1}^{d}(1+q^{h})/(1-q^{h})$.
Each new factor $(1+q^{d})/(1-q^{d})$ exceeds $1$, so the sequence is
strictly increasing; both infinite products converge to positive limits
because $\sum_{h}q^{h}<\infty$, so the limit is finite.  The error bound is
the triangle inequality applied to
$\widetilde{\RvachevUp}(x)-\RvachevUp(x)=\sum_{j}\Wrow{d}{j}e_{j}$, with
equality for the displayed choice of signs.  The numerical value is the
partial product with $400$ factors; the omitted factors change it by a
relative amount below $2\sum_{h>400}4^{-h}<4^{-399}$.
\end{proof}

\begin{table}[htbp]
\centering
\caption{Lebesgue factors $\Lambda_{d}(1/4)=\sum_{j}\AbsoluteValue{\Wrow{d}{j}}$
of the quarter-base rows.  The same numbers appear in the Exponents volume
\cite{proveit-exponents} as \texttt{p3:eq:richardson-\allowbreak stability} with $M-1=d$.}
\label{due:tab:stab-lebesgue}
\begin{tabular}{c@{\qquad}l@{\qquad}l}
\toprule
$d$ & $\Lambda_{d}(1/4)$ exact & decimal\\
\midrule
0 & $1$ & $1$\\
1 & $5/3$ & $1.66667$\\
2 & $17/9$ & $1.88889$\\
3 & $1105/567$ & $1.94885$\\
4 & $3341/1701$ & $1.96414$\\
5 & $3424525/1740123$ & $1.96798$\\
$\infty$ & $\QPochhammer{-1/4}{1/4}{\infty}/\QPochhammer{1/4}{1/4}{\infty}$ & $1.96926$\\
\bottomrule
\end{tabular}
\end{table}

The bounded Lebesgue factor is a property of the geometric nodes, not of
extrapolation in general: the nodes $1,\tfrac14,\tfrac1{16},\ldots$ are
equally spaced on a logarithmic scale and the target $0$ lies at their
accumulation point, whereas polynomial extrapolation from nearly coincident
nodes has weights that grow without bound.  The base matters: at the half
base $q=1/2$, the natural base for samples taken one dyadic level apart when
the modes decay like $2^{-n}$, the same formula gives
$\Lambda_{\infty}(\tfrac12)=8.25598793577825\ldots$, four times larger; the
quarter base is available here precisely because only even derivatives occur
in \cref{due:eq:ext-model}, so consecutive levels are two powers of $1/2$
apart in every mode.  For general $q\in(0,1)$ everything in this proposition
holds verbatim; as $q\to1^{-}$ the factor $\Lambda_{d}(q)$ tends to
$\binom{2d}{d}$-type growth, since each factor $(1+q^{h})/(1-q^{h})$ diverges.

\subsection{Exact rational arithmetic versus floating point}

\begin{proposition}[Denominator bookkeeping]
\label{due:prop:stab-denominator}
Let $x$ be dyadic of depth $s$, $N\ge s$, and let $\Delta$ be the least
common denominator of $\qn{N},\ldots,\qn{N+d}$.  Then
\begin{equation}
 \RvachevUp(x)=\frac{\text{integer}}{P_{d}\,\Delta},
 \qquad P_{d}=\prod_{r=1}^{d}(4^{r}-1)\ \text{odd},
 \label{due:eq:stab-denominator}
\end{equation}
so the denominator of $\RvachevUp(x)$ divides $P_{d}\Delta$; in particular
the extraction row never introduces a new power of $2$, and the $2$-adic
valuation of $\RvachevUp(x)$ is at least $-\TwoAdicValuation(\Delta)$.
\end{proposition}

\begin{proof}
Write \cref{due:eq:ext-base4} over the common denominator $P_{d}\Delta$; the
numerators $(-1)^{d-j}4^{\binom{j+1}{2}}\GaussianBinomial{d}{j}{4}$ are
integers and $\qn{N+j}\Delta$ is an integer.  Every $4^{r}-1$ is odd.
\end{proof}

\begin{remark}[Where floating point actually fails]
\label{due:rmk:stab-float}
% ed.: the cancellation estimate is an upper bound on the possible loss; the
% ed.: sources (S3 sec. 11.3, S4 sec. 11.4) only say "exact input remains
% ed.: preferable".  The Thue--Morse representation is quoted through
% ed.: cor:ext-explicit, which cites thm:set-finite-spline.
Suppose the samples are available only as floating-point numbers with
absolute errors $\AbsoluteValue{e_{j}}\le\varepsilon$, and the dot product
\cref{due:eq:ext-extraction} is evaluated in IEEE arithmetic with unit
roundoff $u$ ($u=2^{-53}$ in double precision) from correctly rounded weights.
The standard forward error bound for a $(d+1)$-term dot product, together with
\cref{due:prop:stab-lebesgue}, gives a total error at most
\[
 \Lambda_{d}(\tfrac14)\Bigl(\varepsilon+\gamma_{d+2}\max_{j}\AbsoluteValue{\tilde q_{N+j}}\Bigr),
 \qquad\gamma_{k}=\frac{ku}{1-ku},
\]
that is, essentially $2\varepsilon$ plus a few units of roundoff.  The row is
therefore not the problem.  The problem is $\varepsilon$ itself when the
samples are computed from the Thue--Morse truncated-power representation
(\cref{due:cor:ext-explicit}) in floating point: in the scaled form
$\upn{n}(x)=\bigl(2^{\binom{n+1}{2}}n!\bigr)^{-1}\sum_{0\le k\le y}
\ThueMorseSign{k}(2y-2k)^{n}$ with $y=2^{n}(x+1-2^{-(n+1)})$, the summands are
integers of size up to $(2y)^{n}\le2^{n(n+2)}$ while their alternating sum is
at most $2^{\binom{n+1}{2}}n!\,\max\AbsoluteValue{\upn{n}}$, so up to
$n(n+2)-\binom{n+1}{2}-\log_{2}n!=\tfrac12n(n+3)-\log_{2}n!$ bits can cancel:
about $23$ bits at $n=7$ and about $43$ bits at $n=10$, close to the entire
double-precision mantissa.  Whether this worst case is attained at a given
point is beside the point; the remedy is free.  With $2y$ an odd integer
(midpoint alignment at $n\ge s$) the sum is an exact integer computation
followed by one rational normalization, and then the extraction is exact by
\cref{due:thm:ext-extraction}.  If, conversely, only a high-precision
floating approximation $\tilde v$ of $\RvachevUp(x)$ is at hand, the exact
value can still be recovered: \cref{due:prop:stab-denominator} bounds the
denominator by $D=P_{d}\Delta$, and a rational $a/b$ with $b\le D$ is
determined uniquely by any approximation with
$\AbsoluteValue{\tilde v-a/b}<1/(2D^{2})$, the continued-fraction convergents
of $\tilde v$ producing it.
\end{remark}

\subsection{Conditioning of the full Vandermonde recovery}

The row is the first line of $V^{-1}$; the other lines recover the decaying
modes and are necessarily much larger, because they must undo the factors
$\Qq^{mN}$.

\begin{proposition}[Norms of the geometric Vandermonde inverse]
\label{due:prop:stab-vandermonde}
Let $V=(\Qq^{ij})_{0\le i,j\le d}$ and
$(V^{-1})_{mj}=[z^{m}]\ell^{(d)}_{j}(z)$ as in \cref{due:thm:rec-inverse}.
Write $\Lambda^{(m)}_{d}\DefinitionEquals\sum_{j=0}^{d}\AbsoluteValue{(V^{-1})_{mj}}$
for the absolute row sums.  Then:
\begin{enumerate}
\item $\Norm{V}_{\infty}=d+1$ (the first row), and
 $\Lambda^{(0)}_{d}=\Lambda_{d}(\tfrac14)$.
\item The last row sums exactly:
 \begin{equation}
  \Lambda^{(d)}_{d}=\sum_{j=0}^{d}\frac{1}{\AbsoluteValue{D_{j}}}
  =4^{\binom{d}{2}}\,\frac{\QPochhammer{-1}{\Qq}{d}}{\QPochhammer{\Qq}{\Qq}{d}}
  =2\cdot4^{\binom{d}{2}}\,\frac{\QPochhammer{-\Qq}{\Qq}{d-1}}{\QPochhammer{\Qq}{\Qq}{d}}
  \qquad(d\ge1).
  \label{due:eq:stab-lastrow}
 \end{equation}
\item The absolute column sums are
 \begin{equation}
  \sum_{m=0}^{d}\AbsoluteValue{(V^{-1})_{mj}}
  =\AbsoluteValue{\ell^{(d)}_{j}(-1)}
  =\frac{\QPochhammer{-1}{\Qq}{d+1}}{(1+\Qq^{j})\AbsoluteValue{D_{j}}},
  \qquad
  \Norm{V^{-1}}_{1}=\max_{j}\frac{\QPochhammer{-1}{\Qq}{d+1}}{(1+\Qq^{j})\AbsoluteValue{D_{j}}}.
  \label{due:eq:stab-columns}
 \end{equation}
\item Hence
 $\kappa_{\infty}(V)=\Norm{V}_{\infty}\Norm{V^{-1}}_{\infty}
 \ge(d+1)\Lambda^{(d)}_{d}\ge2(d+1)4^{\binom d2}$: the recovery of all
 modes is ill conditioned in absolute terms, at least like $4^{d^{2}/2}$.
\end{enumerate}
If the samples carry absolute errors at most $\varepsilon$, the recovered
$B_{m}$ carries an absolute error at most $\Lambda^{(m)}_{d}\varepsilon$
(sharp), and the mode coefficient $\Cr{m}{x}=4^{mN}B_{m}$ an absolute error
at most $4^{mN}\Lambda^{(m)}_{d}\varepsilon$.
\end{proposition}

% ed.: the sources only state the l-infinity factor of the first row; the
% ed.: closed forms (2)-(3) and the table are new, verified in exact
% ed.: arithmetic for d <= 5.
\begin{proof}
(1) The $i$-th row sum of $V$ is $\sum_{j}\Qq^{ij}\le d+1$ with equality only
for $i=0$; the first row of $V^{-1}$ is the extraction row.

(2) $(V^{-1})_{dj}=[z^{d}]\ell^{(d)}_{j}=1/D_{j}$.  By \cref{due:eq:rec-Dj},
$\AbsoluteValue{D_{j}}=\Qq^{\binom{j}{2}+j(d-j)}\QPochhammer{\Qq}{\Qq}{j}\QPochhammer{\Qq}{\Qq}{d-j}$,
and $\binom{j}{2}+j(d-j)=\tfrac12j(2d-j-1)=\binom{d}{2}-\binom{d-j}{2}$.  Hence
\[
 \sum_{j}\frac{1}{\AbsoluteValue{D_{j}}}
 =\frac{\Qq^{-\binom d2}}{\QPochhammer{\Qq}{\Qq}{d}}
  \sum_{j=0}^{d}\Qq^{\binom{d-j}{2}}\GaussianBinomial{d}{j}{\Qq}
 =\frac{\Qq^{-\binom d2}}{\QPochhammer{\Qq}{\Qq}{d}}\,\QPochhammer{-1}{\Qq}{d},
\]
by \cref{due:lem:ext-qbinomial} at $z=-1$ (with $k=d-j$).  Finally
$\QPochhammer{-1}{\Qq}{d}=(1+1)\prod_{h=1}^{d-1}(1+\Qq^{h})
=2\QPochhammer{-\Qq}{\Qq}{d-1}$ and $\Qq^{-\binom d2}=4^{\binom d2}$.

(3) $\ell^{(d)}_{j}(z)=D_{j}^{-1}\prod_{k\ne j}(z-\Qq^{k})$, and the
coefficients of $\prod_{k\ne j}(z-\Qq^{k})$ are
$(-1)^{d-m}e^{(j)}_{d-m}$ with $e^{(j)}_{r}>0$ (elementary symmetric functions
of positive numbers).  Therefore
$\sum_{m}\AbsoluteValue{[z^{m}]\ell^{(d)}_{j}}
=\AbsoluteValue{D_{j}}^{-1}\sum_{r}e^{(j)}_{r}
=\AbsoluteValue{D_{j}}^{-1}\prod_{k\ne j}(1+\Qq^{k})
=\AbsoluteValue{\ell^{(d)}_{j}(-1)}$, and
$\prod_{k\ne j}(1+\Qq^{k})=\QPochhammer{-1}{\Qq}{d+1}/(1+\Qq^{j})$.  The
$1$-norm is the largest column sum.

(4) $\Norm{V^{-1}}_{\infty}=\max_{m}\Lambda^{(m)}_{d}\ge\Lambda^{(d)}_{d}$,
and $\QPochhammer{-1}{\Qq}{d}/\QPochhammer{\Qq}{\Qq}{d}\ge2$.  The error
statements are the triangle inequality on $B=V^{-1}q$, row by row, with
equality for $e_{j}=\varepsilon\operatorname{sign}(V^{-1})_{mj}$.
\end{proof}

\begin{table}[htbp]
\centering
\caption{Absolute row sums $\Lambda^{(m)}_{d}$ of $V^{-1}$ for the quarter
base, exact, and the condition number
$\kappa_{\infty}(V)=(d+1)\max_{m}\Lambda^{(m)}_{d}$.  Row $m=0$ is the
Lebesgue factor of \cref{due:tab:stab-lebesgue}; row $m=d$ is
\cref{due:eq:stab-lastrow}.  Note that for $d\ge4$ the largest row is
$m=d-1$, not $m=d$.}
\label{due:tab:stab-condition}
\begin{tabular}{c@{\quad}l@{\quad}l}
\toprule
$d$ & $\Lambda^{(0)}_{d},\Lambda^{(1)}_{d},\ldots,\Lambda^{(d)}_{d}$ & $\kappa_{\infty}(V)$\\
\midrule
1 & $\tfrac53,\ \tfrac83$ & $\tfrac{16}{3}\approx5.33$\\[2pt]
2 & $\tfrac{17}{9},\ \tfrac{136}{9},\ \tfrac{128}{9}$ & $\tfrac{136}{3}\approx45.3$\\[2pt]
3 & $\tfrac{1105}{567},\ \tfrac{1768}{27},\ \tfrac{8320}{27},\ \tfrac{139264}{567}$
  & $\tfrac{33280}{27}\approx1.23\cdot10^{3}$\\[2pt]
4 & $\tfrac{3341}{1701},\ \tfrac{454376}{1701},\ \tfrac{427648}{81},\ \tfrac{35790848}{1701},\ \tfrac{27262976}{1701}$
  & $\tfrac{178954240}{1701}\approx1.05\cdot10^{5}$\\[2pt]
5 & $\tfrac{3424525}{1740123},\ \tfrac{5479240}{5103},\ \tfrac{438339200}{5103},\ \tfrac{7337123840}{5103},\ \tfrac{27944550400}{5103},\ \tfrac{7174742867968}{1740123}$
  & $\tfrac{55889100800}{1701}\approx3.29\cdot10^{7}$\\
\bottomrule
\end{tabular}
\end{table}

The growth is inherent, not an artefact of the method: the mode
$B_{m}=\Cr{m}{x}4^{-mN}$ is exponentially small in $N$, so any procedure that
reads it off samples of absolute accuracy $\varepsilon$ needs
$\varepsilon\ll\AbsoluteValue{\Cr{m}{x}}4^{-mN}$, and the higher modes are
recoverable in floating point only from samples at the smallest admissible
level $N=s$ and with high working precision.  With exact rational samples the
question does not arise: $V^{-1}$ is a fixed rational matrix, and
\cref{due:eq:rec-modes} returns each $\Cr{m}{x}$ as an exact rational number,
as in \cref{due:ex:rec-sixteenth}.  In terms of cost, the closed row needs
$O(d)$ rational operations once the $P_{k}$ are tabulated, the tableau and
the closed inverse $O(d^{2})$, and a dense solve of
\cref{due:eq:rec-vandermonde} $O(d^{3})$; all are negligible against the
evaluation of the samples themselves, whose truncated-power sums have
$2^{n+1}$ terms.

\begin{summarybox}[title={Section summary}]
The quarter-base row has Lebesgue factor
$\Lambda_{d}(1/4)=\QPochhammer{-1/4}{1/4}{d}/\QPochhammer{1/4}{1/4}{d}
<1.969260353668\ldots$, so inexact samples are amplified by less than a
factor $2$ at every depth (\cref{due:prop:stab-lebesgue}); the extraction
adds only the odd factor $P_{d}$ to the sample denominators
(\cref{due:prop:stab-denominator}); the floating-point hazard is the
alternating Thue--Morse sum that produces the samples, not the row
(\cref{due:rmk:stab-float}); and recovering the decaying modes themselves is
exponentially ill conditioned in absolute terms, with the exact last-row sum
$2\cdot4^{\binom d2}\QPochhammer{-1/4}{1/4}{d-1}/\QPochhammer{1/4}{1/4}{d}$
(\cref{due:prop:stab-vandermonde}), which is why the mode recovery of
\cref{due:sec:recovery} is meant for exact rational input.
\end{summarybox}

% ---- fragment 05-examples-register ----
% ==========================================================================
% Cluster 05: worked examples, the exact algorithm and its verification,
% consequences and extensions, corpus relation and formalization register
% ==========================================================================

\section{Worked exact examples}
\label{due:sec:examples}

The six sources between them work out about twenty dyadic points.  This
section consolidates them into one graded set, in the volume's indexing
($\upn{n}$ has $n+1$ factors, $\qn{n}=\upn{n}(x)$), with every number an
exact rational.  Every displayed value was recomputed for this volume in
exact arithmetic by an independent evaluator (Thue--Morse truncated powers
for the samples, the closed row for the extraction, an exact Vandermonde
solve for the modes, and the derivative tiling of \cref{due:cor:mech-tiling}
for the jet, the last two agreeing term by term), and cross-checked
against the captured verification outputs of the six packages.  Where a
source's text disagreed with its own program output, the output was
followed and the repair recorded in \cref{due:sec:repairs}.
% ed.: S1's tables print the mode coefficients in the (n+1)-convention
% ed.: (its A_ell = C_ell 4^ell); S4's table prints C_m as here; S6 prints
% ed.: A_k = C_k 4^k with the argument folded to a = 1 - |x|.  All numbers
% ed.: below are C_m in the volume's convention, with the source form noted
% ed.: once at x = 1/16.

\subsection{The extraction rows}
\label{due:sec:exr-rows}

The rows of \cref{due:thm:ext-weights}, $\Qq=1/4$, for the orders that
occur through depth $11$:
\begin{center}
\begin{tabular}{@{}cl@{}}
\toprule
$d$ & $(\Wrow{d}{0},\ldots,\Wrow{d}{d})$\\
\midrule
$0$ & $(1)$\\
$1$ & $\bigl(-\tfrac13,\ \tfrac43\bigr)$\\
$2$ & $\bigl(\tfrac1{45},\ -\tfrac49,\ \tfrac{64}{45}\bigr)$\\
$3$ & $\bigl(-\tfrac1{2835},\ \tfrac4{135},\ -\tfrac{64}{135},\ \tfrac{4096}{2835}\bigr)$\\
$4$ & $\bigl(\tfrac1{722925},\ -\tfrac4{8505},\ \tfrac{64}{2025},\ -\tfrac{4096}{8505},\ \tfrac{1048576}{722925}\bigr)$\\
$5$ & $\bigl(-\tfrac1{739552275},\ \tfrac4{2168775},\ -\tfrac{64}{127575},\ \tfrac{4096}{127575},\ -\tfrac{1048576}{2168775},\ \tfrac{1073741824}{739552275}\bigr)$\\
\bottomrule
\end{tabular}
\end{center}
Each row sums to $1$ and annihilates $\Qq^{n},\ldots,\Qq^{dn}$
(\cref{due:lem:ext-moments}); the last entry is
$4^{\binom{d+1}{2}}/P_{d}$ with $P_{d}=\prod_{k\le d}(4^{k}-1)=3,45,2835,
722925,739552275$, and the alternating pattern $4^{\binom{j+1}{2}}$ in the
numerators is \cref{due:eq:ext-integer-product}.  The rows for $d\le3$ are
the Exponents volume's printed weight vectors for $M=d+1$
(\cref{due:sec:extraction}).

\subsection{Symmetry: half the points suffice}
\label{due:sec:exr-symmetry}

\begin{proposition}[Reflection]
\label{due:prop:exr-symmetry}
For $n\ge1$ and $0\le x\le1$,
\begin{equation}
 \upn{n}(x)+\upn{n}(1-x)=1,
 \qquad
 \RvachevUp(x)+\RvachevUp(1-x)=1 .
 \label{due:eq:exr-reflection}
\end{equation}
Consequently, for dyadic $x\in(0,1)$ of depth $s\ge1$: $\dep{1-x}=s$,
$\Cr{m}{1-x}=-\Cr{m}{x}$ for $1\le m\le d$, $\Delta_{s}(1-x)=-\Delta_{s}(x)$,
and $\RvachevUp^{(2m)}(1-x)=-\RvachevUp^{(2m)}(x)$ for $m\ge1$.  Together
with evenness, every quantity at $\pm x$ and $\pm(1-x)$ is determined by
its value at $x$.
\end{proposition}

\begin{proof}
By \cref{due:eq:sharp-finite-fde}, $\upn{n}(x)=\int_{2x-1}^{2x+1}\upn{n-1}
=P_{n-1}(2x+1)-P_{n-1}(2x-1)=1-P_{n-1}(2x-1)$ for $x\ge0$, since
$2x+1\ge1$ is beyond the support; and
$\upn{n}(1-x)=1-P_{n-1}(1-2x)=1-\bigl(1-P_{n-1}(2x-1)\bigr)=P_{n-1}(2x-1)$
by the evenness of $\upn{n-1}$.  The sum is $1$.  The same computation with
Rvachev's equation, or the limit $n\to\infty$, gives the second identity.
The depth is invariant under $x\mapsto1-x$ by
\cref{due:lem:set-depth-representation}.  Differentiating
\cref{due:eq:exr-reflection} $2m$ times gives the derivative statement,
and \cref{due:eq:set-main-law} at $x$ and $1-x$ then gives the modes and
the defects.
\end{proof}

% ed.: S4's "Symmetry" subsection states the reflection for its examples;
% ed.: the proof through the finite functional equation is supplied here.
Thus the pairs $(\tfrac14,\tfrac34)$, $(\tfrac18,\tfrac78)$,
$(\tfrac1{16},\tfrac{15}{16})$, $(\tfrac5{16},\tfrac{11}{16})$,
$(\tfrac1{64},\tfrac{63}{64})$, $(\tfrac1{128},\tfrac{127}{128})$ that the
sources treat separately are one example each.

\subsection{Depth two and three: one mode}
\label{due:sec:exr-one-mode}

\begin{example}[$x=\tfrac14$ and $x=\tfrac34$; $s=2$, $d=1$]
\label{due:ex:exr-quarter}
The point $\tfrac14$ is worked out in \cref{due:ex:set-quarter}:
$\qn{2}=\tfrac{15}{16}$, $\qn{3}=\tfrac{179}{192}$,
$\RvachevUp(\tfrac14)=-\tfrac13\cdot\tfrac{15}{16}+\tfrac43\cdot\tfrac{179}{192}
=\tfrac{67}{72}$, $\Cr{1}{\tfrac14}=\tfrac19$,
$\RvachevUp''(\tfrac14)=-8$.  By reflection, at $x=\tfrac34$:
$\qn{2}=\tfrac1{16}$, $\qn{3}=\tfrac{13}{192}$,
$\RvachevUp(\tfrac34)=\tfrac5{72}$, $\Cr{1}{\tfrac34}=-\tfrac19$,
$\RvachevUp''(\tfrac34)=8$, so that
$\qn{n}=\tfrac5{72}-\tfrac19\,4^{-n}$ for all $n\ge2$ (S4's and S6's
example).  At the knot level $n=1$: $\qn{1}=\upn{1}(\tfrac34)=0$ (the
trapezoid $\upn{1}$ vanishes for $\AbsoluteValue{x}\ge\tfrac34$) while the
law would give $\tfrac5{72}-\tfrac1{36}=\tfrac1{24}$; the defect $-\tfrac1{24}$
is \cref{due:thm:sharp-knot-defect} with $k=3$, $\ThueMorseSign{3}=+1$.
\end{example}

\begin{example}[$x=\tfrac18,\tfrac38,\tfrac78$; $s=3$, $d=1$]
\label{due:ex:exr-eighth}
Three points of depth $3$, all with the single mode ratio $\tfrac14$ and
the row $(-\tfrac13,\tfrac43)$ applied at $N=3$:
\begin{center}
\begin{tabular}{@{}cccccc@{}}
\toprule
$x$ & $\qn{3}$ & $\qn{4}$ & $\RvachevUp(x)$ & $\Cr{1}{x}$ & $\RvachevUp''(x)$\\
\midrule
$\tfrac18$ & $\tfrac{383}{384}$ & $\tfrac{1531}{1536}$ & $\tfrac{287}{288}$ & $\tfrac1{18}$ & $-4$\\[0.2em]
$\tfrac38$ & $\tfrac{287}{384}$ & $\tfrac{1147}{1536}$ & $\tfrac{215}{288}$ & $\tfrac1{18}$ & $-4$\\[0.2em]
$\tfrac78$ & $\tfrac1{384}$ & $\tfrac5{1536}$ & $\tfrac1{288}$ & $-\tfrac1{18}$ & $4$\\
\bottomrule
\end{tabular}
\end{center}
The full jets $(\RvachevUp,\RvachevUp',\RvachevUp'',\RvachevUp''')$ are: at
$\tfrac18$, $(\tfrac{287}{288},-\tfrac5{36},-4,-64)$; at $\tfrac38$,
$(\tfrac{215}{288},-\tfrac{67}{36},-4,64)$; at $\tfrac78$,
$(\tfrac1{288},-\tfrac5{36},4,-64)$; and $\RvachevUp^{(4)}=0$ at all
three, as \cref{due:thm:mech-jet-termination} requires at depth $3$.  The
top derivative $\pm64=\pm2^{\binom42}$ is \cref{due:eq:mech-top-derivative}.
Odd depth: the law already holds at $n=s-1=2$
(\cref{due:prop:mech-odd-onset}): $\qn{2}(\tfrac18)=1$, which is
$\tfrac{287}{288}+\tfrac1{18}\cdot\tfrac1{16}$;
$\qn{2}(\tfrac38)=\tfrac34$, $\qn{2}(\tfrac78)=0$, each equal to the
predicted value.  At $n=s-2=1$ it fails by $\mp\tfrac1{96}$ at all three
points ($\qn{1}=1,\tfrac34,0$ against $\tfrac{97}{96},\tfrac{73}{96},
-\tfrac1{96}$).
\end{example}

\subsection{Depth four and five: two modes}
\label{due:sec:exr-two-modes}

\begin{example}[$x=\tfrac1{16}$; $s=4$, $d=2$]
\label{due:ex:exr-sixteenth}
Three samples at $N=4$ and the row $(\tfrac1{45},-\tfrac49,\tfrac{64}{45})$:
\[
 \qn{4}=\frac{24575}{24576},\qquad
 \qn{5}=\frac{1965959}{1966080},\qquad
 \qn{6}=\frac{94365509}{94371840},
\]
\[
 \RvachevUp(\tfrac1{16})
 =\frac{1}{45}\,\qn{4}-\frac49\,\qn{5}+\frac{64}{45}\,\qn{6}
 =\frac{2073457}{2073600},
\]
and the modes, from the exact Vandermonde solve and, independently, from
the formula $\Cr{m}{x}=(-1)^{m}\am{m}4^{-m}\RvachevUp^{(2m)}(x)$, with the
tiled derivatives $\RvachevUp''(\tfrac1{16})=-\tfrac59$ and
$\RvachevUp^{(4)}(\tfrac1{16})=-1024$:
\[
 \Cr{1}{\tfrac1{16}}=\frac{5}{648},\qquad
 \Cr{2}{\tfrac1{16}}=-\frac{248}{2025},\qquad
 \qn{n}=\frac{2073457}{2073600}+\frac{5}{648}\,4^{-n}-\frac{248}{2025}\,4^{-2n}
 \quad(n\ge4).
\]
In S1's $(n+1)$-convention the same modes read
$A_{1}=\tfrac5{162}$, $A_{2}=-\tfrac{3968}{2025}$, and its output line for
this point, ``extraction row $[1,-20,64]/45$, modes $5/162$, $-3968/2025$'',
is the display above.  The odd jet is $\RvachevUp'(\tfrac1{16})=-\tfrac1{144}$,
$\RvachevUp'''(\tfrac1{16})=-32$; $\RvachevUp^{(5)}=0$.  Even depth: at the
knot level $n=3$, $\qn{3}=\upn{3}(\tfrac1{16})=1$ while the law predicts
$\tfrac{46081}{46080}$; the defect $-\tfrac1{46080}$ is
\cref{due:thm:sharp-knot-defect} with $k=8$, $\ThueMorseSign{8}=-1$,
$\beta_{4}=-\tfrac1{720}$, $2^{-6}$.  The annihilating recurrence
$-\qn{n}+21\qn{n+1}-84\qn{n+2}+64\qn{n+3}=0$ holds for $n\ge4$; its residual
at $n=3$ is $-\tfrac1{46080}$ times the leading coefficient, and at $n=2$
it is nonzero as well ($\qn{2}=1$ against a predicted
$\tfrac{15359}{15360}$).
\end{example}

\begin{example}[$x=\tfrac5{16}$ and $x=\tfrac{15}{16}$; $s=4$, $d=2$]
\label{due:ex:exr-five-sixteenths}
At $x=\tfrac5{16}$: $\qn{4}=\tfrac{6987}{8192}$,
$\qn{5}=\tfrac{1676281}{1966080}$, $\qn{6}=\tfrac{80454331}{94371840}$,
\[
 \RvachevUp(\tfrac5{16})=\frac{1767743}{2073600},\qquad
 \Cr{1}{\tfrac5{16}}=\frac{67}{648},\qquad
 \Cr{2}{\tfrac5{16}}=\frac{248}{2025},
\]
with $\RvachevUp''=-\tfrac{67}9$, $\RvachevUp^{(4)}=1024$,
$\RvachevUp'=-\tfrac{215}{144}$, $\RvachevUp'''=32$.  At the knot level
$n=3$, $\qn{3}=\tfrac{41}{48}=\tfrac{39360}{46080}$ against a predicted
$\tfrac{39359}{46080}$: the defect is $+\tfrac1{46080}$, with $k=10$ and
$\ThueMorseSign{10}=+1$.  By reflection, $x=\tfrac{15}{16}$ (S4's example)
has $\RvachevUp(\tfrac{15}{16})=\tfrac{143}{2073600}$,
$\Cr{1}{\tfrac{15}{16}}=-\tfrac5{648}$, $\Cr{2}{\tfrac{15}{16}}=\tfrac{248}{2025}$,
and $x=\tfrac{11}{16}$ has $\tfrac{305857}{2073600}$, $-\tfrac{67}{648}$,
$-\tfrac{248}{2025}$.  Note that $\tfrac1{16}$ and $\tfrac5{16}$, which are
not reflections of each other, share $\AbsoluteValue{\Cr{2}{x}}=\tfrac{248}{2025}
=\am{2}\,2^{\binom52-4}$, as \cref{due:eq:sharp-top-mode} says every point
of depth $4$ must.
\end{example}

\begin{example}[$x=\tfrac1{32}$ and $x=\tfrac3{32}$; $s=5$, $d=2$]
\label{due:ex:exr-thirty-second}
Odd depth $5$ with two modes.  At $x=\tfrac1{32}$:
\[
 \qn{5}=\frac{3932159}{3932160},\qquad
 \qn{6}=\frac{188743589}{188743680},\qquad
 \qn{7}=\frac{1006632407}{1006632960},
\]
\[
 \RvachevUp(\tfrac1{32})=\frac{33177581}{33177600},\qquad
 \Cr{1}{\tfrac1{32}}=\frac{1}{2592},\qquad
 \Cr{2}{\tfrac1{32}}=-\frac{124}{2025},
\]
with $\RvachevUp''=-\tfrac1{36}$, $\RvachevUp^{(4)}=-512$, the odd jet
$\RvachevUp'=-\tfrac{143}{1036800}$, $\RvachevUp'''=-\tfrac{40}9$,
$\RvachevUp^{(5)}=-32768=-2^{\binom62}$, and $\RvachevUp^{(6)}=0$.  The law
holds already at $n=4$: $\qn{4}=1$ equals the predicted value exactly
(odd onset), and fails at $n=3$ by $\tfrac7{737280}$.  At $x=\tfrac3{32}$:
$\qn{5}=\tfrac{3929279}{3932160}$, $\qn{6}=\tfrac{188601509}{188743680}$,
$\qn{7}=\tfrac{1005869527}{1006632960}$,
$\RvachevUp(\tfrac3{32})=\tfrac{33152381}{33177600}$,
$\Cr{1}{\tfrac3{32}}=\tfrac{73}{2592}$, $\Cr{2}{\tfrac3{32}}=-\tfrac{124}{2025}$,
and again $\qn{4}=\tfrac{1535}{1536}$ is exactly the value the law predicts.
The top mode $\AbsoluteValue{\Cr{2}{x}}=\tfrac{124}{2025}=\am{2}2^{\binom52-5}$
is the odd-depth line of \cref{due:eq:sharp-top-mode}.
\end{example}

\subsection{Depth six and seven: three modes}
\label{due:sec:exr-three-modes}

\begin{example}[$x=\tfrac1{64}$; $s=6$, $d=3$]
\label{due:ex:exr-sixty-fourth}
Four samples at $N=6$ and the row
$(-\tfrac1{2835},\tfrac4{135},-\tfrac{64}{135},\tfrac{4096}{2835})$:
\[
 \qn{6}=\frac{1509949439}{1509949440},\quad
 \qn{7}=\frac{676457348027}{676457349120},\quad
 \qn{8}=\frac{43293270259811}{43293270343680},\quad
 \qn{9}=\frac{1662461577834037}{1662461581197312},
\]
\[
 \RvachevUp(\tfrac1{64})=\frac{561842748287}{561842749440},\qquad
 \Cr{1}{\tfrac1{64}}=\frac{143}{18662400},\qquad
 \Cr{2}{\tfrac1{64}}=-\frac{31}{3645},\qquad
 \Cr{3}{\tfrac1{64}}=\frac{4806656}{2679075},
\]
with $\RvachevUp''=-\tfrac{143}{259200}$, $\RvachevUp^{(4)}=-\tfrac{640}9$,
$\RvachevUp^{(6)}=-2097152=-2^{\binom72}$, and the odd jet
$\RvachevUp'=-\tfrac{19}{16588800}$, $\RvachevUp'''=-\tfrac29$,
$\RvachevUp^{(5)}=-16384$.  Even depth: $\qn{5}=1$ against a predicted
$\tfrac{990904319}{990904320}$, a defect of $+\tfrac1{990904320}
=2^{-15}\cdot\tfrac1{30240}$ with $\ThueMorseSign{k}=-1$ ($k=32$).  By
reflection, $x=\tfrac{63}{64}$ has $\RvachevUp=\tfrac{1153}{561842749440}$
and the negated modes, and its knot-level sample $\qn{5}=0$ falls short of
the prediction by the same amount.
\end{example}

\begin{example}[$x=\tfrac1{128}$ and $x=\tfrac3{128}$; $s=7$, $d=3$]
\label{due:ex:exr-one-twenty-eighth}
At $x=\tfrac1{128}$, four samples at $N=7$:
\[
 \qn{7}=\frac{1352914698239}{1352914698240},\qquad
 \qn{8}=\frac{17317308137431}{17317308137472},
\]
\[
 \qn{9}=\frac{3324923162384629}{3324923162394624},\qquad
 \qn{10}=\frac{212795082392578741}{212795082393255936},
\]
\[
 \RvachevUp(\tfrac1{128})=\frac{179789679820217}{179789679820800},\qquad
 \Cr{1}{\tfrac1{128}}=\frac{19}{298598400},
\]
\[
 \Cr{2}{\tfrac1{128}}=-\frac{31}{72900},\qquad
 \Cr{3}{\tfrac1{128}}=\frac{2403328}{2679075},
\]
with $\RvachevUp''=-\tfrac{19}{4147200}$, $\RvachevUp^{(4)}=-\tfrac{32}9$,
$\RvachevUp^{(6)}=-1048576=-2^{\binom72-1}$ (the odd-depth top even
derivative of \cref{due:eq:mech-subtop-derivative}), and
$\RvachevUp^{(7)}=-268435456=-2^{\binom82}$.  Odd onset: $\qn{6}=1$ equals
the predicted value; at $n=5$ the law fails by $\tfrac{31}{63417876480}$.
At $x=\tfrac3{128}$: $\qn{7}=\tfrac{1352914622303}{1352914698240}$,
$\qn{8}=\tfrac{86586535502003}{86586540687360}$,
$\qn{9}=\tfrac{3324922960085749}{3324923162394624}$,
$\qn{10}=\tfrac{1063975346970795401}{1063975411966279680}$,
$\RvachevUp(\tfrac3{128})=\tfrac{179789668823441}{179789679820800}$,
$\Cr{1}{\tfrac3{128}}=\tfrac{25219}{298598400}$,
$\Cr{2}{\tfrac3{128}}=-\tfrac{2263}{72900}$,
$\Cr{3}{\tfrac3{128}}=\tfrac{2403328}{2679075}$, and $\qn{6}=\tfrac{23592959}{23592960}$
is exactly the predicted value.
\end{example}

\subsection{A table of exact laws}
\label{due:sec:exr-table}

\begin{center}
\small
\begin{longtable}{@{}c c c l l@{}}
\caption{Exact laws $\qn{n}=\RvachevUp(x)+\sum_{m=1}^{d}\Cr{m}{x}4^{-mn}$, valid for all $n\ge s$ (and for $n\ge s-1$ at odd $s$).  Every entry is an exact rational, recomputed for this volume.}
\label{due:tab:exr-laws}\\
\toprule
$x$ & $s$ & $d$ & $\RvachevUp(x)$ & $\Cr{1}{x},\ldots,\Cr{d}{x}$\\
\midrule
\endfirsthead
\toprule
$x$ & $s$ & $d$ & $\RvachevUp(x)$ & $\Cr{1}{x},\ldots,\Cr{d}{x}$\\
\midrule
\endhead
$0$ & $0$ & $0$ & $1$ & ---\\
$\pm1$ & $0$ & $0$ & $0$ & ---\\
$\tfrac12$ & $1$ & $0$ & $\tfrac12$ & ---\\
$\tfrac14$ & $2$ & $1$ & $\tfrac{67}{72}$ & $\tfrac19$\\
$\tfrac34$ & $2$ & $1$ & $\tfrac5{72}$ & $-\tfrac19$\\
$\tfrac18$ & $3$ & $1$ & $\tfrac{287}{288}$ & $\tfrac1{18}$\\
$\tfrac38$ & $3$ & $1$ & $\tfrac{215}{288}$ & $\tfrac1{18}$\\
$\tfrac78$ & $3$ & $1$ & $\tfrac1{288}$ & $-\tfrac1{18}$\\
$\tfrac1{16}$ & $4$ & $2$ & $\tfrac{2073457}{2073600}$ & $\tfrac5{648},\ -\tfrac{248}{2025}$\\
$\tfrac3{16}$ & $4$ & $2$ & $\tfrac{2026943}{2073600}$ & $\tfrac{67}{648},\ \tfrac{248}{2025}$\\
$\tfrac5{16}$ & $4$ & $2$ & $\tfrac{1767743}{2073600}$ & $\tfrac{67}{648},\ \tfrac{248}{2025}$\\
$\tfrac{11}{16}$ & $4$ & $2$ & $\tfrac{305857}{2073600}$ & $-\tfrac{67}{648},\ -\tfrac{248}{2025}$\\
$\tfrac{15}{16}$ & $4$ & $2$ & $\tfrac{143}{2073600}$ & $-\tfrac5{648},\ \tfrac{248}{2025}$\\
$\tfrac1{32}$ & $5$ & $2$ & $\tfrac{33177581}{33177600}$ & $\tfrac1{2592},\ -\tfrac{124}{2025}$\\
$\tfrac3{32}$ & $5$ & $2$ & $\tfrac{33152381}{33177600}$ & $\tfrac{73}{2592},\ -\tfrac{124}{2025}$\\
$\tfrac5{32}$ & $5$ & $2$ & $\tfrac{32842819}{33177600}$ & $\tfrac{215}{2592},\ \tfrac{124}{2025}$\\
$\tfrac1{64}$ & $6$ & $3$ & $\tfrac{561842748287}{561842749440}$ & $\tfrac{143}{18662400},\ -\tfrac{31}{3645},\ \tfrac{4806656}{2679075}$\\
$\tfrac{37}{64}$ & $6$ & $3$ & $\tfrac{966420578371}{2809213747200}$ & $-\tfrac{305857}{18662400},\ \tfrac{2077}{18225},\ \tfrac{4806656}{2679075}$\\
$\tfrac1{128}$ & $7$ & $3$ & $\tfrac{179789679820217}{179789679820800}$ & $\tfrac{19}{298598400},\ -\tfrac{31}{72900},\ \tfrac{2403328}{2679075}$\\
$\tfrac3{128}$ & $7$ & $3$ & $\tfrac{179789668823441}{179789679820800}$ & $\tfrac{25219}{298598400},\ -\tfrac{2263}{72900},\ \tfrac{2403328}{2679075}$\\
\bottomrule
\end{longtable}
\end{center}

Two regularities are visible and both are theorems: within one depth the
top coefficient $\Cr{d}{x}$ has a constant absolute value
(\cref{due:eq:sharp-top-mode}), and reflection $x\mapsto1-x$ negates every
mode (\cref{due:prop:exr-symmetry}).  The denominators of
$\RvachevUp(x)$ are $72=2^{3}\cdot9$, $288=2^{5}\cdot9$,
$2073600=2^{10}\cdot2025$, $33177600=2^{14}\cdot2025$,
$561842749440=2^{21}\cdot267915\cdot\ldots$; their growth is the subject of
\cref{due:prop:stab-denominator}.
% ed.: 561842749440 = 2^21 * 3^5 * 5 * 7^2 * ... is not worth printing in
% ed.: full; the prose stops at the power of two.

\begin{summarybox}[title={Section summary}]
Twenty dyadic points through depth $7$, every number exact and
independently recomputed: the samples, the extraction, the value, every
mode, and the jet.  The examples exhibit the law holding from the onset,
holding one level early at odd depth (\cref{due:ex:exr-eighth,due:ex:exr-thirty-second,due:ex:exr-one-twenty-eighth}),
failing at the knot level at even depth by exactly the Bernoulli amount of
\cref{due:thm:sharp-knot-defect}
(\cref{due:ex:exr-quarter,due:ex:exr-sixteenth,due:ex:exr-five-sixteenths,due:ex:exr-sixty-fourth}),
and the reflection symmetry that halves the work
(\cref{due:prop:exr-symmetry}).
\end{summarybox}

% ==========================================================================
\section{An exact algorithm and its verification}
\label{due:sec:algorithm}

Each source ends with a procedure and a program.  The procedures are the
same procedure and the programs check the same identities; this section
states the procedure once, with its preconditions and postconditions, and
specifies what the single verifier retained with this volume checks.

\subsection{The algorithm}
\label{due:sec:alg-algorithm}

\begin{algorithm}[Exact dyadic extraction]
\label{due:alg:exr-extraction}
\emph{Input.}  A dyadic rational $x\in[-1,1]$, given either in lowest terms
or as a pair $(m,n)$ with $x=m/2^{n}$; the arithmetic is that of
$\RationalNumbers$ with unbounded integers.

\emph{Output.}  The exact rational $\RvachevUp(x)$; optionally every mode
$\Cr{1}{x},\ldots,\Cr{d}{x}$, every even derivative
$\RvachevUp^{(2)}(x),\ldots,\RvachevUp^{(2d)}(x)$, and a certificate.

\begin{enumerate}[label=\textup{\arabic*.}]
\item\label{due:it:alg-depth}
 (\emph{Depth.})  Compute $s=\dep{x}=\max(0,\,n-\TwoAdicValuation(m))$
 (\cref{due:lem:set-depth-representation}; for $x=0$ take $s=0$) and
 $d=\Floor{s/2}$.  If $\AbsoluteValue{x}=1$ return $0$; if $x=0$ return $1$
 (\cref{due:ex:set-low-depth}).
\item\label{due:it:alg-level}
 (\emph{Window.})  Choose $N=s$; if $s\ge3$ is odd, $N=s-1$ is also
 admissible (\cref{due:prop:mech-odd-onset}), and any $N$ above these is
 admissible but wasteful.
\item\label{due:it:alg-samples}
 (\emph{Samples.})  For $j=0,\ldots,d$ compute $\qn{N+j}=\upn{N+j}(x)$
 exactly from the scaled cell formula \cref{due:eq:set-scaled-cell}: with
 $y=2^{n}(x+1)-\tfrac12$ and $n=N+j$,
 \[
  \qn{n}=\frac{2^{-\binom n2}}{n!}\sum_{k=0}^{\Floor y}\ThueMorseSign{k}\,(y-k)^{n},
 \]
 where, since $n\ge s$, $2y$ is an odd integer and the sum involves only
 integer powers of half-integers, so that $2^{\binom{n+1}{2}}n!\,\qn{n}$ is
 an integer (\cref{due:cor:set-scaled-cell}).
\item\label{due:it:alg-weights}
 (\emph{Row.})  Compute $\Wrow{d}{j}$, $0\le j\le d$, from
 \cref{due:eq:ext-integer-product},
 $\Wrow{d}{j}=(-1)^{d-j}4^{\binom{j+1}{2}}/(P_{j}P_{d-j})$ with
 $P_{i}=\prod_{k\le i}(4^{k}-1)$ accumulated once, or from the
 $q$-Pochhammer form \cref{due:eq:ext-extraction}; the two agree
 (\cref{due:thm:ext-weights}).
\item\label{due:it:alg-value}
 (\emph{Value.})  Return $\RvachevUp(x)=\sum_{j=0}^{d}\Wrow{d}{j}\,\qn{N+j}$
 (\cref{due:thm:ext-extraction}); equivalently run the triangular
 tableau of \cref{due:thm:rec-tableau}, whose last entry is the same
 rational.
\item\label{due:it:alg-modes}
 (\emph{Modes and jet, optional.})  Compute
 $\Cr{m}{x}=4^{mN}\sum_{j}\qn{N+j}(-1)^{d-m}e^{(j)}_{d-m}/D_{j}$
 from \cref{due:eq:rec-modes}, or solve the $(d+1)\times(d+1)$ geometric
 Vandermonde system \cref{due:eq:rec-vandermonde} exactly; then
 $\RvachevUp^{(2m)}(x)=(-1)^{m}4^{m}\Cr{m}{x}/\am{m}$ (\cref{due:cor:rec-jet}),
 with $\am{m}$ from the recurrence \cref{due:eq:mech-recurrence}.
\item\label{due:it:alg-certify}
 (\emph{Certificate, optional.})  Any of: (a) repeat step
 \ref{due:it:alg-value} from the window $N+1$ and compare; (b) compute one
 further sample $\qn{N+d+1}$ and compare with the law; (c) evaluate the
 residual of the integer recurrence \cref{due:eq:rec-recurrence-int} at
 $n=N$.  Each must give exact equality
 (\cref{due:prop:rec-diagnostics}\ref{due:it:diag-exact}); a nonzero
 residual proves an indexing or arithmetic error, or that $N$ is below the
 onset.
\end{enumerate}

\emph{Postcondition.}  The returned rational is $\RvachevUp(x)$, with no
limit, tolerance or rational reconstruction; the modes and the jet are
exact; and, for $s\ge2$, every returned mode is nonzero
(\cref{due:thm:sharp-modes}).
\end{algorithm}

\begin{proof}[Correctness]
Steps \ref{due:it:alg-depth}--\ref{due:it:alg-level} are
\cref{due:lem:set-depth-representation,due:thm:set-main,due:prop:mech-odd-onset};
step \ref{due:it:alg-samples} is \cref{due:cor:set-scaled-cell}; steps
\ref{due:it:alg-weights}--\ref{due:it:alg-value} are
\cref{due:thm:ext-weights,due:thm:ext-extraction,due:thm:rec-tableau};
step \ref{due:it:alg-modes} is \cref{due:thm:rec-inverse,due:cor:rec-jet};
step \ref{due:it:alg-certify} is \cref{due:prop:rec-diagnostics}.
\end{proof}

% ed.: the six pseudocode listings (S1, S3, S4 in Python; S2, S5, S6 in
% ed.: prose) are the same seven lines; one listing is kept.
\begin{lstlisting}[style=pseudo,caption={The core of \cref{due:alg:exr-extraction} in exact rational arithmetic; \texttt{up\_n(x, n)} is the sample $\upn{n}(x)$ of step 3 and \texttt{P} the products $P_{i}$.},label={due:lst:exr-core}]
s = depth(x); d = s // 2; N = s
samples = [up_n(x, N + j) for j in range(d + 1)]
P = [1]
for k in range(1, d + 1): P.append(P[-1] * (4**k - 1))
weights = [(-1)**(d-j) * Fraction(4**((j*(j+1))//2), P[j]*P[d-j])
           for j in range(d + 1)]
value = sum(w * q for w, q in zip(weights, samples))   # = up(x), exactly
\end{lstlisting}

\subsection{Cost}
\label{due:sec:alg-cost}

The extraction itself is negligible: $O(d)$ rational operations for the
value once the $P_{i}$ are cached, $O(d^{2})$ for the tableau or the closed
inverse, $O(d^{3})$ for a dense solve, with $d=\Floor{s/2}$.  The cost is
in the samples: the truncated-power sum for $\qn{n}$ has up to $2^{n+1}$
terms, each a power of a half-integer of about $n\log_{2}(2^{n})=n^{2}$
bits, so the reference evaluator is exponential in the depth and is
intended for transparent verification at moderate depth (the exhaustive
runs below reach depth $7$ or $8$ in seconds).  Faster exact evaluators for
the samples---binary block recurrences, repeated integration of the
piecewise polynomial, or the dyadic arithmetic of \cite{arias2018} and of
the Lean corpus \cite{proveit-primary}---can be substituted without
changing any theorem, since the extraction sees only the $d+1$ rationals.

In floating point the picture reverses: the samples are the hazard (an
alternating Thue--Morse sum, \cref{due:rmk:stab-float}) and the row is
benign, amplifying uniform absolute error by less than $1.97$
(\cref{due:prop:stab-lebesgue}); if a denominator bound is known
(\cref{due:prop:stab-denominator}), rational reconstruction from a
high-precision result is possible but never necessary, since exact input
is available.

\subsection{The verification harness}
\label{due:sec:alg-verifier}

Each of the six packages shipped a Python program using only the standard
library and exact fractions (\texttt{fractions.Fraction}), together with its
captured output.
Consolidated, they checked the following, all as equalities of reduced
fractions and never to a tolerance:
\begin{itemize}
\item S1: the $q$-Pochhammer and Gaussian-binomial forms of the row agree;
 interpolation of the mode polynomial and extraction of all modes;
 Bernoulli numbers, cumulants, the coefficients $\gamma^{\mathrm{Rv}}_{2m}$
 and the dyadic derivatives by the Thue--Morse stencil; interpolated and
 differential-operator coefficients coincide; the expansion and the
 recurrence have zero residual;
\item S2: the truncated-power formula, the weights, an exact Gauss--Jordan
 solve, prediction of unused levels, the recurrence; exhaustively at all
 $129$ signed dyadic points of depth at most $6$;
\item S3: the evaluator, the weights, the Bell--Bernoulli recurrence for
 $\am{m}$, exact values and derivatives; $2376$ identities at the $255$
 interior points of depth at most $7$, including sliding-window
 invariance and the derivative formula at several consecutive levels;
\item S4: extraction from four windows agrees; fitted coefficients
 reproduce six further levels; fitted coefficients equal the analytic
 $(-1)^{m}\am{m}4^{-m}\RvachevUp^{(2m)}(x)$ with tiled derivatives; the
 recurrence has zero residual at six consecutive levels; $2570$
 sequence and window identities, $941$ coefficient entries and $1542$
 recurrence identities at all $257$ points of depth at most $7$;
\item S5: the evaluator, the weights, the triangular table, an exact solve,
 overlapping-block and out-of-sample checks at $129$ points of depth at
 most $6$;
\item S6: closed-row and recursive extraction, an exact solve, the reciprocal
 moments from Bernoulli numbers, assertions for two worked examples, and an
 optional exhaustive test through a chosen depth.
\end{itemize}

The volume retains one verifier,
\texttt{verify\_\allowbreak dyadic\_\allowbreak up\_\allowbreak extraction.\allowbreak py}, adapted from S4's (the most
exhaustive) and extended by the claims that are specific to this volume.
Its specification:
\begin{enumerate}[label=\textup{(V\arabic*)}]
\item it evaluates the finite splines by the Thue--Morse truncated-power
 formula in the $(n+1)$-factor indexing, computes the quarter-base row in
 both closed forms, solves the geometric Vandermonde system exactly, and
 evaluates the tiled derivatives;
\item at every reduced dyadic point of depth at most $7$ ($257$ points,
 including $0,\pm1$) it checks the law from four windows, six further
 levels, the analytic-versus-fitted coefficients, and the recurrence
 residuals: $2570$, $941$ and $1542$ identities;
\item it checks $\am{0},\ldots,\am{5}$ against the printed table
 (\cref{due:thm:mech-am-rational-positive});
\item at the $252$ interior points of depth $2$ to $7$ it checks that all
 $684$ mode coefficients are nonzero (\cref{due:thm:sharp-modes}), that the
 reversed row \cref{due:eq:ext-reversed} equals the forward row, and that
 the law holds at $n=s-1$ at all $168$ points of odd depth
 (\cref{due:prop:mech-odd-onset});
\item at the same $252$ points it checks the closed form of the knot-level
 defect, \cref{due:thm:sharp-knot-defect}: the difference between
 $\upn{s-1}(x)$ and the law's prediction equals
 $-\ThueMorseSign{k}2^{-\binom s2}\beta_{s}$, which is nonzero at all $84$
 points of even depth and zero at the $168$ of odd depth;
\item it checks $\sum_{j}\AbsoluteValue{\Wrow{d}{j}}=\QPochhammer{-\Qq}{\Qq}{d}/\QPochhammer{\Qq}{\Qq}{d}$
 for $d\le6$ (\cref{due:prop:stab-lebesgue});
\item it writes nothing unless asked (\texttt{--out-dir}), runs in about
 twelve seconds, and accepts \texttt{--max-depth} for deeper runs.
\end{enumerate}
The knot-defect formula was additionally checked through depth $8$ ($508$
interior points) with a separate script during the consolidation.  All
these computations are evidence that the formulas are transcribed and
indexed correctly; the proofs do not depend on them.

\begin{summarybox}[title={Section summary}]
One algorithm (\cref{due:alg:exr-extraction}): depth, window, $d+1$ exact
samples, one row, one dot product; optionally all modes and the jet, and a
certificate that is a finite exact test.  The samples are the only
expensive part.  One retained verifier checks, at every reduced dyadic
point of depth at most $7$, the law, the coefficients, the recurrence, the
nonvanishing of all modes, both forms of the row, the odd-depth onset, the
Bernoulli formula for the knot-level defect, and the Lebesgue factor.
\end{summarybox}

% ==========================================================================
\section{Consequences and extensions}
\label{due:sec:consequences}

\subsection{What is universal and what is dyadic}
\label{due:sec:cons-universal}

Two logically distinct structures meet in \cref{due:thm:set-main}.  The
algebra---the finite deconvolution by a reciprocal power series, the
geometric Lagrange row, its moment identities, the annihilating
recurrence---is universal: it holds for any sequence of the form
$L+\sum_{r\le d}A_{r}\Qq^{rn}$ and for any base.  The exactness is
arithmetic: it needs the point to become the centre of a polynomial cell
and the jet to terminate, and both are properties of dyadic points.  The
following proposition isolates the universal part.

% ed.: S1's "reusable local principle", S2's "general local-deconvolution
% ed.: principle" and S4's "general geometric scales" are the same
% ed.: statement; one proposition with a proof.
\begin{proposition}[Local deconvolution principle]
\label{due:prop:cons-principle}
Let $K$ be an even probability density on $\RealNumbers$ with compact
support, all of whose moments $\mu_{2j}=\int u^{2j}K$ exist (automatic),
and let $B_{K}(z)=\sum_{j}\mu_{2j}z^{2j}/(2j)!$ be its moment generating
function, $1/B_{K}(z)=\sum_{j}b_{j}z^{2j}$ its formal reciprocal.  Let
$f$ be a function, $x$ a point, $\delta>0$, and $p$ a polynomial such that
\begin{enumerate}[label=\textup{(H\arabic*)}]
\item $f=g\Convolution K_{\delta}$ with $K_{\delta}=\delta^{-1}K(\cdot/\delta)$,
 for some locally integrable $g$;
\item $g=p$ on $[x-\delta',x+\delta']$ for some $\delta'>\delta$ with
 $\operatorname{supp}K_{\delta}\subseteq[-\delta',\delta']$;
\item the jet of $f$ at $x$ vanishes above order $R$: $f^{(r)}(x)=0$ for
 $r>R$.
\end{enumerate}
Then
\begin{equation}
 p(x)=\sum_{j=0}^{\Floor{R/2}}b_{j}\,\delta^{2j}\,f^{(2j)}(x),
 \label{due:eq:cons-principle}
\end{equation}
a finite identity.  If moreover $\delta=\delta_{n}=c\alpha^{n}$ along a
sequence of pairs $(g_{n},p_{n})$ satisfying (H1)--(H2), then
$n\mapsto p_{n}(x)$ is $f(x)$ plus a finite sum of geometric modes
$\alpha^{2jn}$, $1\le j\le\Floor{R/2}$, and the geometric Lagrange row at
base $q=\alpha^{2}$ of order $\Floor{R/2}$ returns $f(x)$ exactly from
$\Floor{R/2}+1$ consecutive values of $p_{n}(x)$.
\end{proposition}

\begin{proof}
On polynomials, convolution with $K_{\delta}$ is the finite operator
$B_{K}(\delta D)$, $D=\Differential/\Differential t$: for a monomial,
$(t^{k}\Convolution K_{\delta})(x)=\int(x-y)^{k}K_{\delta}(y)\,\Differential y
=\sum_{j}\binom{k}{2j}\mu_{2j}\delta^{2j}x^{k-2j}$ by evenness, which is
$B_{K}(\delta D)t^{k}$ evaluated at $x$.  By (H1)--(H2), $f$ and all its
derivatives at $x$ are those of $p\Convolution K_{\delta}=B_{K}(\delta D)p$:
$f^{(r)}(x)=(B_{K}(\delta D)p)^{(r)}(x)$ for every $r$.  The operator
$B_{K}(\delta D)$ is invertible on polynomials with inverse
$(1/B_{K})(\delta D)=\sum_{j}b_{j}\delta^{2j}D^{2j}$, a finite sum on any
polynomial, so $p(x)=\sum_{j}b_{j}\delta^{2j}(B_{K}(\delta D)p)^{(2j)}(x)
=\sum_{j}b_{j}\delta^{2j}f^{(2j)}(x)$, the sum over $2j\le\deg p$; by (H3)
the terms with $2j>R$ vanish, and if $\deg p<R$ the missing terms vanish on
both sides.  The geometric statement is
\cref{due:eq:cons-principle} with $\delta_{n}^{2j}=c^{2j}\alpha^{2jn}$, and
the extraction is the moment identity of the row
(\cref{due:thm:ext-weights,due:lem:ext-moments} hold for every base
$q\in(0,1)$, indeed for every $q$ with $\QPochhammer{q}{q}{d}\ne0$).
\end{proof}

The Rvachev case is the cleanest instance: $K=\RvachevUp$ itself (the
tail is a scaled copy of the whole density), $\delta_{n}=2^{-(n+1)}$ so
$\alpha=\tfrac12$ and $q=\tfrac14$, $g=\upn{n}$ with the knot lattice
aligned to the dyadic denominators, and the jet terminates at the order
$\dep{x}$.  For a general base $b>1$ of sinc products the same principle
gives $q=b^{-2}$; what would be missing is a terminating jet at the natural
lattice, which is a property of the base-$2$ Rvachev function
(\cref{due:sec:jet}) and is not known for other bases.

\subsection{Other windows, strides and functionals}
\label{due:sec:cons-windows}

\begin{proposition}[Arbitrary levels]
\label{due:prop:cons-levels}
Let $n_{0}<n_{1}<\cdots<n_{d}$ be any $d+1$ levels with $n_{0}\ge s$.
Then
\begin{equation}
 \RvachevUp(x)=\sum_{j=0}^{d}\qn{n_{j}}\prod_{\substack{0\le l\le d\\ l\ne j}}
 \frac{-\Qq^{n_{l}}}{\Qq^{n_{j}}-\Qq^{n_{l}}} .
 \label{due:eq:cons-levels}
\end{equation}
In particular, for samples at a fixed stride $t\ge1$, $n_{j}=N+jt$, the row
is \cref{due:eq:ext-extraction} with $\Qq^{t}$ in place of $\Qq$.
\end{proposition}

\begin{proof}
By \cref{due:thm:set-main} the samples are the values at the distinct nodes
$\Qq^{n_{j}}$ of the polynomial $L+\sum_{r}A_{r}z^{r}$ of degree at most
$d$, whose constant term is $\RvachevUp(x)$; \cref{due:eq:cons-levels} is
Lagrange interpolation at those nodes evaluated at $z=0$.  For
$n_{j}=N+jt$ the common factor $\Qq^{N}$ cancels and the nodes are
$1,\Qq^{t},\ldots,\Qq^{dt}$.
\end{proof}

The consecutive window is preferable only because the products collapse
into the $q$-Pochhammer form and the Lebesgue factor is then smallest for
the given $d$.  More generally, any linear functional of the local jet of
$\upn{n}$ at $x$---a derivative, a finite combination of values and
derivatives at points of the same cell---inherits a finite geometric law
by applying the functional to \cref{due:eq:mech-local-deconvolution}
(\cref{due:cor:main-derivative} is the case of one derivative), and the
same rows extract its limit.

\subsection{The distribution-function version}
\label{due:sec:cons-cdf}

The Lean corpus works with distribution functions (Fabius's function) rather
than densities; the law transports.

\begin{proposition}[CDF version]
\label{due:prop:cons-cdf}
Let $P_{n}(x)=\int_{-\infty}^{x}\upn{n}$ and $P(x)=\int_{-\infty}^{x}\RvachevUp$.
For $x\in[-1,1]$ and $n\ge0$,
\begin{equation}
 P_{n}(x)=1-\upn{n+1}\Bigl(\frac{x+1}{2}\Bigr),
 \qquad
 P(x)=1-\RvachevUp\Bigl(\frac{x+1}{2}\Bigr).
 \label{due:eq:cons-cdf-transport}
\end{equation}
Hence, if $x$ is dyadic of depth $s$, the point $y=(x+1)/2$ is dyadic of
depth $s+1$ (depth $1$ if $x=0$), and for every $n+1\ge\dep{y}$,
\begin{equation}
 P_{n}(x)=P(x)-\sum_{m=1}^{\Floor{\dep{y}/2}}\Cr{m}{y}\,4^{-m(n+1)}:
 \label{due:eq:cons-cdf-law}
\end{equation}
the distribution-function sequence at a dyadic point of depth $s$ has
exactly $\Floor{(s+1)/2}$ modes, all nonzero for $s\ge1$, holds from
$n\ge s$ on (from $n\ge s-1$ on when $s$ is even, $s\ge2$), and is
extracted by the row of order $\Floor{(s+1)/2}$.
\end{proposition}

\begin{proof}
\Cref{due:eq:sharp-finite-fde} at the point $y=(x+1)/2\in[0,1]$ gives
$\upn{n+1}(y)=P_{n}(2y+1)-P_{n}(2y-1)=P_{n}(x+2)-P_{n}(x)=1-P_{n}(x)$,
since $x+2\ge1$; the same for $P$.  The depth of $y=(x+1)/2$ is $s+1$ for
$s\ge1$ by \cref{due:lem:set-depth-representation}, and $\dep{\tfrac12}=1$.
Then \cref{due:thm:set-main} at $y$, with its onset $n+1\ge\dep{y}$ and its
odd-depth improvement (\cref{due:cor:sharp-onset}), transports through
\cref{due:eq:cons-cdf-transport}; the number and nonvanishing of the modes
are \cref{due:thm:sharp-modes} at $y$.
\end{proof}

At $x=\tfrac12$ this is \cref{due:prop:sharp-cdf}: $y=\tfrac34$, one mode,
$P_{n}(\tfrac12)-P(\tfrac12)=-\Cr{1}{\tfrac34}4^{-(n+1)}=\tfrac19 4^{-(n+1)}$.
The Lean statement \texttt{Fabius.\allowbreak fabiusUniformSpline\_\allowbreak quarter\_\allowbreak eq\_\allowbreak quadratic}
(\cref{due:sec:register}) is this one-mode CDF law at the quarter point of
the Fabius interval, which is $y=\tfrac34$ in the present coordinates
(\cref{due:sec:set-bridge}).

\subsection{Non-dyadic points}
\label{due:sec:cons-nondyadic}

At a non-dyadic point nothing of the exact structure survives
(\cref{due:rmk:sharp-nondyadic}): the tail window always contains a knot,
and the jet does not terminate.  The Exponents volume's all-orders
expansion and its sharp $C^{r}$ law are the statements that hold there,
and the signed Richardson combination accelerates with the asymptotic
remainder of its \texttt{p3:eq:richardson-\allowbreak leading}.  The contrast is not a
matter of degree: at a dyadic point the remainder is identically zero from
a computable level on; at every other point of $(-1,1)$ it is nonzero for
infinitely many $n$, because a sequence $\qn{n}$ that eventually satisfied
a finite geometric law would make the tail generating function rational
and the jet at $x$ finite by \cref{due:prop:main-generating} read
backwards through \cref{due:eq:mech-local-deconvolution}---which requires
$x$ to be a cell centre for all large $n$, hence dyadic.
% ed.: the last clause is the converse direction; the sources assert it
% ed.: without argument.  The argument: a cell centre of up_n for all
% ed.: large n is an odd multiple of 2^{-(n+1)} for all large n, hence
% ed.: dyadic.  It is given here in one sentence and not as a theorem,
% ed.: because the step "finite law implies cell centre" is only sketched.

\subsection{Open problems}
\label{due:sec:cons-open}

\begin{enumerate}[label=\textup{(O\arabic*)}]
\item\label{due:it:open-two-knots}
 \emph{The levels below the knot level.}  \Cref{due:thm:sharp-knot-defect}
 gives the defect at $n=s-1$ in closed form.  At $n=s-2$ the defect is
 observed to be independent of the point within a depth
 (\cref{due:rmk:sharp-two-knots}); a closed form, presumably in terms of
 partial moments of $\RvachevUp$ over $[-1,\pm\tfrac12]$, would prove that
 the onset at odd depth is exactly $s-1$ for all $s$, which is now known
 through depth $7$.
\item\label{due:it:open-lean}
 \emph{Formalization.}  The general theorem is absent from the Lean corpus;
 the register in \cref{due:sec:register} lists what exists and what a
 proof would have to build.
\item\label{due:it:open-fast}
 \emph{Fast exact samples.}  The extraction is linear in $d$; the
 reference evaluator for the samples is exponential in $n$.  An exact
 evaluator polynomial in $n$ for $\upn{n}$ at dyadic points, in the spirit
 of the dyadic arithmetic of \cite{arias2018}, would make the method
 practical at large depth.
\item\label{due:it:open-bases}
 \emph{Other kernels.}  \Cref{due:prop:cons-principle} applies to any
 self-similar even kernel with a terminating jet on a lattice.  Which
 kernels beyond $\RvachevUp$ (other Rvachev-type solutions of
 functional-differential equations, other bases) have this property is
 not known to the authors of the six sources or to this volume.
\item\label{due:it:open-extremizers}
 \emph{Extremizers of the sharp law.}  \Cref{due:prop:sharp-attained}
 shows that the sup norm in the Exponents volume's sharp $C^{r}$ law is
 attained at every dyadic point of depth $r+2$.  Whether these are all the
 extremizers is not settled here.
\end{enumerate}

\begin{summarybox}[title={Section summary}]
The algebra of the method is universal (\cref{due:prop:cons-principle});
its exactness is arithmetic and belongs to dyadic points.  Arbitrary
levels and strides, derivatives and other jet functionals inherit the
finite law (\cref{due:prop:cons-levels}); the distribution-function
version has $\Floor{(s+1)/2}$ modes and is the form the Lean corpus uses
(\cref{due:prop:cons-cdf}); away from dyadic points only the asymptotic
statements of the Exponents volume remain.  Five open problems are
recorded.
\end{summarybox}

% ==========================================================================
\section{Relation to the corpus and formalization register}
\label{due:sec:register}

\subsection{The Exponents volume}
\label{due:sec:reg-exponents}

Part~III of the Exponents volume \cite{proveit-exponents} studies the same
finite splines at a general point; its objects and this volume's are
related by the bridge $\upn{n}=p_{n+1}$, $\Phin{n}=\Phi_{n+1}$,
$\delta_{n}^{\text{here}}=\delta_{n+1}^{\text{there}}$ of
\cref{due:eq:set-bridge-exponents}, recorded in its Part~0 dictionary
(its table \texttt{p0:tab:dictionary}, row ``finite products'').  What it
proves, and what becomes of it at a dyadic point:

\begin{center}
\small
\begin{tabular}{@{}>{\raggedright\arraybackslash}p{0.47\textwidth}>{\raggedright\arraybackslash}p{0.47\textwidth}@{}}
\toprule
Exponents volume (label cited by text) & At a dyadic point of depth $s$ (this volume)\\
\midrule
\texttt{p3:eq:prefix-product}: the prefix product $\Phi_{n}=\prod_{j<n}\SincRad(\pi\xi/2^{j})$ and $p_{n}$ its inverse transform, of degree at most $n-1$
& \cref{due:def:set-finite-spline}: the same object as $\upn{n-1}$; the Thue--Morse formula \cref{due:thm:set-finite-spline} with its genuine knots\\[0.3em]
\texttt{p3:thm:sharp-prefix-law}: $\Norm{p_{n}^{(r)}-\RvachevUp^{(r)}}_{\infty}=2^{\binom{r+3}{2}-1}/(9\cdot4^{n})$ for $n\ge r+3$
& attained at every dyadic point of depth $r+2$, for every $n\ge r+3$, with the one-mode exact law (\cref{due:prop:sharp-attained})\\[0.3em]
\texttt{p3:cor:kolmogorov}: the CDF discrepancy is exactly $1/(9\cdot4^{n})$, at $x=\tfrac12$
& the $s=2$ density law at $x=\tfrac14$ transported by the finite functional equation (\cref{due:prop:sharp-cdf})\\[0.3em]
\texttt{p3:eq:richardson-\allowbreak combination}, \texttt{-weights}, \texttt{-moments}: $R_{n,M}=\sum_{j<M}w_{M,j}p_{n+j}$ with the quarter-base Lagrange weights
& the same weights, $w_{M,j}=\Wrow{M-1}{j}$ (\cref{due:thm:ext-weights}); the same moment identities (\cref{due:lem:ext-moments})\\[0.3em]
\texttt{p3:eq:richardson-\allowbreak leading}: $R^{(r)}_{n,M}=\RvachevUp^{(r)}-\am{M}q^{\binom M2}\delta_{n}^{2M}\RvachevUp^{(r+2M)}+O(\delta_{n}^{2M+2})$, asymptotic at a general point
& exact, with zero remainder, for $n\ge s+1$ and $M\ge\Floor{s/2}+1$ (\cref{due:cor:sharp-richardson}); the leading coefficient vanishes because the jet terminates\\[0.3em]
\texttt{p3:eq:richardson-\allowbreak stability}: $\sum_{j}\AbsoluteValue{w_{M,j}}<1.9693$
& the same Lebesgue factor, in closed form $\QPochhammer{-\Qq}{\Qq}{M-1}/\QPochhammer{\Qq}{\Qq}{M-1}$ (\cref{due:prop:stab-lebesgue})\\
\bottomrule
\end{tabular}
\end{center}

The contribution of this volume relative to Part~III is therefore a single
sentence: \emph{at a dyadic point the asymptotic expansion is finite and
exact from the level $\dep{x}$ on, so the same Richardson row is an exact
reconstruction operator}, together with the exact knowledge of when the
finite law starts (\cref{due:cor:sharp-onset}).

\subsection{The half-base row}
\label{due:sec:reg-half-base}

% ed.: S6's closing section, kept as the one place where the letter clash
% ed.: (Q as a translation parameter elsewhere in the corpus, Q = 1/4 here)
% ed.: is recorded.
The corpus also contains a general half-base Gaussian-binomial row for
shifted finite splines, with row polynomial
$\QPochhammer{z/2}{1/2}{n}/\QPochhammer{1/2}{1/2}{n}$, which annihilates
modes $2^{-rn}$.  The quarter base is the natural and minimal one for the
centred sequence of this volume because the tail is even: only even
derivative orders $2k$ survive the deconvolution, with the scale factor
$\delta_{n}^{2k}=4^{-k(n+1)}$, so the modes are $4^{-kn}$ and the
quarter-base row of order $\Floor{s/2}$ uses $\Floor{s/2}+1$ samples, where
a half-base row prepared for the absent odd modes would need about twice as
many.  A notational caution: in parts of the corpus the letter $Q$ denotes
a complex translation parameter of a shifted spline; the $\Qq=\tfrac14$ of
this volume is the geometric decay ratio of the centred sequence, and
\cref{due:sec:notation} keeps the two apart.

\subsection{Lean register}
\label{due:sec:reg-lean}

Every name below was checked to resolve in the repository
\cite{proveit-primary} on the day of consolidation, namespace-aware
(the module is given in parentheses).  Three statuses are used:
\textsc{crosswalked}, a declaration states the same fact (possibly in the
corpus's CDF or Appell normalization, with the translation stated);
\textsc{neighbouring}, a declaration a reader could mistake for the
result, with the difference stated; \textsc{absent}, no declaration found.
The crosswalks stated in the body of the volume
(\cref{due:sec:set-bridge,due:sec:coefficients,due:sec:jet}, and
\cref{due:rmk:ext-lean}) are collected here.

\begin{center}
\small
\begin{longtable}{@{}>{\raggedright\arraybackslash}p{0.27\textwidth}>{\raggedright\arraybackslash}p{0.47\textwidth}>{\raggedright\arraybackslash}p{0.20\textwidth}@{}}
\caption{Formalization register.}
\label{due:tab:reg-lean}\\
\toprule
Result in this volume & Lean declaration(s) & Status\\
\midrule
\endfirsthead
\toprule
Result in this volume & Lean declaration(s) & Status\\
\midrule
\endhead
\cref{due:def:set-finite-spline}, \cref{due:thm:set-finite-spline}: the finite spline and its Thue--Morse truncated-power form
& \texttt{Fabius.\allowbreak fabiusUniformSpline} (\texttt{FabiusUniformSpline}), defined \emph{as} the truncated-power sum in the centred-CDF variable; \texttt{Fabius.ProbabilityRepresentation.\allowbreak fabiusUniformSpline\_eq\_centeredPartialCDF} identifies it with the partial-sum distribution function
& \textsc{crosswalked} (CDF form; the density $\upn{n}$ is its derivative, \cref{due:sec:set-bridge})\\[0.3em]
\cref{due:lem:set-box-transform}, \cref{due:prop:set-convolution-model}: the transform of the finite product and the random-sum model
& \texttt{Fabius.\allowbreak centeredSincPartialProduct\_\allowbreak dyadic\_\allowbreak eq\_\allowbreak thueMorse} (\texttt{CenteredRvachevThueMorseFourier}): the denominator-cleared centred sinc prefix equals the Thue--Morse sine polynomial; \texttt{Fabius.ProbabilityRepresentation.\allowbreak rvachevUp\_eq\_weightedSum\_probability} and \texttt{\ldots geometricUniformDensity\_\allowbreak one\_\allowbreak half\_\allowbreak eq\_\allowbreak rvachevUp}: $\RvachevUp$ is the density of $\sum_{k}2^{-k}V_{k}$
& \textsc{crosswalked} (the infinite sum); the finite head--tail split $Z=Z_{n}+\delta_{n}Z'$ of \cref{due:prop:set-head-tail}: \textsc{absent}\\[0.3em]
\cref{due:lem:set-knots}, \cref{due:lem:set-midpoint}: knot lattice, genuine knots, cell centres
& --- & \textsc{absent}\\[0.3em]
\cref{due:def:mech-am}, \cref{due:thm:mech-am-rational-positive}(a),(b): the coefficients $\am{m}$, their recurrence and Bell form
& \texttt{Fabius.\allowbreak rvachevReciprocalMomentRat}, \texttt{\ldots\_\allowbreak eq\_\allowbreak completeBellPolynomial}, \texttt{Fabius.\allowbreak binomialConv\_\allowbreak rvachevRawMomentRat\_\allowbreak reciprocal} (\texttt{RvachevMomentAppell}): $\am{m}=(-1)^{m}\gamma^{\mathrm{Rv}}_{2m}/(2m)!$ is Lean-computable
& \textsc{crosswalked}; positivity (c) and growth (d): \textsc{absent}\\[0.3em]
\cref{due:lem:mech-finite-inverse}, \cref{due:thm:mech-local-deconvolution}: finite polynomial deconvolution
& \texttt{Fabius.\allowbreak rvachevDeconvolvedPolynomial}, \texttt{Fabius.\allowbreak rvachevAppellPolynomialRat}, \texttt{Fabius.\allowbreak integral\_\allowbreak eval\_\allowbreak rvachevAppellPolynomial\_\allowbreak sub\_\allowbreak mul\_\allowbreak rvachev} (\texttt{RvachevMomentAppell}), in the Appell normalization at scale $1$
& \textsc{crosswalked} at scale $1$; the dilation to $\delta_{n}$ and the cell identity: \textsc{absent}\\[0.3em]
\cref{due:thm:mech-stencil}: Thue--Morse derivative stencil
& \texttt{Fabius.\allowbreak iteratedDeriv\_\allowbreak rvachevUp\_\allowbreak eq\_\allowbreak thueMorse\_\allowbreak sum} (\texttt{Differential}); Rvachev's equation itself is \texttt{Fabius.deriv\_rvachevUp} (same module)
& \textsc{crosswalked}\\[0.3em]
\cref{due:thm:mech-jet-termination}(i),(ii): the jet stops at the depth; the top derivative
& \texttt{Fabius.\allowbreak iteratedDeriv\_\allowbreak rvachevUp\_\allowbreak dyadic\_\allowbreak eq\_\allowbreak zero}, \texttt{\ldots\_dyadic\_critical}, \texttt{Fabius.\allowbreak dyadic\_\allowbreak depth\_\allowbreak eq\_\allowbreak max\_\allowbreak nonzero\_\allowbreak iteratedDeriv} (\texttt{DyadicDerivativeFiltration}); \texttt{Fabius.\allowbreak iteratedDeriv\_\allowbreak rvachev\_\allowbreak centeredDyadic\_\allowbreak eq\_\allowbreak zero}, \texttt{\ldots\_ne\_zero} (\texttt{OriginalPaperSupplement}); \texttt{Fabius.\allowbreak rvachevDyadicTaylorPolynomial\_\allowbreak natDegree\_\allowbreak eq}
& \textsc{crosswalked}; part (iii) for $0<r<s$: \textsc{absent} as a separate statement\\[0.3em]
\cref{due:cor:mech-even-jet}, \cref{due:eq:mech-top-even}: the even jet and the sharp amplitude on the matched grid
& \texttt{Fabius.\allowbreak iteratedDeriv\_\allowbreak fabiusReal\_\allowbreak dyadicGrid\_\allowbreak eq\_\allowbreak ite}, \texttt{\ldots\_eq\_zero\_iff}, \texttt{Fabius.\allowbreak abs\_\allowbreak iteratedDeriv\_\allowbreak fabiusReal\_\allowbreak dyadicGrid\_\allowbreak of\_\allowbreak odd} (\texttt{BoundedDerivatives}): on the mesh $m/2^{k}$ the $k$-th derivative vanishes at even $m$ and has absolute value $2^{\binom{k+1}{2}}$ at odd $m$
& \textsc{crosswalked} (CDF form)\\[0.3em]
\cref{due:cor:main-derivative} at $r=s$: $\upn{n}^{(s)}(x)=\RvachevUp^{(s)}(x)$ for all $n\ge s$
& \texttt{Fabius.\allowbreak iteratedDeriv\_\allowbreak fabiusUniformSpline\_\allowbreak dyadic\_\allowbreak center} (\texttt{PlateauLimit}): the $r$-th derivative of the level-$p$ spline at the anchor $2^{-r}$ equals $2^{\binom{r+1}{2}}$ for every $p\ge r$
& \textsc{crosswalked} (one anchor per depth, CDF form)\\[0.3em]
\cref{due:thm:set-main} at $x=\tfrac34$ ($s=2$, one mode), CDF form
& \texttt{Fabius.\allowbreak fabiusUniformSpline\_\allowbreak quarter\_\allowbreak eq\_\allowbreak quadratic}, \texttt{Fabius.\allowbreak reportFiniteFabiusApproximant\_\allowbreak quarter\_\allowbreak eq\_\allowbreak quadratic} (\texttt{QuarterSplineLocalPolynomial}): the level-$n$ centred CDF near its quarter point is $\tfrac5{72}+z+4z^{2}-\tfrac49 4^{-n}$
& \textsc{crosswalked instance}: kernel-verified for one point and one mode; the statement for general depth: \textsc{absent}\\[0.3em]
\cref{due:thm:set-main}, \cref{due:cor:mech-exact-modes}: the exact eventual law at general depth
& --- & \textsc{absent}\\[0.3em]
\cref{due:thm:sharp-knot-defect}, \cref{due:lem:sharp-half-mgf}, \cref{due:cor:sharp-onset}
& --- & \textsc{absent}\\[0.3em]
\cref{due:def:ext-weights}, \cref{due:thm:ext-weights}: the geometric Lagrange weights and their closed forms
& \texttt{Fabius.\allowbreak geometricLagrangeWeight}, \texttt{\ldots\_eq\_product} (\texttt{GeometricLagrangeWeights}); \texttt{\ldots\_\allowbreak eq\_\allowbreak geometricQPochhammer}, \texttt{Fabius.\allowbreak quarterLagrangeWeight\_\allowbreak eq\_\allowbreak qPochhammer}, \texttt{\ldots\_eq\_qBinomial}, \texttt{Fabius.\allowbreak quarter\_\allowbreak qPochhammer\_\allowbreak self\_\allowbreak ne\_\allowbreak zero} (\texttt{GeometricLagrangeQBinomial})
& \textsc{crosswalked}\\[0.3em]
\cref{due:lem:ext-moments}, \cref{due:eq:ext-uncancelled}, \cref{due:prop:ext-rowpoly}: moment identities, uncancelled moments, row polynomial
& \texttt{Fabius.\allowbreak sum\_\allowbreak quarterLagrangeWeight\_\allowbreak eq\_\allowbreak one}, \texttt{\ldots\_mul\_pow\_eq\_zero}, \texttt{Fabius.\allowbreak sum\_\allowbreak geometricLagrangeWeight\_\allowbreak mul\_\allowbreak pow\_\allowbreak card\_\allowbreak add}, \texttt{Fabius.\allowbreak sum\_\allowbreak quarterLagrangeWeight\_\allowbreak firstUncancelled}, \texttt{Fabius.\allowbreak quarterLagrangeWeightPolynomial\_\allowbreak eq\_\allowbreak forwardRichardson} (\texttt{GeometricLagrangeWeights})
& \textsc{crosswalked}\\[0.3em]
\cref{due:thm:ext-extraction}: the row applied to the up-function sequence
& \texttt{Fabius.\allowbreak quarterLagrange\_\allowbreak rvachevFourierProduct\_\allowbreak eq\_\allowbreak gaussian\_\allowbreak tsum} (\texttt{RvachevQBinomialFilter}): the row applied to the \emph{infinite} product on the Fourier side, cancelling the first $d$ even moments
& \textsc{neighbouring} (Fourier side, asymptotic content); the theorem itself: \textsc{absent}, pending the law\\[0.3em]
\cref{due:cor:sharp-richardson}: exactness of the Richardson combination
& \texttt{Fabius.\allowbreak geometricRichardsonPowerSeriesFilter\_\allowbreak coeff\_\allowbreak eq\_\allowbreak qBinomial} (\texttt{QuarterCatalanRichardson}): a formal power-series filter with the same coefficients; its header disclaims any asymptotic meaning
& \textsc{neighbouring}\\[0.3em]
\cref{due:thm:rec-inverse}, \cref{due:thm:rec-recurrence}, \cref{due:prop:stab-lebesgue}: closed inverse, annihilating recurrence, Lebesgue factor
& --- (the general-$q$ Toeplitz form of the row exists in \texttt{GeometricLagrangeQBinomial}; no statement of the inverse, the recurrence, or the absolute row sum was found)
& \textsc{absent}\\
\bottomrule
\end{longtable}
\end{center}

\subsection{What a formalization would build, in order}
\label{due:sec:reg-plan}

The general theorem is absent, and the register shows why: the corpus has
the algebraic layer (the row and its identities), the arithmetic layer
(the jet termination and the derivative amplitudes, in several forms), the
coefficients, and the deconvolution at scale $1$; what it lacks is the
analytic bridge between them.  A proof of \cref{due:thm:set-main} would
add, in this order:
\begin{sloppypar}
\begin{enumerate}[label=\textup{(F\arabic*)}]
\item\label{due:it:F-density}
 the density $\upn{n}$ on $\RealNumbers$ as the derivative of
 \texttt{fabiusUniformSpline}, or directly as the $(n+1)$-fold box
 convolution, with its knot lattice and the cell-centre lemma
 \cref{due:lem:set-midpoint} (the genuine-knot statement is the nonvanishing
 of the Thue--Morse sign, which is elementary);
\item\label{due:it:F-headtail}
 the head--tail factorization $\RvachevUp=\upn{n}\Convolution T_{\delta_{n}}$
 as an identity of measures, from
 \texttt{ProbabilityRepresentation.\allowbreak weightedSumDistribution} by splitting the
 weighted sum at index $n$;
\item\label{due:it:F-deconv}
 the finite deconvolution at scale $\delta$ on polynomials, by dilating
 \texttt{rvachevDeconvolvedPolynomial}, and the cell identity
 \cref{due:thm:mech-local-deconvolution} (a finite computation once
 \ref{due:it:F-density}--\ref{due:it:F-headtail} are available);
\item\label{due:it:F-assemble}
 the assembly with \texttt{dyadic\_\allowbreak depth\_\allowbreak eq\_\allowbreak max\_\allowbreak nonzero\_\allowbreak iteratedDeriv}
 (\cref{due:sec:main-proof}), which gives the law, and with the existing
 row identities, which gives \cref{due:thm:ext-extraction} with no new
 analysis;
\item\label{due:it:F-knot}
 for \cref{due:thm:sharp-knot-defect}: the half-range identity
 \cref{due:lem:sharp-half-mgf} from \texttt{deriv\_rvachevUp} by one
 integration by parts, and the two-piece transfer, which is
 \ref{due:it:F-deconv} with one truncated power added.
\end{enumerate}
\end{sloppypar}
Keeping \ref{due:it:F-density}--\ref{due:it:F-headtail} (measure theory)
separate from \ref{due:it:F-deconv}--\ref{due:it:F-assemble} (finite
algebra) is what makes the statements reusable for the other kernels of
\cref{due:it:open-bases}.

\begin{summarybox}[title={Section summary}]
Relative to the Exponents volume the contribution is exactness at dyadic
points, with the onset known; the quarter base is forced by the even tail.
In Lean the row and its identities, the jet termination and its amplitudes,
the coefficients and the deconvolution at scale $1$ are kernel-verified,
and the one-mode law at the quarter point is a verified instance; the
finite splines' knot geometry, the head--tail split at finite level, the
general law, and the knot-level defect are absent, and the five-step plan
of \cref{due:sec:reg-plan} is the shortest route to them.
\end{summarybox}

\clearpage
\appendix
\part*{Appendices}
\addcontentsline{toc}{part}{Appendices}

% ---- fragment 19-notation ----
\section{Notation, and a concordance with the six sources}
\label{due:sec:notation}

\subsection{The normative source}

\begin{sloppypar}
Every symbol here is defined normatively in the corpus notation catalogue:
the entries \texttt{rvachev-up}, \texttt{up-fourier-product}, \texttt{sinc},
\texttt{fabius-dyadic-cells}, \texttt{q-pochhammer},
\texttt{gaussian-binomial}, \texttt{full-moments}, \texttt{splines},
\texttt{thue-morse-block-sum} and \texttt{rvachev-\allowbreak appell-\allowbreak deconvolution}
cover the material of this volume.  The preamble defines a few document-local
shorthands, listed below; each expands to catalogue notation or names an
object this volume defines in \cref{due:sec:setting}, and none carries a
meaning the text does not state.
\end{sloppypar}

\subsection{Symbols used throughout}

\begin{longtable}{@{}l l p{0.58\textwidth}@{}}
\caption{Principal notation.}\label{due:tab:notation}\\
\toprule
Symbol & Alias & Meaning\\
\midrule
\endfirsthead
\toprule
Symbol & Alias & Meaning\\
\midrule
\endhead
\(\RvachevUp\) & --- & Rvachev's up-function: even, \(\int\RvachevUp=1\),
  support \([-1,1]\), \(\RvachevUp(x)=\FabiusBounded(1-\AbsoluteValue{x})\)\\
\(\FourierTwoPi{f}\) & --- & the \(2\pi\)-convention transform
  \(\int f(x)\EulerE^{-2\pi\ImaginaryUnit tx}\,\Differential x\)\\
\(\SincRad(z)\) & --- & \(\sin z/z\), value \(1\) at \(0\)\\
\(\PhiUp(t)\) & \verb|\PhiUp| & \(\FourierTwoPi{\RvachevUp}(t)
  =\prod_{k\ge0}\SincRad(\pi t/2^{k})\)\\
\(\Phin{n}(t)\) & \verb|\Phin| & the prefix product with factors
  \(k=0,\ldots,n\) (\(n+1\) factors)\\
\(\upn{n}\) & \verb|\upn| & the finite sinc-product spline
  \(\FourierTwoPiOperator^{-1}[\Phin{n}]\); equals \(p_{n+1}\) of the
  Exponents volume and corresponds to \texttt{fabiusUniformSpline}~\(n\)
  in Lean\\
\(\dep{x}\) & \verb|\dep| & dyadic depth: least \(s\ge0\) with
  \(2^{s}x\in\IntegerNumbers\)\\
\(d\) & --- & \(\Floor{\dep{x}/2}\), the number of geometric modes\\
\(\Qq\) & \verb|\Qq| & the ratio base \(1/4\)\\
\(\am{m}\) & \verb|\am| & coefficients of \(1/\PhiUp(z)=\sum_{m}\am{m}(2\pi
  z)^{2m}\); positive rationals\\
\(\Cr{r}{x}\) & \verb|\Cr| & the \(r\)-th mode coefficient,
  \((-1)^{r}\am{r}4^{-r}\RvachevUp^{(2r)}(x)\)\\
\(\qn{n}\) & \verb|\qn| & the sample \(\upn{n}(x)\) at the fixed point \(x\)\\
\(\Wrow{d}{j}\) & \verb|\Wrow| & the extraction weight at node \(\Qq^{j}\)\\
\(\QPochhammer{a}{q}{n}\) & --- & \(q\)-Pochhammer symbol
  \(\prod_{k=0}^{n-1}(1-aq^{k})\)\\
\(\GaussianBinomial{n}{k}{q}\) & --- & Gaussian binomial coefficient\\
\(\ThueMorseSign{k}\) & --- & the signed Thue--Morse sequence
  \((-1)^{\text{binary digit sum of }k}\)\\
\(\Kk\) & \verb|\Kk| & the coefficient field, \(\RationalNumbers\)
  throughout\\
\(\delta_{n}\) & --- & \(2^{-(n+1)}\), the half-width of the knot cells of
  \(\upn{n}\) and the scale of the omitted tail \(T_{\delta_n}\)\\
\(P_{d}\) & --- & \(\prod_{k=1}^{d}(4^{k}-1)\), the odd common denominator of the
  row (\cref{due:sec:extraction})\\
\(\rho_{n}\) & --- & residual of the annihilating recurrence at level \(n\)
  (\cref{due:sec:recovery})\\
\(M(z)\), \(M_{+}(z)\) & --- & the moment generating function
  \(\int\EulerE^{zu}\RvachevUp(u)\,\Differential u\) and its half-range
  companion \(\int_{0}^{1}\EulerE^{zu}\RvachevUp(u)\,\Differential u\)
  (\cref{due:sec:coefficients,due:sec:sharpness})\\
\(M_{j}\) & --- & the absolute moment \(\int\AbsoluteValue{u}^{j}\RvachevUp(u)\,\Differential u\)\\
\(\Delta_{s}(x)\) & --- & the knot-level defect: \(\upn{s-1}(x)\) minus the value the
  law predicts at \(n=s-1\) (\cref{due:sec:sharpness})\\
\(\beta_{s}\) & --- & \(B_{s}/s!\) for even \(s\), \(0\) for odd \(s\); \(B_{s}\) the
  Bernoulli numbers\\
\bottomrule
\end{longtable}

\subsection{Concordance with the absorbed sources}

The six reports used four letters for the depth, four for the mode count,
three for the base, three normalizations of the reciprocal coefficients, and
three names for the up-function itself.

\begin{longtable}{@{}l l l l l l@{}}
\caption{Each source's choice, and this volume's.  S1--S6 are as identified in
\cref{due:tab:provenance}.}
\label{due:tab:concordance}\\
\toprule
& up-function & depth & modes & base & reciprocal coefficients\\
\midrule
\endfirsthead
\toprule
& up-function & depth & modes & base & reciprocal coefficients\\
\midrule
\endhead
S1 & \verb|\up|  & \(r\) & \(m\) & \(\rho\) & \(\gamma_{2\ell}/(2\ell)!\)
  from \(1/M(z)\), \(M\) the MGF\\
S2 & \verb|\upf| & \(d\) & \(M\) & \(Q\) & \(a_m\)\\
S3 & \verb|\up|  & \(m\) & \(d\) & \(4^{-1}\) & \(a_j\), Bell--Bernoulli\\
S4 & \verb|\upf| & \(s\) & \(d\) & \(\rho\) & \(a_m\)\\
S5 & \verb|\rv|  & \(s\) & \(R\) & \(Q\) & \(\alpha_r\)\\
S6 & \verb|\upf| & \(s\) & \(d\) & \(\rho\) & reciprocal moments\\
\midrule
here & \(\RvachevUp\) & \(\dep{x}\) & \(d\) & \(\Qq\) & \(\am{m}\)\\
\bottomrule
\end{longtable}

Two of these differences change formulas rather than glyphs.

\begin{enumerate}
\item \textbf{Where the factor \(4^{-r}\) lives.}  Writing the law as
  \(\upn{n}(x)=\RvachevUp(x)+\sum_{r}\Cr{r}{x}4^{-rn}\) (S4, S5, S6) or as
  \(\sum_{r}(-1)^{r}\am{r}\RvachevUp^{(2r)}(x)4^{-r(n+1)}\) (S1, S2, S3) is
  the same statement; the coefficient \(\Cr{r}{x}\) absorbs one factor
  \(4^{-r}\).  This volume prints both forms once, in \cref{due:sec:main},
  and uses the first elsewhere.
\item \textbf{The sign of the coefficients.}  S1's \(\gamma_{2\ell}/(2\ell)!\)
  is \((-1)^{\ell}\am{\ell}\), so S1's mode coefficients carry no explicit
  sign while everyone else's carry \((-1)^{r}\); \cref{due:sec:coefficients}
  proves the identity.  S6's ``reciprocal moments'' are the catalogue's
  binomial-convolution reciprocal \(\gamma^{\mathrm{Rv}}_{2m}\), which the
  same section identifies with \(\am{m}\) up to the factor
  \((2\pi)^{2m}(2m)!\) and a sign.
\end{enumerate}

\begin{warningbox}[title={Index conventions that do not agree across the corpus}]
\label{due:rmk:notation-indexing}
All six sources count the finite spline by its largest factor index:
\(\upn{n}\) has the \(n+1\) factors \(k=0,\ldots,n\).  The Exponents volume
counts factors: its \(p_{n}\) has \(n\) factors, so \(\upn{n}=p_{n+1}\).  The
Lean corpus's \texttt{fabiusUniformSpline}~\(n\) is the centred CDF
counterpart of \(\upn{n}\) (its argument is the number of summands minus
one).  Every statement transported from the Exponents volume into this one
shifts its index by one; \cref{due:sec:setting} records the bridge once, and
the Exponents volume's Part~0 dictionary table (\texttt{p0:tab:dictionary})
records it from the other side.
\end{warningbox}

% ---- fragment 20-repairs ----
\section{Ledger of editorial repairs}
\label{due:sec:repairs}

Every correction this consolidation made to its sources is marked at the
point of repair in this document's own \LaTeX{} source with an
\texttt{\%~ed.:} comment.  The table below is generated from those marks, so it
cannot drift out of step with the text.  Repairs of pure notation --- a
symbol renamed to its catalogue command --- are not listed; those are recorded
once, in the concordance of \cref{due:tab:concordance}.  What is listed is
every place where a source's mathematical statement, proof, hypothesis, index
range, or cost claim was wrong, incomplete, or overstated.

Three kinds of defect recur often enough to name.  The first is an
\emph{omitted hypothesis}: a statement true under conditions the source did not
print, most often a missing empty-product convention, a missing guard-precision
requirement, or a missing restriction to \(n\ge1\).  The second is
\emph{overclaiming}: an analytic conclusion drawn from formal algebra alone,
typically a convergence or uniformity assertion that the argument does not
support.  The third is a \emph{cost or sharpness claim} asserted without the
count that would establish it.

\begin{longtable}{@{}p{0.20\textwidth} p{0.72\textwidth}@{}}
\caption{Editorial repairs, in document order.}\label{due:tab:repairs}\\
\toprule
Location & Repair\\
\midrule
\endfirsthead
\toprule
Location & Repair\\
\midrule
\endhead
\cref{due:thm:set-extraction} & the proof of the main theorem is distributed over the volume; the section references below were attached at assembly.\\
\cref{due:rmk:set-d-zero} & this section proved the row in full by geometric Lagrange interpolation, and the extraction section proves the identical theorem by the identical argument. One proof: the statement stays here, the proof lives there.\\
The result in one page & section titles below are descriptive; the orchestrator should replace them by \textbackslash cref's to the assembled parts.\\
Fourier normalization and the finite spl & the sources say "integrate n+1 times from -infinity"; the closure (a polynomial vanishing on a half-line is zero) is supplied here, as is the genuine-knot statement (3), which the sources do not state and which \textbackslash cref\{due:lem:set-midpoint\} needs for "x is a knot".\\
\cref{due:rmk:set-prouhet} & the denominator bound is new; it is what makes the phrase "exact rational arithmetic" quantitative for the verifiers.\\
Fourier normalization and the finite spl & a breakable tcolorbox placed directly after a heading inherits the heading's \textbackslash nobreak and is moved whole to the next page, leaving the heading orphaned; the lead-in paragraph above prevents that.\\
Fourier normalization and the finite spl & consolidates S2 "Index conversion to the repository convention", S4 "Factor-count and user indexing" and S6 "Exact indexing bridge"; the Lean correspondence was checked against the definition of fabiusUniformSpline and the module docstring of QuarterSplineLocalPolynomial.lean.\\
\cref{due:lem:set-knots} & the catalogue's audit rule demands invariance under (m,n) -> (2m,n+1) for any exact evaluator; none of the sources states it.\\
Dyadic depth and cell geometry & (i) merges S1 "Cell-center stabilization", S2 "A dyadic point becomes a cell midpoint", S4 "Exact midpoint alignment" and S5 "Dyadic midpoint lemma" (the latter two in factor-count indexing, N >= s+1, which is n >= s here); (ii) is S1's item 3; (iii) and the "if and only if" are added so that the onset is explained for every level below s, not only for s-1.\\
\cref{due:rmk:set-tail-scope} & S1's "Exact self-similar tail" and S4's "the tail sees only one polynomial piece" state the r = 0 case; S6 states the jet form "in the distributional sense". The differentiation under the integral for r <= n-1 and the absolute-continuity argument at r = n are supplied here, since the deconvolution section uses r = n.\\
\cref{due:rmk:mech-cumulants-probabilistic} & S1 derives the cumulants probabilistically (cumulants add under convolution and scale homogeneously); that argument is recorded here as a remark and the analytic derivation above is the proof, because it also supplies the domain of validity of every series used later.\\
\cref{due:prop:mech-qpoch} & (d) is not in any source; it quantifies the size of the modes and is the analytic reason the extraction weights of the later sections need no growth hypothesis on the coefficients.\\
\cref{due:rmk:mech-positivity-routes} & the sources disagree on which argument proves positivity; the orchestrating question was "which proof establishes positivity, not just nonvanishing". This remark records the answer exactly.\\
\cref{due:rmk:mech-six-forms} & S1, S2 and S6 differentiate the spline distributionally and restrict to r <= n; S4 and S5 differentiate the smooth kernel instead. The kernel route is adopted because it needs no distribution theory and holds for all r, which is what makes the second proof of jet termination in \textbackslash cref\{due:rmk:mech-second-proof\} available.\\
\cref{due:prop:mech-odd-onset} & S3's odd-level onset is the one genuinely stronger statement among the six; the endpoint case s = 1 is separated from it because it depends on a convention for the value of a jump.\\
\cref{due:thm:mech-jet-termination} & the two values are derived from the first mechanism at d = 0, as S3 does, rather than from a separate argument; positivity follows S4.\\
\cref{due:cor:mech-even-jet} & S6 states the cutoff for even derivatives in terms of half the depth; this corollary is the exact reconciliation, and it also carries S1's "highest even derivative is nonzero" with its magnitude.\\
\cref{due:cor:mech-exact-modes} & this equivalence is editorial. It is what makes S3's Prouhet cancellation and the other five sources' jet-termination arguments two proofs of one fact rather than two different theorems.\\
Termination of the dyadic derivative jet & rationality by the Vandermonde argument of S6, which needs no separate arithmetic theorem about up at dyadic points.\\
\cref{due:rmk:main-ingredients} & none of the six sources isolates the rationality step from the classical arithmetic of Arias de Reyna; the volume derives it from the law itself, so the volume is self-contained on this point.\\
\cref{due:rmk:main-derivative-parity} & S2 states the derivative law with the depth letter in the summation bound (its d is this volume's s) and the onset "n >= d"; S4 states it for p\_\{n+1\}; S3 states it for the centred splines. All three are the same statement; the mode count floor((s-r)/2) is S4's.\\
\cref{due:rmk:sharp-modes-sources} & S6 is the source that proves all modes nonzero; S1's "sharp maximal mode" proves only the top one, and S3's "sharp mode" corollary is the same top-mode statement. The volume's proof (cor:mech-even-jet) is S6's argument, made exact by the tiling identity of cor:mech-tiling.\\
\cref{due:lem:sharp-half-mgf} & this lemma is new to the volume; it is the one-line consequence of Rvachev's equation up'(x) = 2up(2x+1) - 2up(2x-1) that the sources' numerics were tracking without naming.\\
\cref{due:rmk:sharp-depth-one} & the s = 1 case is the convention question of prop:mech-odd-onset seen once more: the same computation at n = 0 gives -J/2.\\
\cref{due:rmk:sharp-two-knots} & the sources' onset statements (S1 "cell-center level", S2/S4/S5 n >= s or N >= s+1, S3 the odd-level sharpening) are all consistent with the corollary; none of them claimed exactness at even depth, which is new here.\\
\cref{due:rmk:sharp-half} & the cluster brief asked whether the CDF law 1/(9 4\textasciicircum n) and the density law at x = 1/2 are the same object; they are not (at x = 1/2 the density law has no mode at all), but the CDF law at 1/2 IS the density law at 1/4 through the finite functional equation, which is the precise statement.\\
\cref{due:def:ext-weights} & cross-reference to the main theorem attached at assembly.\\
\cref{due:lem:ext-qbinomial} & the six sources prove this identically (S1 App. B, S2 App. A, S3 App. A, S4 App. A, S5 App. A, S5 Thm "Universal q-Lagrange weights"); one proof.\\
The quarter-base Gaussian-binomial extra & S1 Thm 8.1, S2 Thm 7.1, S3 Thm 1.3, S4 Thm 8.1, S5 Thm 8.1, S6 Thm 1.2 are all this statement; the forward/reversed/integer/base-4 displays are reconciled here once.\\
\cref{due:ex:ext-hand} & the Thue--Morse truncated-power representation is proved in the setting part of this volume; it is quoted here in the (n+1)-factor indexing, and was re-validated in exact arithmetic (up(1/2)=1/2, up(3/4)=5/72, up(7/8)=1/288) before being printed.\\
Recovering every mode and the Romberg ta & S1 Thm "Full coefficient reconstruction", S2 Prop "Closed inverse for all modes", S3 Thm "Mode projector", S4 sec. "The Vandermonde system", S5 sec. 9, S6 sec. 5.1 -- merged into one statement. The symmetry [z\textasciicircum m] l\_j = [z\textasciicircum j] l\_m (which reconciles S1/S2's row formula with S3's operator formula) is not remarked in any source; it is proved here.\\
\cref{due:prop:rec-minimal} & all rationals below recomputed exactly from the truncated-power representation and Gauss--Jordan on the 3x3 system; they agree with S1 verification\_output.txt lines 16-17, S2 experiment\_output.txt lines 25-27, and S5's table (5/648, -248/2025).\\
\cref{due:thm:rec-tableau} & S4 Prop "Why d+1 samples are generically necessary", S6 sec. 5.5, S1 sec. 11.3. The sources prove (2) only for consecutive levels and only for linear rules; the polynomial construction below covers both.\\
\cref{due:thm:rec-tableau} & the nonvanishing of C\_d(x) for s >= 2 is S6's theorem (S3 proves the top mode only); it is proved in the jet section and used here only as a hypothesis.\\
\cref{due:ex:rec-tableau} & S1 Alg "Exact quarter-base tableau", S5 Alg "Exact q-Richardson table", S6 sec. 5.4 -- one statement; the Neville identification (2) and the explicit mode formula (3) are supplied.\\
\cref{due:thm:rec-recurrence} & recomputed exactly; row m=2 was checked constant for n = 4,5,6, which uses samples through q\_8.\\
\cref{due:rmk:rec-below-onset} & S4 sec. "Annihilating recurrences and onset diagnostics", S2 Thm "Eventual annihilating recurrence" + diagnostics, S1 Cor "Annihilating recurrence" + Remark "Why initial values can fail", S3 Thm 10.1. The sources state onset detection informally ("holds on several successive blocks"); (1) makes it a finite, provably sufficient test.\\
Recovering every mode and the Romberg ta & exact residuals computed for this remark: x = 7/8, 3/8, 1/8 (s = 3), 3/32, 1/32 (s = 5): rho\_\{s-1\} = 0, rho\_\{s-2\} != 0; x = 15/16, 1/16 (s = 4): rho\_\{s-1\} != 0. Cross-references attached at assembly; the closed form of rho\_\{s-1\} is the sharpness section's knot-defect theorem.\\
Stability of the extraction & S1 Prop "Exact l-infinity noise amplification", S2 sec. 9, S3 Prop "Lebesgue factor", S4 Prop "Uniform weight bound", S5 sec. 9.2, S6 sec. 5.5 -- one statement. Monotonicity and the limit are proved (the sources assert them); the constant is recomputed to 46 digits from the partial products (400 factors; the tail is below 4\textasciicircum \{-400\}).\\
\cref{due:prop:stab-vandermonde} & the cancellation estimate is an upper bound on the possible loss; the sources (S3 sec. 11.3, S4 sec. 11.4) only say "exact input remains preferable". The Thue--Morse representation is quoted through cor:ext-explicit, which cites thm:set-finite-spline.\\
Stability of the extraction & the sources only state the l-infinity factor of the first row; the closed forms (2)-(3) and the table are new, verified in exact arithmetic for d <= 5.\\
\cref{due:prop:exr-symmetry} & S1's tables print the mode coefficients in the (n+1)-convention (its A\_ell = C\_ell 4\textasciicircum ell); S4's table prints C\_m as here; S6 prints A\_k = C\_k 4\textasciicircum k with the argument folded to a = 1 - |x|. All numbers below are C\_m in the volume's convention, with the source form noted once at x = 1/16.\\
\cref{due:ex:exr-quarter} & S4's "Symmetry" subsection states the reflection for its examples; the proof through the finite functional equation is supplied here.\\
Worked exact examples & 561842749440 = 2\textasciicircum 21 * 3\textasciicircum 5 * 5 * 7\textasciicircum 2 * ... is not worth printing in full; the prose stops at the power of two.\\
An exact algorithm and its verification & the six pseudocode listings (S1, S3, S4 in Python; S2, S5, S6 in prose) are the same seven lines; one listing is kept.\\
\cref{due:prop:cons-principle} & S1's "reusable local principle", S2's "general local-deconvolution principle" and S4's "general geometric scales" are the same statement; one proposition with a proof.\\
Consequences and extensions & the last clause is the converse direction; the sources assert it without argument. The argument: a cell centre of up\_n for all large n is an odd multiple of 2\textasciicircum \{-(n+1)\} for all large n, hence dyadic. It is given here in one sentence and not as a theorem, because the step "finite law implies cell centre" is only sketched.\\
Relation to the corpus and formalization & S6's closing section, kept as the one place where the letter clash (Q as a translation parameter elsewhere in the corpus, Q = 1/4 here) is recorded.\\
\bottomrule
\end{longtable}

\begin{remark}[What is not in this ledger]
\label{due:rmk:repairs-scope}
Three classes of change are deliberately excluded.  Deduplication is not a
repair: where several sources proved the same result, the strongest statement
and the best proof were kept and the rest collapsed, which the provenance
section records in aggregate rather than line by line.  Restructuring is not a
repair.  And a source claim that this volume could neither prove nor refute has
not been silently dropped: it is stated here as an explicitly labelled open
problem or conjecture, with the source's assertion quoted as an assertion.
\end{remark}

% ---- fragment 98-provenance ----
\section{Provenance of the consolidation}
\label{due:sec:provenance}

\subsection{What was merged}

This volume is an editorial consolidation, carried out on 3~September 2026, of
six independently written reports that arrived as bare directories on
2~September 2026 in the intake commit
\texttt{8f822212d} and were filed the same day under
\texttt{inverse-\allowbreak and-\allowbreak sampling/dyadic-\allowbreak up-\allowbreak extraction/} as standalone quick-gate
receipts.  No archive or checksum ledger accompanied them.  Every one of the
six proves the same theorem --- that at a dyadic point the finite
sinc-product spline is, after finitely many levels, \emph{exactly} the
up-function value plus finitely many geometric modes of ratio \(4^{-1},
4^{-2},\ldots\) --- and every one derives the same quarter-base
Gaussian-binomial extraction row from it.  They differ in normalization,
indexing, the letters used for the two key integers, and in the layers each
adds around the shared core.

\emph{Editorial} merge means: one statement of each result, in the strongest
form the six jointly support and that we can prove; the best proof kept and
gaps filled; the six normalizations and four index letters reconciled once and
recorded; errors corrected and marked in the source with \verb|% ed.:|
comments; and every layer that is not redundant preserved.  Repetition was
collapsed, content was not.

\subsection{Absorbed sources and their receipts}

\begin{table}[htbp]
\centering
\small
\caption{The six absorbed reports as filed on 2~September 2026 and as merged
(no byte changed between filing and merging).  Page counts are of the retained
arrival PDFs, which were not rebuilt.}
\label{due:tab:provenance}
\begin{tabularx}{\textwidth}{@{}lrrX@{}}
\toprule
Source & Lines & Bytes & Source SHA-256 \\
\midrule
\multicolumn{4}{@{}>{\raggedright\arraybackslash}p{\dimexpr\textwidth-2\tabcolsep\relax}@{}}{\textbf{S1} \(=\) \texttt{Dyadic-\allowbreak Up-\allowbreak Extraction/dyadic\_\allowbreak up\_\allowbreak extraction.\allowbreak tex} (22 A4 pp.)}\\
& 1{,}725 & 58{,}695 & \texttt{9ab6bd6d04f2725f\allowbreak af1b96a94c8b64d7\allowbreak 1e78c053665196ad\allowbreak 7a4d4cdd126177ff}\\
\addlinespace
\multicolumn{4}{@{}>{\raggedright\arraybackslash}p{\dimexpr\textwidth-2\tabcolsep\relax}@{}}{\textbf{S2} \(=\) \texttt{Exact\_\allowbreak Dyadic\_\allowbreak Up\_\allowbreak Extraction/dyadic\_\allowbreak up\_\allowbreak extraction.\allowbreak tex} (20 Letter pp.)}\\
& 1{,}277 & 46{,}386 & \texttt{80d2feedf09875a2\allowbreak 278ad0310edac7e6\allowbreak 3c65d483ccace42a\allowbreak 8fb2f28528381e46}\\
\addlinespace
\multicolumn{4}{@{}>{\raggedright\arraybackslash}p{\dimexpr\textwidth-2\tabcolsep\relax}@{}}{\textbf{S3} \(=\) \texttt{Exact\_\allowbreak Geometric\_\allowbreak Tails\_\allowbreak Rvachev\_\allowbreak Up/up\_\allowbreak dyadic\_\allowbreak extraction.\allowbreak tex} (21 A4 pp.)}\\
& 1{,}530 & 54{,}622 & \texttt{8669456e7848e9a1\allowbreak 91daf358f716afff\allowbreak fa02747895f45654\allowbreak 3478e7a1dafd8581}\\
\addlinespace
\multicolumn{4}{@{}>{\raggedright\arraybackslash}p{\dimexpr\textwidth-2\tabcolsep\relax}@{}}{\textbf{S4} \(=\) \texttt{dyadic\_\allowbreak up\_\allowbreak extraction\_\allowbreak package/dyadic\_\allowbreak up\_\allowbreak extraction.\allowbreak tex} (26 A4 pp.)}\\
& 1{,}801 & 60{,}526 & \texttt{9840d5381484887474\allowbreak 6e7f82f6a9dc2c3d0\allowbreak aedfb542d1f1b8ac30\allowbreak d81a86251819}\\
\addlinespace
\multicolumn{4}{@{}>{\raggedright\arraybackslash}p{\dimexpr\textwidth-2\tabcolsep\relax}@{}}{\textbf{S5} \(=\) \texttt{rvachev\_\allowbreak q\_\allowbreak extrapolation/rvachev\_\allowbreak q\_\allowbreak extrapolation.\allowbreak tex} (19 A4 pp.)}\\
& 1{,}334 & 41{,}936 & \texttt{c692ae4d4b2db6bc\allowbreak 18fa1e0be43b3e21\allowbreak 4dee0c6f193ab6be\allowbreak b9536bbe92839108}\\
\addlinespace
\multicolumn{4}{@{}>{\raggedright\arraybackslash}p{\dimexpr\textwidth-2\tabcolsep\relax}@{}}{\textbf{S6} \(=\) \texttt{rvachev\_\allowbreak up\_\allowbreak dyadic\_\allowbreak extrapolation\_\allowbreak package/rvachev\_\allowbreak up\_\allowbreak dyadic\_\allowbreak extrapolation.\allowbreak tex} (15 A4 pp.)}\\
& 1{,}143 & 36{,}098 & \texttt{e100dec4adc34acd\allowbreak ce1b71af1446e236\allowbreak 96ef75efd6e98251\allowbreak cd46123045e912dc}\\
\bottomrule
\end{tabularx}
\end{table}

Each package also shipped an exact-arithmetic verifier in Python and its
captured output; S4 shipped two figures and two CSV tables as well.  The six
verifiers were consolidated into the single script named in
\cref{due:sec:algorithm}; the figures were not retained, since every quantity
they plot is printed exactly in \cref{due:sec:examples}.  Git history is the
archive for all of it.

\subsection{What each source contributed}

No source was a superset of the others.

\begin{itemize}
\item \textbf{S1} contributed the knot-geometry explanation of the onset
  (\(x\) becomes a cell centre exactly at level \(\dep{x}\), and is a knot one
  level earlier), the Thue--Morse derivative stencil, the
  \(q\)-Pochhammer closed form of the reciprocal-product coefficients, the
  exact \(\ell^{\infty}\) noise-amplification factor, and the reconstruction
  of every mode coefficient.  It alone used the moment-generating-function
  normalization \(\gamma_{2\ell}/(2\ell)!\), reconciled in
  \cref{due:sec:coefficients}.
\item \textbf{S2} contributed the cleanest statement and the depth
  \(\dep{x}\) as the organizing integer, the terminating Taylor jet, the
  closed inverse for all modes, and the explicit index conversion to the
  repository convention.
\item \textbf{S3} contributed the Prouhet-cancellation route to the local
  degree collapse, the Bell--Bernoulli form of the coefficients, the
  \emph{odd-level onset} (\(n=\dep{x}-1\) already suffices when the depth is
  odd, with the symmetric-rectangle convention in the exceptional case
  \(\dep{x}=1\)), the Lebesgue factor, the mode projector with its even-jet
  recovery, and a verification of 2{,}376 identities at 255 interior dyadic
  points through level 7.
\item \textbf{S4} contributed the factor-count versus user-indexing
  discussion, head--tail self-similarity with strict positivity, the
  derivative-tiling proof of jet termination, the integer base-4 and
  common-denominator forms of the extraction row, the uniform weight bound,
  the argument that \(d+1\) samples are generically necessary, onset
  diagnostics through the annihilating recurrence, and the largest
  verification (2{,}570 window identities at 257 points, 941 fitted-versus-
  analytic coefficient checks, 1{,}542 recurrence checks).
\item \textbf{S5} contributed the one-page statement, the jet transfer at a
  cell centre, the rational tail generating function, the universality of the
  \(q\)-Lagrange weights, the Romberg-style \(q\)-Richardson tableau, and the
  coefficient table through \(\am{5}\).
\item \textbf{S6} contributed the proof that \emph{all} \(d\) modes are
  nonzero for depth at least~2, so that \(d\) is the exact number of geometric
  modes, together with the indexing bridge to the centred Fabius splines and
  the relation to the repository's half-base \(q\)-row.
\end{itemize}

\subsection{What the consolidation added}

Beyond the merge, the following results are new to the volume; none of them
is in any of the six sources.
\begin{itemize}
\item The closed form of the defect at the knot level,
  \cref{due:thm:sharp-knot-defect}: at every dyadic point of depth \(s\ge2\)
  the sample at \(n=s-1\) differs from the law's prediction by exactly
  \(-\ThueMorseSign{k}2^{-\binom s2}B_{s}/s!\) (zero for odd \(s\)).  It rests
  on the half-range identity \(zM_{+}(z)=\EulerE^{z/2}M(z/2)-1\) of
  \cref{due:lem:sharp-half-mgf}, a one-line consequence of Rvachev's
  equation, and it settles what the sources had only verified: the onset is
  exact at even depth (\cref{due:cor:sharp-onset}).  The formula was checked
  in exact arithmetic at all \(508\) interior points of depth \(2\) to \(8\), and
  the retained verifier checks it through depth \(7\).
\item The reflection identity \(\upn{n}(x)+\upn{n}(1-x)=1\) with its proof
  through the finite functional equation
  \(\upn{n+1}(x)=\int_{2x-1}^{2x+1}\upn{n}\) (\cref{due:prop:exr-symmetry},
  \cref{due:prop:sharp-cdf}); the sources used the symmetry of their examples
  without proving it.
\item The identification of the Exponents volume's exact CDF law at
  \(x=\tfrac12\) with the \(s=2\) density law at \(x=\tfrac14\), and the
  attainment of its sharp \(C^{r}\) law at every dyadic point of depth
  \(r+2\) (\cref{due:prop:sharp-cdf,due:prop:sharp-attained}).
\item Rationality of the whole even jet derived from the law itself rather
  than from the classical arithmetic (\cref{due:cor:main-rational}), and a
  proof of the local deconvolution principle the sources stated as a
  heuristic (\cref{due:prop:cons-principle}).
\end{itemize}

\subsection{Conventions reconciled}

Four different letters were used for the depth (\(r\), \(d\), \(m\), \(s\)),
four for the mode count (\(m\), \(M\), \(d\), \(R\)), three for the ratio base
(\(\rho\), \(Q\), \(\rho\)), and three normalizations for the reciprocal
coefficients.  The concordance in \cref{due:sec:notation} maps each source's
choice to this volume's \(\dep{x}\), \(d\), \(\Qq\), \(\am{m}\).  Two
reconciliations change formulas rather than glyphs: the mode coefficient is
written here as \(\Cr{r}{x}\) with the geometric factor \(4^{-rn}\), so that a
source's \(4^{-r(n+1)}\) form corresponds to absorbing one factor \(4^{-r}\)
into the coefficient; and S1's \(\gamma_{2\ell}/(2\ell)!\) equals
\((-1)^{\ell}\am{\ell}\), which \cref{due:sec:coefficients} proves.  All six
sources agree on the indexing of the finite spline itself (\(n+1\) factors,
\(k=0,\ldots,n\)); the volume keeps it and states its relation to the corpus's
\(p_{n}\) once, in \cref{due:sec:setting}.

\subsection{Status of the claims}

Every theorem, proposition, lemma, and corollary carries a proof.  Where a
source asserted without proof, the proof is supplied or the assertion is
demoted to a labelled open problem.  Every rational number printed was checked
against the six packages' captured verification outputs, which were treated as
ground truth over the prose; each disagreement is recorded in
\cref{due:sec:repairs}.  The plan was to draft each cluster of sections and
then have it reviewed by an independent adversarial pass; the drafting of
the setting, mechanism and extraction sections was completed that way, but
the adversarial passes and the remaining drafts were cut off by a resource
limit, and the main-theorem, sharpness, examples, algorithm, consequences and
register sections were written by the consolidating editor.  In place of the
adversarial passes, the editor's own refutation pass recomputed every printed
rational with an independent exact evaluator, ran the retained verifier,
resolved every cited Lean name namespace-aware against the repository, and
audited every cross-reference by script; the reader should weigh the proofs
accordingly.

None of the proofs is machine-checked.  The formalization register in
\cref{due:sec:register} records which statements have kernel-verified
counterparts in the Lean corpus --- the extraction row's weights and their
moment cancellation do; the exact cell identity at \(x=1/4\) is a verified
instance of the main law in CDF form --- and which do not: the general
dyadic-depth theorem is absent from Lean.

\clearpage
\begin{thebibliography}{99}
\addcontentsline{toc}{section}{References}

\bibitem{arias1982}
J.~Arias de Reyna,
``Definici\'on y estudio de una funci\'on indefinidamente diferenciable de
soporte compacto,''
\emph{Revista de la Real Academia de Ciencias Exactas, F\'isicas y Naturales
de Madrid} \textbf{76} (1982), no.~1, 21--38;
English translation, ``An infinitely differentiable function with compact
support: definition and properties,'' arXiv:1702.05442 (2017).

\bibitem{arias2018}
J.~Arias de Reyna,
``Arithmetic of the Fabius function,''
\emph{Integers} \textbf{18} (2018), article A51;
arXiv:1702.06487.

\bibitem{brezinski-redivo-zaglia}
C.~Brezinski and M.~Redivo-Zaglia,
\emph{Extrapolation Methods: Theory and Practice},
Studies in Computational Mathematics, vol.~2, North-Holland, Amsterdam, 1991.

\bibitem{deboor2001}
C.~de Boor,
\emph{A Practical Guide to Splines}, revised edition,
Applied Mathematical Sciences, vol.~27, Springer, New York, 2001.

\bibitem{dlmf-q}
NIST Digital Library of Mathematical Functions,
Chapter~17, ``\(q\)-Hypergeometric and related functions,'' in particular
\S17.2 on \(q\)-Pochhammer symbols, Gaussian binomial coefficients, and the
finite \(q\)-binomial theorem,
F.~W.~J. Olver et al., eds., \url{https://dlmf.nist.gov/17.2}.

\bibitem{fabius1966}
J.~Fabius,
``A probabilistic example of a nowhere analytic \(C^{\infty}\)-function,''
\emph{Zeitschrift f\"ur Wahrscheinlichkeitstheorie und Verwandte Gebiete}
\textbf{5} (1966), 173--174.

\bibitem{gasper-rahman}
G.~Gasper and M.~Rahman,
\emph{Basic Hypergeometric Series}, second edition,
Encyclopedia of Mathematics and its Applications, vol.~96,
Cambridge University Press, Cambridge, 2004.

\bibitem{haugland2016}
J.~K. Haugland,
``Evaluating the Fabius function,''
arXiv:1609.07999 (2016).

\bibitem{proveit-exponents}
V.~Reshetnikov,
\emph{Geometric \(q\)-Fabius Frontiers} (the consolidated exponents and
\(q\)-series volume; Part~III: finite dyadic sinc products and
piecewise-polynomial approximants, with the sharp \(C^{r}\) prefix law and the
signed quarter-base Richardson acceleration),
\texttt{ProveIt} repository,
\url{https://github.com/VladimirReshetnikov/ProveIt/tree/main/Analysis/FabiusFunction/docs/semi-formalized-research-frontiers/drafts/exponents-and-q-series}.

\bibitem{proveit-primary}
V.~Reshetnikov,
\emph{The Fabius Function and Rvachev's Up-Function: Exact Arithmetic,
Probability, Fourier Analysis, and Corrected Sharp Asymptotics},
\texttt{ProveIt} repository documentation,
\url{https://github.com/VladimirReshetnikov/ProveIt/tree/main/Analysis/FabiusFunction/docs/Fabius_Function_and_Rvachev_Up}.

\bibitem{richardson1911}
L.~F. Richardson,
``The approximate arithmetical solution by finite differences of physical
problems involving differential equations, with an application to the
stresses in a masonry dam,''
\emph{Philosophical Transactions of the Royal Society of London, Series~A}
\textbf{210} (1911), 307--357.

\bibitem{rvachev1990}
V.~A. Rvachev,
``Compactly supported solutions of functional-differential equations and
their applications,''
\emph{Russian Mathematical Surveys} \textbf{45} (1990), no.~1, 87--120.

\end{thebibliography}

\end{document}
